% !TEX program = latexmk
\documentclass[uplatex]{jsarticle}
\usepackage{amsthm,amssymb,amsmath,url,mathrsfs,bm}
\usepackage{xparse}
\usepackage{ascmac}
\usepackage[dvipdfmx]{graphicx}
\usepackage[dvipdfmx]{xcolor}
\usepackage{tikz}
\usetikzlibrary{arrows, arrows.meta, positioning, calc, fit, shapes, backgrounds}
\usepackage{float}
\usepackage{array}
\usepackage{tabularx}
\usepackage{longtable}
\usepackage{multirow}
\usepackage[dvipdfmx,unicode=true,pdfencoding=auto]{hyperref}
\usepackage{pxjahyper}

\pdfstringdefDisableCommands{%
  \def\ref#1{ref}%
  \def\label#1{}%
}

\newcounter{thm}[section]
\renewcommand{\thethm}{\thesection.\arabic{thm}}
\newcommand{\boxedthmtitle}[3]{%
  \IfNoValueTF{#3}{#1~#2}{#1~#2（#3）}%
}
\makeatletter
\newcommand{\beginboxedthm}[2]{%
  \refstepcounter{thm}%
  \IfNoValueTF{#2}{\def\@currentlabelname{}}{\def\@currentlabelname{#2}}%
  \par\medskip\noindent
  \begin{itembox}[l]{\boxedthmtitle{#1}{\thethm}{#2}}%
}
\makeatother
\newcommand{\finishboxedthm}{%
  \end{itembox}%
  \par\medskip
}
\NewDocumentEnvironment{thm}{o}{\beginboxedthm{定理}{#1}}{\finishboxedthm}
\NewDocumentEnvironment{prop}{o}{\beginboxedthm{命題}{#1}}{\finishboxedthm}
\NewDocumentEnvironment{dfn}{o}{\beginboxedthm{定義}{#1}}{\finishboxedthm}
\NewDocumentEnvironment{lem}{o}{\beginboxedthm{補題}{#1}}{\finishboxedthm}
\NewDocumentEnvironment{exa}{o}{\beginboxedthm{例}{#1}}{\finishboxedthm}
\NewDocumentEnvironment{quiz}{o}{\beginboxedthm{問}{#1}}{\finishboxedthm}
\NewDocumentEnvironment{nota}{o}{\beginboxedthm{記号}{#1}}{\finishboxedthm}
\NewDocumentEnvironment{cor}{o}{\beginboxedthm{系}{#1}}{\finishboxedthm}
\NewDocumentEnvironment{rem}{o}{\beginboxedthm{注意}{#1}}{\finishboxedthm}
\renewcommand{\labelenumi}{(\arabic{enumi})}
\setcounter{tocdepth}{3}

\title{はじめてのデジタルアイデンティティ ノート\\[0.5em]\large 認証・認可・ID管理の定義と主張の整理}
\author{\thanks{この資料は、いとう りょう『はじめてのデジタルアイデンティティ』（技術評論社, 2026）の学習目的のために作成されたものです。許可なく他の場所での再配布・転載などはしないでください。}}
\date{\today}
\begin{document}

%% =======================================================
%%  part: front
%% =======================================================
\maketitle
\tableofcontents
\newpage

%% =======================================================
%%  はじめに（原著の前書き）
%% =======================================================
\section*{このノートについて}
\addcontentsline{toc}{section}{このノートについて}

このノートは、いとう りょう『はじめてのデジタルアイデンティティ―Webサービスに欠かせない認証・認可・ID管理』\cite{Ito2026} の内容を、数学のノートのように「定義」「命題」「例」「注意」に整理し直したものである。
原著の第1章から第6章までの本文とコラムをすべて含め、章立て（第 $N$ 章 $\to$ \verb|\section|、$N$-$M$ $\to$ \verb|\subsection|）も原著に合わせている。
原著の図は画像としては収録せず、図が表している内容は地の文・表・図式（TikZ）に書き直した。
原著の注意書き、目次、謝辞、索引、奥付は収録していない。

以下の「はじめに」「想定する読者層」「本書の構成と各章の概要」は原著の前書きの内容であり、文中の「筆者」は原著者を指す。

\subsection*{はじめに}

現在、オンラインサービスでは様々な商品が購入可能となり、オフライン店舗の決済に使える決済アプリも普及している。
決済の実績に応じた報酬、C2C（個人間取引）サービスやコンテンツ配信の報酬、個人間送金の残高などを現金化できるサービスも増えている。
オンラインサービスが提供する内容の多様化、高度化に伴い、ユーザーの身元確認や当人認証といった本人確認の仕組みやユーザー情報の管理に注目が集まっている。

金融機関における口座開設や取引、携帯電話や不動産の契約を扱うサービスでは、法制度によって身分証を用いた身元確認が義務付けられている。
これまで店舗で行っていた手続きがオンラインで可能となることに伴い、身元確認もオンラインのみで完結できる仕組みが導入されている。
オンラインサービスにログインするために実施される当人認証の方式についても、脅威と対策を繰り返す形で変化が続いており、最近は長年使われてきたパスワード認証を置き換える方式として登場したパスキー認証に注目が集まっている。

オンラインサービスでは様々な種類の情報がユーザーに関連づけられており、このような属性情報の集合は\underline{デジタルアイデンティティ}（digital identity）と呼ばれている（正式な定義は第1章で与える）。
上述の本人確認の仕組みが高度化された結果、それにより取得した情報も含めたユーザー情報が、悪意のある第三者の攻撃により外部に漏洩したという事案も報告されている。
このような情報は様々な形で悪用される可能性があるため、ますますユーザー情報の管理の重要性が高まっている。

筆者は15年以上、コンシューマー向けサービスにてこれらのデジタルアイデンティティに関わる業務に従事してきた。
その中で、サービスが提供する内容に応じた本人確認やユーザー情報の管理についての要件を深く理解し、適切な設計、実装、運用を行うことの重要性、やりがい、同時に難しさを感じてきた。
デジタルアイデンティティの専門家にはセキュリティだけではなくプライバシー、関連する法制度にも深い造詣が求められる。
しかし、そのような専門家による組織を用意できるサービスはごく一部の大規模なところに限られ、そうではない規模の企業やスタートアップ企業ではデジタルアイデンティティの専任が不在だったり、経験豊富なシニアエンジニアやセキュリティの専門家が担当することもあるがデジタルアイデンティティに精通していないケースも見られる。
独学で高い専門性を身につけるためのハードルはとても高く、体系的に情報を得られるようにはなっていないのが現状である。

本書はこのような課題を持つ開発者にデジタルアイデンティティについての基本的かつ体系的な知識を提供することを目指している。
「デジタルアイデンティティとは」という概念から「実際のアプリケーションの設計、実装のポイント」までをカバーしており、非開発者やマネジメント層にも読める内容になっている。
個人の能力向上、組織内の共通認識、ジュニアエンジニアの教育などのために活用されることが期待されている。

\subsection*{想定する読者層、本書を読むことで得られること}

本書の読者層として、デジタルアイデンティティに興味がある人、業務で触れる可能性がある人を想定している。
よって、ユーザー情報を扱う全てのオンラインサービスの関係者が対象である。

\begin{itemize}
  \item 開発者は、概念、設計、実装までの一連の流れを押さえることで、現場のサービスの品質と開発速度の向上が期待できる。
  \item 開発者とともにサービスを担当するプロダクトマネージャー、企画、デザイン担当者は、概念と実際のアプリケーションの機能の関係をチーム内の共通認識にしておくことで、サービスの品質向上、安定した運用が期待できる。
  \item 経営層、マネジメント層には、サービスにおけるデジタルアイデンティティの立ち位置と重要性を理解し、自らが関わる事業の成長につなげることが期待されている。
\end{itemize}

\subsection*{本書の構成と各章の概要}

本書はこの分野についての情報を体系的にまとめることを目的としているため、冒頭から順に読み進めることが勧められている。
プロダクトマネージャー、企画、営業、デザイン、QA といった非開発者にも、第1章から第4章までの内容は重点的に読んでほしい、というのが原著者の意図である。
第5章、第6章は実際に開発する場面でとくに役立つ内容になっている。
各章の概要は次のとおりである。

\begin{description}
  \item[第1章 デジタルアイデンティティとID管理の概要]
  現実世界において私たちが身近に感じられる「アイデンティティ」の概念から始め、それがデジタル世界でどのように「デジタルアイデンティティ」として扱われるのかを解説する。
  さらに、デジタルアイデンティティとプライバシーの関係や想定しうるリスクを明らかにする。
  その上で、オンラインサービスがデジタルアイデンティティを管理する「ID管理」の基本的な考え方と構成要素について概観する。

  \item[第2章 ID管理の構成要素]
  ID管理の構成要素について、現実世界における実例と、デジタル世界における定義と概要を解説する。
  現実世界における実例でそれぞれの要素のイメージをつかみつつ、デジタル世界における定義を理解することで、オンラインサービスなどの機能との関連づけを意識する。

  \item[第3章 単一サービスにおけるID管理機能]
  あるサービス単体で提供するID管理機能について、処理内容と実装例を解説する。
  ID管理におけるアイデンティティの状態と、その状態遷移に求められるアクションを体系的に整理したアイデンティティライフサイクルを念頭に置きながら、実際に何が行われているか、どのように実現されているかを整理する。

  \item[第4章 複数サービスが関わるID管理機能]
  複数のサービスが連携して提供するID管理機能について、機能単位で関連するID管理の構成要素を解説する。
  サービス間でアイデンティティ情報をやり取りする「ID連携」を導入することで、ID管理機能がどのように変化するかを解説する。

  \item[第5章 ID管理機能の設計時に意識したいポイント]
  これまで触れたID管理機能の設計において、意識したいポイントを解説する。
  設計における重要な4つの観点、関連する各種ガイドラインや原則、広く採用されている標準化仕様を意識することで、設計精度の向上を目指す。

  \item[第6章 ID管理機能を支える技術と開発スタイル]
  ID管理機能の実装にあたり参考にできる情報を整理する。
  開発の現場で必要となる技術をあらかじめ押さえておき、いくつかの開発スタイルの特徴を意識することで、実際の開発業務に関わる際の心構えができる。
\end{description}

\newpage

%% =======================================================
%%  part: ch1
%% =======================================================
\section{デジタルアイデンティティとID管理の概要}

この章では、\cite{Ito2026} の第1章に沿って、本書全体の基礎となる概念を扱う。
前半（デジタルアイデンティティとは）では、まず現実社会の例を用いてエンティティ、属性情報、アイデンティティ、自観と他観を定義し、
それを踏まえてデジタルアイデンティティを定義し、その構成要素となる属性情報と、プライバシー侵害のリスクを現実世界と比較しながら論じる。
後半（ID管理とその構成要素）では、デジタルアイデンティティを安全に運用するための ID 管理を定義し、オンラインサービスの種類による ID 管理の違いを述べたのち、
ID 管理を構成する 5 つの要素（身元確認、当人認証、ID 連携、アクセス制御、セッション管理）の概要を定義する。
これらの構成要素は次章でより精密に扱う。

\subsection{デジタルアイデンティティとは}

デジタルアイデンティティと聞くと、オンラインサービスのアカウントやログインの仕組みを想像する人が多いであろう。
しかし、「デジタル」という修飾を外した「アイデンティティ」そのものが何であるかを先に整理しておかないと、デジタルアイデンティティの意味も曖昧になる。
そこでこの項では、まず現実社会の例を挙げながらアイデンティティとその周辺の概念を定義し、次にそれをデジタル世界に移したデジタルアイデンティティを定義し、
最後にアイデンティティとプライバシーの関係を、現実世界とデジタル世界を比較しながら論じる。

\subsubsection{現実社会におけるエンティティとアイデンティティ}

この項では、\emph{デジタル}がつかない\emph{アイデンティティ}とは何かを、現実社会の例を挙げながら整理する。
以後の章で用いるエンティティ、属性情報、アイデンティティ、自観、他観という基本用語は、すべてここで導入する。

出発点として、自己紹介という身近な場面を考える。

\begin{exa}[自己紹介]\label{exa:ch1-self-introduction}
私たちは自己紹介をする際、氏名、年齢、職業、趣味などの情報を用いて自分自身のことを説明する。
自己紹介を受けた相手は、その内容に加えて、態度や話し方の印象や容姿などの情報をもとに、その人のことを認識する。
\end{exa}

この例には、「説明する主体」「主体についての情報」「主体を認識してもらうために使われる情報のまとまり」という三つのものが現れている。
これらにそれぞれ名前を付ける。

\begin{dfn}[エンティティ]\label{dfn:entity}
例 \ref{exa:ch1-self-introduction} において自己紹介をする主体（例えば「私」）のように、
アイデンティティによって表現・認識される存在そのものを \underline{エンティティ}（実体、entity）という。
\end{dfn}

\begin{dfn}[属性情報]\label{dfn:attribute}
エンティティに関する性質や特徴を表す情報を \underline{属性情報}（attribute）という。
\end{dfn}

\begin{dfn}[アイデンティティ]\label{dfn:identity}
エンティティが「こういう人間である」と他者に提示し、認識してもらうために利用する属性情報の集合を \underline{アイデンティティ}（identity）という。
すなわち、エンティティ $e$ の属性情報全体の集合を $A_e$ とするとき、アイデンティティとは、$e$ を提示・認識するために用いられる部分集合
\[
I \subseteq A_e
\]
のことである。
\end{dfn}

エンティティとアイデンティティの関係をまとめると、エンティティという一つの実体があり、その性質や特徴を表す属性情報がいくつもあって、
そのうち他者への提示・認識に使われる属性情報を集めたものがアイデンティティである、という関係になる。図 \ref{fig:ch1-entity-identity} では、例 \ref{exa:ch1-self-introduction} の氏名、年齢、職業、趣味がそれぞれエンティティ $e$ の性質や特徴を表す属性情報であり、それらを集めたアイデンティティ $I$ が $e$ を提示・認識するのに用いられる様子を示す。

\begin{figure}[H]
\centering
\begin{tikzpicture}[
  box/.style={draw, rounded corners, align=center, minimum height=8mm, inner sep=4pt},
  >=Stealth]
\node[box] (e) {エンティティ $e$\\（例：「私」）};
\node[box, right=22mm of e] (a1) {属性情報\\（氏名）};
\node[box, below=3mm of a1] (a2) {属性情報\\（職業）};
\node[box, above=3mm of a1] (a0) {属性情報\\（年齢）};
\node[box, below=3mm of a2] (a3) {属性情報\\（趣味）};
\node[draw, dashed, rounded corners, fit=(a0)(a1)(a2)(a3), inner sep=6pt,
  label={[align=center]above:アイデンティティ $I\subseteq A_e$}] (I) {};
\draw[->] (I.west) -- node[above, align=center, font=\small]{提示・認識に\\用いる} (e);
\end{tikzpicture}
\caption{エンティティとアイデンティティの関係}\label{fig:ch1-entity-identity}
\end{figure}

現実世界のアイデンティティでは、具体的に次のような属性情報が利用される。

\begin{exa}[現実世界のアイデンティティで利用される属性情報]\label{exa:ch1-real-world-attributes}
\begin{enumerate}
\item \underline{基本4情報：}氏名、住所、生年月日、性別。
\item \underline{各種身分証明書の情報：}保険証の保険者番号、運転免許証の番号、顔写真。
\item \underline{身体的特徴：}顔の表情や特徴、身長、体型、服装、声のトーンや話し方、姿勢、ジェスチャー。
\item \underline{経歴：}学歴、職歴。
\item \underline{周囲との関わりに関する情報：}家族構成、友人関係、職場での印象など。
\end{enumerate}
\end{exa}

次に、一つのエンティティに対してアイデンティティが一つに定まるかどうかを考える。
私たちは家族、友人、職場の同僚といった様々な人間関係を持っていることに注目する。

\begin{exa}[人間関係による情報の使い分け]\label{exa:ch1-relationships}
家族や友人には話せることでも職場の同僚には控える話題がある一方で、職場の同僚と共有している情報を家族や友人には明かさない、といった違いがある。
\end{exa}

\begin{prop}[一つのエンティティは複数のアイデンティティを持ちうる]\label{prop:ch1-multiple-identities}
一つのエンティティ $e$ は、複数のアイデンティティ $I_1, I_2, \dots, I_n \subseteq A_e$ を持ちうる。
\end{prop}

\begin{proof}
アイデンティティが提示の相手ごとに異なりうることを示せばよい。
例 \ref{exa:ch1-relationships} のように、人は相手や状況に応じて会話の内容や共有する情報を変える。
すなわち、家族や友人に提示する属性情報の集合を $I_1$、職場の同僚に提示する属性情報の集合を $I_2$ とすると、
家族や友人には話すが同僚には控える情報は $I_1$ に属して $I_2$ に属さず、同僚とは共有するが家族や友人には明かさない情報は $I_2$ に属して $I_1$ に属さない。
よって一般に $I_1 \neq I_2$ である。
定義 \ref{dfn:identity} より $I_1$ も $I_2$ も同じエンティティのアイデンティティであるから、一つのエンティティが複数のアイデンティティを持ちうる。
\end{proof}

ここまでのアイデンティティは、エンティティ自身が「私はこういう人間です」と認識して提示するものであった。
しかし、例 \ref{exa:ch1-self-introduction} で見たように、相手は提示された内容だけでなく態度や話し方の印象、容姿などからもその人を認識している。
この「誰の認識に基づくか」の違いを区別するために、次の用語を導入する。

\begin{dfn}[自観と他観]\label{dfn:self-other-view}
エンティティ $e$ について、
\begin{enumerate}
\item 例 \ref{exa:ch1-self-introduction} の自己紹介のように、「私はこういう人間です」とエンティティ自身が認識して他者に提示する属性情報群、すなわち自らが捉えるアイデンティティを \underline{自観} という。
\item 他者（相手）がそのエンティティをどのように捉えているか、という認識に基づく属性情報群もまたアイデンティティであり、これを \underline{他観} という。
\end{enumerate}
以後、$e$ の自観を $I^{\mathrm{self}}_e$、ある他者から見た $e$ の他観を $I^{\mathrm{other}}_e$ と書くことがある。
\end{dfn}

\begin{rem}[自観と他観は一致するとは限らない]\label{rem:ch1-self-other-gap}
自観と他観は必ずしも一致するわけではない。すなわち一般に $I^{\mathrm{self}}_e \neq I^{\mathrm{other}}_e$ でありうる。
「自分はこう思っているが相手はそうは思っていない」といった自観と他観の間に生じる差異が、時には悩みや人間関係のトラブルに発展することもある（図 \ref{fig:ch1-self-other-view}）。
\end{rem}

\begin{figure}[H]
\centering
\begin{tikzpicture}[
  box/.style={draw, rounded corners, align=center, minimum height=8mm, inner sep=4pt},
  >=Stealth]
\node[box] (e) {エンティティ $e$\\（「私」）};
\node[box, right=40mm of e] (o) {他者\\（相手）};
\node[box, below=10mm of e] (s) {自観 $I^{\mathrm{self}}_e$\\「私はこういう人間です」};
\node[box, below=10mm of o] (t) {他観 $I^{\mathrm{other}}_e$\\相手が捉える「私」};
\draw[->] (e) -- node[above, font=\small]{提示} (o);
\draw[->] (e) -- node[left, font=\small]{自ら認識} (s);
\draw[->] (o) -- node[right, font=\small]{認識} (t);
\draw[<->, dashed] (s) -- node[below, font=\small]{一致するとは限らない} (t);
\end{tikzpicture}
\caption{自観と他観}\label{fig:ch1-self-other-view}
\end{figure}

最後に、ここまでの定義が人間に限られないことを注意しておく。

\begin{rem}[エンティティ・アイデンティティは人間以外にも適用できる]\label{rem:ch1-non-human-entity}
ここまで紹介したエンティティ、アイデンティティの関係は、人間だけに適用できるものではない。
動物、植物、人工物を問わず、全ての存在するものがエンティティであり、それを認識するための属性情報群がアイデンティティである。
したがって定義 \ref{dfn:entity} と定義 \ref{dfn:identity} の「人間」「人」は、任意の存在に読み替えてよい。
\end{rem}

\subsubsection{デジタルアイデンティティとそれを構成する属性情報}

前項の現実世界でのアイデンティティの考え方を踏まえ、この項では、オンラインサービスにおいて「自分自身」を表現する情報、すなわちデジタルアイデンティティについて詳しく見ていく。
まず定義を与え、次にそれが現代において重要視されるようになった背景をユーザーコミュニケーションの変化から説明し、
最後にデジタルアイデンティティを構成する具体的な属性情報とその利用目的を整理する。

\begin{dfn}[デジタルアイデンティティ]\label{dfn:digital-identity}
デジタル世界で特定のエンティティを表現、認識するための属性情報の集合を \underline{デジタルアイデンティティ}（digital identity）という。
すなわち、オンラインサービスにおいて「自分自身」を表現する情報のことであり、定義 \ref{dfn:identity} のアイデンティティをデジタル世界に移したものである。
\end{dfn}

このデジタルアイデンティティが現代において非常に重要視されるようになった背景には、オンラインサービスの多様化がある。
その中の一例として、ユーザーコミュニケーションの変化に注目する。

\begin{exa}[オンラインサービスのユーザーコミュニケーションの変化]\label{exa:ch1-communication-change}
かつての匿名性が高かった掲示板（BBS）では、「特定の話題」がコミュニケーションの中心であった。
ところがブログや SNS が主流になると、「誰が」発信しているのか、そして「誰に」発信しているのかというユーザー同士の関係性がより重視されるようになった。
\end{exa}

この「人物や関係性」を重視する流れの中で、オンラインサービス上で使われるアイデンティティも次のように変化した。

\begin{exa}[オンラインサービス上で使われるアイデンティティの変化]\label{exa:ch1-online-identity-change}
\begin{enumerate}
\item 当初は、完全な匿名やハンドルネームが主流であった。
\item 次第に、ニックネームやプロフィール画像など、より個人を識別しやすいデジタル上のアイデンティティが用いられるようになっていった。
\item さらには、実名、顔写真、所属組織といった現実世界のアイデンティティに近い情報を利用するサービスも増えてきた。
\end{enumerate}
\end{exa}

そして、現実世界で私たちが状況に応じてアイデンティティを使い分ける（命題 \ref{prop:ch1-multiple-identities}）のと同様に、
デジタル世界においても、利用するサービスや目的に応じて、複数のデジタルアイデンティティを使い分けることが一般的になっている。

\begin{exa}[デジタル世界でのアイデンティティの使い分け]\label{exa:ch1-digital-multiple-identities}
仕事用 SNS と趣味のコミュニティで、表示するユーザー名やプロフィール画像を変えるといった使い方がこれにあたる。
同じエンティティが、仕事用 SNS 向けのデジタルアイデンティティと、趣味のコミュニティ向けのデジタルアイデンティティを別々に持っていることになる。
\end{exa}

それでは、具体的にデジタルアイデンティティはどのような属性情報から構成されているのであろうか。主なものとして、以下が挙げられる。

\begin{exa}[デジタルアイデンティティを構成する主な属性情報]\label{exa:ch1-digital-attributes}
\begin{enumerate}
\item \underline{ユーザー識別子：}ユーザー ID、アカウント名。
\item \underline{プロフィール情報：}ニックネーム、プロフィール画像、自己紹介文。
\item \underline{連絡先：}メールアドレスや電話番号。
\item \underline{ユーザー認証に利用する情報：}パスワードやデバイス情報、生体情報（指紋認証、顔認証などに利用する情報）。
\item \underline{外部サービスとの連携情報：}他の SNS やプラットフォームのアカウントなど、外部サービスのアカウント情報。
\item \underline{行動属性、サービスの利用履歴：}購入履歴、ウェブサイト閲覧履歴、位置情報、投稿内容、「いいね」の履歴など。
\end{enumerate}
\end{exa}

\begin{rem}[属性情報の利用目的]\label{rem:ch1-attribute-uses}
例 \ref{exa:ch1-digital-attributes} の属性情報は、単に個人を識別するだけでなく、
サービス利用時のユーザー認証、コンテンツのパーソナライズ、ユーザー間の信頼関係構築、ターゲティング広告など、様々な目的で利用されている。
\end{rem}

オンラインサービスの内容によっては、デジタルアイデンティティが、現実世界のアイデンティティに含まれる情報群（例 \ref{exa:ch1-real-world-attributes}）を含む場合がある。

\begin{exa}[現実世界の属性情報を含むデジタルアイデンティティ]\label{exa:ch1-real-attributes-online}
次の現実世界の属性情報を含むデジタルアイデンティティがある。
\begin{enumerate}
\item 基本4情報（氏名、住所、生年月日、性別）。
\item 経歴。
\item 周囲との関わりに関する情報。
\end{enumerate}
これら現実世界の属性情報を使うサービスとして、金融機関や行政機関が提供するサービスがある。
\end{exa}

\begin{prop}[信頼できるデジタルアイデンティティの必要性]\label{prop:ch1-trusted-digital-identity}
金融機関や行政機関が提供するサービスでは、現実世界ですでに確認されているアイデンティティを、そのままデジタル世界でも活用するための仕組み、
すなわち信頼できるデジタルアイデンティティが必要になる。
\end{prop}

\begin{proof}
必要性の根拠を、手続きのオンライン化と本人確認の厳密さの二点から述べる。
まず、これまで窓口で行われていた手続きをオンラインで完結させる DX（デジタルトランスフォーメーション）の動きが加速している。
次に、金融機関や行政機関の手続きでは、窓口で行われていたときと同様に厳密な本人確認が必要である。
よって、窓口で確認していたのと同等の確からしさで、現実世界ですでに確認されているアイデンティティをデジタル世界でもそのまま活用できる仕組み、
すなわち信頼できるデジタルアイデンティティが必要になる。
\end{proof}

以上をまとめると、デジタルアイデンティティを構成する属性情報には、例 \ref{exa:ch1-digital-attributes} のようなデジタル世界に固有の属性情報と、
例 \ref{exa:ch1-real-attributes-online} のような現実世界のアイデンティティに由来する属性情報とがあり、サービスの内容に応じてそれらが組み合わされる。

\subsubsection{アイデンティティとプライバシー}

前項までで、デジタルアイデンティティを構成する様々な属性情報とその利用について見てきた。
この項では特に、「他者からの認識（他観）とのずれ」と「プライバシー保護」という二つの観点から、
現実世界のアイデンティティとデジタルアイデンティティの違いと注意点を掘り下げる。
前者は定義 \ref{dfn:self-other-view} の自観と他観の差異の問題であり、後者は命題 \ref{prop:ch1-multiple-identities} の複数のアイデンティティの使い分けに伴う情報の漏洩・拡散の問題である。

\underline{アイデンティティの認識差：}
まず、アイデンティティの認識差について述べる。
現実世界では、コミュニケーションを通じて相手に提示したアイデンティティ（自観）と、提示された側が受け取ったアイデンティティ（他観）の間に何らかの差異が生じることは珍しくない。
自分が意図した印象と相手の受け取り方が異なり、それが悩みの種となることもある（注意 \ref{rem:ch1-self-other-gap}）。
デジタルアイデンティティについては、次が成り立つ。

\begin{prop}[自観と他観のずれはデジタルアイデンティティでも起こりうる]\label{prop:ch1-digital-view-gap}
デジタルアイデンティティにおいても、提示された情報に対する受け取り方（他観）のずれは、現実世界と同様に、あるいはそれ以上に起こりうる。
\end{prop}

\begin{proof}
一見すると差異が生じにくいように思える理由を確認したうえで、それでもずれが生じる要因を挙げる。
まず、デジタルアイデンティティの場合、プロフィール情報などのデータ自体はそのまま伝達される。
そのため、提示した属性情報そのものが途中で変わることはなく、一見すると自観と他観の差異は生じにくいように思える。
しかし、他観は伝達されたデータそのものではなく、それに対する受け手の受け取り方に基づく。
次に、テキスト中心のコミュニケーションでは、ニュアンスの誤解が生じる。
また、リッチコンテンツを用いた場合でも、相手が得るのは断片的な情報であり、そこから相手が抱く印象には違いが生じる。
よって、提示された情報に対する受け取り方（他観）のずれは、現実世界と同様に、あるいはそれ以上に起こりうる。
\end{proof}

\underline{プライバシーの問題：}
アイデンティティの認識差に加え、もう一つ重要な論点がプライバシーの問題である。
特にデジタルアイデンティティにおいては、複数のアイデンティティを使い分ける際に、意図せずプライベートな情報が漏洩・拡散してしまうリスクが伴う。
この種の漏洩は、現実世界でも次のように起こりうる。

\begin{exa}[コミュニティをまたぐ情報の漏洩]\label{exa:ch1-community-leak}
現実世界でも、あるコミュニティ（家族、友人、職場など）の中だけで共有していたはずの情報が別のコミュニティに漏れてしまい、
「知られたくない相手に知られてしまった」という状況は起こりうる。
噂話のように情報が意図せず拡散していく可能性は常に存在する。
\end{exa}

\begin{rem}[意図しない漏洩・拡散はプライバシー侵害にあたる]\label{rem:ch1-privacy-violation}
例 \ref{exa:ch1-community-leak} のような、意図しない情報の漏洩や拡散は \underline{プライバシー侵害} にあたる。
\end{rem}

プライバシー侵害は現実世界とデジタル世界の双方で起こりうるが、その性質には、主に以下の四つの点で顕著な違いが見られる。

\underline{プライバシー侵害による被害：}
現実世界では、ストーカー被害のような身体的危害、住所の漏洩からの物理的な侵入・盗難、情報公開による近隣や職場での評判への悪影響などが考えられる。
一方、デジタルアイデンティティにはオンラインサービスの利用履歴なども含まれるため（例 \ref{exa:ch1-digital-attributes}）、その侵害内容は多岐にわたる。
例えば、決済情報の悪用による経済的被害、なりすましアカウント作成などのアイデンティティ盗用、個人情報が不特定多数に晒されることによる深刻なプライバシー侵害、
そして広範囲への拡散による社会的信用の失墜などが挙げられる。

\underline{影響範囲、スピード、持続性：}
現実世界での情報拡散は、その範囲やスピードにおいて比較的限定的である。
物理的な書類や写真の拡散には限界があり、情報の削除や回収も不可能ではない。
それに対し、デジタル世界、特にインターネット上では、情報は瞬時に、かつ広範囲（場合によっては世界規模）に拡散する。
一度流出した情報を完全に削除することは極めて困難であり（いわゆる \underline{デジタルタトゥー}）、影響が半永久的に持続する可能性も秘めている。
サービスやデータベースからの情報漏えいの場合、データ量が膨大となる可能性もある。

\underline{属性情報の取得手段：}
現実世界では、盗撮や盗聴、ストーキングや尾行、書類の盗難などがあり、物理的な行動を伴う。
デジタル世界では、ユーザーを狙うフィッシング、サービスを狙うハッキングや、両者をターゲットにするマルウェアといった攻撃や、
トラッキング技術の悪用、公開情報の悪用など、物理的な行動を伴わない手段が多く、対策も変わってくる。

\underline{加害者の特定：}
匿名性が高く特定が難しいデジタル世界と、物理的な行動を伴う現実世界とでは、加害者の特定方法についても違いがある。

以上の四つの観点を表 \ref{tab:ch1-privacy-comparison} にまとめる。

\begin{table}[H]
\centering
\caption{プライバシー侵害の性質：現実世界とデジタル世界の比較}\label{tab:ch1-privacy-comparison}
\begin{tabularx}{\linewidth}{|p{0.17\linewidth}|X|X|}
\hline
観点 & 現実世界 & デジタル世界 \\
\hline
被害 &
ストーカー被害のような身体的危害、住所の漏洩からの物理的な侵入・盗難、情報公開による近隣や職場での評判への悪影響 &
利用履歴なども含むため多岐にわたる。決済情報の悪用による経済的被害、なりすましアカウント作成などのアイデンティティ盗用、個人情報が不特定多数に晒されることによる深刻なプライバシー侵害、広範囲への拡散による社会的信用の失墜 \\
\hline
影響範囲、スピード、持続性 &
範囲・スピードは比較的限定的。物理的な書類や写真の拡散には限界があり、削除や回収も不可能ではない &
瞬時に、かつ広範囲（場合によっては世界規模）に拡散。完全な削除は極めて困難（デジタルタトゥー）で、影響が半永久的に持続しうる。漏えいデータ量が膨大になりうる \\
\hline
属性情報の取得手段 &
盗撮、盗聴、ストーキング、尾行、書類の盗難など、物理的な行動を伴う &
フィッシング（ユーザーを狙う）、ハッキング（サービスを狙う）、マルウェア（両者を狙う）、トラッキング技術の悪用、公開情報の悪用など、物理的な行動を伴わない手段が多く、対策も変わる \\
\hline
加害者の特定 &
物理的な行動を伴う &
匿名性が高く、特定が難しい \\
\hline
\end{tabularx}
\end{table}

この比較から、プライバシー侵害に対する責任について次の主張が導かれる。

\begin{prop}[オンラインリスクの管理にはユーザーとサービスの双方の取り組みが不可欠である]\label{prop:ch1-shared-responsibility}
デジタル世界では、ユーザーは情報公開やプライバシー設定に主体的な責任を持つことが求められ、
サービスは、ユーザーが情報を容易に管理・保護できるよう、分かりやすい機能と透明性の高い設計を提供する責務を負う。
オンラインリスクの管理には、この両者の取り組みが不可欠である。
\end{prop}

\begin{proof}
表 \ref{tab:ch1-privacy-comparison} の比較から、デジタル情報の性質を取り出し、そこから両者の責任を導く。
まず、プライバシー侵害は現実世界とデジタル世界の両方で起こりうるが、その手段、影響の範囲、スピード、持続性には大きな違いがある。
とくに「影響範囲、スピード、持続性」の観点から、デジタル情報は瞬時に広範囲へ拡散し、一度流出すると完全な削除が極めて困難である。
すなわちデジタル情報は「拡散しやすく削除困難」である。
したがって、流出した後に回収することに頼ることはできず、そもそも何を公開するかを決める段階での管理が重要になる。
よって、情報を公開する当事者であるユーザーは、情報公開やプライバシー設定に主体的な責任を持つことが求められる。
次に、ユーザーがその責任を果たすには、情報を管理・保護する手段がサービス側に用意されていなければならない。
よって、サービスもまた、ユーザーが情報を容易に管理・保護できるよう、分かりやすい機能と透明性の高い設計を提供する責務を負う。
以上より、オンラインリスクの管理には、ユーザーとサービスの両者の取り組みが不可欠である。
\end{proof}

\subsection{ID管理とその構成要素}

前項では、デジタルアイデンティティの概念と、その活用に伴うプライバシー侵害などのリスクへの向き合い方について見てきた。
この項では、そのリスクを低減してデジタルアイデンティティを安全に運用するための仕組みである ID 管理を定義する。
次に、オンラインサービスの種類によって求められる ID 管理のあり方がどう異なるかを述べ、
最後に、ID 管理を構成する 5 つの主要な要素の概要を定義する。

\subsubsection{アイデンティティ管理（ID管理）とは}

この項では、オンラインサービス提供者に厳格な管理体制が求められる理由を述べ、そのための一連のプロセスや技術として ID 管理を定義する。

\begin{prop}[デジタルアイデンティティの不適切な管理はリスクを高める]\label{prop:ch1-mismanagement-risk}
デジタルアイデンティティが不適切に管理されると、情報漏洩、不正アクセス、なりすましといったインシデントが発生し、
深刻なプライバシー侵害や経済的損害につながるリスクが高まる。
\end{prop}

\begin{proof}
デジタルアイデンティティに含まれる情報の性質から説明する。
デジタルアイデンティティには、氏名や連絡先といった個人情報だけでなく、サービスへのログインに必要な認証情報（パスワードなど）や、
購買履歴・閲覧履歴といった機密性の高い（センシティブな）データが含まれている（例 \ref{exa:ch1-digital-attributes}）。
まず、これらの情報が不適切に管理されると、個人情報やセンシティブなデータが外部に漏れる情報漏洩が起こりうる。
次に、認証情報が漏れれば、第三者がサービスに不正アクセスしたり、本人になりすましたりすることが可能になる。
よって、こうしたインシデントの結果として、深刻なプライバシー侵害や経済的損害につながるリスクが高まる。
\end{proof}

命題 \ref{prop:ch1-mismanagement-risk} のリスクを低減し、デジタルアイデンティティを安全に運用するために、オンラインサービス提供者には厳格な管理体制が求められる。
具体的には、ユーザーが誰であるかを正確に識別・検証し、そのユーザーに与えられた権限に基づいて、許可された情報や機能にのみアクセスできるように制御する必要がある。

\begin{dfn}[アイデンティティ管理（ID管理）]\label{dfn:identity-management}
デジタルアイデンティティ、および関連づけられたエンティティによるシステムやサービスへのアクセス権限を、適切に管理・運用するための一連のプロセスや技術を
\underline{アイデンティティ管理}（identity management）、略して \underline{ID管理} という。
\end{dfn}

\begin{rem}[適切な ID 管理の意義]\label{rem:ch1-identity-management-value}
適切な ID 管理は、ユーザーの利便性を損なうことなく、デジタルアイデンティティの安全性とプライバシー保護を両立させ、
ひいては信頼性の高いデジタル社会を支える基盤となる。
\end{rem}

\subsubsection{オンラインサービスの特性によるID管理の違い}

私たちが現実世界で人間関係によって複数のアイデンティティを使い分けているように（命題 \ref{prop:ch1-multiple-identities}）、
デジタル世界でもオンラインサービスの内容に応じて複数のデジタルアイデンティティを使い分けている（例 \ref{exa:ch1-digital-multiple-identities}）。
この項では、オンラインサービスをその目的や対象ユーザーによって分類し、それぞれで求められる ID 管理のあり方がどう異なるかを見る。

\begin{dfn}[オンラインサービスの種類]\label{dfn:ch1-service-types}
オンラインサービスは、その目的や対象ユーザーによって、大きく以下の種類に分類できる。
\begin{enumerate}
\item \underline{コンシューマー向けサービス}（B2C）：個人ユーザーを対象としたサービス。SNS や EC サイトなど。
\item \underline{エンタープライズ向けサービス}（B2B）：企業ユーザーを対象としたサービス。業務アプリケーション、グループウェア、クラウド基盤など。
\item \underline{政府提供のサービス}（G2C/G2B）：個人／企業ユーザーを対象とした公共サービス。マイナンバーポータルや確定申告用サイトなど。
\end{enumerate}
\end{dfn}

これらのサービスでは、それぞれ求められる ID 管理のあり方も異なる。その主な違いを具体的に見ていく。

\begin{exa}[コンシューマー向けサービス（B2C）の ID 管理]\label{exa:ch1-b2c}
コンシューマー向けサービス（B2C）では、基本的にユーザー（個人）自身がアカウントのライフサイクル（作成、変更、削除）を管理する。
サービス利用開始時にユーザーが自らアカウントを作成し、不要になれば自身で削除することが可能である。
近年では、プライバシー保護の観点から、ユーザー自身による詳細なプライバシー設定や、\underline{データポータビリティ}（自身のデータの持ち出し）への対応なども重要視されている。
\end{exa}

\begin{exa}[エンタープライズ向けサービス（B2B）の ID 管理]\label{exa:ch1-b2b}
エンタープライズ向けサービス（B2B）では、通常、サービスを利用する企業（組織）が、従業員のアカウントを一元的に管理する。
人事部門や情報システム部門などが、従業員の入社、異動、退職といったライフサイクルに合わせて、アカウントの発行、権限変更、停止、削除を行う。
さらに、役職や所属部署、担当業務に応じて、アクセスできる情報システムやデータ、利用可能な機能を細かく制御する厳格な権限管理が不可欠である。
\end{exa}

\begin{exa}[政府提供のサービス（G2C/G2B）の ID 管理]\label{exa:ch1-g2c}
行政機関、政府提供のサービス（G2C/G2B）では、多くの場合、既存の公的な情報基盤（住民基本台帳ネットワークシステムなど）と連携し、厳格な本人確認が求められる。
マイナンバーカードを用いた公的個人認証（JPKI）などが利用されるケースも増えている。
また、取り扱う情報の機密性が極めて高いため、非常に高度なセキュリティ基準に基づいた ID 管理が実装されている。
\end{exa}

以上の違いを表 \ref{tab:ch1-idm-by-service} にまとめる。

\begin{table}[H]
\centering
\caption{オンラインサービスの種類による ID 管理の違い}\label{tab:ch1-idm-by-service}
\begin{tabularx}{\linewidth}{|p{0.2\linewidth}|X|X|}
\hline
種類 & 対象と例 & ID 管理の主な特徴 \\
\hline
コンシューマー向け（B2C） &
個人ユーザー。SNS、EC サイトなど &
ユーザー自身がアカウントのライフサイクル（作成、変更、削除）を管理する。詳細なプライバシー設定、データポータビリティへの対応が重要視される \\
\hline
エンタープライズ向け（B2B） &
企業ユーザー。業務アプリケーション、グループウェア、クラウド基盤など &
利用企業が従業員のアカウントを一元管理する。入社・異動・退職に合わせた発行、権限変更、停止、削除。役職・所属部署・担当業務に応じた厳格な権限管理 \\
\hline
政府提供（G2C/G2B） &
個人／企業ユーザー。マイナンバーポータル、確定申告用サイトなど &
公的な情報基盤（住民基本台帳ネットワークシステムなど）と連携した厳格な本人確認。公的個人認証（JPKI）の利用。非常に高度なセキュリティ基準 \\
\hline
\end{tabularx}
\end{table}

\begin{rem}[ID 管理システムの設計にあたって]\label{rem:ch1-idm-design}
このように、オンラインサービスの種類によって ID 管理のあり方は大きく異なる。
効果的で安全な ID 管理システムを設計、開発、運用するためには、まず ID 管理を構成する基本的な要素を理解した上で、
それぞれのサービスの特性や要件（対象ユーザー、セキュリティレベル、適用される法規制など）を考慮に入れることが極めて重要になる。
\end{rem}

\subsubsection{ID管理の構成要素}

ID 管理は、ユーザーのデジタルアイデンティティ情報を安全に取り扱い、適切な人に適切なタイミングでアクセスを許可するための重要な仕組みである。
多くの場合、ID 管理は以下の要素が連携しあって機能している。
この項では、注意 \ref{rem:ch1-idm-design} で述べた「ID 管理を構成する基本的な要素」として、5 つの主要な構成要素の概要を定義する。
各要素の定義やオンラインサービスでの使い所、実装方法は次章で精密化する。

\begin{dfn}[身元確認]\label{dfn:identity-proofing}
アカウント登録時などに、利用者が申告した氏名や住所といった属性情報が本当に本人のものかを、公的な身分証明書などを用いて検証するプロセスを
\underline{身元確認}（identity proofing）という。
\end{dfn}

\begin{rem}[身元確認は信頼の基点である]\label{rem:ch1-proofing-anchor}
身元確認は、信頼できるデジタルアイデンティティ（命題 \ref{prop:ch1-trusted-digital-identity}）の基点となる。
\end{rem}

\begin{dfn}[当人認証]\label{dfn:authentication}
サービスへログインする際に、「あなたは本当に登録されている本人ですか？」を確認するプロセスを \underline{当人認証}（authentication）という。
\end{dfn}

\begin{exa}[当人認証の方法]\label{exa:ch1-authentication-methods}
当人認証では、ID とパスワードの組み合わせだけでなく、SMS 認証コード、指紋や顔などの生体情報を用いた \underline{多要素認証}（MFA）など、
様々な方法で正規のユーザーであることを確かめる。多要素認証の精密な定義は第2章で扱う（定義 \ref{dfn:mfa}）。
\end{exa}

\begin{dfn}[ID連携]\label{dfn:identity-federation}
一つの ID とパスワード（あるいは認証方法）で、連携している複数の異なるサービスにログインできるようにする仕組みを
\underline{ID 連携}（identity federation）という。\underline{シングルサインオン}（SSO）とも呼ばれる。
\end{dfn}

\begin{rem}[ID 連携の利点]\label{rem:ch1-federation-benefit}
ID 連携により、ユーザーはサービスごとにログインし直す手間が省け、利便性が大きく向上する。
\end{rem}

\begin{dfn}[アクセス制御]\label{dfn:access-control}
当人認証（定義 \ref{dfn:authentication}）によりユーザーの検証が済んだ後、「その人がサービスのどの情報にアクセスし、何を行えるか」を、
予め定められた役割や権限に基づいて制御する仕組みを \underline{アクセス制御}（authorization / access control）という。
\end{dfn}

\begin{rem}[アクセス制御の効果]\label{rem:ch1-access-control-benefit}
アクセス制御によって必要な人に必要な権限だけを与えることで、セキュリティを高め、プライバシーリスクを軽減する。
\end{rem}

\begin{dfn}[セッション管理]\label{dfn:session-management}
ユーザーがログインしてからログアウトするまでの一連の通信状態を \underline{セッション}（session）という。
セッションを安全に管理する仕組みを \underline{セッション管理}（session management）という。
\end{dfn}

\begin{exa}[セッション管理の働き]\label{exa:ch1-session-management}
セッション管理は、一定時間操作がない場合に自動でログアウトさせたり、不正なアクセス（\underline{セッションハイジャック}など）に気づくきっかけをユーザーに提供したりすることで、
サービス利用の安全性向上に寄与する。
\end{exa}

以上の 5 つの構成要素を、ID 管理の流れに沿って並べると表 \ref{tab:ch1-idm-components} のようになる。

\begin{table}[H]
\centering
\caption{ID 管理の 5 つの主要な構成要素}\label{tab:ch1-idm-components}
\begin{tabularx}{\linewidth}{|p{0.32\linewidth}|X|}
\hline
構成要素 & 概要 \\
\hline
身元確認（Identity Proofing） & 申告された属性情報が本人のものかを公的な身分証明書などで検証する。信頼できるデジタルアイデンティティの基点 \\
\hline
当人認証（Authentication） & ログイン時に登録されている本人かを確認する。ID とパスワード、SMS 認証コード、生体情報を用いた多要素認証など \\
\hline
ID 連携（Identity Federation） & 一つの ID とパスワード（認証方法）で連携する複数サービスにログインできるようにする。シングルサインオン（SSO） \\
\hline
アクセス制御（Authorization / Access Control） & 認証後、役割や権限に基づいてアクセスできる情報と行える操作を制御する \\
\hline
セッション管理（Session Management） & ログインからログアウトまでの通信状態を安全に管理する。自動ログアウト、不正アクセスに気づくきっかけの提供 \\
\hline
\end{tabularx}
\end{table}

本章では、デジタルアイデンティティの基本概念から、それがもたらすプライバシーに関する課題、そしてそれらを管理・保護するための ID 管理の概要までを解説した。
これらの基礎知識は、私たちが日々利用するオンラインサービスを安全かつ効果的に活用し、信頼できるデジタル社会を築いていく上で不可欠なものである。
次章では、これらの構成要素のそれぞれについてより深く理解するために、定義やオンラインサービスでの使い所、実装方法に踏み込んでいく。

%% =======================================================
%%  part: ch2a
%% =======================================================
\section{ID管理の構成要素}

この章では、\cite{Ito2026} の第2章に沿って、第1章で概要を述べた ID 管理の構成要素を一つずつ詳しく見ていく。
前半では、身元確認（定義 \ref{dfn:identity-proofing}）と当人認証（定義 \ref{dfn:authentication}）を現実世界とデジタル世界の両面から精密化し、デジタル世界で使われている当人認証の方式（パスワード認証、ワンタイムパスワード、TOTP、プッシュ通知、バックアップコード、FIDO 認証とパスキー）とその変遷を、共通の枠組みで整理する。
あわせて、身分証明書のデジタル化と、ユーザー認証に対する脅威分析の 2 つの軸を、コラムとして紹介する。
後半では、ID 連携、アクセス制御、セッション管理を順に扱う。

\subsection{身元確認}\label{subsec:ch2-identity-proofing}

第1章では、身元確認（定義 \ref{dfn:identity-proofing}）を ID 管理の構成要素の一つとして概要だけ述べた。
この節では、まず現実世界の窓口での手続きを例に身元確認が何を確かめる行為なのかを確認し、次に NIST やデジタル庁のガイドラインにおける定義を参照して、デジタル世界の身元確認を「収集・検証・照合」の 3 つのプロセスとして精密化する。
最後に、オンラインサービスで身元確認が必要になる場面と、その代表的な実現方法（メールアドレスの確認、eKYC）を見る。

\subsubsection{現実世界における身元確認}

身元確認が何を確かめる行為なのかを理解するには、私たちが日常的に経験している窓口での手続きを観察するのがよい。
そこで確認されているのは「申請者がどんな資格を持っているか」ではなく、「申請者が申告した情報が正しく、かつ申請者自身のものであるか」である。

\begin{exa}[窓口での身分証明書による身元確認]\label{exa:ch2-counter-proofing}
銀行などの金融機関、市役所などの行政機関の窓口で書類提出などを伴う手続きを行う際、身元確認として免許証やマイナンバーカードといった身分証明書の提示が求められる。
このとき、窓口の担当者は、申請者が運転免許を持っているかどうかを気にしているわけでは\emph{ない}。
担当者は次のことを確認する。
\begin{enumerate}
\item 身分証明書が有効であり、偽造されたものではないこと。
\item 身分証明書に記載されている情報と、手続きの中で申告された情報が一致すること。
\item 身分証明書にある顔写真が、手続きを行う人物と一致すること。
\end{enumerate}
この一連の流れが、現実世界における\underline{身元確認}（Identity Proofing）である。
\end{exa}

現実世界では、さまざまな場面で法律や制度に基づく厳密な身元確認が義務付けられている。
これらは主に、犯罪の防止、取引の安全確保、行政サービスの適正な提供などを目的としている。
表 \ref{tab:ch2-real-world-proofing} に、関連する法律ごとに目的・場面・確認方法をまとめる。

{\small
\begin{longtable}{|p{0.19\textwidth}|p{0.16\textwidth}|p{0.29\textwidth}|p{0.22\textwidth}|}
\caption{現実世界におけるさまざまな身元確認}\label{tab:ch2-real-world-proofing}\\
\hline
関連する法律 & 目的 & 場面 & 確認方法 \\ \hline
\endfirsthead
\hline
関連する法律 & 目的 & 場面 & 確認方法 \\ \hline
\endhead
\hline
\endfoot
犯罪による収益の移転防止に関する法律（犯罪収益移転防止法）
& マネーロンダリング（資金洗浄）やテロ資金供与を防止するため
& 金融機関（銀行、証券会社、保険会社など）での口座開設、大口の現金取引、為替取引、融資など。クレジットカードの契約。ファイナンスリース契約。宅地建物の売買契約またはその代理・媒介。宝石・貴金属等の売買契約（一定額以上）。士業（弁護士、司法書士、行政書士など）が特定の業務を行う場合
& 公的な身分証明書（マイナンバーカード、運転免許証、パスポート、在留カードなど）の提示を受け、氏名、住所、生年月日などを確認・記録する。金融機関などでは、取引を行う目的や職業なども確認される（取引時確認） \\ \hline
携帯音声通信事業者による契約者等の本人確認等及び携帯音声通信役務の不正な利用の防止に関する法律（携帯電話不正利用防止法）
& 振り込み詐欺など、携帯電話や PHS が犯罪に不正利用されることを防止するため
& 携帯電話・PHS の新規契約、譲渡、名義変更。SIM カードの契約
& 公的な身分証明書により、契約者の氏名、住所、生年月日等を確認・記録する \\ \hline
古物営業法
& 盗品の流通を防止・発見するため
& 古物商（リサイクルショップ、中古品買取店など）が古物を買い受ける、または交換する場合（対価が 1 万円未満の場合などを除く）
& 相手方の身分証明書（運転免許証、マイナンバーカードなど）により、氏名、住所、年齢等を確認する \\ \hline
公職選挙法
& なりすましによる不正投票を防止し、選挙の公正を確保するため
& 投票所での投票時（ただし、厳密な「身分証明書の提示義務」ではなく、選挙人名簿との照合が基本。忘れた場合でも、本人と確認できれば投票できる運用が多い）。期日前投票や不在者投票
& 投票所入場整理券（持参が推奨される）や口頭での確認が行われ、それでも本人確認が難しい場合に身分証明書の提示が求められることがある。期日前投票・不在者投票ではより厳格な確認が行われることがある \\ \hline
各種行政手続き（戸籍法、住民基本台帳法、出入国管理及び難民認定法、道路交通法など）
& 行政サービスの適正な提供、公的な記録の正確性維持、資格の付与など
& 戸籍謄本・抄本、住民票の写しなどの請求。印鑑登録、印鑑証明書の発行。マイナンバーカードの申請・交付。パスポート（旅券）の申請・交付。運転免許証の申請・交付・更新。在留カードの交付・更新など（外国人住民の場合）
& それぞれの手続きで定められた公的な身分証明書（マイナンバーカード、運転免許証、パスポート、在留カード、健康保険証など）の提示が必要。複数の書類が必要な場合もある \\ \hline
\end{longtable}
}

このように、ある手続きを行うにあたって、信頼を構築するプロセスとして身元確認が行われている。
この観察は、申請された情報の「状態」という観点から次のようにまとめられる。

\begin{prop}[身元確認は申請情報を「主張」から「確認済み」に移す]\label{prop:ch2-claimed-to-verified}
手続きにおいて申請された情報は、身元確認が済むまでは「申請者が主張している」状態にすぎず、身元確認によってその正当性が検証されて初めて「確認済み」の状態となる。
\end{prop}
\begin{proof}
表 \ref{tab:ch2-real-world-proofing} の各手続きを見ると、申請者が申告した氏名・住所・生年月日などの情報は、それだけでは申請者本人の言明にすぎず、窓口がそれを信じる根拠はない。
まず、身元確認では公的な身分証明書の有効性と真正性を確かめ、次に申告内容と記載内容の一致を確かめ、さらに顔写真と申請者の一致を確かめる（例 \ref{exa:ch2-counter-proofing}）。
よって、この一連の手順を経た後で初めて、申請情報は第三者の発行した書類によって裏付けられ、正当性が検証された「確認済み」の情報として扱える。
\end{proof}

\begin{exa}[報道における「自称」]\label{exa:ch2-self-proclaimed}
命題 \ref{prop:ch2-claimed-to-verified} を象徴するのが、刑事事件の報道における「自称」の扱いである。
いわゆる「現行犯逮捕」「緊急逮捕」が行われた場合、警察による身元確認が済むまでは「自称 会社員」「自称（氏名）」といった表現がされることがある。
これは、被疑者が名乗った職業や氏名が、まだ「主張」の状態にとどまっていることを表している。
\end{exa}

\subsubsection{デジタル世界における身元確認}

デジタル世界の身元確認の定義も、現実世界に即したものとなっている。
この項では、代表的なガイドラインにおける身元確認の定義を参照し、そこに共通する「収集・検証・照合」の 3 つのプロセスとして定義 \ref{dfn:identity-proofing} を精密化する。
そのうえで、オンラインサービスで身元確認が必要になる場面と、その実現方法を見ていく。

デジタル世界での身元確認については、米国国立標準技術研究所（NIST）が発行している「SP 800-63 Digital Identity Guidelines」のようなガイドラインが参考になる。
これは米国の連邦政府機関向けに策定された、電子的な本人確認に関するデジタルアイデンティティのガイドラインであるが、各国の政府機関や業界標準など、さまざまなガイドラインからも広く参照されている。

\begin{rem}[NIST SP 800-63 Revision 4 における身元確認]\label{rem:ch2-nist-proofing}
2025 年 7 月に発行された NIST SP 800-63 Revision 4 では、身元確認は次のように定義されている。
\begin{quote}
The processes used to collect, validate, and verify information about a subject to establish assurance in the subject's claimed identity.
\end{quote}
すなわち、「対象者が主張するアイデンティティの信頼性を確立するために、その情報を収集、検証、照合するプロセス」である。
\end{rem}

\begin{rem}[デジタル庁 DS-511 における身元確認]\label{rem:ch2-ds511-proofing}
令和 7 年 9 月 30 日にデジタル庁から発行された「DS-511 行政手続等での本人確認におけるデジタルアイデンティティの取扱いに関するガイドライン」という資料では、身元確認を次のように定義している。
\begin{quote}
申請者を一意に識別するとともに、その実在性を確認すること。
\end{quote}
具体的には、申請者の属性情報を収集することで申請者を一意に識別するとともに、収集した属性情報が真正かつ申請者自身のものであることを本人確認書類により検証することで、申請者が実在かつ生存する人物であることを確認する。
\end{rem}

どちらのドキュメントにおいても、「収集、検証、照合」の 3 つのプロセスに触れられている。
そこで、身元確認（定義 \ref{dfn:identity-proofing}）を、この 3 つのプロセスの合成として精密化しておく。

\begin{dfn}[身元確認の 3 つのプロセス]\label{dfn:ch2-proofing-processes}
身元確認は、次の 3 つのプロセスからなる。
\begin{enumerate}
\item \underline{収集}（Collection）：情報の提出、証明書類の提示を受けること。氏名、生年月日、住所などの属性情報を収集し、母集団の中で一意に識別する。
\item \underline{検証}（Validation）：提出された情報や証明書類の正当性を検証すること。提示された本人確認書類が偽装・改ざん・複製された不正なものでないこと、有効なものであることを、物理的または電子的に検証する。
\item \underline{照合}（Verification）：提供された身分証明書や情報と対象者との結びつき、すなわち当人性を確認すること。顔写真や暗証番号などを用いて、提出された本人確認書類が確かに対象者自身のものであることを検証する。
\end{enumerate}
\end{dfn}

例 \ref{exa:ch2-counter-proofing} の窓口の手続きに当てはめれば、申告と身分証明書の提示が収集、身分証明書の有効性・非偽造性の確認が検証、顔写真と申請者の一致の確認が照合にあたる。

\underline{身元確認が必要になる場面：}
オンラインサービスにおいて、身元確認が必要になる場面がいくつかある。
\begin{enumerate}
\item 新規アカウント登録時：サービス提供者が利用者の正当性を確認する。
\item アカウントリカバリー：利用者の正当性を確認して状態を復旧する。
\item 属性情報の確認が必要な場面：年齢・住所・職業など、サービス条件を満たすかを確認する。
\item 高リスク取引の前：不正アクセスやなりすまし防止のために確認する。
\item 規制業種での利用者登録：法令を遵守するために確認する。
\end{enumerate}

\underline{メールアドレスの確認：}
オンラインサービスの新規アカウント登録時によく使われている方法が、メールアドレスの確認である。
特にコンシューマー向けサービスにおいて、メールアドレスは連絡先、ユーザー識別子といった用途の属性情報として利用される。

\begin{exa}[メールアドレスの確認]\label{exa:ch2-email-verification}
サービスの利用開始時、ユーザーは利用したいメールアドレスを送信する。
オンラインサービスはメールアドレスの形式を確認し、ワンタイムパスワードや URL を送信して有効性を確認する。
送信されたワンタイムパスワードの検証や URL へのアクセスにより、そのユーザーがメールを受信できたこととみなす。
こうしてメールアドレスの提示、有効性の検証、受信による所有確認をした上で、確認済みの属性情報として保持する。
\end{exa}

\begin{prop}[メールアドレスの確認も身元確認である]\label{prop:ch2-email-is-proofing}
例 \ref{exa:ch2-email-verification} のフローには身分証明書などは一切出てこないが、この一連のフローも身元確認である。
\end{prop}
\begin{proof}
定義 \ref{dfn:ch2-proofing-processes} の 3 つのプロセスがそろっていることを確かめればよい。
まず、ユーザーがメールアドレスを提示することは、属性情報の収集にあたる。
次に、サービスがメールアドレスの形式を確認し、ワンタイムパスワードや URL を送信してその有効性を確かめることは、検証にあたる。
さらに、ワンタイムパスワードの検証や URL へのアクセスによって、そのユーザーがそのメールを受信できたこと、すなわちそのメールアドレスを所有していることを確かめることは、属性情報と対象者の結びつきを確かめる照合にあたる。
よって、このフローは収集・検証・照合を経て属性情報を「確認済み」にしており（命題 \ref{prop:ch2-claimed-to-verified}）、身元確認である。
\end{proof}

\underline{アカウントリカバリーにおける身元確認：}
アカウントリカバリー（第3章で扱う。定義 \ref{dfn:account-recovery}）は、なんらかの理由でログインできなくなったり、サービス側から無効化されたりといった状態から復旧するためのプロセスである。
ユーザー自身の操作、もしくはカスタマーサービスとのコミュニケーションにおいて身元確認を行うことで正当性を確認し、ログインできない状態の復旧を行う。
サービス自体の内容や復旧したい状態により、メールアドレスの確認から、現実世界と同等の身分証明書を用いた身元確認まで、求められる方法が異なる。

\underline{eKYC：}
身分証明書を利用した身元確認を必要とするケースのために、次の仕組みがある。

\begin{dfn}[eKYC]\label{dfn:ch2-ekyc}
実店舗や行政機関へ出向くことなく、身分証明書を利用した身元確認をオンラインで行う仕組みを \underline{eKYC}（electronic Know Your Customer）という。
\end{dfn}

eKYC には、身分証の特性、ユーザーが利用可能なデバイスなどにより、いくつかの種類がある。
\begin{enumerate}
\item 身分証明書とユーザーのセルフィー（自撮り）を画像で送信するもの：利便性は高いものの、偽造リスクを考慮する必要がある。
\item 動画を利用するもの：静止画よりは確認能力の向上が期待できるものの、ユーザーの利便性が低下し、カメラ、マイクといったデバイスに依存する部分が増える。
\item IC チップに保存された情報を読み取るもの：モバイル端末の進化も伴い、利便性と安全性の両方を享受できるユーザーが増えたといえる。
\item これら以外にも、各種データベースとの照合を行うものなどがある。
\end{enumerate}

\begin{rem}[eKYC の方式の選択]\label{rem:ch2-ekyc-choice}
eKYC のどの方式を利用するかは、サービスや機能単位の内容、ユーザー層、制度による制約などの外的要因などから決まる。
\end{rem}

\subsubsection{コラム：身分証明書のデジタル化（スマートフォン搭載）}

近年、身分証明書のデジタル化、特にスマートフォンへの搭載が世界的な潮流となっている。
この動きは、物理的なカードを持ち歩く不便さを解消し、利便性を高めることを目的としており、身元確認とも密接に関係している。
このコラムでは最近のトレンドとして、日本の「マイナンバーカードのスマートフォン搭載」と、米国を中心に進む「モバイル運転免許証（mDL）」の取り組みを例に、その概要を紹介する。

\underline{日本のマイナンバーカード：スマートフォン搭載で拡がる利便性}

\begin{exa}[マイナンバーカードのスマートフォン搭載]\label{exa:ch2-mynumber-smartphone}
日本では、マイナンバーカードの機能をスマートフォンに取り込む取り組みが進行中である。
実現方法の一つとして、マイナンバーカードの IC チップに格納された電子証明書（公的個人認証サービス JPKI の署名用・利用者証明用電子証明書など）を、スマートフォンの\underline{セキュアエレメント}（SE）と呼ばれる安全な領域に保存するものがある。
これにより、利用者は物理カードなしに、スマートフォンの生体認証やパスコード認証を通じて、マイナポータルへのログイン、オンラインの行政手続き、民間のオンラインサービスを利用できるようになる。
対面では、医療機関でのマイナ保険証としての利用が始まっており、将来的には役所窓口や民間での身元確認において利用場面の拡大が期待されている。
\end{exa}

\begin{rem}[個人番号はやり取りされない]\label{rem:ch2-no-individual-number}
例 \ref{exa:ch2-mynumber-smartphone} のような利用では、「個人番号」のやり取りは行われない。
\end{rem}

\underline{米国のモバイル運転免許証（mDL）：プライバシー重視の国際標準}

\begin{exa}[米国のモバイル運転免許証（mDL）]\label{exa:ch2-mdl}
米国では、運転免許証をデジタル化しスマートフォンに格納する\underline{モバイル運転免許証}（mDL）の導入が進んでいる。
日本のマイナンバーカードの取り組みとコンセプトは似ているが、mDL は国際規格（ISO/IEC 18013-5）への準拠や、提示する情報をユーザーが選択できる\underline{選択的開示}機能に特徴がある。
利用場面としては、酒類購入時の年齢確認、空港の保安検査、レンタカー利用、法執行機関への提示などが想定されている。
その際、免許証の情報を全て提示するのではなく、用途に応じた必要最低限の情報に絞ることができる。
\end{exa}

\subsection{当人認証}\label{subsec:ch2-authentication}

私たちがさまざまなサービスを利用する際、多くの場合、その利用資格を持つ本人であることを示す必要がある。
この「利用しようとしている人が正当な人物であること」を確認する手続きが当人認証（定義 \ref{dfn:authentication}）である。
この節では、まず身近な現実世界の例を通して当人認証がどのようなものかを見て、認証の 3 要素と多要素認証を定義する。
次に、しばしば混同される「本人確認」という言葉が身元確認と当人認証のどちらを指すのかを整理し、最後に NIST とデジタル庁のガイドラインを参照して、デジタル世界の当人認証を 4 つの局面からなるライフサイクルとして精密化する。

\subsubsection{現実世界における当人認証}

当人認証が何をしているのかを、最も身近な例である ATM から見ていく。
そこから当人認証に用いられる情報や手段（認証要素）を取り出して分類し、それらを組み合わせることで認証の強度を高める考え方へ進む。

\underline{ATM での現金引き出し：}

\begin{exa}[ATM での現金引き出し]\label{exa:ch2-atm}
ATM で銀行のキャッシュカードを使って現金を引き出す際には、\underline{暗証番号}（PIN：Personal Identification Number）の入力が求められる。
これは、操作している人物がキャッシュカードの正当な所有者であることを、原則として本人だけが知っている情報（暗証番号）によって確認する行為である。
キャッシュカードを持っているだけでは不十分で、正しい暗証番号を入力して初めて本人であると認められる。
\end{exa}

この例のように、当人認証とは、何らかのサービスや権利を利用しようとする人物が、その資格を持つ本人であるかを確認するプロセスである。
現実世界の当人認証では、利用者（例 \ref{exa:ch2-atm} では ATM の利用者）が、サービス（銀行）に対して、あらかじめ持っている情報や手段（キャッシュカードと暗証番号）を示し、サービスがそれを確かめる。

\underline{認証の 3 要素：知識情報、所持情報、生体情報}

例 \ref{exa:ch2-atm} における暗証番号やキャッシュカード、あるいは生年月日や後述する指紋など、当人認証に用いられる情報や手段にはいくつかの種類がある。
これらを性質によって分類しておくと、認証の強度を議論するときの共通の言葉になる。

\begin{dfn}[認証要素と認証の 3 要素]\label{dfn:authentication-factors}
当人認証に用いられる情報や手段を\underline{認証要素}（Authentication Factor）と呼ぶ。
認証要素は、その性質によって大きく次の 3 種類に分類され、\underline{認証の 3 要素}として知られている。
\begin{enumerate}
\item \underline{知識情報}（Something You Know）：本人だけが知っている情報。
\item \underline{所持情報}（Something You Have）：本人だけが持っている物理的なモノ。
\item \underline{生体情報}（Something You Are）：本人固有の身体的特徴や行動的特徴。
\end{enumerate}
\end{dfn}

\begin{nota}[認証要素のカテゴリ]\label{nota:ch2-factor-category}
以後、認証の 3 要素のカテゴリの集合を
\[
C := \{\text{知識},\ \text{所持},\ \text{生体}\}
\]
と書く。
認証要素全体の集合を $F$ とし、各認証要素 $f\in F$ に対してそれが属するカテゴリを $c(f)\in C$ と書く。
すなわち $c:F\to C$ は、認証要素をその種類に対応させる写像である。
認証要素の組合せ $S\subseteq F$ に対して、$c(S) := \{c(f)\mid f\in S\}$ はその組合せに現れるカテゴリの集合である。
\end{nota}

\begin{exa}[認証の 3 要素の具体例]\label{exa:ch2-factor-examples}
\begin{enumerate}
\item 知識情報の例：暗証番号（PIN）、パスワード、合言葉、特定の記憶（母親の旧姓など。秘密の質問の答えに使われることがある）。
\item 所持情報の例：キャッシュカード、診察券、社員証、学生証、家の鍵、車の鍵、スマートフォン、ハードウェアトークン、印鑑（日本特有）。
\item 生体情報の例：指紋、顔、声、虹彩、静脈パターン、署名（筆跡）。
\end{enumerate}
\end{exa}

\underline{多要素認証（MFA）：}

これら 3 種類の認証要素は、それぞれ単独で用いられることもあるが、各要素には固有のリスクが存在する。

\begin{prop}[単一の認証要素に固有のリスク]\label{prop:ch2-single-factor-risk}
認証の 3 要素のいずれも、単独では突破されるリスクを持つ。
\end{prop}
\begin{proof}
各カテゴリについて、そのリスクを挙げればよい。
まず、知識情報は、忘れたり他人に知られたりする可能性がある。
次に、所持情報は、盗まれたり紛失したりする可能性がある。
さらに、生体情報も、状況によっては読み取りが困難だったり、模倣されたりするリスクがゼロではない。
よって、どのカテゴリの認証要素も単独ではリスクを免れない。
\end{proof}

そこで、単一の認証要素だけに頼るリスクを軽減し、認証の強度を高めるためのアプローチとして、多要素認証がある。

\begin{dfn}[多要素認証]\label{dfn:mfa}
認証の 3 要素のうち、複数の\emph{異なる}認証要素を組み合わせた当人認証を\underline{多要素認証}（Multi-Factor Authentication：MFA）という。
記号 \ref{nota:ch2-factor-category} を用いれば、認証要素の組合せ $S\subseteq F$ による当人認証が多要素認証であるとは、
\[
|c(S)| \geq 2
\]
が成り立つこと、すなわち $S$ に 2 つ以上の異なるカテゴリの認証要素が含まれることである。
\end{dfn}

\begin{prop}[多要素認証は安全性を向上させる]\label{prop:ch2-mfa-improves}
多要素認証では、仮に一つの認証要素が突破されても、他の認証要素が不正利用を防ぐ別の壁となるので、安全性を向上できる。
\end{prop}
\begin{proof}
多要素認証 $S$ では $|c(S)|\geq 2$ であるから、$S$ にはカテゴリの異なる認証要素 $f_1, f_2$ が含まれる。
命題 \ref{prop:ch2-single-factor-risk} の証明で見たように、各カテゴリのリスクは性質が異なる（知識情報は漏洩、所持情報は盗難・紛失、生体情報は読み取り困難・模倣）。
よって、例えば $f_1$ が知られたり盗まれたりして突破されても、性質の異なる $f_2$ が同時に突破されているとは限らず、$f_2$ が不正利用を防ぐ別の壁として残る。
したがって、異なる認証要素の認証方式を組み合わせることで、不正ログインのハードルを上げることができる。
\end{proof}

\begin{exa}[ATM における多要素認証]\label{exa:ch2-atm-mfa}
例 \ref{exa:ch2-atm} の ATM は、「キャッシュカード」（所持情報）と「暗証番号」（知識情報）という 2 つの異なる認証要素を組み合わせた、典型的な多要素認証である。
また、銀行によっては「指紋認証」を利用できる ATM がある。
この場合も、「キャッシュカード」（所持情報）と「指紋」（生体情報）を組み合わせた多要素認証である。
\end{exa}

\underline{同一要素の組合せ：}
多要素認証とは異なり、\emph{同じ}種類の認証要素を複数組み合わせることで安全性を高める方法もある。
記号 \ref{nota:ch2-factor-category} でいえば、$|S|\geq 2$ だが $|c(S)|=1$ となる組合せである。
これは多要素認証（異なる種類の要素の組合せ）とは区別されるが、特定の状況下では有効な対策となり得る。

\begin{exa}[2 つの鍵穴のあるドア]\label{exa:ch2-two-keys}
家のドアに鍵穴が 2 つあり、それぞれ異なる鍵（どちらも所持情報）が必要な場合がある。
この「鍵と鍵」の組合せは所持情報という同一要素の組合せであるが、侵入するには両方の鍵を入手するか、2 つの鍵をピッキングで解除する必要がある。
単一の鍵よりも時間がかかるため、不正な侵入のリスクを低減できる。
\end{exa}

\begin{rem}[同一要素の組合せでもハードルは上がる]\label{rem:ch2-same-factor}
例 \ref{exa:ch2-two-keys} のように、認証要素は同じでも、管理を分けるなどにより、不正ログインのハードルを上げられる。
ただし、これは定義 \ref{dfn:mfa} の意味の多要素認証ではない。
\end{rem}

\subsubsection{「本人確認」は身元確認か、当人認証か}

先に説明した身元確認と、この節で扱っている当人認証は、密接に関係している。
しかし日常語の「本人確認」はその両方を指しうるため、ここで両者の関係を整理しておく。

\begin{rem}[「本人確認」は身元確認か、当人認証か]\label{rem:honnin-kakunin}
一般的に「本人確認」という言葉は、文脈によって「身元確認」と「当人認証」のどちらか、または両方を指す場合がある。
例えば、銀行で口座を開設する際の「本人確認」は主に身元確認を指し、ATM で現金を引き出す際の「本人確認」は当人認証を指す。
この二つを区別して理解することが重要である。
\end{rem}

身元確認と当人認証は、次のような関係にある。
\begin{enumerate}
\item まず、公的書類などで「その人が誰であるか」を確認し、その情報を記録する。同時に、当人認証に必要な情報（例：暗証番号の設定、キャッシュカードの発行）を準備する。$\Rightarrow$ 身元確認
\item 以降、サービスを利用する際には、登録・発行された認証要素を用いて「登録された本人であるか」を都度確認する。$\Rightarrow$ 当人認証
\item キャッシュカードを紛失したり、暗証番号を忘れたりして当人認証が困難になった場合、再度身元確認を行い、サービスが利用可能な状態に戻す（例：カード再発行、暗証番号再設定）。$\Rightarrow$ アカウントリカバリー
\end{enumerate}

表 \ref{tab:ch2-proofing-vs-authn} に、身元確認と当人認証の違いをまとめる。

\begin{table}[H]
\centering
\caption{身元確認と当人認証の関係}\label{tab:ch2-proofing-vs-authn}
\begin{tabular}{|p{0.13\textwidth}|p{0.38\textwidth}|p{0.38\textwidth}|}
\hline
 & 身元確認（Identity Proofing） & 当人認証（Authentication） \\ \hline
目的 & 対象者が「誰であるか」を特定し、信頼性を確立する & 利用者が「登録された本人であるか」を確認する \\ \hline
タイミング & 主に初回登録時、重要な契約・手続き開始時 & サービス利用時、施設入館時など、アクセスの都度 \\ \hline
主な手段 & 公的証明書（運転免許証、パスポート等）、書類確認 & 暗証番号、カードキー、社員証、生体特徴（指紋等） \\ \hline
確認内容 & 実在性、氏名、住所、生年月日などの属性情報 & 登録された認証要素との一致 \\ \hline
例 & 銀行口座開設時の本人確認書類提出 & ATM での暗証番号入力、オフィスへの入館証提示 \\ \hline
\end{tabular}
\end{table}

この関係は、現実世界だけでなく、デジタル世界のオンラインサービスにも適用されている。

\subsubsection{デジタル世界における当人認証}

当人認証と身元確認との関係は、デジタル世界においてより明確に定義されている。
この項では、NIST とデジタル庁のガイドラインにおける当人認証の定義を参照して「認証器」という概念を導入し、当人認証を 4 つの局面からなるライフサイクルとして定義 \ref{dfn:authentication} を精密化する。
この 4 つの局面は、次の節で各認証方式を整理するときの共通の枠組みになる。

\begin{rem}[NIST SP 800-63 Rev.4 における当人認証]\label{rem:ch2-nist-authn}
NIST（米国国立標準技術研究所）のガイドライン（SP 800-63 Rev.4）では、当人認証を次のように定義している。
\begin{quote}
\textit{The process by which a claimant proves possession and control of one or more authenticators bound to a subscriber account to demonstrate that they are the subscriber associated with that account.}
\end{quote}
\end{rem}

\begin{rem}[デジタル庁のガイドラインにおける当人認証]\label{rem:ch2-digital-agency-authn}
デジタル庁の「本人確認ガイドライン改定方針 令和 6 年度とりまとめ」（令和 7 年 4 月発行）では、当人認証を次のように定義している。
\begin{quote}
申請者の当人性を確認すること。
\end{quote}
具体的には、対象手続を利用しようとする者が、身元確認時に登録された者と同一の人物であることを、申請者と紐づけて登録した認証器を用いて確認する。
\end{rem}

いずれの定義も、「利用者がサービスに登録済みのアカウントに紐づけられた認証器を用い、そのアカウントの正当な利用者本人であることを確認するプロセス」を指している。
これらの定義に出てくる「認証器」という言葉を、次のように理解しておく。

\begin{dfn}[認証器]\label{dfn:ch2-authenticator}
認証要素（定義 \ref{dfn:authentication-factors}）を実現するための具体的な手段やツールを\underline{認証器}（Authenticator）という。
例えば、パスワード、スマートフォンアプリ、セキュリティキーと呼ばれるデバイスなどである。
\end{dfn}

\begin{prop}[デジタル世界の当人認証]\label{prop:ch2-digital-authn}
デジタル世界の当人認証（定義 \ref{dfn:authentication}）は、利用者がサービスに登録済みのアカウントに紐づけられた認証器を用い、そのアカウントの正当な利用者本人であることを確認するプロセスである。
\end{prop}
\begin{proof}
注意 \ref{rem:ch2-nist-authn} と注意 \ref{rem:ch2-digital-agency-authn} の二つの定義の共通部分を取り出せばよい。
まず、NIST の定義では、claimant（主張者）が subscriber account（登録済みのアカウント）に bound された（紐づけられた）一つ以上の authenticator（認証器）の所持と制御を示すことで、そのアカウントに対応する subscriber（登録者）であることを示す。
次に、デジタル庁の定義では、利用しようとする者が身元確認時に登録された者と同一であることを、申請者と紐づけて登録した認証器を用いて確認する。
よって、どちらも「登録済みのアカウントに紐づけられた認証器を用いて、そのアカウントの正当な利用者本人であることを確認する」ことを述べている。
\end{proof}

さらに、デジタル庁のガイドラインでは、当人認証に関連するプロセスとして次の 4 つを挙げている。

\begin{dfn}[当人認証の 4 つの局面]\label{dfn:authenticator-phases}
当人認証に関連する次の 4 つのプロセスを、\underline{当人認証の 4 つの局面}と呼ぶ。
\begin{enumerate}
\item \underline{認証器の登録}：サービス利用開始時などに、パスワードを設定したり、認証アプリを連携させたりして、認証に使う手段を用意するプロセス。アカウント登録時等の身元確認プロセスにおいて認証器を登録するなどして、当人認証に用いる認証器を利用者と紐づけて登録する。
\item \underline{当人認証の実施}：サービス利用時に、登録された認証器を用いて本人であることを確認するプロセス。手続きやサービスを利用しようとする申請者が、あらかじめ登録されている利用者と同一の人物であることを、認証器によって確認する。
\item \underline{盗難・紛失等への対応}：利用者から認証器が利用できなくなった等の報告を受け、認証器の無効化やアカウントの一時停止といった対応を行うプロセス。利用者から認証器の盗難や紛失の報告を受けた際に、認証器の無効化やアカウントの停止等の対応を行う。
\item \underline{アカウントの回復}：ログインできない状態になった利用者が、再びサービスを利用できるように回復させるプロセス。認証器の盗難・紛失、故障による交換、パスワードの忘失などによって利用者がアカウントにログインできなくなった状態を回復する。
\end{enumerate}
以下、これらを略して (1) 登録、(2) 認証、(3) 対応、(4) 回復とも呼ぶ。
\end{dfn}

\begin{rem}[当人認証はライフサイクルである]\label{rem:ch2-authn-lifecycle}
定義 \ref{dfn:authenticator-phases} のように、デジタル世界の当人認証は、単に認証を行う（局面 (2)）だけでなく、その準備（局面 (1)）、問題発生時の対応（局面 (3)）、そして復旧（局面 (4)）までを含む一連のライフサイクルとして捉えられている。
特に、認証器の登録（局面 (1)）やアカウントの回復（局面 (4)）のプロセスでは、身元確認が必要となる場合がある。
これは、現実世界における身元確認と当人認証の関係（前項の (1)--(3)）に対応している。
\end{rem}

\subsection{デジタル世界で使われている当人認証の方式とその変遷}\label{subsec:ch2-authn-methods}

デジタル世界では、さまざまな当人認証の方式（以下、\underline{認証方式}）が利用されている。
この節では主要な認証方式について、その概要と、定義 \ref{dfn:authenticator-phases} の 4 つの局面（登録、認証、対応、回復）がどのように行われるか、そして認証要素（定義 \ref{dfn:authentication-factors}）の観点から整理する。
各項は、認証方式の定義、4 つの局面の記述、認証要素の分類と特徴、の順に統一した形で述べる。
そのうえで、多要素認証とパスワードレス認証という組合せ方の違い、フィッシング攻撃の進化とそれに耐性を持つ FIDO 認証、そしてパスキー認証の登場までの変遷をたどる。

\subsubsection{パスワード認証}

パスワード認証は、オンラインサービスで最も古くから使われ、広く知られている認証方式である。
他の認証方式は、その多くがパスワード認証の弱点を補う、あるいは置き換えるために登場したので、まずこれを基準として整理する。

\begin{dfn}[パスワード認証]\label{dfn:password-auth}
ユーザーが自身で考えた文字列（\underline{パスワード}）をあらかじめサービスに登録しておき、認証時にユーザーが送信したパスワードが登録されたものと一致するかをサービスが確認する認証方式を、\underline{パスワード認証}という。
\end{dfn}

パスワード認証の 4 つの局面（定義 \ref{dfn:authenticator-phases}）は次のとおりである。
\begin{enumerate}
\item \underline{登録：}アカウント登録時や設定変更機能を通じて、ユーザーは自身で考えたパスワード（文字列）をサービスに登録する。
\item \underline{認証：}ユーザーは ID 等と共にパスワードをサービスに送信し、サービスは登録されたパスワードと一致するかを確認する。
\item \underline{対応：}パスワード漏洩の疑いがある場合など、必要に応じてサービス側でパスワードの強制変更やアカウントの無効化を行う。
\item \underline{回復：}パスワードを忘れた場合、メールアドレス等を用いたリセット手続き（簡易的な身元確認を含む場合がある）や、より厳格な身元確認を経て、パスワードを再設定する。
\end{enumerate}

\begin{prop}[パスワード認証の認証要素と特徴]\label{prop:ch2-password-factor}
パスワード認証の認証要素は知識情報にあたる。
パスワード認証は利用環境の自由度が高いという利点を持つ一方、不正ログインのリスクが高いことが指摘されている。
\end{prop}
\begin{proof}
まず、パスワード認証はユーザーがパスワードを記憶することが前提であり、本人だけが知っている情報によって認証するので、認証要素は知識情報（定義 \ref{dfn:authentication-factors}）にあたる。
次に、記憶した文字列を送信するだけで認証でき、特定のデバイスを必要としないため、利用環境の自由度が高い。
一方で、ユーザー側の課題として推測されやすいパスワードの使用やパスワードの使い回し、フィッシング詐欺による漏洩があり、サービス側の課題としてサーバーへの侵入による漏洩がある。
よって、ユーザー側・サービス側双方の課題により、不正ログインのリスクが高い。
\end{proof}

\subsubsection{メールや SMS を用いたワンタイムパスワード（OTP）認証}

パスワードを記憶させる代わりに、その都度サービスから一度だけ有効なパスワードを送る方式がある。
パスワード認証と組み合わせて、あるいは単体で利用されることもある。

\begin{dfn}[メール・SMS による OTP 認証]\label{dfn:otp-auth}
一度だけ有効なパスワードを\underline{ワンタイムパスワード}（OTP：One-Time Password）という。
メールアドレスやスマートフォンの SMS（ショートメッセージサービス）宛に OTP を送信し、それをユーザーに入力させる認証方式を、\underline{メール・SMS による OTP 認証}という。
\end{dfn}

メール・SMS による OTP 認証の 4 つの局面は次のとおりである。
\begin{enumerate}
\item \underline{登録：}ユーザーは OTP を受け取るための連絡先（メールアドレスや電話番号）をサービスに登録する。
\item \underline{認証：}サービスは登録された連絡先に OTP を送信する。ユーザーは受け取った OTP をサービスに入力し、サービスはそれが正しいかを確認する。
\item \underline{対応：}連絡先として登録したメールアカウントや電話番号が利用できなくなった場合、サービス側でその連絡先を無効化する手続きが必要である。
\item \underline{回復：}身元確認等を経て、新しい連絡先を登録し直す。
\end{enumerate}

\begin{prop}[OTP 認証の認証要素]\label{prop:ch2-otp-factor}
メール・SMS による OTP 認証の認証要素は所持情報に分類される。
\end{prop}
\begin{proof}
この方式は、「登録されたメールアドレスや電話番号に送られた OTP にアクセスできること」をもって本人性を確認する方式である。
OTP にアクセスできることは、それらのメールアドレスや電話番号が利用可能なデバイスやアカウントを持っていることにほぼ等しい。
よって、確認しているのは本人が持っているモノであり、認証要素は所持情報（定義 \ref{dfn:authentication-factors}）に分類される。
\end{proof}

\subsubsection{時間ベースのワンタイムパスワード（TOTP）認証}

OTP をサービスから送る代わりに、ユーザーの手元のアプリやデバイスで生成させる方式がある。
主にパスワード認証と組み合わせて、多要素認証として利用される。

\begin{dfn}[TOTP 認証]\label{dfn:totp-auth}
専用のアプリ（\underline{ソフトウェアトークン}）や物理的なデバイス（\underline{ハードウェアトークン}）が、一定時間ごと（例：30 秒、60 秒）に変化するワンタイムパスワードを生成し、それをユーザーが入力する認証方式を、\underline{時間ベースのワンタイムパスワード認証}（\underline{TOTP 認証}）という。
\end{dfn}

TOTP 認証の 4 つの局面は次のとおりである。
\begin{enumerate}
\item \underline{登録：}サービスから提示される QR コード等を認証アプリで読み取るか、シークレットキーを手動で設定する（ハードウェアトークンの場合は事前配布・登録が必要な場合もある）。
\item \underline{認証：}ユーザーは認証アプリやハードウェアトークンが表示する OTP をサービスに入力し、サービスは時刻とシークレットキーに基づいて期待される OTP と一致するかを検証する。
\item \underline{対応：}認証アプリをインストールしたデバイスやハードウェアトークンが利用できなくなった場合、サービス側でこの認証方式を無効化する手続きが必要である。
\item \underline{回復：}身元確認等を経て、新しいデバイスで再度 TOTP 認証を設定する。
\end{enumerate}

\begin{rem}[TOTP 認証の変遷]\label{rem:ch2-totp-history}
古くは金融機関などでハードウェアトークンが用いられていたが、スマートフォンの普及に伴い、Google Authenticator や Microsoft Authenticator などの認証アプリ（ソフトウェアトークン）の利用が一般的になった。
\end{rem}

\begin{prop}[TOTP 認証の認証要素]\label{prop:ch2-totp-factor}
TOTP 認証の認証要素は所持情報にあたる。
\end{prop}
\begin{proof}
TOTP 認証で入力される OTP は、シークレットキーを設定したアプリやデバイスだけが、時刻に応じて生成できるものである。
よって、正しい OTP を入力できることは「OTP を生成するアプリやデバイスを持っていること」を確認しているのであり、認証要素は所持情報にあたる。
\end{proof}

\subsubsection{モバイル認証アプリ（プッシュ通知）を用いた認証}

OTP を手で入力させる代わりに、ユーザーのスマートフォンに通知を送り、承認操作をしてもらう方式がある。
これも主にパスワード認証と組み合わせて、多要素認証として利用される。

\begin{dfn}[モバイル認証アプリ（プッシュ通知）による認証]\label{dfn:push-auth}
ログイン要求があった際に、事前に連携させたスマートフォンの認証アプリにプッシュ通知を送り、アプリ上でユーザーが承認操作（例：「はい、私です」ボタンをタップ）を行うことで認証する方式を、\underline{モバイル認証アプリ（プッシュ通知）を用いた認証}という。
\end{dfn}

この方式の 4 つの局面は次のとおりである。
\begin{enumerate}
\item \underline{登録：}ユーザーは事前に、サービス指定のモバイルアプリをインストールし、自分のアカウントと連携させる。
\item \underline{認証：}サービスへのログイン試行時、連携されたモバイルアプリに通知が送られる。ユーザーがアプリ上で承認操作を行うと、その結果がサービスに伝わり認証が完了する。場合によっては、アプリ側で PIN 入力や生体認証が追加で求められることもある。
\item \underline{対応：}連携したスマートフォンが利用できなくなった場合、ユーザー自身がサービス上で連携を解除する手続きが必要である。
\item \underline{回復：}身元確認等を経て、新しいデバイスで再度アプリ連携を設定する。
\end{enumerate}

\begin{prop}[プッシュ通知による認証の認証要素]\label{prop:ch2-push-factor}
モバイル認証アプリ（プッシュ通知）を用いた認証の認証要素は所持情報にあたる。
\end{prop}
\begin{proof}
通知はアカウントと連携した特定のモバイルアプリにだけ送られ、承認操作はそのアプリ上でしか行えない。
よって、この方式は「アカウントと連携した特定のモバイルアプリ（が動作するデバイス）を持っていること」を確認する方式であり、認証要素は所持情報にあたる。
\end{proof}

\begin{dfn}[Decoupled Authentication]\label{dfn:ch2-decoupled}
ログイン操作を行うデバイスと、認証の承認操作を行うデバイスが異なる認証を、\underline{Decoupled Authentication} と呼ぶことがある。
モバイル認証アプリ（プッシュ通知）を用いた認証は、この意味で Decoupled Authentication と呼ばれることもある。
\end{dfn}

\subsubsection{バックアップコードを用いた認証}

多要素認証を設定すると、普段使っている認証手段を失ったときにログインできなくなるおそれがある。
これに備える予備の認証手段として、次の方式がある。

\begin{dfn}[バックアップコード]\label{dfn:backup-code}
サービスが生成してユーザーに提示する、一つまたは複数の使い切りのコードを\underline{バックアップコード}という。
バックアップコードを用いた認証は、主に多要素認証を設定しているアカウントで、普段使っている認証手段（例：認証アプリ、スマートフォン）が利用できなくなった場合の予備の認証手段として提供されることが多い方式である。
\end{dfn}

バックアップコードを用いた認証の 4 つの局面は次のとおりである。
\begin{enumerate}
\item \underline{登録：}サービスが一つまたは複数の使い切りのバックアップコードを生成し、ユーザーに提示する。ユーザーはこれを安全な場所に保管する（印刷、メモ、パスワードマネージャーなど）。
\item \underline{認証：}他の認証手段が使えない場合に、ユーザーは保管しておいたバックアップコードの一つを入力し、サービスはそれが有効なコードかを確認する。
\item \underline{対応：}バックアップコードが第三者に知られた可能性がある場合、サービス側で既存のコードを無効化し、再発行する機能がある場合がある。
\item \underline{回復：}全ての認証手段とバックアップコードを失った場合は、通常の身元確認等によるアカウント回復プロセスが必要になる。なお、バックアップコード自体が回復手段であるため、この局面は他の方式とは少し意味合いが異なる。
\end{enumerate}

\begin{prop}[バックアップコードの認証要素と課題]\label{prop:ch2-backup-factor}
バックアップコードを用いた認証の認証要素は所持情報にあたる。
ただし、ユーザーがコードを安全かつ紛失しないように保管できるかが課題となる。
\end{prop}
\begin{proof}
まず、バックアップコードはユーザーが記憶するものではなく、印刷やメモ、パスワードマネージャーなどに記録して保管するものである。
よって、この方式は「バックアップコードそのもの（を記録した媒体やデータ）を持っていること」を確認する方式であり、認証要素は所持情報にあたる。
次に、所持情報は盗まれたり紛失したりする可能性があるので（命題 \ref{prop:ch2-single-factor-risk}）、コードを記録した媒体やデータを第三者に知られないよう安全に、かつ紛失しないように保管できるかが課題となる。
\end{proof}

\subsubsection{多要素認証（MFA）とパスワードレス認証}

これまで見てきた認証方式は、単独で使うこともできるが、組み合わせて使うことも多い。
この項では、認証方式の組合せ方として、パスワード認証に他の方式を追加する多要素認証と、そもそもパスワードを使わないパスワードレス認証を対比する。

\begin{prop}[パスワード認証の強化としての多要素認証]\label{prop:ch2-password-plus-mfa}
パスワード認証（知識情報）は単体ではセキュリティリスクが高いとされている。
そこで、パスワード認証に、認証要素の異なる他の認証方式を組み合わせることでセキュリティを強化するのが、多要素認証（定義 \ref{dfn:mfa}）である。
\end{prop}
\begin{proof}
まず、命題 \ref{prop:ch2-password-factor} より、パスワード認証は推測・使い回し・漏洩などにより不正ログインのリスクが高い。
次に、パスワード認証の認証要素は知識情報であり、OTP 認証、TOTP 認証、プッシュ通知による認証、バックアップコードはいずれも所持情報である（命題 \ref{prop:ch2-otp-factor}、命題 \ref{prop:ch2-totp-factor}、命題 \ref{prop:ch2-push-factor}、命題 \ref{prop:ch2-backup-factor}）。
よって、パスワード認証にこれらのいずれかを加えた組合せ $S$ は $|c(S)|\geq 2$ を満たす多要素認証となり、命題 \ref{prop:ch2-mfa-improves} により、パスワードが突破されても所持情報の認証要素が別の壁として残る。
\end{proof}

\begin{exa}[典型的な MFA の構成]\label{exa:ch2-mfa-config}
「パスワード（知識情報）＋ TOTP（所持情報）」などが、典型的な MFA の構成である。
\end{exa}

\begin{dfn}[2 段階認証]\label{dfn:ch2-two-step}
「パスワード＋秘密の質問」のように、\emph{同一の}認証要素を用いる認証方式を組み合わせたものを\underline{2 段階認証}と呼ぶ。
記号 \ref{nota:ch2-factor-category} でいえば、$|S|\geq 2$ かつ $|c(S)|=1$ の組合せである。
\end{dfn}

\begin{rem}[2 段階認証は多要素認証ではない]\label{rem:ch2-two-step-not-mfa}
定義 \ref{dfn:ch2-two-step} の 2 段階認証は、多要素認証では\emph{ない}。
例えば「パスワード＋秘密の質問」は、どちらも知識情報（秘密の質問の答えは例 \ref{exa:ch2-factor-examples} の「特定の記憶」）なので、$|c(S)|=1$ であり、定義 \ref{dfn:mfa} の条件 $|c(S)|\geq 2$ を満たさない。
\end{rem}

\underline{MFA の導入と要求方法：}
2010 年頃から、コンシューマー向けサービスは、パスワードに関する不正ログイン被害の深刻化を受け、既存のパスワード認証に別の要素を追加する形で MFA の導入が進んだ。
MFA の要求方法には、いくつかのパターンがある。
\begin{enumerate}
\item ユーザー自身の設定に基づいて要求される場合。
\item サービス側がセキュリティポリシーとして強制的に要求する場合。
\item ログイン時の環境やユーザーの行動パターンを分析し、リスクが高いと判断された場合にのみ要求する場合。
\end{enumerate}

\begin{dfn}[リスクベース認証]\label{dfn:ch2-risk-based}
ログイン時の環境やユーザーの行動パターンを分析し、リスクが高いと判断された場合にのみ MFA を要求する方式を、\underline{リスクベース認証}という。
\end{dfn}

\underline{パスワードレス認証：}
一方、パスワード自体の管理の煩雑さや漏洩リスクといった根本的な課題を解決するため、最初からパスワードを使わない認証方式も登場している。

\begin{dfn}[パスワードレス認証]\label{dfn:passwordless}
最初からパスワードを使わない認証方式を、\underline{パスワードレス認証}という。
\end{dfn}

\begin{exa}[パスワードレス認証の例]\label{exa:ch2-passwordless-examples}
「メールや SMS を用いた OTP 認証のみ」（定義 \ref{dfn:otp-auth}）、「モバイル認証アプリのみ」（定義 \ref{dfn:push-auth}）などが、パスワードレス認証に該当する。
\end{exa}

\begin{rem}[MFA とパスワードレス認証の関係]\label{rem:ch2-mfa-vs-passwordless}
シンプルなパスワードレス認証は、見方を変えれば「MFA の構成からパスワード認証を取り除いたもの」と捉えることもできる。
例えば MFA の構成「パスワード＋ OTP」からパスワード認証を取り除くと、例 \ref{exa:ch2-passwordless-examples} の「OTP 認証のみ」になる。
\end{rem}

\subsubsection{フィッシング耐性をもつ FIDO 認証}

多要素認証の普及は、攻撃者の側にも手法の変化を促した。
この項では、当人認証に対する脅威であるフィッシング攻撃がどのように進化したかを述べ、その進化した攻撃にも対抗できる性質（フィッシング耐性）を持つ FIDO 認証を、他の方式と同じ形で整理する。

\begin{dfn}[フィッシング攻撃]\label{dfn:ch2-phishing}
当人認証に対する脅威として、\underline{フィッシング}（Phishing）攻撃がある。
古典的なフィッシング攻撃は、偽のログインページを用意してユーザーに ID とパスワードを入力させ、それを盗み取ることを目的としていた。
\end{dfn}

しかし、多要素認証の普及により、単にパスワードを盗むだけでは不正ログインが難しくなった（命題 \ref{prop:ch2-password-plus-mfa}）。
そこで攻撃手法も進化し、次のような攻撃が登場した。

\begin{dfn}[リアルタイムフィッシング]\label{dfn:ch2-realtime-phishing}
ユーザーと正規サービスの間に攻撃者が介在して認証プロセス全体を中継し、最終的にログイン後のセッション情報（Cookie など）を奪い取ることを目的とするフィッシング攻撃を、\underline{リアルタイムフィッシング}（中間者攻撃型フィッシング、Adversary-in-the-Middle Phishing）という。
\end{dfn}

\begin{prop}[リアルタイムフィッシングには MFA でも防御が困難]\label{prop:ch2-realtime-phishing-mfa}
MFA を導入していても、パスワードや OTP の入力、あるいは認証アプリの同意など、偽サイト上で認証処理が進められる認証方式では、リアルタイムフィッシングへの防御は困難である。
\end{prop}
\begin{proof}
リアルタイムフィッシングでは、攻撃者がユーザーと正規サービスの間で認証プロセス全体を中継する（定義 \ref{dfn:ch2-realtime-phishing}）。
まず、ユーザーが偽サイトに入力したパスワードや OTP は、攻撃者によってそのまま正規サービスへ転送されるので、正規サービスから見れば正しい認証要素が提示されたことになる。
次に、認証アプリの同意のように入力を伴わない方式でも、偽サイトでのログイン操作に合わせて攻撃者が正規サービスへのログインを試み、ユーザーがその通知に同意すれば、認証は正規サービス上で完了する。
よって、MFA であっても認証は成功し、攻撃者はログイン後のセッション情報を奪い取ることができる。
したがって、偽サイト上で認証処理が進められる認証方式では、この攻撃への防御は困難である。
\end{proof}

\begin{dfn}[フィッシング耐性]\label{dfn:ch2-phishing-resistance}
リアルタイムフィッシングのような巧妙なフィッシング攻撃にも対抗できる特性を、\underline{フィッシング耐性}（Phishing Resistance）と呼ぶ。
\end{dfn}

FIDO（Fast Identity Online）認証、特に Web ブラウザでの利用を想定した FIDO2 規格に基づく認証方式は、このフィッシング耐性を持つことが大きな特徴である。

\begin{dfn}[FIDO 認証]\label{dfn:fido-auth}
ユーザーのデバイス上でサービスごとに公開鍵と秘密鍵のペア（\underline{FIDO クレデンシャル}）を生成して公開鍵をサービスに登録し、認証時にはデバイス内の秘密鍵でチャレンジに対するデジタル署名を生成してサービスが公開鍵で検証する認証方式を、\underline{FIDO}（Fast Identity Online）\underline{認証}という。
このとき用いるデバイス（PC、スマートフォン、FIDO セキュリティキー）を \underline{FIDO 認証器}という。
\end{dfn}

FIDO 認証の 4 つの局面は次のとおりである。
\begin{enumerate}
\item \underline{登録：}ユーザーは自身のデバイス（PC、スマートフォン、または FIDO セキュリティキー）上で、サービスごとに、FIDO 認証で利用する情報である公開鍵と秘密鍵のペア（FIDO クレデンシャル）を生成する。秘密鍵はデバイス内の安全な領域（セキュアエレメント等）に保管され、公開鍵がサービス側に登録される。
\item \underline{認証：}サービス（Web サイト）は、自身のドメイン情報と\underline{チャレンジ}（ランダムなデータ）をブラウザに送る。ブラウザは、ユーザーに認証操作（例：指紋認証、顔認証、PIN 入力）を要求する。認証操作が成功すると、デバイス内の秘密鍵でチャレンジに対するデジタル署名が生成される。この際、ブラウザはアクセス先のドメインが登録時のドメインと一致することも検証する。生成されたデジタル署名がサービスに送られ、サービスは登録済みの公開鍵を使って署名を検証する。検証が成功すれば認証完了である。
\item \underline{対応：}FIDO 認証器（デバイス）を紛失した場合、サービス側でそのデバイスに紐づく FIDO クレデンシャル（公開鍵）を無効化する必要がある。
\item \underline{回復：}身元確認等を経て、新しいデバイスで改めて FIDO クレデンシャルを生成・登録する。
\end{enumerate}

\begin{prop}[FIDO 認証のフィッシング耐性]\label{prop:ch2-fido-phishing-resistance}
FIDO 認証はフィッシング耐性（定義 \ref{dfn:ch2-phishing-resistance}）を持つ。
\end{prop}
\begin{proof}
FIDO 認証の認証の局面 (2) では、ブラウザがアクセス先のドメインが登録時のドメインと一致することを検証する。
すなわち、認証プロセスにおいてサービスのドメイン情報が検証される。
よって、ユーザーが誤ってフィッシングサイトにアクセスしても、そのドメインは正規サービスの登録時のドメインと一致しないので、正規のサービスに対する認証プロセスは実施されない。
したがって、攻撃者が偽サイトを介して正規サービスへの認証を中継しようとしても（定義 \ref{dfn:ch2-realtime-phishing}）、正規サービスに対する署名は得られない。
これがフィッシング耐性の仕組みである。
\end{proof}

\begin{rem}[FIDO 認証の課題]\label{rem:ch2-fido-issue}
FIDO 認証では、秘密鍵がデバイスに固定され、原則として取り出せないことがセキュリティの高さにつながる。
その一方で、デバイスの紛失や機種変更時の再登録が面倒であるという課題があった。
\end{rem}

\subsubsection{パスキー認証の登場}

注意 \ref{rem:ch2-fido-issue} で述べた、デバイスの紛失や機種変更時の再登録という FIDO 認証の課題を解決し、利便性を向上させるために登場したのが、パスキーによる認証である。

\begin{dfn}[パスキー]\label{dfn:passkey}
FIDO クレデンシャル（定義 \ref{dfn:fido-auth}）の一種で、その秘密鍵をデバイス間やクラウドサービス（例：Apple iCloud キーチェーン、Google パスワードマネージャー）経由で同期・バックアップできるようにしたものを、\underline{パスキー}（Passkeys）という。
\end{dfn}

\begin{rem}[「パスキー」という用語]\label{rem:ch2-passkey-term}
「パスキー」という用語は、文脈により別のものを指し示すことがある。
以下では定義 \ref{dfn:passkey} の意味で用いる。
\end{rem}

\begin{prop}[パスキーの利便性とフィッシング耐性]\label{prop:ch2-passkey-properties}
パスキーを用いると、ユーザーは複数のデバイスで同じパスキーを利用したり、デバイスを買い替えても簡単に引き継いだりできる。
また、パスキーはフィッシング耐性を備えている。
\end{prop}
\begin{proof}
まず、パスキーでは秘密鍵をデバイス間やクラウドサービス経由で同期・バックアップできる（定義 \ref{dfn:passkey}）ので、秘密鍵が一つのデバイスに固定されず、複数のデバイスで同じパスキーを使ったり、新しいデバイスへ引き継いだりできる。
よって、注意 \ref{rem:ch2-fido-issue} の再登録の課題が解消される。
次に、パスキーも FIDO 認証の規格に基づいているので、認証プロセスでドメイン情報が検証され、命題 \ref{prop:ch2-fido-phishing-resistance} と同じ理由によりフィッシング耐性を備えている。
\end{proof}

\underline{認証方式の変遷のまとめ：}
このように、デジタル世界の当人認証方式は、パスワード認証から始まり、多要素認証による強化、そしてフィッシング耐性を持つ FIDO 認証やパスキー認証へと、脅威に対抗しながら利便性も考慮しつつ進化を続けている。

\begin{rem}[認証方式を理解するための観点]\label{rem:ch2-evaluation-viewpoints}
認証方式の特徴を理解する際には、認証の 3 要素（知識・所持・生体）だけでなく、フィッシング耐性といったセキュリティ特性や利便性も重要な観点となる。
\end{rem}

\subsubsection{コラム：ユーザー認証における脅威分析のための 2 つの軸}

この節で紹介したように、世の中には多種多様な認証方式が存在する。
このコラムでは、これらのユーザー認証に対する脅威を 2 つの軸で分析する方法を紹介する。
このフレームワークを用いることで、どの脅威が特に危険で早急な対策が必要なのか、あるいはどのような条件下で危険度が増すのかを、体系的に整理できる。

\begin{dfn}[脅威分析の 2 つの軸]\label{dfn:ch2-threat-axes}
ユーザー認証に対する脅威を分析するための、次の 2 つの軸を考える。
\begin{enumerate}
\item \underline{ターゲットとなるユーザーのアクションが必要かどうか}：攻撃を成功させるために、正規ユーザー自身が何らかの操作（例：偽サイトでの情報入力、不審なリンクのクリック、マルウェアの実行許可など）を自発的に行う必要があるか否かを判断する軸。
\item \underline{ターゲットとなるユーザーのデバイスや環境にアクセスする必要があるか}：攻撃者がユーザーのデバイスやその周辺環境（メールアカウント、ブラウザのセッション情報など）への物理的または論理的（リモートアクセスを含む）なアクセス権を必要とするか否かを判断する軸。
\end{enumerate}
\end{dfn}

\begin{nota}[脅威の 4 象限]\label{nota:ch2-threat-quadrants}
各脅威 $t$ に対して、軸 (1) の値を $a(t)\in\{\text{不要},\ \text{必要}\}$、軸 (2) の値を $d(t)\in\{\text{不要},\ \text{必要}\}$ とする。
脅威は組 $(a(t), d(t))$ によって 4 つの象限に分類される。
表 \ref{tab:ch2-threat-quadrants} のように、列を軸 (1)（正規ユーザーのアクション）、行を軸 (2)（デバイスや環境へのアクセス）にとって配置し、左上・右上・左下・右下の象限と呼ぶ。
\end{nota}

\begin{table}[H]
\centering
\caption{脅威の分類と特徴}\label{tab:ch2-threat-quadrants}
\begin{tabular}{|p{0.1\textwidth}|p{0.06\textwidth}|p{0.36\textwidth}|p{0.36\textwidth}|}
\hline
\multicolumn{2}{|c|}{} & \multicolumn{2}{c|}{正規ユーザーのアクション} \\ \cline{3-4}
\multicolumn{2}{|c|}{} & 不要 & 必要 \\ \hline
\multirow{6}{=}{デバイスや環境へのアクセス}
& 不要 & （左上）ユーザーの知らないところで行われる。\newline 例：パスワードリスト攻撃、SIM スワップやメールアカウントを奪取された状態からの MFA 実行 & （右上）正規のサービスを装った誘導などを伴う。\newline 例：リアルタイムフィッシング攻撃 \\ \cline{2-4}
& 必要 & （左下）デバイス盗難などの物理的なアクセス、マルウェアなどの論理的なアクセスを伴う。\newline 例：セキュリティキーやスマートフォンを奪取し、ロック解除済みの状態からの MFA 施行 & （右下）正規ユーザーと同等の権限を持った状態での脅威。\newline 例：物理的な脅迫を伴う認証突破 \\ \hline
\end{tabular}
\end{table}

以下、各象限の脅威の特徴と、求められる対策を述べる。

\underline{左上：正規ユーザーのアクション、デバイスや環境へのアクセスがどちらも不要}

\begin{prop}[左上の象限の脅威には早急な対策が必要]\label{prop:ch2-quadrant-upper-left}
左上の象限に分類される脅威は、攻撃者がリモートから、ユーザーの行動やデバイスへの直接的なアクセスなしに不正ログインを試みるものである。
代表的な例としてパスワードリスト攻撃が挙げられる。
この象限の脅威は、ユーザーが何も知らずに被害に遭う可能性が高いので、サービス側はここに当てはまる脅威を残さないようにすべく、早急な対策が必要である。
\end{prop}
\begin{proof}
この象限の脅威は、ユーザーのアクションもデバイスや環境へのアクセスも必要としない。
よって、攻撃はユーザーの関与なしに、ユーザーの知らないところで進行し、ユーザーが異変に気づいて防ぐ機会がない。
したがって、ユーザーが何も知らずに被害に遭う可能性が高く、ユーザー側の注意では防げないので、サービス側が早急に対策してこの象限の脅威を残さないようにする必要がある。
\end{proof}

\underline{右上：正規ユーザーのアクションが必要／デバイスや環境へのアクセス不要}

\begin{prop}[右上の象限の脅威にはシステムと啓発の両面が必要]\label{prop:ch2-quadrant-upper-right}
右上の象限の脅威は、攻撃者がリモートからユーザーを誘導し、ユーザー自身の行動によって被害が発生するものである。
リアルタイムフィッシング攻撃が典型例であり、攻撃者はユーザーのデバイスに直接アクセスする必要はない。
この象限の脅威に対しては、サービスがフィッシング耐性のある認証方式を用意して使ってもらうシステム的な防御と、ユーザーへの啓発の両面からのアプローチが重要である。
\end{prop}
\begin{proof}
この象限の脅威は、ユーザー自身の行動によって被害が発生するので、ユーザーの「騙されやすさ」に依存する。
まず、システム的には、フィッシング耐性のある認証方式（命題 \ref{prop:ch2-fido-phishing-resistance}、命題 \ref{prop:ch2-passkey-properties}）を用意して使ってもらえば、ユーザーが騙されて偽サイトにアクセスしても正規サービスへの認証は実施されない。
次に、ユーザーが騙されること自体を減らすには、ユーザーへの啓発が必要である。
よって、両面からのアプローチが重要である。
\end{proof}

\underline{左下：正規ユーザーのアクション不要／デバイスや環境へのアクセスが必要}

\begin{prop}[左下の象限の脅威とデバイスのセキュリティ管理]\label{prop:ch2-quadrant-lower-left}
左下の象限の脅威は、攻撃者がユーザーのデバイスやその周辺環境へのアクセス権を既に持っている状態（例：デバイスの盗難、マルウェア感染によるリモートアクセス）で、ユーザーの明示的なアクションなしに不正ログインを試みるものである。
リモートからの攻撃に比べると発生頻度は低いかもしれないが、一度アクセスを許すと深刻な被害につながる可能性がある。
この象限の脅威は、デバイスのセキュリティ管理が鍵となる。
\end{prop}
\begin{proof}
まず、この象限の攻撃はデバイスへのアクセスを前提とするので、リモートからの攻撃に比べて攻撃の前提条件が厳しく、発生頻度は低いかもしれない。
一方で、攻撃者がデバイスや環境へのアクセス権を持ってしまえば、ユーザーのアクションなしに不正ログインを試みられるので、一度アクセスを許すと深刻な被害につながる可能性がある。
よって、攻撃者にデバイスや環境へのアクセスを許さないこと、すなわちデバイスのセキュリティ管理が鍵となる。
\end{proof}

\begin{rem}[所持情報の普及とデバイス管理の重要性]\label{rem:ch2-device-management}
所持情報を用いる認証方式の普及とともに、デバイスや環境の管理の重要性は増している。
\end{rem}

\underline{右下：正規ユーザーのアクション、デバイスや環境へのアクセスがどちらも必要}

\begin{prop}[右下の象限の脅威の特徴]\label{prop:ch2-quadrant-lower-right}
右下の象限は、攻撃を成功させるためにユーザーのアクションとデバイスや環境へのアクセスの両方が必要となる、最も条件が厳しい脅威である。
この象限の脅威は、他の象限の脅威に比べると発生頻度は極めて低いが、万一発生した場合はユーザーにとって精神的・物理的に大きな負担となる。
また、サービスにとっては、正規のユーザーが自身の意思に基づいて行った処理との区別が難しいという側面もある。
\end{prop}
\begin{proof}
まず、この象限の攻撃は 2 つの軸の条件を両方とも満たす必要があるので、最も条件が厳しく、発生頻度は極めて低い。
次に、例えばユーザーを物理的に脅迫し、目の前でデバイスを操作させてログインさせるようなケースが考えられ、このような事態はユーザーにとって精神的・物理的に大きな負担となる。
さらに、このとき認証操作は正規ユーザー本人が自分のデバイスで行っているので、サービスから見ると、正規のユーザーが自身の意思に基づいて行った処理との区別が難しい。
\end{proof}

\begin{exa}[右下の象限の脅威の例]\label{exa:ch2-lower-right-example}
例えば、金融系サービスが管理する資産情報を元に、実際の自宅に強盗が押し入るといった犯罪が行われる可能性がある。
\end{exa}

\underline{特定の認証における脅威の整理：}
このフレームワークは、ざっくりとした分類以外に、認証方式の評価にも使える。
例として、パスワード認証（定義 \ref{dfn:password-auth}）において発生しうる脅威を整理すると、表 \ref{tab:ch2-password-threats} のようになる。

\begin{table}[H]
\centering
\caption{パスワード認証の脅威の分類}\label{tab:ch2-password-threats}
\begin{tabular}{|p{0.1\textwidth}|p{0.06\textwidth}|p{0.36\textwidth}|p{0.36\textwidth}|}
\hline
\multicolumn{2}{|c|}{} & \multicolumn{2}{c|}{正規ユーザーのアクション} \\ \cline{3-4}
\multicolumn{2}{|c|}{} & 不要 & 必要 \\ \hline
\multirow{6}{=}{デバイスや環境へのアクセス}
& 不要 & 複数のサービスでパスワードを使い回している。\newline 対策：パスワードマネージャーの使用 & 使い回していないが、パスワードマネージャーを使っていない。\newline 対策：パスワードマネージャーの使用 \\ \cline{2-4}
& 必要 & パスワードマネージャーが動いているスマートフォンが盗まれ、アンロックされた。\newline 対策：適切なデバイス管理 & 言われるがままに正規のサービスにログインをしなければここまでは…… \\ \hline
\end{tabular}
\end{table}

同様に、メールや SMS によるワンタイムパスワード（OTP）認証やパスキー認証など、それぞれの認証方式における脅威を分類できる。

\begin{prop}[脅威の重心の変化]\label{prop:ch2-threat-shift}
これまでの認証セキュリティは、主にリモートからの攻撃、特にパスワードリスト攻撃やフィッシング攻撃といった左上や右上の象限に焦点が当てられてきた。
しかし、今後は脅威の重心が変化し、左下の象限の脅威、すなわちユーザーのデバイスや環境自体が攻撃のターゲットとなるケースについて、より深く検討し、対策を講じる必要が出てくると考えられる。
\end{prop}
\begin{proof}
まず、多要素認証で所持情報を利用する認証方式が加えられるようになったことで、パスワードを盗むだけでは不正ログインが難しくなった（命題 \ref{prop:ch2-password-plus-mfa}）。
次に、パスキーのようなフィッシング耐性の高い認証方式が普及するにつれて、右上の象限のリアルタイムフィッシングも成功しにくくなる（命題 \ref{prop:ch2-passkey-properties}）。
こうして左上・右上の象限の脅威への対策が進むと、所持情報を持つデバイスや環境そのものへのアクセスが攻撃者にとって残された手段となる（注意 \ref{rem:ch2-device-management}）。
よって、脅威の重心は変化していき、「正規ユーザーのアクション不要、デバイスや環境へのアクセスが必要」という左下の象限の脅威について、より深く検討し対策を講じる必要が出てくると考えられる。
\end{proof}

\begin{rem}[今後重要になる対策]\label{rem:ch2-future-measures}
命題 \ref{prop:ch2-threat-shift} により、デバイスのセキュリティ、マルウェア対策、そしてデバイス紛失時の対応などが、これまで以上に重要になってくる。
\end{rem}

%% =======================================================
%%  part: ch2b
%% =======================================================
%% ===== ch2b: 2-4 ID連携, 2-5 アクセス制御, 2-6 セッション管理 =====

\subsection{ID連携}\label{subsec:ch2b-identity-federation}

ここまでで、身元確認と当人認証の具体的な方法を見てきた。しかし、利用者が新しいサービスを使うたびに、それらを一から行うのは利用者にとってもサービス提供者にとっても負担が大きい。この項では、ある主体が行った確認の結果を別の主体が信頼して利用する仕組みである ID 連携（定義 \ref{dfn:identity-federation}）を、現実世界の例から出発して精密化する。まず現実世界の「信頼の連携」を例で確認し、次に NIST とデジタル庁による定義、登場人物（申請者／利用者、IdP、RP）とアサーションを定義する。その上で、ID 連携を身元確認に使う場合と当人認証に使う場合を順に論じ、最後に ID 連携を実現するフェデレーションプロトコル（SAML と OIDC）を紹介する。

\subsubsection{現実世界のID連携}

ID 連携の考え方は、デジタル技術に固有のものではない。この項では、身の回りにある「ある組織の保証を別の組織が信頼する」場面を観察し、身元確認との違いを明確にしておく。

\begin{prop}[確認結果の再利用によるコストと手間の削減]\label{prop:ch2b-reuse-of-verification}
世の中のあらゆるサービスを利用するたびに、身元確認や当人認証を一から行うのは大変である。あるサービスで行った確認結果を、別のサービスが信頼して利用できれば、サービス提供側のコストも、利用者の手間も大きく減らせる。
\end{prop}

\begin{proof}
確認の作業がどこで発生するかを数えればよい。まず、確認結果を共有しない場合、利用者はサービスを一つ利用するたびに身元確認や当人認証の手続きを行い、各サービス提供者もそれぞれ確認の仕組みを用意して運用しなければならない。次に、あるサービスで行った確認結果を別のサービスが信頼して利用できるならば、後者のサービスは自ら確認を行う必要がなくなり、利用者も同じ手続きを繰り返さずに済む。よって、サービス提供側のコストと利用者の手間の双方が大きく減る。
\end{proof}

この「信頼の連携」に近い考え方は、私たちの身の回りにも存在する。普段使っている身分証明書は、発行元の組織による身元確認の結果である。

\begin{rem}[身分証明書を使う身元確認とID連携の違い]\label{rem:ch2b-proofing-vs-federation}
同じ身分証明書を使う場面でも、何を信頼しているかによって二つを区別する。
\begin{enumerate}
\item 身分証明書を使って、利用者が\emph{申告した}情報が正しいかを確認するのが身元確認（定義 \ref{dfn:identity-proofing}）である。
\item 証明書が持つ「この人は確かに本人である」「この人は特定の資格を持っている」といった\emph{発行元の保証そのもの}を信頼し、別のサービスで利用するのが ID 連携である。
\end{enumerate}
ID 連携では、時には、それによって本人であることの証明（当人認証，定義 \ref{dfn:authentication}）も兼ねることがある。
\end{rem}

\underline{現実世界におけるID連携の例：}以下に、現実世界で発行元の保証が別の組織に受け入れられている例を挙げる。

\begin{exa}[学生証による学生料金]\label{exa:ch2b-student-id}
映画館や美術館などで学生証を提示すると、学生料金でサービスを利用できる。これは、施設側が「学生証の発行元（学校）が『この人は学生である』と保証している」ことを信頼し、その情報を受け入れているからである。ここで、保証を発行する組織は学校であり、保証を信頼して受け入れる組織は映画館や美術館である。
\end{exa}

\begin{exa}[運転免許証とマイナンバーカードの利用]\label{exa:ch2b-license-mynumber}
この他にも、次のような場面は ID 連携である。
\begin{enumerate}
\item レンタカーを借りる際に運転免許証を提示して、運転資格を証明すること。
\item マイナンバーカードを健康保険証として利用すること。
\end{enumerate}
\end{exa}

このように、現実世界でも、ある組織が発行した証明や保証を、別の組織が信頼して受け入れることで、手続きがスムーズになったり、特定のサービスが利用できたりする場面がある（原書の図「現実世界のID連携」はこの関係を描いたものである）。

\subsubsection{デジタル世界におけるID連携}

現実世界の例で見たような「信頼の連携」の考え方は、デジタル世界でどのように実現されているのだろうか。この項では、NIST とデジタル庁のガイドラインによる ID 連携の定義を確認し、ID 連携の主要な登場人物と、連携される情報であるアサーションを定義する。

NIST は、ID 連携に関するガイドライン（SP 800-63 Rev.4）で、ID 連携を次のように定義している。

\begin{quote}
\textit{Federation is a process that enables the subscriber account defined in SP800-63A to be used with an RP that does not verify one of the authenticators bound to the subscriber account.}

\textit{Instead, a service known as an identity provider, or IdP, makes the subscriber account available through a federation protocol to the relying party, or RP.}
\end{quote}

これを日本語で整理すると、次の定義になる。これは第1章の定義 \ref{dfn:identity-federation} を、デジタル世界について精密化したものである。

\begin{dfn}[ID連携（NIST SP 800-63 Rev.4 による定義）]\label{dfn:ch2b-federation-nist}
\underline{ID 連携}（federation）とは、加入者アカウント（ユーザーアカウント）に紐づく認証器（パスワードなど）を直接検証しない RP（定義 \ref{dfn:rp}）に対して、そのアカウントを使えるようにするプロセスである。その代わりに、IdP（定義 \ref{dfn:idp}）と呼ばれるサービスが、フェデレーションプロトコル（定義 \ref{dfn:federation-protocol}）を通じて加入者アカウントを RP で利用可能にする。
\end{dfn}

ここでいう認証器の例であるパスワードについては、定義 \ref{dfn:password-auth} を参照せよ。この定義にあるように、ID 連携は IdP（Identity Provider）と呼ばれる主体がユーザーの認証を行い、その結果を RP（Relying Party）と呼ばれる別の主体に引き継ぐことで、RP 側での当人認証などを可能にする仕組みである。RP は自分では認証器を検証しないが、IdP が検証した結果を引き継ぐことで、当人認証を実現する。

また、デジタル庁の「DS-511 行政手続等での本人確認におけるデジタルアイデンティティの取扱いに関するガイドライン」（令和7年9月30日発行）では、ID 連携（文脈上、当人認証における連携を指す）を次のように定義している。

\begin{dfn}[ID連携（デジタル庁 DS-511 による定義）]\label{dfn:ch2b-federation-ds511}
\underline{ID 連携}とは、身元確認や当人認証を、他者に依拠して実現することである。具体的には、信頼できる ID プロバイダと連携し、ID プロバイダによって行われた身元確認や当人認証の結果に関する情報を入手することで、対象手続における本人確認を実現する。
\end{dfn}

「本人確認」という語が身元確認と当人認証のどちらを指しうるかについては、注意 \ref{rem:honnin-kakunin} を参照せよ。

ここで、ID 連携の仕組みを理解するために、主要な登場人物を整理しておく。

\begin{dfn}[申請者／利用者]\label{dfn:ch2b-applicant-user}
\underline{申請者／利用者}とは、RP（依拠当事者，定義 \ref{dfn:rp}）のサービスを利用したい主体である。
\end{dfn}

\begin{dfn}[IDプロバイダ（IdP）]\label{dfn:idp}
\underline{ID プロバイダ}（IdP: Identity Provider）とは、申請者／利用者に対する身元確認や当人認証を行い、その結果を証明する主体である。
\end{dfn}

\begin{dfn}[依拠当事者（RP）]\label{dfn:rp}
\underline{依拠当事者}（RP: Relying Party）とは、自身のサービスにおける身元確認や当人認証を、IdP に依拠（信頼）する主体である。
\end{dfn}

ID 連携では、IdP が申請者／利用者に対して行った身元確認や当人認証の結果を RP に連携する。この連携される情報に名前を付けておく。

\begin{dfn}[アサーション]\label{dfn:ch2b-assertion}
ID 連携において、IdP が申請者／利用者に対して行った身元確認や当人認証の結果として RP に連携される情報を、\underline{アサーション}（assertion）という。
\end{dfn}

\begin{dfn}[フェデレーションプロトコル]\label{dfn:federation-protocol}
ID 連携において、IdP と RP の間でアサーションをやり取りする際に使用される通信規格の総称を、\underline{フェデレーションプロトコル}（federation protocol）という。代表的なものに SAML（定義 \ref{dfn:ch2b-saml}）と OIDC（定義 \ref{dfn:ch2b-oidc}）があり、本項の最後の小項で扱う。
\end{dfn}

以上の登場人物とアサーションの関係（原書の図「ID連携の概念」）を図 \ref{fig:ch2b-federation-concept} にまとめる。IdP は申請者／利用者に対して身元確認や当人認証を行い、その結果をアサーションとして RP に連携する。RP はそのアサーションを信頼して、申請者／利用者にサービスを提供する。

\begin{figure}[H]
\centering
\begin{tikzpicture}[
  node distance=14mm and 20mm,
  box/.style={draw, rounded corners, align=center, minimum width=34mm, minimum height=12mm},
  >={Stealth[length=2.5mm]}
]
\node[box] (user) {申請者／利用者};
\node[box, above right=of user] (idp) {IdP\\（ID プロバイダ）};
\node[box, below right=of idp] (rp) {RP\\（依拠当事者）};
\draw[->] (idp) -- node[sloped, above, font=\small]{身元確認・当人認証} (user);
\draw[->] (idp) -- node[sloped, above, font=\small]{アサーション} (rp);
\draw[->] (user) -- node[below, font=\small]{サービスの利用} (rp);
\end{tikzpicture}
\caption{ID連携の概念：IdP が行った身元確認・当人認証の結果をアサーションとして RP に連携し、RP はそれに依拠する。}
\label{fig:ch2b-federation-concept}
\end{figure}

\subsubsection{身元確認としての利用}

ID 連携は、RP が身元確認（定義 \ref{dfn:identity-proofing}）を IdP に依拠するために使うことができる。この項では、オンラインサービスで身元確認が必要となる場面を確認し、ID 連携を使うと何が省略できるのか、またどこに注意が必要なのかを論じる。

オンラインサービスで身元確認が必要となるのは、主に次のような場面である。
\begin{enumerate}
\item サービスの初回登録時（新規登録の詳細は第3章で扱う。定義 \ref{dfn:signup}）。
\item 当人認証に用いる認証器を紛失した際のアカウントリカバリー（当人認証の局面としては定義 \ref{dfn:authenticator-phases} の「アカウントの回復」にあたり、詳細は第3章で扱う。定義 \ref{dfn:account-recovery}）。
\end{enumerate}

\begin{prop}[ID連携による身元確認プロセスの簡略化]\label{prop:ch2b-federation-proofing}
ID 連携を利用する場合、RP は身元確認を IdP に依拠することで、RP 自身の確認プロセスを簡略化または省略できる。
\end{prop}

\begin{proof}
ID 連携を利用しない場合と利用する場合を比較する。まず、ID 連携を利用しない場合は、サービス提供者（RP）が自ら利用者の身元確認プロセスを実施しなければならない。次に、ID 連携を利用する場合は、RP がこの身元確認を IdP に依拠する。すなわち、信頼できる IdP が行った身元確認の結果（確認済みの属性情報など。属性情報については定義 \ref{dfn:attribute}）をアサーション（定義 \ref{dfn:ch2b-assertion}）として受け取り、それを信頼する。RP はその結果を自分で確認し直す必要がないので、RP 自身の確認プロセスを簡略化または省略できる。
\end{proof}

\begin{exa}[ID連携によるメールアドレス確認の省略]\label{exa:ch2b-verified-email}
アカウント登録時にメールアドレスの有効性確認を行うサービス（RP）があるとする。このサービスが ID 連携を利用する場合、すでにメールアドレス確認を済ませている IdP から「確認済みメールアドレスである」という情報を含むアサーションを受け取る。RP はこのアサーションを信頼することで、自らメール確認を行う手間を省き、アカウント登録を完了させることができる（原書の図「ID連携を用いた身元確認」）。
\end{exa}

\begin{figure}[H]
\centering
\begin{tikzpicture}[
  box/.style={draw, rounded corners, align=center, minimum width=36mm, minimum height=12mm},
  >={Stealth[length=2.5mm]}
]
\node[box] (idp) {IdP\\（メールアドレス確認済み）};
\node[box, right=40mm of idp] (rp) {RP\\（アカウント登録を完了）};
\draw[->] (idp) -- node[above, font=\small, align=center]{アサーション\\「確認済みメールアドレス」} (rp);
\end{tikzpicture}
\caption{ID連携を用いた身元確認：RP は IdP のアサーションを信頼し、自らのメール確認を省く。}
\label{fig:ch2b-federation-proofing}
\end{figure}

さらに、厳格な身元確認が必要な場合にも ID 連携は効果を持ちうる。

\begin{prop}[厳格な身元確認におけるID連携の効果]\label{prop:ch2b-strict-proofing}
運転免許証やマイナンバーカードなどを用いた厳格な身元確認が必要な場合、ID 連携を利用すれば、氏名、住所といった基本情報の入力や、本人確認書類画像のアップロードや IC チップを用いたそれらの確認作業といったプロセスを省略できる可能性がある。これにより、サービスの開発・運用コストの削減や、ユーザーの利便性向上が期待できる。
\end{prop}

\begin{proof}
命題 \ref{prop:ch2b-federation-proofing} を厳格な身元確認に当てはめる。まず、厳格な身元確認では、利用者は基本情報（氏名、住所など）を入力し、本人確認書類の画像をアップロードするか IC チップを読み取らせ、サービス側はそれらを確認する。次に、IdP がすでにこれらの確認を行っているならば、RP はその結果をアサーションとして受け取って依拠でき、これらのプロセスを自ら行わずに済む可能性がある。よって、RP はこれらの確認の仕組みを開発・運用するコストを削減でき、利用者は入力やアップロードの手間を省けるので利便性が向上する。
\end{proof}

\begin{rem}[ID連携で代替できる身元確認のレベルは規制による]\label{rem:ch2b-proofing-regulation}
実際にどのレベルの身元確認を ID 連携で代替できるかは、そのサービスが準拠すべき法律や規制（例：犯罪収益移転防止法など）によって要件が異なるため、注意が必要である。命題 \ref{prop:ch2b-strict-proofing} が「省略できる\emph{可能性がある}」と述べているのはこのためである。
\end{rem}

\subsubsection{当人認証としての利用}

ID 連携のもう一つの使い方は、RP が当人認証（定義 \ref{dfn:authentication}）を IdP に依拠することである。この項では、当人認証が必要となる場面と、ID 連携を使わない場合の問題点を確認し、ID 連携によってシングルサインオンのような利便性が得られることを示す。

オンラインサービスで当人認証（ログインなど）が必要となるのは、主に次のような場面である。
\begin{enumerate}
\item 登録済みサービスへのログイン（ログインの詳細は第3章で扱う。定義 \ref{dfn:login}）。
\item 重要な操作（例：登録情報の変更、決済など）の実行時。
\end{enumerate}
などである。

\begin{prop}[ID連携を使わない場合のクレデンシャル管理の煩雑さ]\label{prop:ch2b-credential-burden}
ID 連携を利用しない場合、サービス提供者（RP）は自らログイン機能などを開発・運用しなければならず、ユーザーはサービスごとに ID やパスワード（\underline{クレデンシャル}，credential）を管理する必要がある。利用するサービスが増えるほど、クレデンシャルの管理は煩雑になりがちである。
\end{prop}

\begin{proof}
まず、ID 連携を利用しない場合、各 RP は当人認証を他者に依拠できないので、自らログイン機能などを開発・運用するしかない。次に、そのため各サービスはそれぞれ独自のクレデンシャルを発行することになり、ユーザーはサービスごとに ID やパスワードを管理しなければならない。よって、ユーザーが管理すべきクレデンシャルの数は利用するサービスの数とともに増え、管理は煩雑になりがちである。
\end{proof}

ID 連携を利用する場合、RP は当人認証プロセスを IdP に依拠する。その流れは次のとおりである（原書の図「ID連携を用いた当人認証」）。
\begin{enumerate}
\item ユーザーが IdP で認証に成功する。
\item IdP はその結果を示すアサーションを RP に送る。
\item RP はそのアサーションを信頼し、「IdP で認証されたユーザー」に対応する自社サービスのアカウントへのログインを許可する。
\end{enumerate}

\begin{dfn}[シングルサインオン（SSO）]\label{dfn:ch2b-sso}
ユーザーが IdP のアカウント一つで複数の RP のサービスにログインできることを、\underline{シングルサインオン}（SSO: Single Sign-On）という。
\end{dfn}

\begin{prop}[ID連携による当人認証の集約]\label{prop:ch2b-federation-sso}
ID 連携を当人認証に利用すると、ユーザーは IdP のアカウント一つで複数の RP のサービスにログインできる、いわゆるシングルサインオン（定義 \ref{dfn:ch2b-sso}）のような利便性を享受できる。
\end{prop}

\begin{proof}
上の流れ (1)--(3) において、ユーザーが認証を行うのは IdP だけであり、各 RP は IdP から受け取ったアサーションを信頼して自社サービスのアカウントへのログインを許可する。この流れは RP がいくつあっても同じ IdP に対して成り立つ。よって、ユーザーは IdP のアカウント一つを管理するだけで、複数の RP のサービスにログインできる。これは命題 \ref{prop:ch2b-credential-burden} で述べたサービスごとのクレデンシャル管理の煩雑さを解消する。
\end{proof}

\begin{figure}[H]
\centering
\begin{tikzpicture}[
  box/.style={draw, rounded corners, align=center, minimum width=30mm, minimum height=11mm},
  >={Stealth[length=2.5mm]}
]
\node[box] (user) {ユーザー};
\node[box, right=36mm of user] (idp) {IdP};
\node[box, below=16mm of idp] (rp) {RP\\（自社サービスのアカウント）};
\draw[->] (user) -- node[above, font=\small]{(1) IdP で認証} (idp);
\draw[->] (idp) -- node[right, font=\small]{(2) アサーション} (rp);
\draw[->] (rp) -- node[sloped, below, font=\small]{(3) ログインを許可} (user);
\end{tikzpicture}
\caption{ID連携を用いた当人認証の流れ}
\label{fig:ch2b-federation-authn}
\end{figure}

\subsubsection{ID連携を実現するためのフェデレーションプロトコル}

ID 連携が実際に安全に機能するためには、IdP と RP がアサーションをどのような形式と手順でやり取りするかを取り決めておく必要がある。この項では、フェデレーションプロトコル（定義 \ref{dfn:federation-protocol}）の役割を確認し、その代表例である SAML と OIDC を定義する。

IdP と RP の間で安全かつ標準化された方法でアサーションをやり取りするために、フェデレーションプロトコルが用いられる。代表的なものに、SAML と OIDC（OpenID Connect）がある。

\begin{dfn}[SAML]\label{dfn:ch2b-saml}
\underline{SAML}（Security Assertion Markup Language）は、OASIS（構造化情報標準促進協会）によって 2000 年代初頭に策定された標準規格である。XML というデータ形式をベースにしており、認証情報やユーザー属性情報を安全に交換することに重点が置かれている。
\end{dfn}

SAML は、主に企業内システムやクラウドサービス（SaaS）間の認証連携（シングルサインオン）で広く利用されている。

\begin{dfn}[OIDC（OpenID Connect）]\label{dfn:ch2b-oidc}
\underline{OIDC}（OpenID Connect）は、OpenID Foundation によって 2014 年に策定された比較的新しいプロトコルである。IETF により同時期に策定された OAuth 2.0 というアクセス許可の仕組みをベースとして、当人認証、身元確認を実現するための仕様である。JSON という軽量なデータ形式と Web API（REST API）をベースにしている。
\end{dfn}

OIDC は JSON と Web API をベースにしているため、Web アプリケーションやスマートフォンアプリでのログイン実装に適している。

\begin{exa}[ソーシャルログイン]\label{exa:ch2b-social-login}
皆さんもよく利用する「Google でログイン」「LINE でログイン」といったソーシャルログイン機能の多くは、OIDC（定義 \ref{dfn:ch2b-oidc}）を基盤としている。この場合、Google や LINE が IdP、ログインボタンを置いているサービスが RP にあたる。
\end{exa}

\begin{rem}[OIDCの仕様拡張]\label{rem:ch2b-oidc-extension}
OpenID ファウンデーションでは、現在も新しい利用シーンに対応するための仕様拡張などが活発に進められている。
\end{rem}

以上の SAML と OIDC についての記述を、表 \ref{tab:ch2b-saml-oidc} に対比してまとめておく（内容は上の記述の再整理である）。

\begin{table}[H]
\centering
\caption{SAML と OIDC の対比}
\label{tab:ch2b-saml-oidc}
\begin{tabularx}{\linewidth}{|l|X|X|}
\hline
 & SAML & OIDC \\
\hline
策定 & OASIS（2000 年代初頭） & OpenID Foundation（2014 年） \\
\hline
基盤 & --- & OAuth 2.0（IETF により同時期に策定されたアクセス許可の仕組み） \\
\hline
データ形式 & XML & JSON と Web API（REST API） \\
\hline
主な利用場面 & 企業内システムやクラウドサービス（SaaS）間の認証連携（シングルサインオン） & Web アプリケーションやスマートフォンアプリでのログイン、ソーシャルログイン \\
\hline
\end{tabularx}
\end{table}

\subsection{アクセス制御}\label{subsec:ch2b-access-control}

身元確認と当人認証を経てサービスの利用が許可されたとしても、それだけで利用者に何をさせてよいかが決まるわけではない。この項では、第1章で概要を述べたアクセス制御（定義 \ref{dfn:access-control}）を、「主体・対象・操作」という三つの要素を使って精密化する。まず現実世界の入退室管理を例に考え方を確認し、次にデジタル世界でのアクセス制御を定義して、代表的な五つの方式（DAC, MAC, RBAC, ABAC, PBAC）を定義・比較する。

\subsubsection{現実世界のアクセス制御}

アクセス制御が何を制限するものなのかを、まず現実世界で確かめる。この項では、アクセス制御の三要素（主体・対象・操作）を言葉で導入し、社員証による入退室管理を例として、それがデジタル世界でも重要であることを論じる。

これまで紹介した身元確認、当人認証を経てサービスの利用が許可されたとしても、誰もがそのサービスの全ての機能を利用できるわけではない。誰が（どの主体が）、何に対して（どの対象に）、どのような操作（読み取り、書き込み、入室など）を行えるかを適切に制限・管理する仕組み、それがアクセス制御である。以後、アクセス制御を行う者が制限する相手を\underline{主体}（subject）、制限される資源を\underline{対象}（object）、主体が対象に対して行う行為を\underline{操作}（operation）とよぶ。

私たちの身近な現実世界にも、アクセス制御の例はたくさんある。

\begin{exa}[社員証による入退室管理]\label{exa:ch2b-office-badge}
オフィスビルでの入退室管理を考える。
\begin{enumerate}
\item まず、社員証をセキュリティゲートにかざしてビルに入館する。
\item さらに、自分の部署があるフロアや執務室に入る際にも社員証が必要になる。
\item しかし、同じ会社の社員であっても、所属部署や役職によっては立ち入りが許可されていないエリア（例えば、サーバールームや役員室など）がある。
\end{enumerate}
これは、社員という主体に対し、場所という対象への入室操作を、その人の属性（所属部署、役職）に基づいて制御している例である。
\end{exa}

\begin{prop}[デジタル世界におけるアクセス制御の重要性]\label{prop:ch2b-access-control-importance}
このようなアクセス制御の考え方は、デジタル世界においても極めて重要である。
\end{prop}

\begin{proof}
デジタル世界にも保護すべき対象が多数あることから示す。まず、Web サイトの会員ページ、社内システム、クラウド上のファイル、データベースの情報など、デジタル化された様々なリソースが存在する。次に、これらのリソースに対して「誰が」「何に」「何をしてよいか」を適切に管理しなければ、権限のない主体がリソースを読み取ったり操作したりできてしまい、情報漏洩や不正利用といった深刻な問題につながりかねない。よって、デジタル世界でもアクセス制御は極めて重要である。
\end{proof}

\subsubsection{デジタル世界のアクセス制御}

現実世界での「主体・対象・操作」の考え方を、情報システムの言葉に置き換えて定義する。この項では、デジタル世界のアクセス制御を定義したのち、システムの要件やセキュリティポリシーに応じて使い分けられる代表的な方式（モデル）を一つずつ定義し、最後に表で比較する。

現実世界の例で見たように、アクセス制御は「誰が（主体）、何に対して（対象）、どのような操作（権限）を行えるか」を管理する仕組みである。デジタル世界においては、これは次のように言い換えることができる。

\begin{dfn}[デジタル世界のアクセス制御]\label{dfn:ch2b-digital-access-control}
デジタル世界における\underline{アクセス制御}とは、認証されたユーザーやプロセス（主体）が、情報システム上のリソース（データ、機能、アプリケーションなど＝対象）に対して、許可された操作（読み取り、書き込み、実行など＝\underline{権限}）のみを行えるように制限・管理することである。

主体の集合を $S$、対象の集合を $O$、操作の集合を $A$ とし、許可された組の集合を
\[
P\subseteq S\times O\times A
\]
と書けば、アクセス制御とは、主体 $s\in S$ が対象 $o\in O$ に操作 $a\in A$ を行おうとしたとき、$(s,o,a)\in P$ である場合に限りそれを許すことである。
\end{dfn}

これは第1章の定義 \ref{dfn:access-control} の精密化である。以下に述べる各方式は、許可される組の集合 $P$ を\emph{何を基準に、誰が}定めるかの違いとして理解できる。

デジタル世界で利用されているアクセス制御の方式（モデル）はいくつかあり、システムの要件やセキュリティポリシーに応じて使い分けられている。代表的なものを紹介する。

\begin{dfn}[DAC（任意アクセス制御）]\label{dfn:dac}
\underline{DAC}（Discretionary Access Control）とは、ユーザーが自身の所有するリソースに対するアクセス権を自由に設定できる柔軟なアクセス制御のモデルである。すなわち、対象 $o$ に関する許可（$P$ のうち $o$ を含む組）は、$o$ の所有者が定める。
\end{dfn}

DAC は、UNIX や Windows のファイルシステムなど、ユーザーごとの個別設定が求められる環境で使われる。

\begin{dfn}[MAC（強制アクセス制御）]\label{dfn:mac}
\underline{MAC}（Mandatory Access Control）とは、組織や管理者が定義した機密ラベルやポリシーに基づき、厳格にアクセスを制限するアクセス制御のモデルである。ユーザーは設定を変更できない。すなわち、$P$ は組織や管理者がラベルとポリシーから定め、ユーザーはそれを変更できない。
\end{dfn}

MAC は、軍事・政府機関など、情報の機密性が最優先される環境で採用される。

\begin{dfn}[RBAC（ロールベースアクセス制御）]\label{dfn:rbac}
\underline{RBAC}（Role-Based Access Control）とは、ユーザーに\underline{ロール}（役割）を割り当て、ロールに紐づく権限でアクセスを制御するモデルである。ロールの集合を $R$ とし、ユーザーへのロールの割り当てを $U\!A\subseteq S\times R$、ロールに紐づく権限を $P\!A\subseteq R\times (O\times A)$ とすれば、
\[
(s,o,a)\in P \iff \text{ある } r\in R \text{ が存在して } (s,r)\in U\!A \text{ かつ } (r,(o,a))\in P\!A
\]
である。
\end{dfn}

\begin{prop}[RBACによる管理の簡素化]\label{prop:ch2b-rbac-simplification}
RBAC では、ポリシーの作成、保守、監査が簡素化される。
\end{prop}

\begin{proof}
RBAC で管理すべきものが何かを見る。定義 \ref{dfn:rbac} により、権限はユーザーに直接ではなくロールに紐づけられ、ユーザーにはロールが割り当てられるだけである。したがって、権限についての管理はロールのみの管理となり、個々のユーザーごとに権限を設定・確認する必要がない。よって、ポリシーの作成、保守、監査はロール単位で行えばよく、簡素化される。
\end{proof}

RBAC は、企業の業務システムや病院、金融機関など、職務ごとに権限を整理したい場面でよく利用される。

\begin{dfn}[ABAC（属性ベースアクセス制御）]\label{dfn:abac}
\underline{ABAC}（Attribute-Based Access Control）とは、ユーザー、リソース、操作、環境などの複数の属性に基づき、柔軟かつ動的にアクセス制御を行うモデルである。すなわち、主体・対象・操作・環境の属性の組に対して許可か拒否かを返す判定によって、$(s,o,a)\in P$ かどうかがアクセスのたびに決まる。
\end{dfn}

ABAC が用いる属性のうちユーザーの属性は、第1章の属性情報（定義 \ref{dfn:attribute}）にあたる。ABAC は、ゼロトラストネットワークやクラウド環境など、高い柔軟性ときめ細かい制御が求められる場面に適している。

\begin{dfn}[PBAC（ポリシーベースアクセス制御）]\label{dfn:pbac}
\underline{PBAC}（Policy-Based Access Control）とは、属性や条件に基づいたルール（\underline{ポリシー}）を用いてアクセス制御を行うモデルである。
\end{dfn}

\begin{rem}[PBACとABACの関係]\label{rem:ch2b-pbac-abac}
PBAC は、ABAC（定義 \ref{dfn:abac}）をポリシー中心に実装した形式ともいえる。PBAC は、XACML や Open Policy Agent（OPA）といったポリシーエンジンなどと連携し、クラウドサービスやマイクロサービスでの外部ポリシー管理に利用される。
\end{rem}

以上の五つの方式を、制御の基準、柔軟性、管理者制御の強さ、適用例、特徴の観点から比較したものが表 \ref{tab:ch2b-access-control-models} である。

\begin{table}[H]
\centering
\caption{アクセス制御方式の比較}
\label{tab:ch2b-access-control-models}
\small
\begin{tabularx}{\linewidth}{|l|>{\raggedright\arraybackslash}p{24mm}|>{\raggedright\arraybackslash}p{10mm}|>{\raggedright\arraybackslash}p{15mm}|>{\raggedright\arraybackslash}p{22mm}|>{\raggedright\arraybackslash}X|}
\hline
モデル名 & 制御の基準 & 柔軟性 & 管理者制御の強さ & 適用例 & 特徴 \\
\hline
DAC & オブジェクトの所有者 & 高い & 弱い & UNIX ファイルシステム & ユーザーが自分のリソースに対するアクセス権を自由に設定可能 \\
\hline
MAC & セキュリティラベルとポリシー & 低い & 非常に強い & 軍事・政府システム & 組織が一貫して強制的にアクセス制御を行う。ユーザーは設定を変更できない \\
\hline
RBAC & ユーザーのロール（役割） & 中程度 & 強い & 業務システム、病院、金融 & ロール単位で管理しやすく、セキュリティと運用のバランスが良い \\
\hline
ABAC & ユーザー・リソース・環境などの属性 & 非常に高い & 高い（実装が複雑） & クラウド、ゼロトラスト環境 & 動的かつ詳細な制御が可能だが、設計・運用が難しい \\
\hline
PBAC & 明示的に定義されたポリシー（通常は属性ベース） & 非常に高い & 高い & クラウドサービス、マイクロサービス & ABAC の実装スタイルの 1 つ。外部ポリシーエンジンで一元管理でき、拡張性とガバナンスが高い \\
\hline
\end{tabularx}
\end{table}

\begin{rem}[アクセス制御方式の選択]\label{rem:ch2b-choosing-model}
オンラインサービスにアクセス制御を導入する際は、どのような情報や機能を保護したいのか、誰にどのような権限を与える必要があるのかといった要件を明確にし、将来的な拡張性なども考慮した上で、これらのアクセス制御モデルの中から最適な方式を選択、あるいは組み合わせて利用することが重要である。
\end{rem}

\subsection{セッション管理}\label{subsec:ch2b-session-management}

当人認証とアクセス制御によって利用者の正当性と許可範囲が定まったとしても、それを操作のたびに確かめ直すのでは使いものにならない。この項では、第1章で概要を述べたセッション管理（定義 \ref{dfn:session-management}）を、ライフサイクル（開始・保持・終了）とログインセッションの概念を使って精密化する。まず現実世界のホテルのカードキーを例にセッション管理の考え方とリスクを確認し、次にデジタル世界のログインセッション、それが保持する情報、セッションハイジャックのリスクとその対策（トークンの種類とハンドリング、ユーザーに提供すべき機能）を論じる。

\subsubsection{現実世界のセッション管理}

セッション管理がなぜ必要なのかを、まず確認の繰り返しという問題から説明する。この項では、セッションのライフサイクルを定義し、ホテルのカードキーの例でそれを具体化したうえで、セッション中に生じるリスクを論じる。

\begin{prop}[確認を操作のたびに繰り返すことの問題]\label{prop:ch2b-repeated-checks}
現実世界とデジタル世界のサービス、どちらにおいても、まず「当人認証」によって利用者の正当性を確認し、「アクセス制御」によって許可された範囲の操作のみを許可することが基本である。しかし、サービスをある程度の時間使い続ける状況で、操作のたびにこれらの確認を繰り返すのは非常に煩雑で、利便性を損なう。
\end{prop}

\begin{proof}
確認の回数を数える。サービスをある程度の時間使い続ける状況では、利用者は多数の操作を行う。操作のたびに当人認証（定義 \ref{dfn:authentication}）とアクセス制御（定義 \ref{dfn:access-control}）の確認を繰り返すならば、利用者は操作の回数だけ当人認証の手続きを行うことになる。よって、利用者にとって非常に煩雑であり、利便性を損なう。
\end{proof}

そこで登場するのが「セッション管理」の考え方である。第1章の定義 \ref{dfn:session-management} を、ここでは次のように言い表しておく。

\begin{dfn}[セッション管理（認証済み状態の維持）]\label{dfn:ch2b-session-management-refined}
\underline{セッション管理}とは、一度、当人認証によって本人であることが確認され、アクセスが許可された状態を、一定期間維持する仕組みである。この維持される状態（およびその期間）を\underline{セッション}（session）という。
\end{dfn}

\begin{prop}[セッション管理による利便性の回復]\label{prop:ch2b-session-convenience}
セッション管理により、利用者は何度も当人認証の手続きを繰り返すことなく、スムーズにサービスを使い続けられる。
\end{prop}

\begin{proof}
定義 \ref{dfn:ch2b-session-management-refined} により、セッション管理のもとでは、一度当人認証を経てアクセスが許可された状態が一定期間維持される。したがって、その期間内の操作では当人認証の手続きを改めて行う必要がない。よって、命題 \ref{prop:ch2b-repeated-checks} で述べた繰り返しの煩雑さが解消され、利用者はスムーズにサービスを使い続けられる。
\end{proof}

まずは、このセッション管理の考え方に近い仕組みを、現実世界の例で見てみよう。その前に、セッションがたどる段階に名前を付けておく。

\begin{dfn}[セッションのライフサイクル]\label{dfn:ch2b-session-lifecycle}
セッション管理には、一般的に次の三つの段階からなる\underline{ライフサイクル}がある。
\begin{enumerate}
\item \underline{セッション開始}：確認を行い、アクセスが許可された状態を作る。
\item \underline{セッション保持}：その状態を維持し、その間は確認を繰り返さずにサービスを利用させる。
\item \underline{セッション終了}：その状態を無効にし、以後はサービスを利用させない。
\end{enumerate}
\end{dfn}

各段階の説明は、次の例（ホテルのカードキー）での具体的な内容を一般化したものである。ライフサイクルを状態の遷移として描くと、図 \ref{fig:ch2b-session-lifecycle} のようになる。

\begin{figure}[H]
\centering
\begin{tikzpicture}[
  st/.style={draw, rounded corners, align=center, minimum width=28mm, minimum height=11mm},
  >={Stealth[length=2.5mm]}
]
\node[st] (start) {セッション開始};
\node[st, right=22mm of start] (keep) {セッション保持};
\node[st, right=22mm of keep] (end) {セッション終了};
\draw[->] (start) -- (keep);
\draw[->] (keep) -- (end);
\draw[->] (keep.north east) .. controls +(10mm,12mm) and +(-10mm,12mm) .. node[above, font=\small]{繰り返し利用} (keep.north west);
\end{tikzpicture}
\caption{セッションのライフサイクル}
\label{fig:ch2b-session-lifecycle}
\end{figure}

\underline{ホテルのカードキー利用：}セッション管理の考え方に近い現実世界の仕組みとして、ホテルのカードキーを取り上げる。

\begin{exa}[ホテルのカードキー]\label{exa:ch2b-hotel-cardkey}
ホテルのカードキーの利用は、定義 \ref{dfn:ch2b-session-lifecycle} の三つの段階に次のように対応する。
\begin{enumerate}
\item \underline{セッション開始：}ホテルにチェックインする際、予約者本人かの確認（簡単な身元確認や当人認証）が行われ、部屋のカードキーが発行される。この時点で、あなたはその部屋へのアクセスが許可された状態になる。
\item \underline{セッション保持：}チェックインからチェックアウトまでの滞在期間中、そのカードキーを使って何度も部屋に出入りできる。毎回フロントで本人確認をする必要はない。
\item \underline{セッション終了：}チェックアウト時刻を迎えるか、手続きを行うと、カードキーは無効化され、部屋に入れなくなる。これでセッションは終了となる。
\end{enumerate}
この例では、最初に一度だけ確認（チェックイン手続き）を行えば、セッションが終了するまでの間は、利用者は「信頼された状態」として扱われ、繰り返しサービスを利用（部屋へ入室）できる。
\end{exa}

この例の各段階で利用者が誰とやり取りするか（原書の図「現実世界のセッション管理」）を図 \ref{fig:ch2b-hotel-session} に示す。

\begin{figure}[H]
\centering
\begin{tikzpicture}[
  box/.style={draw, rounded corners, align=center, minimum width=26mm, minimum height=10mm},
  lab/.style={align=center, font=\small},
  >={Stealth[length=2.5mm]}
]
\node[box] (u1) {利用者};
\node[box, right=58mm of u1] (f1) {ホテルのフロント};
\draw[<->] (u1) -- node[lab, above]{\textbf{セッション開始}\\身元確認などを行い、カードキーを発行} (f1);
\node[box, below=20mm of u1] (u2) {利用者};
\node[box, right=58mm of u2] (f2) {ホテルの玄関、居室};
\draw[<->] (u2) -- node[lab, above]{\textbf{セッション保持}\\カードキーを所持していれば\\ホテルの玄関、居室の解錠が可能} (f2);
\node[box, below=20mm of u2] (u3) {利用者};
\node[box, right=58mm of u3] (f3) {ホテルのフロント};
\draw[<->] (u3) -- node[lab, above]{\textbf{セッション終了}\\チェックアウトしてカードキーを返却} (f3);
\end{tikzpicture}
\caption{現実世界のセッション管理（ホテルのカードキー）}
\label{fig:ch2b-hotel-session}
\end{figure}

しかし、ここで注意点がある。

\begin{prop}[セッション中の紛失・盗難のリスク]\label{prop:ch2b-session-theft-risk}
セッション中にカードキーや乗車券を紛失したり、盗まれたりして第三者の手に渡ってしまった場合、セッションが有効な限り、その第三者があなたになり代わってサービスを利用できてしまう。したがって、物理的なカードキーや乗車券を安全に管理すること、そしてセッションの有効期間を適切に設定し、必要ならば適切に身元確認や当人認証を要求することが重要である。
\end{prop}

\begin{proof}
まず、例 \ref{exa:ch2b-hotel-cardkey} のセッション保持の段階では、カードキーを所持していれば、毎回の本人確認なしに入室できる。すなわち、セッション保持中に確かめられているのは「カードキーを持っていること」だけであり、持っているのが本人かどうかは確かめられていない。よって、カードキーや乗車券が第三者の手に渡れば、セッションが有効な限り、その第三者が本人になり代わってサービスを利用できる。
次に、このリスクを抑える方法を考える。第三者の手に渡ること自体を防ぐには、物理的なカードキーや乗車券を安全に管理すればよい。また、渡ってしまった場合の被害の期間を限るには、セッションの有効期間を適切に設定すればよい。さらに、必要な場面では改めて身元確認や当人認証を要求すれば、カードキーを持っているだけの第三者を排除できる。以上より、これらが重要であることが分かる。
\end{proof}

このようなセッション管理の考え方と、それに伴う利便性やセキュリティ上の考慮点は、デジタル世界においても同様に重要である。

\subsubsection{デジタル世界のセッション管理}

デジタル世界のセッション管理として代表的な例は、オンラインサービスにおけるログイン状態の保持である。この項では、HTTP がステートレスであることからログインセッションの管理が必要になることを示し、ログインセッションを実現する手段、ログインセッションが保持する情報を定義する。続いて、セッションハイジャックのリスクと、それを軽減するトークンの種類とハンドリングの手法、さらにサービスがユーザーに提供すべき機能を論じる。

\begin{dfn}[ステートレスなプロトコル]\label{dfn:ch2b-stateless}
通信の各リクエストが独立しており、それに応じたデータを返すサーバー側が前回のリクエストとの関連性を保持しないプロトコルを、\underline{ステートレス}（stateless，状態を保持しない）なプロトコルという。
\end{dfn}

Web 通信の基本である HTTP はステートレスなプロトコルである。つまり、ブラウザやモバイルアプリから送られる各リクエストは独立しており、それに応じたデータを返すサーバー側は、前回のリクエストとの関連性を保持しない。

\begin{dfn}[ログインセッション]\label{dfn:login-session}
オンラインサービスにおいて、ログイン状態を表現し、維持するためのセッション（定義 \ref{dfn:ch2b-session-management-refined}）を\underline{ログインセッション}（login session）という。
\end{dfn}

\begin{prop}[ログインセッション管理の必要性]\label{prop:ch2b-need-login-session}
会員制サービスなどのアプリケーションでは、ログインセッションの管理が必要である。
\end{prop}

\begin{proof}
まず、HTTP はステートレスなので、サーバーは前回のリクエストとの関連性を保持せず、あるリクエストを送ったのがすでにログインした利用者かどうかを、HTTP の仕組みだけからは知ることができない。次に、会員制サービスなどのアプリケーションでは、ログインした利用者にだけ会員向けの機能を提供するため、ログイン状態を表現し、維持しなければならない。よって、HTTP の上でログイン状態を表現・維持する仕組み、すなわちログインセッション（定義 \ref{dfn:login-session}）の管理が必要である（原書の図「オンラインサービスのセッション管理」）。
\end{proof}

このログインセッションを実現するために、オンラインサービスでは HTTP Cookie を利用したセッション Cookie、ブラウザのローカルストレージ、あるいはモバイルアプリのデータストレージを利用したセッショントークンなどが用いられている。

\begin{dfn}[トークン]\label{dfn:ch2b-token}
ID 管理における\underline{トークン}（token）とは、様々なコンテキストによって「許可証」、「引換券」、「証明書」の役割を果たすデジタルデータのことである。
\end{dfn}

\begin{dfn}[セッションCookieとセッショントークン]\label{dfn:ch2b-session-cookie-token}
ログインセッションを実現するために用いられる次のものを考える。
\begin{enumerate}
\item HTTP Cookie を利用したものを\underline{セッション Cookie} という。
\item ブラウザのローカルストレージ、あるいはモバイルアプリのデータストレージを利用したものを\underline{セッショントークン}という。
\end{enumerate}
\end{dfn}

セッション Cookie やセッショントークンの実装方法はいくつかあるが、ログインセッションでは一般的に次のような情報を保持している点が共通の特徴である。

\begin{dfn}[ログインセッションが保持する情報]\label{dfn:ch2b-session-contents}
ログインセッションは、一般的に次の情報を保持する。
\begin{enumerate}
\item \underline{ユーザー ID}。
\item \underline{セッションの開始日時}：いつから。
\item \underline{セッションの有効期限}：いつまで（セッションタイムアウトの設定も重要）。
\item \underline{当人認証のコンテキスト}：認証日時、認証方式、認証強度など（例：どのような方法・レベルで本人確認が行われたか）。
\item \underline{アクセス制限のコンテキスト}：ロールや属性情報など、ポリシー判定に利用する情報（例：ユーザーにどの操作や情報へのアクセス権限があるか）。
\end{enumerate}
記号的には、ログインセッションを五つ組
\[
s=\bigl(\mathit{uid},\ t_{\mathrm{start}},\ t_{\mathrm{exp}},\ c_{\mathrm{authn}},\ c_{\mathrm{authz}}\bigr)
\]
として表すことができ、時刻 $t$ にそのセッションが有効であるのは、少なくとも $t_{\mathrm{start}}\le t< t_{\mathrm{exp}}$ のときに限る。
\end{dfn}

項目 (4) の「本人確認」は、当人認証の意味である（注意 \ref{rem:honnin-kakunin}）。項目 (5) のロールや属性は、アクセス制御の方式 RBAC（定義 \ref{dfn:rbac}）や ABAC（定義 \ref{dfn:abac}）でポリシー判定に用いられる情報にあたる。

オンラインサービスにおけるセッション管理には、次のリスクがある。

\begin{dfn}[セッションハイジャック]\label{dfn:ch2b-session-hijacking}
セッション情報が第三者の手に渡り悪用されることを、\underline{セッションハイジャック}（session hijacking）などとよぶ。
\end{dfn}

\begin{exa}[インフォスティーラーによるセッション情報の窃取]\label{exa:ch2b-infostealer}
実際に、\underline{インフォスティーラー}と呼ばれる情報窃取型マルウェアにより、セッション Cookie やセッショントークンが盗まれる事案が発生している。
\end{exa}

このリスクを軽減するため、まずはセッション管理の対象となる情報（定義 \ref{dfn:ch2b-session-contents}）を適切に定義して実装する必要がある。その上で、次のようなトークンの種類とハンドリングの手法がある。

\begin{dfn}[ベアラートークン]\label{dfn:ch2b-bearer-token}
提示された HTTP Cookie やトークン文字列そのものを信用するトークンを、\underline{ベアラートークン}（Bearer Token）という。
\end{dfn}

\begin{dfn}[送信者制限つきトークン]\label{dfn:ch2b-sender-constrained-token}
提示された HTTP Cookie やトークン文字列と、その提示者（送信元）の組み合わせが正当であるかを検証するトークンを、\underline{送信者制限つきトークン}（Sender-Constrained Token）という。
\end{dfn}

二つの違いは、現実世界の乗り物の券で例えると分かりやすい。

\begin{exa}[電車の切符と国際線の搭乗チケット]\label{exa:ch2b-ticket-boarding-pass}
\begin{enumerate}
\item 現実世界の電車の切符は、通常、所有者との一致を厳密には検証しないため、第三者がそれを提示して電車に乗ることができる。よって、ベアラートークン（定義 \ref{dfn:ch2b-bearer-token}）の一種と考えることができる。
\item 一方、国際線の航空機の搭乗チケットではパスポートとの検証が必要であり、送信者制限つきトークン（定義 \ref{dfn:ch2b-sender-constrained-token}）の一種と考えることができる。
\end{enumerate}
\end{exa}

\begin{rem}[ベアラートークンとセッションハイジャック]\label{rem:ch2b-bearer-hijack}
命題 \ref{prop:ch2b-session-theft-risk} で、カードキーや乗車券が第三者の手に渡ると、セッションが有効な限り第三者がなり代われることを見た。例 \ref{exa:ch2b-ticket-boarding-pass} のとおり、ベアラートークンはまさに「提示されたものそのもの」を信用するので、盗まれたセッション Cookie やトークン（例 \ref{exa:ch2b-infostealer}）をそのまま提示されれば区別できない。送信者制限つきトークンは提示者との組み合わせを検証するので、この点で対策となる。
\end{rem}

オンラインサービスで送信者制限つきトークンを実現するために、ハードウェアとセッション Cookie を紐づける仕組みの標準化が進められている。

\begin{exa}[Device Bound Session Credentials（DBSC）]\label{exa:ch2b-dbsc}
W3C で仕様策定が進められている \underline{Device Bound Session Credentials}（DBSC）では、公開鍵暗号の鍵管理を証明するためのプロトコルとブラウザインフラストラクチャを導入する。
\end{exa}

\begin{exa}[普段と異なるアクセスパターンの検知]\label{exa:ch2b-anomalous-access}
新しい技術を利用せずに今できることとして、アクセス元の情報（例：IP アドレス、デバイス情報など）を保持しておいて、普段と異なるアクセスパターンが検知された場合に、不正利用を疑って追加の当人認証を要求するといった対応が有効となる。例えば、通常日本からアクセスしているユーザーが、短時間のうちに海外の IP アドレスからアクセスしてきた場合などである（追加の当人認証については、第3章の再認証も参照。定義 \ref{dfn:reauthentication}）。
\end{exa}

また、サービスがユーザーに対して次のような機能を提供することも重要である。

\begin{prop}[ユーザーに提供すべきセッション管理機能]\label{prop:ch2b-user-session-features}
サービスは、ユーザーに対して次の機能を提供することが重要である。
\begin{enumerate}
\item 現状の有効なセッション情報を確認する機能。
\item 不要になったセッション情報を無効化する機能。
\end{enumerate}
\end{prop}

\begin{proof}
各機能が何を可能にするかを示す。まず、(1) の機能があれば、ユーザーは自分のアカウントに現在どのようなセッションが有効になっているかを見られるので、不正利用されていないかの確認をユーザー自身が行えるようになる。次に、(2) の機能があれば、ユーザーが能動的に既存セッションを終了させられるので、不要なセッションや見覚えのないセッションが悪用されるリスクを軽減できる。よって、両機能の提供により、サービス側のリスク判定だけでなくユーザー自身も安全性の確保に参加でき、これは重要である。
\end{proof}

\begin{exa}[有効なセッション一覧の確認と無効化]\label{exa:ch2b-session-list}
Google や GitHub といったサービスでは、現在有効なセッション一覧の確認や無効化が可能である。
\end{exa}

\begin{rem}[ユーザーと協力した安全性の確保]\label{rem:ch2b-user-cooperation}
サービス側のリスク判定（例 \ref{exa:ch2b-anomalous-access}）だけではなく、ユーザーと協力して安全性を保つ取り組み（命題 \ref{prop:ch2b-user-session-features}）を進める必要がある。
\end{rem}

%% =======================================================
%%  part: ch3
%% =======================================================
\section{単一サービスにおけるID管理機能}

この章では、\cite{Ito2026} の第3章に沿って、単一のサービス内で完結する ID 管理機能を扱う。
まず ID 管理の基本構成であるアイデンティティレジスターと、アイデンティティの状態と状態遷移を体系化したアイデンティティライフサイクルを定義する。
そのうえで、新規登録、ログイン、ログアウト、再認証、ユーザー情報の設定変更、セッション管理、本人確認（KYC）、アカウントリカバリー、無効化と復旧、退会と復旧、セキュリティイベントログの各機能を順に定義し、それぞれの主なプロセスを整理する。
各機能がアイデンティティライフサイクルのどの遷移に対応するかも、可能な限り明示する。
最後にコラムとして、不正ログインが疑われるときの対処の流れを扱う。

\subsection{ID管理の基本構成とライフサイクル}

この項では、以降のすべての機能の土台となる二つの概念を導入する。
一つはアイデンティティ情報を保持するデータストアであるアイデンティティレジスター、もう一つはアイデンティティがとりうる状態とその間の遷移を整理したアイデンティティライフサイクルである。
後の項で扱う各機能は、アイデンティティレジスターの内容を読み書きし、アイデンティティの状態を遷移させる操作として理解できる。

\subsubsection{単一サービス内で完結するID管理}

ID 管理（定義 \ref{dfn:identity-management}）をどの範囲で実現するかによって、構成は大きく二つに分かれる。
本章では、その一方である単一サービス内で完結する場合に焦点を当て、その実現方法を述べる。

\begin{dfn}[単一サービス内で完結するID管理]\label{dfn:ch3-single-service-idm}
大きく分けて、ID 管理には次の 2 種類がある。
\begin{enumerate}
\item 単一のサービス内で完結する場合。
\item 複数のサービスで連携する場合。
\end{enumerate}
(1) を \underline{単一サービス内で完結する ID 管理}という。本章では (1) を扱い、(2) は次の第4章で扱う。
\end{dfn}

歴史的には、(1) の構成が出発点であった。

\begin{rem}[単一サービス構成は最も基本的な構成である]\label{rem:ch3-single-service-basic}
インターネットが普及し、会員機能を提供するサービスが増え始めた当初、サービスの利用開始から終了までの全機能は、単一のサービス内で完結していた。
ID 管理の観点でも、すべての機能が一つのサービス内に閉じており、これは最もシンプルで基本的な構成といえる。
\end{rem}

\subsubsection{ユーザーの情報を管理するアイデンティティレジスター}

ID 管理機能を実現するには、アイデンティティ情報をどこかに保持しなければならない。
その保持場所を「アイデンティティレジスター」と呼ぶ。
ここでは、アイデンティティレジスターが主に何を管理するかを定義し、ID 管理の他の構成要素のデータとの関係を述べる。

\begin{dfn}[アイデンティティレジスター]\label{dfn:identity-register}
ID 管理機能を実現するためにアイデンティティ情報を管理するデータストアを、\underline{アイデンティティレジスター}（Identity Register）という。
アイデンティティレジスターは、主に次の情報を管理する。
\begin{enumerate}
\item \underline{識別子}（Identifier）。
\item 当人認証（定義 \ref{dfn:authentication}）のための\underline{クレデンシャル}（Credentials）。
\item 属性情報（Attributes。定義 \ref{dfn:attribute}）。
\end{enumerate}
\end{dfn}

アイデンティティレジスターが管理するのは上の三つが中心であるが、ID 管理にはこれ以外の構成要素もあり、それらもユーザーに関連づけられたデータを必要とする。

\begin{rem}[他の構成要素のデータの置き場所]\label{rem:ch3-register-other-data}
アクセス制御（定義 \ref{dfn:access-control}）、セッション管理（定義 \ref{dfn:session-management}）、ID 連携（定義 \ref{dfn:identity-federation}）といった ID 管理の他の構成要素についても、ユーザーに関連づけられたデータを管理する必要がある。
これらのデータは、アイデンティティレジスターまたは別の専用データストアで管理される。
\end{rem}

本書の図では、アイデンティティレジスターが識別子・クレデンシャル・属性情報を保持し、アクセス制御・セッション管理・ID 連携のデータがアイデンティティレジスターまたは別のデータストアに置かれる様子が描かれている。
記号で書けば、ユーザー $u$ についてアイデンティティレジスターが保持するレコードは、おおよそ
\[
\bigl(\mathrm{id}(u),\ \mathrm{cred}(u),\ A(u)\bigr)
\]
という組である。ここで $\mathrm{id}(u)$ は識別子、$\mathrm{cred}(u)$ は当人認証のためのクレデンシャル、$A(u)$ は属性情報の集合である。

\subsubsection{ユーザーの状態と状態遷移を表すアイデンティティライフサイクル}

アイデンティティレジスターに記録されたユーザーは、登録されてから削除されるまでの間に、利用できる状態、一時停止された状態など、いくつかの状態を経る。
ID 管理におけるアイデンティティの状態と、その状態遷移に求められるアクションを体系的に整理した概念は、国際規格
「ISO/IEC 24760-1:2025 Information security, cybersecurity and privacy protection --- A framework for identity management」
において「アイデンティティライフサイクル」（Identity Lifecycle）として定義されている。
本書は同規格におけるアイデンティティライフサイクルを図で示しており、ここではその内容を、状態の集合と遷移からなる状態遷移系として定義し直す。
まず状態を、次に遷移のためのアクションを定義する。

\begin{dfn}[アイデンティティの状態]\label{dfn:ch3-identity-states}
アイデンティティの状態として、次の 5 つが定義されている。
\begin{enumerate}
\item \underline{未登録}（Unknown）：ユーザー情報が登録されていない、または退会によりデータが完全に消去された状態。
\item \underline{仮登録}（Established）：データは存在するものの、本登録が完了しておらず、サービスを利用できない状態。
\item \underline{本登録}（Active）：本登録が完了し、サービスを完全に利用できる状態。
\item \underline{無効化状態}（Suspended）：アカウントが一時的に停止され、サービスを利用できない状態。
\item \underline{論理削除状態}（Archived）：データは保持されているもののサービスは利用できない、退会済みの状態。
\end{enumerate}
これらの状態の集合を
\[
\mathcal{S}:=\{\mathsf{Unknown},\ \mathsf{Established},\ \mathsf{Active},\ \mathsf{Suspended},\ \mathsf{Archived}\}
\]
とおく。
\end{dfn}

5 つの状態のうち、サービスを利用できるのは本登録だけである。
また、データが存在しないのは未登録だけであり、仮登録・本登録・無効化状態・論理削除状態ではいずれもデータが存在する。
この二つの観点は、以下の遷移の意味を理解するうえで役に立つ。

\begin{dfn}[アイデンティティライフサイクルのアクション]\label{dfn:ch3-lifecycle-actions}
状態遷移のためのアクションとして、次のものが定義されている。
\begin{enumerate}
\item \underline{登録}（Enrolment）：サービスが身元確認（定義 \ref{dfn:identity-proofing}）を行い、その結果をアイデンティティレジスターへ登録するプロセス。
\item \underline{有効化}（Activation）：特定のユーザーがサービスを利用できるよう、アイデンティティ情報をアクティブな状態にするプロセス。
\item \underline{保守}（Maintenance）：アイデンティティレジスターに登録されているユーザー情報を更新するプロセス。
\item \underline{アイデンティティ情報の調整}（Identity adjustment）：新しい情報に基づき、有効化に関する内容が変更され、アイデンティティ情報を更新するプロセス。
\item \underline{一時停止}（Suspension）：ユーザーのアカウントを一時的に利用不可な状態としてマークするプロセス。
\item \underline{再有効化}（Reactivation）：一時的に利用不可な状態を再び利用可能な状態に戻すプロセス。
\item \underline{削除}（Delete）：登録されたアイデンティティ情報を完全に削除するプロセス。
\item \underline{保存}（Archive）：ユーザーのアカウントをサービス利用不可の状態にし、統計処理などの目的で情報を保管するプロセス。
\item \underline{再登録}（Enrolment/restore）：保存されていた情報を利用して、アカウントを復元・再登録するプロセス。
\end{enumerate}
これらのアクションの集合を $\Sigma$ と書く。
\end{dfn}

状態とアクションがそろったので、アイデンティティライフサイクルを状態遷移系として定義する。
ただし、本書の地の文は各アクションの始点と終点を必ずしも明記していない。
そこで以下では、各遷移の始点と終点を本章の記述（3-1 から 3-11 の各項）から読み取り、その根拠と確度を表 \ref{tab:ch3-lifecycle-transitions} にまとめる。

\begin{dfn}[アイデンティティライフサイクル]\label{dfn:identity-lifecycle}
定義 \ref{dfn:ch3-identity-states} の状態集合 $\mathcal{S}$ と定義 \ref{dfn:ch3-lifecycle-actions} のアクション集合 $\Sigma$、および遷移関係
\[
\delta\subseteq \mathcal{S}\times\Sigma\times\mathcal{S}
\]
の組 $(\mathcal{S},\Sigma,\delta)$ を\underline{アイデンティティライフサイクル}（Identity Lifecycle）という。
$(s,a,t)\in\delta$ のとき $s\xrightarrow{a}t$ と書き、「アクション $a$ により状態 $s$ から状態 $t$ へ遷移する」と読む。
本章で用いる $\delta$ は表 \ref{tab:ch3-lifecycle-transitions} で与える。
\end{dfn}

\begin{rem}[形式化について]\label{rem:ch3-lifecycle-formalization}
状態と状態遷移のためのアクションは本書（およびそこで引用される ISO/IEC 24760-1:2025）が定義するものであり、これを三つ組 $(\mathcal{S},\Sigma,\delta)$ として書くのはこのノートでの整理である。
遷移関係 $\delta$ のうち、本書の地の文に始点と終点がはっきり書かれているものは表で「明示」とし、他の記述から読み取ったものは「読み取り」、本書の記述だけでは決まらないものは「曖昧」と注記する。
\end{rem}

\begin{longtable}{p{0.17\textwidth}p{0.12\textwidth}p{0.12\textwidth}p{0.37\textwidth}p{0.08\textwidth}}
\caption{アイデンティティライフサイクルの遷移（始点・終点と本書における根拠）}\label{tab:ch3-lifecycle-transitions}\\
\hline
アクション & 始点 & 終点 & 本書における根拠 & 確度\\
\hline
\endfirsthead
\hline
アクション & 始点 & 終点 & 本書における根拠 & 確度\\
\hline
\endhead
\hline
\endfoot
登録 (Enrolment) & 未登録 & 仮登録 &
3-2 で、新規登録はユーザーを「未登録」から「本登録」へ変え、その途中に「登録」フェーズと「有効化」フェーズを経るとされる。管理者による承認が必要な場合は、まず「登録」が行われ、承認を経て「有効化」される。仮登録の定義は「データは存在するが本登録が完了していない状態」である。以上から、登録の終点はデータが存在しかつ未利用の状態、すなわち仮登録と読める。終点を「仮登録」と明記した文はない。 & 読み取り\\
\hline
有効化 (Activation) & 仮登録 & 本登録 &
3-3 に「もし『仮登録』状態であれば『有効化』プロセスへ……誘導する」とある。3-2 で、有効化を経てサービスを利用できるようになる（本登録になる）とされる。 & 明示\\
\hline
保守 (Maintenance) & 本登録 & 本登録 &
3-6 で、サービスの利用を続ける中でユーザー自身が情報を更新する設定変更機能が「保守」にあたるとされる。保守は情報の更新であり、状態を変える旨の記述はない。本登録以外の状態で保守が行われるかについての記述はない。 & 読み取り\\
\hline
アイデンティティ情報の調整 (Identity adjustment) & 不明 & 不明 &
定義（「新しい情報に基づき、有効化に関する内容が変更され、アイデンティティ情報を更新する」）以外に記述がなく、始点・終点は本書の地の文からは決まらない。図 \ref{fig:ch3-lifecycle} では、情報の更新であることから仮に本登録上の自己遷移として破線で描く。 & 曖昧\\
\hline
一時停止 (Suspension) & 本登録 & 無効化状態 &
3-10 で、利用可能な「本登録」状態から利用が制限された「無効化」状態へ遷移するプロセスである「一時停止」とされる。 & 明示\\
\hline
再有効化 (Reactivation) & 無効化状態 & 本登録 &
3-10 で、「無効化」状態から再び「本登録」状態へ復帰させるための「再有効化」とされる。3-3 でも「無効化」状態であれば「再有効化」プロセスへ誘導するとされる。 & 明示\\
\hline
削除 (Delete) & 仮登録・本登録・無効化状態・論理削除状態 & 未登録 &
3-11 で、物理削除は「アカウントの状態を問わず、データが存在しない『未登録』状態への遷移である『削除』」とされる。未登録の定義にも「退会によりデータが完全に消去された状態」が含まれる。始点は「状態を問わず」なので、データが存在する 4 状態すべてとした。 & 明示\\
\hline
保存 (Archive) & 本登録 & 論理削除状態 &
3-11 で、論理削除は「『本登録』から『論理削除』状態への遷移である『保存』」とされる。本登録以外（例えば無効化状態）からの保存についての記述はない。 & 明示\\
\hline
再登録 (Enrolment/restore) & 論理削除状態 & 仮登録（推定） &
定義は「保存されていた情報を利用して、アカウントを復元・再登録する」であり、保存されていた情報が存在するのは論理削除状態なので、始点は論理削除状態と読める。3-11 で、論理削除からの復旧は「『再登録』と『有効化』の組み合わせ」に相当するとされる。有効化を仮登録から本登録への遷移と読むと再登録の終点は仮登録となるが、終点を明記した文はない。 & 始点は読み取り、終点は曖昧\\
\end{longtable}

表 \ref{tab:ch3-lifecycle-transitions} の遷移を図示したものが図 \ref{fig:ch3-lifecycle} である。
破線は、終点（または始点・終点）が本書の地の文からは確定しない遷移を表す。

\begin{figure}[H]
\centering
\begin{tikzpicture}[
  st/.style={draw, rounded corners, align=center, minimum width=2.3cm, minimum height=1.0cm, font=\small},
  tr/.style={-{Stealth[length=2.5mm]}, thick},
  amb/.style={-{Stealth[length=2.5mm]}, thick, dashed},
  lb/.style={font=\scriptsize, fill=white, inner sep=1pt}
]
\node[st] (U) at (0,0) {未登録\\(Unknown)};
\node[st] (E) at (4.2,0) {仮登録\\(Established)};
\node[st] (A) at (8.4,0) {本登録\\(Active)};
\node[st] (S) at (8.4,3.4) {無効化状態\\(Suspended)};
\node[st] (R) at (8.4,-3.4) {論理削除状態\\(Archived)};

\draw[tr] (U) to[bend left=18] node[lb, above] {登録} (E);
\draw[tr] (E) to[bend left=18] node[lb, below] {削除} (U);
\draw[tr] (E) -- node[lb, above] {有効化} (A);
\draw[tr] ([yshift=3mm]A.east) .. controls +(1.4,0.6) and +(1.4,-0.6) .. node[lb, right] {保守} ([yshift=3mm]A.east);
\draw[amb] ([yshift=-3mm]A.east) .. controls +(1.6,-0.2) and +(1.0,-1.2) .. node[lb, right] {調整} ([xshift=6mm]A.south);
\draw[tr] (A) to[bend left=22] node[lb, left] {一時停止} (S);
\draw[tr] (S) to[bend left=22] node[lb, right] {再有効化} (A);
\draw[tr] (A) to[bend left=22] node[lb, right] {保存} (R);
\draw[amb] (R) to[bend left=10] node[lb, below left] {再登録} (E);
\draw[tr] (A.north west) .. controls +(-1.5,1.6) and +(1.5,1.6) .. node[lb, above] {削除} (U.north east);
\draw[tr] (S.west) .. controls +(-5,0) and +(0,2.5) .. node[lb, above] {削除} (U.north);
\draw[tr] (R.west) .. controls +(-5,0) and +(0,-2.5) .. node[lb, below] {削除} (U.south);
\end{tikzpicture}
\caption{アイデンティティライフサイクルの状態遷移図（表 \ref{tab:ch3-lifecycle-transitions} に基づく。破線は本書の記述からは始点・終点が確定しない遷移）}\label{fig:ch3-lifecycle}
\end{figure}

この状態遷移系は規格上の整理であり、実際のアプリケーションがすべてを実装しているわけではない。

\begin{rem}[状態と状態遷移の簡略化]\label{rem:ch3-lifecycle-simplification}
すべてのアプリケーションで、これらの状態と状態遷移が厳密に実装されているわけではない。
エンタープライズ向けサービスの ID 管理では社員の状況に応じて細かい状態管理が行われるのに対して、コンシューマー向けサービスでは一部の状態と状態遷移を簡略化、省略する場合もある。
例えば次のような場合がある。
\begin{enumerate}
\item 仮登録状態を設けずに直接本登録とする。
\item データを保存せず即時削除する。
\item 一時停止・再有効化機能を持たない。
\end{enumerate}
(1) は $\mathsf{Established}$ を経由しない、(2) は保存（$\mathsf{Active}\to\mathsf{Archived}$）を用いず削除だけを用いる、(3) は $\mathsf{Suspended}$ を用いない、という簡略化にあたる。
\end{rem}

\begin{rem}[ライフサイクルを意識することの意義]\label{rem:ch3-lifecycle-importance}
注意 \ref{rem:ch3-lifecycle-simplification} のように実装が簡略化される場合があるとしても、これらの状態と状態遷移を意識することは、ID 管理の中核であるアイデンティティレジスターのデータ管理をより深く理解するうえで重要である。
\end{rem}

実際、定義 \ref{dfn:ch3-lifecycle-actions} のアクションのうち、登録はレジスターへの登録、保守と調整はレジスター上の情報の更新、保存はデータを保持したままの利用停止、削除はデータの完全な削除であり、いずれもアイデンティティレジスター（定義 \ref{dfn:identity-register}）のデータに対する操作である。
また定義 \ref{dfn:ch3-identity-states} の各状態は、データが存在するか、存在するならサービスを利用できるかによって区別されている。
このため、状態遷移を意識することが、レジスターのデータ管理を体系的に捉えることにつながる。

以下の各項では、各機能がこのライフサイクルのどの遷移に対応するかを、定義 \ref{dfn:identity-lifecycle} を参照しながら述べる。

\subsection{新規登録}

この項では、ユーザーがサービスを利用できるようになるための最初の機能である新規登録を扱う。
新規登録は、アイデンティティライフサイクル（定義 \ref{dfn:identity-lifecycle}）において未登録の状態から本登録の状態へ移る過程であり、ライフサイクル上で非常に重要な役割を担う。
まず新規登録を定義してその実装パターンを整理し、次に新規登録の主なプロセスを順に述べる。

\subsubsection{サービスを利用する最初のステップ}

新規登録は、ユーザーがサービスを利用可能になるための最初のステップである。
この機能により、ユーザーの状態は「未登録」から「本登録」へと変わる。

\begin{dfn}[新規登録]\label{dfn:signup}
ユーザーの状態をアイデンティティライフサイクルの未登録から本登録へ変え、ユーザーがサービスを利用できるようにする機能を\underline{新規登録}という。
新規登録は、システムにユーザー情報を記録する「登録」フェーズと、サービス利用を正式に許可する「有効化」フェーズを経る。
すなわち、表 \ref{tab:ch3-lifecycle-transitions} の遷移で書けば
\[
\mathsf{Unknown}\xrightarrow{\text{登録}}\mathsf{Established}\xrightarrow{\text{有効化}}\mathsf{Active}
\]
に対応する（中間の状態 $\mathsf{Established}$ については表 \ref{tab:ch3-lifecycle-transitions} の「登録」の行を参照。注意 \ref{rem:ch3-lifecycle-simplification} のとおり、仮登録を設けずに直接本登録とする実装もある）。
\end{dfn}

本書の図では、アイデンティティライフサイクルの図の上で、新規登録が未登録から登録・有効化を経て本登録に至る部分として示されている。
この機能の実装方法は、サービスの特性によっていくつかのパターンに分けられる。

\begin{dfn}[新規登録機能の実装パターン]\label{dfn:ch3-signup-patterns}
新規登録機能の実装パターンとして、次の 3 つがある。
\begin{enumerate}
\item \underline{ユーザー自身による登録}：ユーザーが自ら情報を入力して登録する。
\item \underline{管理者による事前登録}：人事担当者などの管理者が事前にユーザー情報を登録し、利用者に通知する。
\item \underline{管理者による招待とユーザー自身による登録}：管理者がユーザーのメールアドレス宛に招待リンクを送信し、ユーザーがそのリンクから登録を完了させる。
\end{enumerate}
\end{dfn}

\begin{exa}[実装パターンが採用される場面]\label{exa:ch3-signup-patterns}
定義 \ref{dfn:ch3-signup-patterns} の各パターンは、次のような場面で採用される。
\begin{enumerate}
\item ユーザー自身による登録は、多くのコンシューマー向けサービスで採用されている。
\item 管理者による事前登録は、企業や教育機関で利用されるシステムに多い。
\item 管理者による招待とユーザー自身による登録は、企業向けサービスなどで利用される。
\end{enumerate}
\end{exa}

\subsubsection{新規登録の主なプロセス}

新規登録は一つの操作に見えるが、内部ではいくつかのプロセスが順に行われる。
ここでは、事前条件の確認、身元確認、クレデンシャルの設定、登録と有効化、ログインセッション作成の 5 つに分けて述べる。

\begin{nota}[新規登録の主なプロセス]\label{nota:ch3-signup-process}
新規登録の主なプロセスは、次の順に行われる。
\begin{enumerate}
\item 事前条件の確認。
\item 身元確認。
\item クレデンシャルの設定。
\item 登録、有効化。
\item ログインセッション作成。
\end{enumerate}
以下、この項ではこれらを「プロセス (1)」から「プロセス (5)」と呼ぶ。
\end{nota}

\underline{(1) 事前条件の確認：}
サービスによっては、新規登録に際してユーザーが満たすべき条件（年齢制限など）が設けられている場合がある。

\begin{dfn}[事前条件の確認]\label{dfn:ch3-precondition-check}
新規登録に際して、ユーザーがサービスの定める条件を満たしているかを確かめるプロセスを\underline{事前条件の確認}という。
事前条件には、年齢制限などのユーザーが満たすべき条件のほか、次のものがこの段階で行われる。
\begin{enumerate}
\item 利用規約への同意取得。
\item 招待制サービスにおける有効な招待リンクの検証。
\end{enumerate}
\end{dfn}

\underline{(2) 身元確認：}
新規登録を試みるユーザーが誰であるかを、サービスはまだ知らない。
そこで、ユーザーが提供した属性情報の正しさを検証する必要がある。
これは身元確認（定義 \ref{dfn:identity-proofing}）にほかならない。

\begin{dfn}[新規登録における身元確認]\label{dfn:ch3-signup-proofing}
新規登録を試みるユーザーは、氏名、メールアドレス、電話番号などの属性情報をアプリケーションに提供する。
サービス側が、送信された情報が
\begin{enumerate}
\item 正しく、
\item 実在するものであり、
\item 送信者のものであること
\end{enumerate}
を検証するプロセスを、\underline{新規登録における身元確認}という。これは身元確認（定義 \ref{dfn:identity-proofing}）を新規登録の場面で行うものである。
\end{dfn}

\begin{exa}[新規登録における身元確認の手法]\label{exa:ch3-signup-proofing-methods}
多くのコンシューマー向けサービスでは、入力されたメールアドレスや電話番号の所有者確認のため、ワンタイムパスワード（OTP）を送信する手法が一般的である（第2章のメール・SMS による OTP 認証、定義 \ref{dfn:otp-auth} と同様の仕組みを用いる）。
サービス内容によっては、この段階で運転免許証などの公的な身分証明書を用いた、より厳格な身元確認を行うこともある。
\end{exa}

\underline{(3) クレデンシャルの設定：}
新規登録の後、ユーザーはログインによってサービスを利用し続ける。
そのためには、ログインで使う認証の手段をこの時点で用意しておかなければならない。

\begin{dfn}[クレデンシャルの設定]\label{dfn:ch3-credential-setup}
新規登録の際に、その後のログインで使用するクレデンシャルを発行・登録するプロセスを\underline{クレデンシャルの設定}という。
\end{dfn}

\begin{exa}[クレデンシャルの設定の例]\label{exa:ch3-credential-setup}
クレデンシャルの設定に相当するものとして、次がある。
\begin{enumerate}
\item 最も代表的な例はパスワード（定義 \ref{dfn:password-auth}）の設定である。
\item 最近のトレンドでは、パスキー（定義 \ref{dfn:passkey}）の設定がこのプロセスに相当する。
\item セキュリティ強化のため多要素認証（MFA。定義 \ref{dfn:mfa}）を必須とするサービスでは、この段階でワンタイムパスワード生成アプリの設定や、そのバックアップコード（定義 \ref{dfn:backup-code}）の発行などが行われる。
\end{enumerate}
\end{exa}

\underline{(4) 登録、有効化：}
新規登録プロセスでは、アイデンティティライフサイクルの「登録」と「有効化」という 2 つの重要な処理が行われる（定義 \ref{dfn:signup}）。
これらの処理をどのタイミングで実行するかは、特に身元確認の方法に依存する。

\begin{prop}[登録と有効化のタイミング]\label{prop:ch3-enrol-activate-timing}
新規登録における「登録」と「有効化」の実行タイミングは、身元確認の方法によって、例えば次のように異なる。
\begin{enumerate}
\item \underline{メールアドレス確認が必須の場合：}ユーザーが入力した情報がシステムに登録され、かつメールアドレスの有効性が確認された時点で、「登録」と「有効化」が完了し、サービスが利用可能になる。
\item \underline{メールアドレス確認が任意の場合：}ユーザーが情報を入力した段階で「登録」と「有効化」が行われ、一部機能が利用可能になることがある。この場合、メールアドレス確認などの追加の身元確認を完了すると、利用可能なサービスの範囲が広がる、といった段階的な権限付与が行われることもある。
\item \underline{管理者による承認が必要な場合：}ユーザー自身による登録手続きが完了した時点でまず「登録」が行われる。その後、システム管理者（例：企業の人事担当者）による承認を経て「有効化」され、初めてサービスが利用可能になる。
\end{enumerate}
\end{prop}

\begin{proof}
有効化はサービスの利用を正式に許可するフェーズである（定義 \ref{dfn:signup}）から、どの時点で有効化するかは、サービスが利用を許可するのに必要な確認がどの時点でそろうかで決まる、という観点から各場合を確かめる。
まず (1) では、メールアドレスの有効性の確認が利用許可の条件であるから、情報の登録とメールアドレスの確認がそろった時点で登録と有効化が完了する。
次に (2) では、メールアドレス確認は利用許可の条件ではないから、入力の段階で登録と有効化を行える。ただし確認済みかどうかで許可する範囲を変えることはでき、そのとき追加の身元確認によって利用範囲が広がる段階的な権限付与となる。
最後に (3) では、管理者の承認が利用許可の条件であるから、ユーザーの手続き完了時点では登録だけが行われ、承認を経て有効化される。
よって、いずれの場合も登録と有効化のタイミングは身元確認（およびそれに付随する承認）の方法に依存する。
\end{proof}

\begin{rem}[ライフサイクル上の位置づけ]\label{rem:ch3-enrol-activate-lifecycle}
命題 \ref{prop:ch3-enrol-activate-timing} の (3) では、登録の後、承認を待つ間はデータが存在するがサービスを利用できない。これは仮登録（$\mathsf{Established}$）の定義に合う。
一方 (1) と (2) では登録と有効化が同時に完了するので、仮登録の状態は（存在するとしても）外からは見えない。
\end{rem}

プロセス (2) と (3) で得られたものは、このプロセス (4) で保存される。

\begin{rem}[身元確認の結果とクレデンシャルの登録時点]\label{rem:ch3-register-timing}
プロセス (2) の身元確認の結果と、プロセス (3) のクレデンシャルは、プロセス (4) の時点でアイデンティティレジスター（定義 \ref{dfn:identity-register}）に登録される。
\end{rem}

\underline{(5) ログインセッション作成：}
新規登録を終えたユーザーにすぐにもう一度ログインを求めるのは不便である。
そのため、新規登録の最後にログインセッション（定義 \ref{dfn:login-session}）を作成することが多い。

\begin{dfn}[新規登録におけるログインセッション作成]\label{dfn:ch3-signup-session}
ユーザー自身が新規登録を完了した場合、多くはそのまま自動的にログイン状態になる。これは、ログイン状態を維持するための「ログインセッション」（定義 \ref{dfn:login-session}）が作成されたことを意味する。
新規登録の最後にこのログインセッションを作成するプロセスを、\underline{新規登録におけるログインセッション作成}という。
\end{dfn}

\begin{rem}[ログインセッションの一般的な実現方法]\label{rem:ch3-session-implementation}
ログインセッションは一般的に、クライアントサイド（Web ブラウザやモバイルアプリ）で\underline{セッショントークン}を保持し、サーバーサイドでそのトークンに紐づくユーザー情報を管理することで実現される。
\end{rem}

\begin{prop}[新規登録で作成するセッションに含めるべき情報]\label{prop:ch3-signup-session-info}
新規登録で作成されるログインセッションには、デバイスやネットワーク情報に加え、「新規登録直後に発行されたセッションである」ことを示す情報を含めることが重要である。
\end{prop}

\begin{proof}
ログインセッションは新規登録のほか、次の項で述べるログインでも作成される。
両者は同じログインセッションという形をとるが、前者はクレデンシャルを設定した直後に作られ、後者は登録済みのクレデンシャルによる当人認証を経て作られる、という違いがある。
この違いを後から把握できるようにするには、セッション自体に発行の経緯を記録しておく必要がある。
よって、デバイスやネットワーク情報に加えて「新規登録直後に発行されたセッションである」ことを示す情報を含めることが重要である。
\end{proof}

本書の図では、新規登録で生成されるログインセッションが、デバイス・ネットワーク情報と「新規登録直後のセッションである」旨の情報を持つものとして示されている。

\subsection{ログイン}

この項では、新規登録を済ませたユーザーがサービスを利用するたびに行う機能であるログインを扱う。
ログインは当人認証（定義 \ref{dfn:authentication}）と同一視されがちであるが、実際にはユーザーの状態の確認とログインセッションの作成を含む一連の手続きである。
まずログインを定義し、次にその主なプロセスを述べる。

\subsubsection{本人であることを証明してサービスを利用する}

新規登録が済んでいても、サービスは、いまアクセスしてきた者が登録済みのユーザー本人であるかを知らない。
その確認と、サービスを利用する権限の付与を行うのがログインである。

\begin{dfn}[ログイン]\label{dfn:login}
新規登録（定義 \ref{dfn:signup}）を済ませたユーザーが、自身が登録済みの本人であることをシステムに証明し、サービスを利用する権限を得るための一連の手続きを\underline{ログイン}という。
このプロセスには、単なる当人認証（定義 \ref{dfn:authentication}）だけでなく、次の重要なプロセスが含まれる。
\begin{enumerate}
\item ユーザーがサービスを利用可能な「本登録」状態であることの確認。
\item その後のログインセッション（定義 \ref{dfn:login-session}）の作成。
\end{enumerate}
\end{dfn}

ログイン自体はアイデンティティライフサイクルの遷移を引き起こさないが、ユーザーの状態が本登録（$\mathsf{Active}$）であることを前提とし、そうでない場合には有効化や再有効化への誘導を伴う（後述のプロセス (2)）。

\subsubsection{ログインの主なプロセス}

ログインの主なプロセスは、当人認証、ユーザーがサービスを利用できる状態であることの確認、ログインセッションの作成の 3 つである。

\begin{nota}[ログインの主なプロセス]\label{nota:ch3-login-process}
ログインの主なプロセスは、次の順に行われる。
\begin{enumerate}
\item 当人認証。
\item ユーザーがサービスを利用できる状態であることの確認。
\item ログインセッションの作成。
\end{enumerate}
以下、この項ではこれらを「プロセス (1)」から「プロセス (3)」と呼ぶ。
\end{nota}

\underline{(1) 当人認証：}
ログイン時にはまず、ユーザーが登録済み本人であることを確認する当人認証（定義 \ref{dfn:authentication}）を行う。
具体的な認証方式はサービスによってさまざまであるが（第2章参照）、実装パターンは主に次の 2 つに分類できる。

\begin{dfn}[当人認証の実装パターン]\label{dfn:ch3-login-auth-patterns}
ログインにおける当人認証の実装パターンとして、次の 2 つがある。
\begin{enumerate}
\item \underline{ユーザーの識別をしてから検証する}方式：ユーザーがメールアドレスやユーザー ID などを先に入力し、その識別子に紐づく認証方法（パスワード入力、パスキー認証など）をシステムが要求する。
\item \underline{ユーザーの識別と検証を同時に行う}方式：ユーザーがユーザー ID とパスワードなどを同時に入力し、システムがその組み合わせの正当性を一度に検証する。
\end{enumerate}
\end{dfn}

二つのパターンは、どのような場合に用いられるかが異なる。

\begin{prop}[各実装パターンが利用される場合]\label{prop:ch3-login-auth-pattern-usage}
\begin{enumerate}
\item 「識別をしてから検証する」方式は、ユーザー毎に認証方法が分かれている場合や、識別子に紐づくユーザーにメールや SMS を送信するようなケースで利用される。
\item 「識別と検証を同時に行う」方式は、パスワードや事前に発行されたバックアップコード（定義 \ref{dfn:backup-code}）のように静的なものを入力する場合に利用される。
\end{enumerate}
\end{prop}

\begin{proof}
各方式で、検証を始める前にシステムが識別子を知っている必要があるかどうかを確かめる。
まず (1) の場合、ユーザー毎に認証方法が分かれているなら、どの認証方法を要求するかを決めるためにユーザーを先に識別しなければならない。同様に、メールや SMS を送るには、送信先となる識別子に紐づくユーザーを先に特定しなければならない。よって識別を先に行う方式が用いられる。
次に (2) の場合、パスワードやバックアップコードは事前に決まっている静的な値であり、サービス側から何かを送ったり認証方法を選んだりする必要がない。よって識別子とまとめて一度に入力させ、組み合わせの正当性を一度に検証できる。
\end{proof}

\begin{exa}[サービスによる実装パターンの違い]\label{exa:ch3-google-facebook}
本書では、Google と Facebook とで当人認証の実装パターンが異なることを図で示している。
すなわち、定義 \ref{dfn:ch3-login-auth-patterns} の 2 つのパターンは、いずれも大規模なサービスで実際に採用されている。
\end{exa}

\underline{(2) ユーザーがサービスを利用できる状態であることの確認：}
当人認証に成功しても、そのユーザーがサービスを利用してよい状態にあるとは限らない。

\begin{dfn}[利用可能状態の確認]\label{dfn:ch3-login-state-check}
アイデンティティレジスター（定義 \ref{dfn:identity-register}）を参照し、ユーザーがサービスを利用可能な状態であるかを確認するプロセスを\underline{利用可能状態の確認}という。
この確認は当人認証と同時、またはその直後に行われる。
アイデンティティライフサイクルにおける「本登録」（$\mathsf{Active}$）状態であれば、次のプロセスに進む。
\end{dfn}

\begin{rem}[状態に応じた誘導]\label{rem:ch3-login-state-guidance}
もし「仮登録」状態であれば「有効化」プロセスへ、「無効化」状態であれば「再有効化」プロセスへ誘導するなど、状態に応じた対応が求められることもある。
ライフサイクルの記号でいえば、$\mathsf{Established}$ のユーザーは有効化 $\mathsf{Established}\to\mathsf{Active}$ へ、$\mathsf{Suspended}$ のユーザーは再有効化 $\mathsf{Suspended}\to\mathsf{Active}$ へ誘導される（表 \ref{tab:ch3-lifecycle-transitions}）。
\end{rem}

\underline{(3) ログインセッションの作成：}
新規登録時（定義 \ref{dfn:ch3-signup-session}）と同様に、ログインセッション（定義 \ref{dfn:login-session}）を作成する。

\begin{prop}[認証コンテキストの保持]\label{prop:ch3-auth-context}
複数の認証方式を提供しているサービスでは、どの認証方式でログインしたか、サービス内の認証強度の基準などの\underline{認証コンテキスト}をセッション情報として保持しておくことが、セキュリティ上特に重要である。
\end{prop}

\begin{proof}
複数の認証方式を提供しているサービスでは、ログインセッションが存在するという事実だけでは、そのセッションがどの方式、どの強度の当人認証を経て作られたかがわからない。
認証方式によって強度は異なるから、後の操作の可否や追加の認証の要否をセキュリティ上適切に判断するには、その情報が必要である（次の項の再認証で、認証方式や認証強度がより強度の高い認証を要求する際の判断材料になることを見る）。
よって、認証コンテキストをセッション情報として保持しておくことがセキュリティ上特に重要である。
\end{proof}

本書は、ログインで作成されるログインセッションも図で示している。新規登録で作成されるもの（命題 \ref{prop:ch3-signup-session-info}）と並べると、両者の違いはセッションに記録される発行の経緯と認証コンテキストに現れる。

\subsection{ログアウト}

この項では、ログインで作成したログインセッションをユーザー自身の意思で終わらせる機能であるログアウトを扱う。
ログインセッションは便利である反面、残り続けることがリスクになる。
まずそのリスクとログアウトの定義を述べ、次にログアウトで行われるプロセスを、サーバーサイドとクライアントサイドに分けて述べる。

\subsubsection{サービス利用を安全に終了する}

ログインセッション（定義 \ref{dfn:login-session}）が有効である限り、その端末からはログイン状態でサービスを利用できる。
これは正規のユーザーにとっては便利であるが、第三者がその端末に触れる場合には問題になる。

\begin{prop}[有効なログインセッションのリスク]\label{prop:ch3-session-risk}
ログインセッションが有効な状態で第三者が端末にアクセスできると、不正利用や個人情報の漏洩といったリスクが生じる。
\end{prop}

\begin{proof}
ログインセッションはログイン状態を維持するためのものであり（定義 \ref{dfn:ch3-signup-session}）、それが有効である間は、改めて当人認証を受けなくてもサービスを利用する権限が保たれる。
よって、その端末にアクセスできる第三者は、正規のユーザーとしてサービスを操作でき、不正利用や、閲覧できる個人情報の漏洩が起こりうる。
\end{proof}

\begin{dfn}[ログアウト]\label{dfn:logout}
ユーザーが自らの意思でログインセッション（定義 \ref{dfn:login-session}）を能動的に破棄し、未ログイン状態に戻すための操作を\underline{ログアウト}という。
\end{dfn}

ログアウトはアイデンティティライフサイクル（定義 \ref{dfn:identity-lifecycle}）の状態を変えない。変わるのはログインセッションの有無であり、ユーザーの状態は本登録（$\mathsf{Active}$）のままである。

\begin{exa}[不特定多数が触れる環境でのログアウト]\label{exa:ch3-shared-pc}
図書館やインターネットカフェの PC、あるいは他者と共有しているデバイスなど、不特定多数が触れる可能性のある環境でサービスを利用した後は、必ずログアウトを行い、命題 \ref{prop:ch3-session-risk} のリスクを防ぐことが強く推奨される。
本書では、共用 PC でログアウトを忘れたことにより、後からその PC を使った第三者に不正利用される様子を図で示している。
\end{exa}

\subsubsection{ログアウトで行われるプロセス}

ログインセッションは、クライアントサイドが保持するセッショントークンと、サーバーサイドが管理するトークンに紐づくユーザー情報とによって実現されている（注意 \ref{rem:ch3-session-implementation}）。
したがって、ログアウト処理もサーバーサイドとクライアントサイド双方での処理を組み合わせて行われる。

\begin{nota}[ログアウトで行われるプロセス]\label{nota:ch3-logout-process}
ログアウトで行われるプロセスは次の 3 つである。
\begin{enumerate}
\item サーバーサイドでのログインセッション無効化。
\item クライアントサイドでの認証情報削除。
\item ユーザーへのフィードバックとリダイレクト。
\end{enumerate}
\end{nota}

\underline{(1) サーバーサイドでのログインセッション無効化：}
サーバー側で、該当ユーザーの現在のログインセッション情報を無効化または破棄する。

\begin{prop}[サーバーサイドでの無効化の効果]\label{prop:ch3-server-side-logout}
サーバーサイドでログインセッション情報を無効化または破棄すれば、たとえクライアント側にセッショントークンが残っていたり第三者がそれを知り得たとしても、サービスへのアクセスは拒否されるようになる。
\end{prop}

\begin{proof}
ログインセッションは、サーバーサイドでセッショントークンに紐づくユーザー情報を管理することで実現される（注意 \ref{rem:ch3-session-implementation}）。
サーバー側でその情報を無効化または破棄すると、提示されたセッショントークンに紐づく有効なセッションがサーバー側に存在しなくなる。
よって、トークンそのものがクライアント側に残っていても、あるいは第三者がトークンを知っていても、そのトークンによるアクセスは拒否される。
\end{proof}

\underline{(2) クライアントサイドでの認証情報削除：}
ユーザーのブラウザやモバイルアプリに保存されているセッショントークンを削除する。
これにより、次回以降サービスにアクセスした際は、未ログイン状態として扱われる。

\underline{(3) ユーザーへのフィードバックとリダイレクト：}
処理が正常に完了したことをユーザーに明確に伝えるため、「ログアウトしました」などのメッセージを表示したり、ログインページやトップページへ自動的に遷移（リダイレクト）させたりする。

本書の図は、ログアウトではサーバーサイドでの処理 (1) とクライアントサイドでの処理 (2) の双方が行われることを示している。

\subsection{再認証}

この項では、ログイン状態のユーザーに改めて本人であることの確認を求める機能である再認証を扱う。
ログアウトがログインセッションを終わらせることでリスクを下げるのに対し、再認証はログインセッションを保ったまま、重要な操作の前に当人認証をやり直すことでリスクを下げる。
まず再認証を定義して再認証が要求される場面を整理し、次に再認証で行われるプロセスを述べる。

\subsubsection{重要な操作の前に本人であることを再確認する}

命題 \ref{prop:ch3-session-risk} のリスクは、共用端末でログアウトを忘れた場合に限らず、ログインセッションが有効なままの端末が第三者に操作されるときに常に生じる。
これに対するもう一つの対策が再認証である。

\begin{dfn}[再認証]\label{dfn:reauthentication}
ログイン状態のユーザーに対し、重要な操作を行う前などに改めて本人であることの確認を求める仕組みを\underline{再認証}という。
再認証は、公共端末におけるログアウト（定義 \ref{dfn:logout}）と同様に、ログインセッションが有効なままの端末を第三者に操作されるリスクを低減するために用いられる。
\end{dfn}

本書は、再認証の概要を図で示している。
再認証もアイデンティティライフサイクルの状態は変えず、本登録（$\mathsf{Active}$）のユーザーのログインセッションに対して行われる。

\begin{exa}[再認証が要求される場合]\label{exa:ch3-reauth-triggers}
サービスは、主に次のような場合に再認証を要求する。
\begin{enumerate}
\item ログインしてから一定の時間が経過した場合（コンシューマー向けサービスの平均利用時間との差異、エンタープライズ向けサービスの業務時間など）。
\item 重要なユーザー情報の変更（連絡先、クレデンシャル、個人情報など）。
\item 決済処理や入出金など、金銭が関わる重要な操作を行う場合。
\item IP アドレスや位置情報などサービスへのアクセス元情報の環境が変更された場合。
\item その他、セキュリティ上重要と判断される操作を行う場合。
\end{enumerate}
\end{exa}

再認証を頻繁に要求すると、ユーザーの負担が増える。この点については次のことがいえる。

\begin{rem}[利便性の高い認証方式による負担の軽減]\label{rem:ch3-reauth-convenience}
端末の画面ロック解除の仕組みを利用するパスキー認証（定義 \ref{dfn:passkey}）のように、利便性の高い認証方式を採用することで、ユーザーの負担を抑えつつ安全性を高められる。
\end{rem}

\subsubsection{再認証で行われるプロセス}

再認証で行われるプロセスは、ログインセッションに紐づいているユーザーに対する当人認証と、ログインセッションの情報の更新の 2 つである。

\begin{nota}[再認証で行われるプロセス]\label{nota:ch3-reauth-process}
再認証で行われるプロセスは、次の順に行われる。
\begin{enumerate}
\item ログインセッションに紐づいているユーザーに対する当人認証。
\item ログインセッションの情報を更新。
\end{enumerate}
\end{nota}

\underline{(1) ログインセッションに紐づいているユーザーに対する当人認証：}
再認証における当人認証は、ログインにおける当人認証（記号 \ref{nota:ch3-login-process} のプロセス (1)）と比べて、識別の段階を省ける点が異なる。

\begin{prop}[再認証ではクレデンシャルの検証のみを行う]\label{prop:ch3-reauth-verify-only}
再認証では、ユーザーの識別は行わず、そのログインセッションに紐づく本人であることを確認するためのクレデンシャル検証のみを行う。
\end{prop}

\begin{proof}
通常のログインでは、まずユーザーを識別し、その後に当人認証を行う（定義 \ref{dfn:ch3-login-auth-patterns}）。
一方、再認証の時点では既にログインセッションが存在し、ログインセッションには紐づくユーザー情報がサーバーサイドで管理されている（注意 \ref{rem:ch3-session-implementation}）。
したがって「誰が」アクセスしているかは明確であり、改めて識別する必要はない。
よって、そのログインセッションに紐づく本人であることを確認するクレデンシャル検証のみを行えばよい。
\end{proof}

\begin{exa}[認証方式ごとの再認証]\label{exa:ch3-reauth-methods}
認証方式ごとの再認証の例は次のとおりである。
\begin{enumerate}
\item \underline{パスワード認証におけるパスワード再確認：}パスワードの入力を要求する（パスワード認証は定義 \ref{dfn:password-auth}）。
\item \underline{ワンタイムパスワード（OTP）を用いる認証：}SMS や認証アプリで生成されたワンタイムパスワードの入力を要求する。
\item \underline{パスキー認証：}指紋認証、顔認証など、利用しているデバイスに搭載されたロック解除機能や、セキュリティキーの操作を要求する（パスキーは定義 \ref{dfn:passkey}）。
\end{enumerate}
サービス提供者は、安全性や利便性を考慮し、これらの認証方法を単独または組み合わせて利用する。
\end{exa}

\underline{(2) ログインセッションの情報を更新：}
当人認証が完了したら、サービスはログインセッションに関連する情報を更新する。
いつ、どの方法で再認証が行われたかを記録し、その後の再認証要求の判定などに利用する。

\begin{dfn}[再認証で更新されるセッション情報]\label{dfn:ch3-reauth-session-update}
再認証によって更新される主なログインセッションの情報には、次のようなものがある。
\begin{enumerate}
\item \underline{最終認証日時}：次回の再認証要求のタイミングを判断するために利用する。
\item \underline{利用した認証方式や認証強度}：操作のリスクレベルに応じて、より強度の高い認証を要求する際の判断材料とする。
\item \underline{再認証が行われた IP アドレスやデバイス情報}：不正アクセス調査における重要な手がかりとなる。
\end{enumerate}
\end{dfn}

本書の図は、再認証によってこれらのログインセッション情報が更新される様子を示している。
(2) の認証方式・認証強度は、ログイン時に保持した認証コンテキスト（命題 \ref{prop:ch3-auth-context}）を再認証の結果で更新するものと見ることができる。

\subsection{ユーザー情報の設定変更}

この項では、登録済みのユーザー情報をユーザー自身が変更・更新する機能である設定変更を扱う。
これはアイデンティティライフサイクルにおける「保守」に対応する。
まず設定変更がなぜ必要かを述べて定義し、次に設定変更の対象となる主な情報を分類する。

\subsubsection{登録した情報を変更・更新する}

サービスの利用を続ける中で、住所や連絡先といった登録情報は変化する。

\begin{prop}[設定変更機能の必要性]\label{prop:ch3-settings-necessity}
ユーザー自身で情報を更新する設定変更機能は不可欠である。
\end{prop}

\begin{proof}
登録情報が古いままの場合に生じる不都合を挙げ、それを防ぐ手段が必要であることを示す。
まず、サービスの利用を続ける中で、住所や連絡先といった登録情報は変化する。
古い情報のままでは、例えば次のようなさまざまな不都合が生じる可能性がある。
\begin{enumerate}
\item 重要なお知らせが届かない。
\item 購入した商品が旧住所に配送される。
\item 連絡先を使ったログインができなくなる。
\end{enumerate}
よって、こうした問題を防ぎ、ユーザーが常に最新の状態でサービスを利用できるようにするため、ユーザー自身で情報を更新する設定変更機能は不可欠である。
\end{proof}

\begin{dfn}[ユーザー情報の設定変更]\label{dfn:ch3-settings-change}
アイデンティティレジスター（定義 \ref{dfn:identity-register}）に登録されたユーザー情報を、ユーザー自身が変更・更新する機能を\underline{ユーザー情報の設定変更}という。
これはアイデンティティライフサイクルにおける「保守」プロセスにあたり、表 \ref{tab:ch3-lifecycle-transitions} の遷移 $\mathsf{Active}\xrightarrow{\text{保守}}\mathsf{Active}$ に対応する。
\end{dfn}

\subsubsection{設定変更の対象となる主な情報}

設定変更の対象となる情報は、サービス特性に依存するものを除き、一般的に次のように分類される。

\begin{dfn}[設定変更の対象となる情報の分類]\label{dfn:ch3-settings-categories}
設定変更の対象となる主な情報は、次の 5 つに分類される。
\begin{enumerate}
\item \underline{プロフィール情報}：ニックネーム、アイコン、自己紹介文など。
\item \underline{連絡先情報}：メールアドレス、電話番号など、通知や身元確認に利用される情報。
\item \underline{基本4情報などの情報群}：氏名、住所、生年月日など、法的な本人確認や配送に用いられる情報。
\item \underline{認証方式とクレデンシャル}：パスワード変更、多要素認証（MFA。定義 \ref{dfn:mfa}）の設定、パスキーの管理など、セキュリティに直結する項目。
\item \underline{外部サービス連携}：ソーシャルログイン（ID 連携。定義 \ref{dfn:identity-federation}）アカウントの接続・解除など。
\end{enumerate}
\end{dfn}

設定変更機能は便利であるが、それゆえに悪用されたときの影響も大きい。

\begin{prop}[重要な情報の変更前には再認証を要求する]\label{prop:ch3-settings-reauth}
連絡先やクレデンシャルといった重要な情報を変更する際には、事前に再認証（定義 \ref{dfn:reauthentication}）を要求するのが一般的である。
\end{prop}

\begin{proof}
設定変更機能は利便性が高い一方、第三者に悪用されるとアカウント乗っ取りにつながるリスクをはらんでいる。
とくに連絡先やクレデンシャルは、通知や身元確認、当人認証に使われる情報である（定義 \ref{dfn:ch3-settings-categories}）から、これを第三者に書き換えられると、正規のユーザーがアカウントにアクセスできなくなり、第三者がアカウントを掌握することになる。
そこで、変更の直前に改めて本人であることを確認すれば、ログインセッションが有効なままの端末を第三者が操作していても変更を防げる（定義 \ref{dfn:reauthentication}）。
よって、重要な情報の変更前には再認証を要求するのが一般的である。これは例 \ref{exa:ch3-reauth-triggers} の (2) とも対応する。
\end{proof}

\subsection{セッション管理}

この項では、ユーザー自身が自分のログインセッションを確認し、不要なものを終了させる機能であるセッション管理を扱う。
第1章ではセッション管理を ID 管理の構成要素として概観した（定義 \ref{dfn:session-management}）。ここでは、それをユーザーに提供する機能として述べる。
まずセッション管理機能を定義し、次に提供される機能と、それを実現するためにログインセッションに求められる特性を述べる。

\subsubsection{ログインセッションをユーザー自身が確認する}

ログアウト（定義 \ref{dfn:logout}）はいま使っている端末のセッションを終わらせる操作であるが、他の端末に残ったセッションには及ばない。

\begin{prop}[残存するログインセッションのリスク]\label{prop:ch3-stale-session-risk}
共用端末でのログアウト忘れや、使用しなくなった端末にログインセッションが残っている場合、不正利用のリスクが高まる。
\end{prop}

\begin{proof}
命題 \ref{prop:ch3-session-risk} より、ログインセッションが有効な状態で第三者が端末にアクセスできると不正利用のリスクが生じる。
共用端末でログアウトを忘れた場合や、使用しなくなった端末にセッションが残っている場合は、正規のユーザーの手元を離れた端末に有効なセッションが存在し、第三者がその端末にアクセスしうる。
よって、不正利用のリスクが高まる。
\end{proof}

\begin{dfn}[セッション管理機能]\label{dfn:ch3-session-management-function}
ユーザー自身が有効なログインセッションの一覧を確認し、不要なセッションを遠隔で強制的に終了させるための機能を\underline{セッション管理}（機能）という。
\end{dfn}

\begin{rem}[不正ログインが疑われる場合のセッション管理]\label{rem:ch3-session-management-incident}
不正ログインが疑われる際には、いつ、どの環境からログインされたかを確認し、不審なセッションを即座に削除することが不可欠である。
ユーザー自身がログイン状況を主体的に管理できるようにすることで、サービスの信頼性と安全性を高める。
この点は本章末尾のコラム（\ref{subsec:ch3-column}）で詳しく述べる。
\end{rem}

本書は、セッション管理機能のインターフェース例を図で示している。後のコラムで言及される「ログインセッション」「Web sessions」の画面がその実例である。

\subsubsection{セッション管理で提供される機能}

セッション管理では、主に次の機能が提供される。

\begin{dfn}[セッション管理で提供される機能]\label{dfn:ch3-session-management-features}
\begin{enumerate}
\item \underline{有効なログインセッション一覧の表示}：各セッションが利用されたデバイス、IP アドレス、ブラウザ、最終利用日時などを一覧で確認できる。
\item \underline{ログインセッションの強制終了}：不要または不審なセッションを選択し、個別に、あるいは一括で強制的に終了（ログアウト）させる。
\end{enumerate}
\end{dfn}

これらの機能を実現するためには、ログインセッションの側にも条件が必要である。

\begin{prop}[ログインセッションに求められる特性]\label{prop:ch3-session-properties}
定義 \ref{dfn:ch3-session-management-features} の機能を実現するため、ログインセッションには次の特性が求められる。
\begin{enumerate}
\item ログインした環境（デバイスやブラウザ）ごとに一意に作成される。
\item ユーザーアカウントに紐づけられ、本人だけが一覧を参照できる。
\item ユーザー自身の操作で、任意のセッションを強制的に終了できる。
\end{enumerate}
\end{prop}

\begin{proof}
各特性が、定義 \ref{dfn:ch3-session-management-features} のどの機能に必要かを確かめる。
まず、一覧表示で各セッションのデバイスやブラウザを区別して示し、不要なものだけを選んで終了させるには、セッションが環境ごとに一意に作成されていなければならない。よって (1) が必要である。
次に、一覧表示は「ユーザー自身が」自分の有効なログインセッションを確認する機能である（定義 \ref{dfn:ch3-session-management-function}）から、セッションがユーザーアカウントに紐づいていなければそのユーザーのセッションの一覧を作れず、またその一覧は本人が参照するためのものである。よって (2) が必要である。
最後に、強制終了の機能はユーザー自身が任意のセッションを終了させるものであるから、(3) が必要である。
\end{proof}

\subsection{本人確認（KYC）}

この項では、新規登録時以外に行われる身元確認である本人確認（KYC）を扱う。
第2章で、「本人確認」という言葉が文脈によって身元確認と当人認証のどちらか、または両方を指しうることを述べた（注意 \ref{rem:honnin-kakunin}）。
ここでの本人確認は、そのうち身元確認の側に重点を置くものである。
まず本人確認を定義し、次に本人確認が行われる典型的な場面を三つに分けて述べる。

\subsubsection{現実世界の身元を証明する}

オンラインサービスでは、新規登録時以外にも身元確認を行う場面がある。

\begin{dfn}[本人確認（KYC）]\label{dfn:kyc}
オンラインサービスにおいて、新規登録時以外に行われる身元確認（定義 \ref{dfn:identity-proofing}）を、一般的に\underline{本人確認}（Know Your Customer: KYC）と呼ぶ。
\end{dfn}

\begin{rem}[再認証と本人確認の違い]\label{rem:ch3-reauth-vs-kyc}
前述のとおり、再認証（定義 \ref{dfn:reauthentication}）が当人認証（定義 \ref{dfn:authentication}）を要求するのに対し、本人確認ではより身元確認（定義 \ref{dfn:identity-proofing}）を要求するのが一般的である。
\end{rem}

本書は、本人確認の設計概念を図で示している。
本人確認が行われる場面は、次の三つに分けられる。

\begin{nota}[本人確認が行われる場面]\label{nota:ch3-kyc-cases}
本人確認が行われる典型的な場面として、次の三つを扱う。
\begin{enumerate}
\item 本人確認を必須とするサービス。
\item 任意で本人確認を求めるサービス。
\item カスタマーサポートにおける本人確認。
\end{enumerate}
\end{nota}

\underline{(1) 本人確認を必須とするサービス：}
法令によって本人確認が義務付けられているサービスがある。

\begin{exa}[本人確認を必須とするサービス]\label{exa:ch3-kyc-mandatory}
金融、通信、古物取引など、法令で本人確認が義務付けられているサービスでは、特定の機能を利用する前に公的身分証などを用いた本人確認が必須となる。
\end{exa}

\begin{rem}[任意のタイミングでの本人確認]\label{rem:ch3-kyc-anytime}
ID 管理の観点では、必須のタイミングだけでなく、ユーザーが任意のタイミングで本人確認を済ませておけるよう、設定変更機能（定義 \ref{dfn:ch3-settings-change}）の一部として提供することが望ましい。
\end{rem}

\underline{(2) 任意で本人確認を求めるサービス：}
法令上の義務がなくても、サービス自身の判断で本人確認を導入する場合がある。

\begin{exa}[任意で本人確認を求めるサービス]\label{exa:ch3-kyc-optional}
オークション・フリマ、クラウドソーシングといった個人間取引（C2C）サービスやマッチングサービスなど、法令上の義務はなくとも、サービスの信頼性・安全性向上のために任意の本人確認を導入するケースがある。
このようなサービスでは、本人確認済みのユーザーにのみ特定の機能を開放したり、特別なバッジを付与したりする。
\end{exa}

\underline{(3) カスタマーサポートにおける本人確認：}
本人確認は、サービスの画面上だけでなく、カスタマーサポートの窓口でも必要になる。

\begin{rem}[カスタマーサポートにおける身元確認]\label{rem:ch3-kyc-support}
ユーザー情報の変更や解約、アカウントリカバリーの補助など、機微な操作を伴うカスタマーサポートへの問い合わせ（お問い合わせフォームからの連絡など）の際に、依頼者が正当なユーザー本人であることを確認するための手段として身元確認（定義 \ref{dfn:identity-proofing}）を行う場合がある。
これにより、第三者による不正な操作を防ぐ。
\end{rem}

\subsection{アカウントリカバリー}

この項では、クレデンシャルが利用できなくなってログインできなくなったユーザーが、アカウントへのアクセスを取り戻すための機能であるアカウントリカバリーを扱う。
これは第2章で述べた当人認証の 4 つの局面（定義 \ref{dfn:authenticator-phases}）のうち「アカウントの回復」にあたる。
まずアカウントリカバリーを定義し、次にそのプロセスを身元確認とクレデンシャル再設定の二段階に分けて述べ、最後に注意点を挙げる。

\subsubsection{ログインできなくなったアカウントを復旧する}

ログイン（定義 \ref{dfn:login}）は登録済みのクレデンシャルによる当人認証を前提とするから、クレデンシャルが使えなくなればログインもできなくなる。

\begin{exa}[ログイン不能に陥る場合]\label{exa:ch3-login-failure}
「パスワードを忘れた」「認証アプリを入れたスマートフォンを紛失した」など、何らかの理由でクレデンシャルが利用できなくなり、ログイン不能に陥る場合がある。
\end{exa}

\begin{dfn}[アカウントリカバリー]\label{dfn:account-recovery}
クレデンシャルが利用できなくなり、利用可能な他の認証方法でもログインできない場合に、アカウントへのアクセスを復旧させるための機能を\underline{アカウントリカバリー}という。
アカウントリカバリーのプロセスは、大きく次の二つのステップで構成され、正当なユーザーであることを確認したうえで、再びサービスにアクセスできるようにする。
\begin{enumerate}
\item 身元確認。
\item （当人認証のための）クレデンシャルの再設定。
\end{enumerate}
\end{dfn}

本書はアカウントリカバリーをアイデンティティライフサイクル（定義 \ref{dfn:identity-lifecycle}）のどの遷移に当たるとも述べていない。ユーザーの状態を変える旨の記述もなく、行われるのはアイデンティティレジスターに登録されたクレデンシャルの再設定であるから、レジスターに登録されている情報の更新という点では保守（定義 \ref{dfn:ch3-lifecycle-actions}）と同じ性格をもつと読める。

\subsubsection{アカウントリカバリーで行われるプロセス}

定義 \ref{dfn:account-recovery} の二つのステップを順に述べる。

\underline{(1) 身元確認：}
ログインできなくなったユーザーは当人認証を受けられないから、正当なユーザーであることを別の方法で確かめる必要がある。
そこで、ログインできなくなったユーザーに対し、登録済みの属性情報などを用いて身元確認（定義 \ref{dfn:identity-proofing}）を行う。
どの方法を要求するかはサービスのポリシーによる。

\begin{dfn}[アカウントリカバリーにおける身元確認の方針]\label{dfn:ch3-recovery-proofing-policy}
アカウントリカバリーにおける身元確認で、どの方法を要求するかの方針として、主に次のものが考えられる。
\begin{enumerate}
\item \underline{新規登録で実施した身元確認を要求する}：新規登録時に実施したメールアドレス確認などを再度要求する。
\item \underline{これまで実施した最も強度の高い方法を要求する}：事前にマイナンバーカードなどでオンラインの本人確認（\underline{eKYC}）を実施済みのユーザーには、同様の方法を要求する。
\item \underline{再設定したいクレデンシャルの強度に合わせる}：パスキー（定義 \ref{dfn:passkey}）のような強度の高い認証情報を再設定する際は、eKYC など同等レベルの身元確認を要求する。
\end{enumerate}
\end{dfn}

\begin{rem}[各方針に必要なポリシー定義]\label{rem:ch3-recovery-policy-definition}
定義 \ref{dfn:ch3-recovery-proofing-policy} の (2) を採用するには、身元確認の強度に対するポリシー定義が必要である。
(3) を採用するには、身元確認、当人認証それぞれの強度と関係に対するポリシー定義が必要である。
\end{rem}

本書は、メールアドレスを用いた身元確認を図で示している。メールアドレス確認は、(1) の方針で例に挙げた方法である。

\underline{(2) 当人認証のためのクレデンシャル再設定：}
身元確認が完了したら、ユーザーはログインのための新しいクレデンシャルを設定する。
このフローは新規登録時の設定フロー（定義 \ref{dfn:ch3-credential-setup}）と類似していることが多く、ユーザーは指示に従って新しいパスワードやパスキーを登録する。

\begin{exa}[再設定用 URL の送信]\label{exa:ch3-reset-url}
アカウントリカバリーでこれまで広く使われてきたものとして、メールアドレスや電話番号への確認が完了した後に、クレデンシャル再設定用の URL を送信する方法がある。
\end{exa}

この方法には、デバイスにまたがる点で注意が要る。

\begin{prop}[開始デバイスと再設定デバイスが異なる場合]\label{prop:ch3-recovery-device-mismatch}
アカウントリカバリーを開始したデバイスと、リンクを開いて再設定を行うデバイスが異なる場合がある。
このとき、開始したデバイス側ではログイン状態が回復しない可能性がある。
したがって、パスキーや所持情報を用いる認証方式を提供する場合、ユーザーへの通知や導線についてシステム設計時の配慮が求められる。
\end{prop}

\begin{proof}
例 \ref{exa:ch3-reset-url} の方法では、再設定用の URL はメールや SMS で届くため、ユーザーはそれを、アカウントリカバリーを開始したデバイスとは別のデバイスで開くことがある。
例えば PC で開始し、スマートフォンでリンクを開いてクレデンシャルの再設定を完了した場合、再設定の操作はスマートフォン側で完了しているので、PC 側ではログイン状態が回復しない可能性がある。
このように再設定を行う環境が開始時と異なりうることを踏まえ、パスキーや所持情報を用いる認証方式を提供する場合には、ユーザーへの通知や導線についてシステム設計時の配慮が求められる。
\end{proof}

本書は、クレデンシャル再設定の環境についてのこの留意事項を図で示している。

\begin{prop}[クレデンシャル再設定の通知]\label{prop:ch3-recovery-notification}
クレデンシャルの再設定が完了した際は、その旨をユーザーに通知し、身に覚えのない操作に気づけるようにすることが重要である。
\end{prop}

\begin{proof}
第三者が正規ユーザーになりすましてアカウントリカバリーを行うリスクが存在する。
そのような不正なリカバリーが成功した場合、正規ユーザーはクレデンシャルが再設定されたことを知らなければ、被害に気づくことができない。
再設定の完了をユーザーに通知すれば、正規ユーザーは自分が行っていない再設定、すなわち身に覚えのない操作に気づける。
よって、再設定完了時の通知は重要である。
\end{proof}

\subsection{無効化と復旧}

この項では、サービス提供者がアカウントの利用を一時的に停止する無効化と、それを解除する復旧を扱う。
これらはアイデンティティライフサイクルの一時停止と再有効化に対応する。
まず無効化を、次に無効化したアカウントの復旧を定義する。

\subsubsection{アカウントの利用を一時的に停止・再開する}

これまでの機能は主にユーザー自身の操作によるものであったが、無効化はサービス提供者の判断で行われる。

\begin{dfn}[アカウントの無効化]\label{dfn:ch3-account-suspension}
利用規約違反や不正ログインの疑い、料金未払いなどの理由から、サービス提供者がユーザーを一時的に利用不可にすることを\underline{アカウントの無効化}という。
これは、アイデンティティライフサイクルにおいて、利用可能な「本登録」状態から、利用が制限された「無効化」状態へ遷移するプロセスである「一時停止」に相当する。すなわち
\[
\mathsf{Active}\xrightarrow{\text{一時停止}}\mathsf{Suspended}
\]
である。
\end{dfn}

\subsubsection{無効化したアカウントを復旧する}

無効化は一時的なものであり、原因が解決されれば元に戻せる場合がある。

\begin{dfn}[無効化したアカウントの復旧]\label{dfn:ch3-account-reactivation}
無効化されたユーザーが、原因となった問題を解決することで元の状態に戻ることを\underline{無効化したアカウントの復旧}という。
このプロセスは、「無効化」状態から再び「本登録」状態へ復帰させるための「再有効化」に相当する。すなわち
\[
\mathsf{Suspended}\xrightarrow{\text{再有効化}}\mathsf{Active}
\]
である。
\end{dfn}

\begin{exa}[無効化の原因の解決]\label{exa:ch3-suspension-resolution}
原因となった問題の解決には、例えば次のものがある。
\begin{enumerate}
\item 不正利用の疑いに対する身元確認。
\item 未払い料金の清算。
\end{enumerate}
\end{exa}

本書の図は、アイデンティティライフサイクルの図の上で、これらの遷移（一時停止と再有効化）を示している（図 \ref{fig:ch3-lifecycle} では本登録と無効化状態の間の二本の矢印）。
なお、ログイン時に無効化状態のユーザーを再有効化プロセスへ誘導することもある（注意 \ref{rem:ch3-login-state-guidance}）。

\subsection{退会と復旧}

この項では、サービスの利用を終了する退会と、退会したアカウントの復旧を扱う。
退会は内部的なデータの扱いによって論理削除と物理削除に分かれ、それぞれがアイデンティティライフサイクルの保存と削除に対応する。
どちらを採用するかによって、退会後に復旧できるかどうかが変わる。

\subsubsection{サービスの利用をやめてアカウントを削除する}

ユーザー自身の意思や契約期間満了により、サービス利用を終了する際に行われるのが退会処理である。
本書は、退会手続きを図で示している。

\begin{dfn}[退会（論理削除・物理削除）]\label{dfn:account-deletion}
ユーザー自身の意思や契約期間満了により、サービス利用を終了する際に行われる処理を\underline{退会}という。
退会は、内部的なデータの扱いによって、大きく次の 2 種類に分けられる。
\begin{enumerate}
\item \underline{論理削除}：ユーザーはサービスを利用できない状態になるが、属性情報や利用履歴などのデータはデータベース上に保持される。このプロセスはアイデンティティライフサイクルにおける「本登録」から「論理削除」状態への遷移である「保存」に相当する。すなわち
\[
\mathsf{Active}\xrightarrow{\text{保存}}\mathsf{Archived}
\]
である。
\item \underline{物理削除}：ユーザーに関連するすべてのデータがデータベースから完全に削除される。このプロセスは、アカウントの状態を問わず、データが存在しない「未登録」状態への遷移である「削除」に相当する。すなわち
\[
s\xrightarrow{\text{削除}}\mathsf{Unknown}\qquad (s\in\{\mathsf{Established},\mathsf{Active},\mathsf{Suspended},\mathsf{Archived}\})
\]
である。
\end{enumerate}
\end{dfn}

\begin{prop}[物理削除からは復旧できない]\label{prop:ch3-physical-deletion-irreversible}
物理削除を一度行うと、復旧は不可能である。
\end{prop}

\begin{proof}
物理削除では、ユーザーに関連するすべてのデータがデータベースから完全に削除される（定義 \ref{dfn:account-deletion}）。
復旧のもとになるデータがどこにも残らないので、元のアカウントを復元することはできない。
ライフサイクルの言葉でいえば、物理削除後の状態は未登録（$\mathsf{Unknown}$）であり、表 \ref{tab:ch3-lifecycle-transitions} において未登録から出る遷移は新たな登録だけである。
\end{proof}

\subsubsection{退会したアカウントを復旧できる場合とできない場合}

退会済みユーザーの復旧可否は、サービスが論理削除と物理削除のどちらの方式を採用しているかに依存する。

\begin{prop}[論理削除からの復旧]\label{prop:ch3-logical-deletion-restore}
論理削除を採用している場合、ユーザーはデータを含めて復旧できる可能性が高い。
これはアイデンティティライフサイクルにおける「再登録」と「有効化」の組み合わせに相当する。
\end{prop}

\begin{proof}
論理削除では、属性情報や利用履歴などのデータがデータベース上に保持されている（定義 \ref{dfn:account-deletion}）。
再登録は保存されていた情報を利用してアカウントを復元・再登録するプロセスであり（定義 \ref{dfn:ch3-lifecycle-actions}）、保持されているデータを用いてアカウントを復元できる。
さらに有効化によりアイデンティティ情報をアクティブな状態にすれば、ユーザーは再びサービスを利用できる。
よって、データを含めて復旧できる可能性が高く、その過程は再登録と有効化の組み合わせ
\[
\mathsf{Archived}\xrightarrow{\text{再登録}}\ \cdot\ \xrightarrow{\text{有効化}}\mathsf{Active}
\]
に相当する（中間の状態は本書に明記されておらず、表 \ref{tab:ch3-lifecycle-transitions} では仮登録と推定した）。
対照的に物理削除では、命題 \ref{prop:ch3-physical-deletion-irreversible} より復旧できない。
\end{proof}

\begin{exa}[論理削除の後に物理削除を行う運用]\label{exa:ch3-grace-period}
誤操作による退会に備え、一定期間は論理削除とし、その期間を過ぎたら物理削除を行う、という運用をしているサービスもある。
ライフサイクルの言葉では、まず $\mathsf{Active}\xrightarrow{\text{保存}}\mathsf{Archived}$ を行い、一定期間の後に $\mathsf{Archived}\xrightarrow{\text{削除}}\mathsf{Unknown}$ を行うことにあたる。
その期間内であれば、命題 \ref{prop:ch3-logical-deletion-restore} により復旧できる可能性が高い。
\end{exa}

本書の図は、アイデンティティライフサイクルの図の上で、退会（保存・削除）と復旧（再登録・有効化）の遷移を示している（図 \ref{fig:ch3-lifecycle} を参照）。

\subsection{セキュリティイベントログ}

この項では、アカウントに対して行われた重要な操作の履歴を記録し、ユーザー自身が確認できるようにする仕組みであるセキュリティイベントログを扱う。
まずセキュリティイベントログを定義してその意義を述べ、次にログに記録すべき情報を 5W1H の観点で整理する。
最後にコラムとして、不正ログインが疑われるときの対処の流れと、それを支える機能を述べる。

\subsubsection{自分のアカウントの履歴を確認する}

これまで紹介したログイン、設定変更、アカウントリカバリーといった機能は、いずれもユーザーのアカウント状態やセキュリティ設定に直接関わる重要な操作である。

\begin{dfn}[セキュリティイベントログ]\label{dfn:security-event-log}
ログイン（定義 \ref{dfn:login}）、設定変更（定義 \ref{dfn:ch3-settings-change}）、アカウントリカバリー（定義 \ref{dfn:account-recovery}）など、アカウント状態やセキュリティ設定に直接関わる操作について、「いつ、誰によって、何が行われたか」を記録し、ユーザー自身がその履歴を確認できるようにする仕組みを\underline{セキュリティイベントログ}という。
\end{dfn}

\begin{rem}[セキュリティイベントログの用途]\label{rem:ch3-security-log-uses}
セキュリティイベントログは、ユーザーへの提示以外にも、カスタマーサポートによる問い合わせ調査、データ解析など、さまざまな目的で利用できる。
\end{rem}

\begin{prop}[セキュリティイベントログの意義]\label{prop:ch3-security-log-value}
セキュリティイベントログは、ユーザー自身がアカウントの安全性を確認するための強力な手段となり、サービスへの信頼感を醸成するうえで非常に重要である。
\end{prop}

\begin{proof}
ユーザーが自身の身に覚えのない操作や不正アクセスの兆候に気づいた際、セキュリティイベントログを確認すれば、「いつ、誰によって、何が行われたか」が記録されているので、何が起こったのかを客観的な事実として把握できる。
よって、ユーザー自身がアカウントの安全性を確認するための強力な手段となる。
そしてそのような手段をサービスがユーザーに提供していること自体が、サービスへの信頼感を醸成するうえで非常に重要である。
\end{proof}

本書は、セキュリティイベントログの画面例を図で示している。後のコラムで言及される「セキュリティアクティビティ」「Security log」の画面がその実例である。

\subsubsection{ログに記録すべき情報}

実用的なログを設計するためには、一般的に「5W1H」の観点で記録すべき項目を定義するのが効果的である。

\begin{dfn}[セキュリティイベントログの記録項目（5W1H）]\label{dfn:ch3-security-log-5w1h}
セキュリティイベントログに記録すべき項目を、5W1H の観点で表 \ref{tab:ch3-security-log-5w1h} のように定める。
\end{dfn}

\begin{table}[H]
\centering
\caption{セキュリティイベントログに記録すべき情報（5W1H）}\label{tab:ch3-security-log-5w1h}
\begin{tabular}{p{0.22\textwidth}p{0.34\textwidth}p{0.32\textwidth}}
\hline
観点 & 記録する情報 & 例\\
\hline
When（いつ） & イベントの発生日時（タイムスタンプ） & \\
Where（どこで） & アクセス元の IP アドレス、デバイス情報、位置情報などの環境情報 & \\
Who（誰が） & 操作を行ったユーザーの識別子など & \\
What（何を） & イベントの種類 & \texttt{login.success}、\texttt{password.change}\\
Why/How（なぜ／どのように） & イベントの詳細情報 & 再認証のトリガーとなった理由（why）、使用された認証方式（how）\\
\hline
\end{tabular}
\end{table}

\begin{rem}[ユーザーに提示する範囲]\label{rem:ch3-security-log-display}
サービスは、これらを含むログの項目をすべてユーザーに提示する必要はない。
重要なのは、ユーザーが、履歴が自分自身のものかそうではないかを判断でき、何が行われたのかを把握できることである。
\end{rem}

\subsubsection{コラム：自分のアカウントに不正ログインされたかも？と思ったら}\label{subsec:ch3-column}

身に覚えのない通知が届くなどして、「自分のアカウントに不正ログインされたかも？」と思ったとき、どうすればよいか。
巷では「パスワードを変更しましょう」「パスキーを設定しましょう」といった対策を耳にするが、それよりも先にやるべきことがある。
それが「ログインセッションの確認」である。
このコラムでは、その理由と、続いて行うべきこと、そしてそれを可能にするためにサービス開発者が用意すべき機能を述べる。

\underline{「家に誰かいないかチェック」、つまり「ログインセッションの確認」：}
まずは現実の例で考える。

\begin{exa}[自宅の鍵が開いていたら]\label{exa:ch3-house-unlocked}
自宅に帰ったとき、自宅の鍵が開いていたとしたら、どうすればよいか。
貴重品が無事か確認する……その前に、まだ空き巣や不審者が室内に潜んでいる可能性を考え、家に誰かいないかをまず確認するべきである。
オンラインサービスの不正ログイン対応も、これとまったく同じ話である。
\end{exa}

オンラインサービスにおける「家に誰かいないかチェック」は「ログインセッションの確認」に該当する。

\begin{prop}[不正ログインが疑われたらまずログインセッションを確認する]\label{prop:ch3-check-sessions-first}
不正ログインが疑われた場合に真っ先に確認すべきは、ログインセッションである。
\end{prop}

\begin{proof}
オンラインサービスで個人情報を閲覧したり、投稿や設定変更を行ったりするには、ログインしている必要がある。
したがって、攻撃者がアカウント上で何かを行えるのは、攻撃者のログインセッションが有効である間である。
もし悪意のある攻撃者と正規のユーザーが同時にログインしていると、正規のユーザーがパスワード変更などの対策を行っても、攻撃者がそれを上書きでき、互いに設定を上書きしあう「設定変更合戦」のような事態になりかねない。
よって、パスワード変更などの対策よりも先に、攻撃者のログインセッションが残っていないかを確認すべきである。これは例 \ref{exa:ch3-house-unlocked} で、貴重品の確認より先に家に誰かいないかを確認するのと同じである。
\end{proof}

\begin{rem}[心当たりのないセッションの強制ログアウト]\label{rem:ch3-force-logout}
不正ログインが疑われる状況は緊急事態である。
現在操作している自身のセッション以外、心当たりのないログインセッションはすべて強制的にログアウトさせる（定義 \ref{dfn:ch3-session-management-features} の強制終了）。
すべての確認・変更作業は、そのあとに行う。
\end{rem}

\underline{可能ならば「何が行われたか」を確認し、設定を変更する：}
不審なログインセッションをすべて無効にしたら、次にアカウントでどのような操作が行われたか（アクティビティ履歴）を確認する。
すべての設定項目を一つずつ確認するのが理想であるが、現実的には難しい場合もある。

\begin{exa}[アクティビティ確認機能]\label{exa:ch3-google-activity}
例えば Google には、ログインや設定変更といったセキュリティ関連の操作履歴（アクティビティ）を確認する機能がある。
これをチェックすることで、不正ログインや意図しない変更が行われたことに気づける。
これはセキュリティイベントログ（定義 \ref{dfn:security-event-log}）の実例である。
\end{exa}

\begin{prop}[ログイン方法の見直し]\label{prop:ch3-review-login-methods}
アクティビティの確認が終わったら、連絡先やログイン方法（パスワード、パスキー、登録済みのメールアドレスや電話番号など）の設定を見直すべきである。
\end{prop}

\begin{proof}
攻撃者はログイン中に設定変更を行えるから、攻撃者によって新たなログイン方法が登録されている可能性がある。
その場合、ログインセッションを強制終了させても、攻撃者は登録したログイン方法で再び侵入できる。
よって、連絡先やログイン方法の設定を見直し、攻撃者が登録したものを取り除く必要がある。
\end{proof}

\begin{rem}[最後にもう一度ログインセッションを確認する]\label{rem:ch3-recheck-sessions}
同時進行で攻撃者も設定変更などをしている可能性があるため、すべての設定変更が終わった後に、念のため再度ログインセッションを確認することも重要である。
\end{rem}

以上の流れをまとめると、次のようになる。
\begin{enumerate}
\item 心当たりのないログインセッションをすべて強制的にログアウトさせる（命題 \ref{prop:ch3-check-sessions-first}、注意 \ref{rem:ch3-force-logout}）。
\item 可能ならアクティビティ履歴で何が行われたかを確認する（例 \ref{exa:ch3-google-activity}）。
\item 連絡先やログイン方法の設定を見直す（命題 \ref{prop:ch3-review-login-methods}）。
\item 念のため再度ログインセッションを確認する（注意 \ref{rem:ch3-recheck-sessions}）。
\end{enumerate}

\underline{開発者は上記の流れを意識して機能を用意することが重要：}
ユーザーが上の流れを実行できるかどうかは、サービスが必要な機能を提供しているかどうかにかかっている。

\begin{prop}[開発者が標準機能として提供すべきもの]\label{prop:ch3-developer-features}
サービス開発者側は、ユーザーがこうした自己防衛を行えるよう、「ログインセッションの管理機能」と「セキュリティイベントの確認機能」を、ID 管理の標準機能として提供することが強く求められる。
\end{prop}

\begin{proof}
上の流れの (1) と (4) はログインセッションの確認と強制終了であり、セッション管理機能（定義 \ref{dfn:ch3-session-management-function}）がなければユーザーには実行できない。
(2) は何が行われたかの確認であり、セキュリティイベントログ（定義 \ref{dfn:security-event-log}）がなければ実行できない。
よって、ユーザーが自己防衛を行えるようにするには、これら二つの機能を ID 管理の標準機能として提供する必要がある。
\end{proof}

\begin{exa}[標準機能の実例]\label{exa:ch3-feature-examples}
本書がセッション管理機能のインターフェース例として示した「ログインセッション」「Web sessions」の画面や、セキュリティイベントログの例として示した「セキュリティアクティビティ」「Security log」の画面は、まさに命題 \ref{prop:ch3-developer-features} の機能の実例である。
\end{exa}

\begin{rem}[実装の現状]\label{rem:ch3-feature-status}
大手サービスは、長年にわたって攻撃と戦ってきた経験から、こうした機能の重要性を認識し、実装しているケースが多い。
一方で、比較的新しいサービスや、パスワード認証（定義 \ref{dfn:password-auth}）のみを提供しているサービスでは、これらの機能が見過ごされがちなのが現状である。
\end{rem}

これらの機能が当たり前に実装されるようになれば、さらにその先にある、ID 連携サービス間でのセキュリティ情報共有（SSF の活用）といった、より高度な施策も視野に入ってくる。
複数のサービスにまたがるこうした施策は、第4章で扱う（セキュリティイベントの通知、定義 \ref{dfn:security-event-notification}）。

%% =======================================================
%%  part: ch4
%% =======================================================
\section{複数サービスが関わるID管理機能}

この章では、\cite{Ito2026} の第4章に沿って、複数のサービスが関わる場合の ID 管理機能を扱う。
第3章では、単一のサービスがすべての ID 管理機能を担うモデルを扱った。本章では、サービス間でアイデンティティ情報をやり取りする ID 連携（定義 \ref{dfn:identity-federation}）を導入すると、第3章の各機能がどのように変化するかを調べる。
構成は次のとおりである。まず ID 連携の登場人物と ID 管理機能との関係、および ID 連携で行われる各種の情報通知（アイデンティティプロビジョニングとセキュリティイベントの通知）を整理する。次に、IdP の側で強化される ID 管理機能を述べる。最後に、RP の側で ID 管理機能がどのように分業化され、第3章の新規登録・ログイン・ログアウト・再認証・アカウントリカバリー・本人確認・退会・セキュリティイベントログがそれぞれどう変わるかを、第3章のプロセスと対比しながら述べる。

%======================================================================
\subsection{ID連携の概要}

第3章で扱った ID 管理は、一つのサービスの中ですべてが閉じていた。ID 連携を導入すると、ID 管理機能の一部を別のサービスに任せたり、別のサービスへ情報を提供したりするようになるため、各サービスが担う ID 管理機能の性格が変わる。この項では、ID 連携の登場人物を固定し、それぞれの立場で ID 管理機能がどう変わるかを二つの観点にまとめた上で、ID 連携の中で行われる情報通知の仕組みを導入する。

%----------------------------------------------------------------------
\subsubsection{ID連携の登場人物とID管理機能の関係}

前章では、単一のサービスがすべての ID 管理機能を担うモデルを解説した。以下ではこれを \underline{自己完結型}のモデルと呼ぶ。本章では、サービス間でアイデンティティ情報をやり取りする「ID 連携」を導入することで、ID 管理機能がどのように変化するのかを解説する。
はじめに、ID 連携における登場人物と、ID 管理機能との関係を整理しておく。IdP と RP そのものの定義は第2章で与えた（定義 \ref{dfn:idp}、定義 \ref{dfn:rp}）。ここでは、本章の議論で必要になる役割分担の観点から、三者をあらためて記号としてまとめる。

\begin{nota}[ID連携の登場人物]\label{nota:ch4-federation-actors}
本章では、ID 連携の登場人物を次の三者とする。
\begin{enumerate}
\item \underline{ユーザー}：Relying Party のサービスを利用したい主体。
\item \underline{Identity Provider}（IdP，定義 \ref{dfn:idp}）：ユーザーに対する身元確認や当人認証を行い、その結果を RP へ提供するサービス。
\item \underline{Relying Party}（RP，定義 \ref{dfn:rp}）：自社サービスの身元確認や当人認証を、IdP に委ねるサービス。
\end{enumerate}
以後、特に断らない限り、IdP と RP はこの意味で用いる。
\end{nota}

\begin{figure}[H]
\centering
\begin{tikzpicture}[
  box/.style={draw, rounded corners, minimum width=3.2cm, minimum height=1.0cm, align=center},
  >=Stealth]
\node[box] (user) at (0,0) {ユーザー};
\node[box] (idp) at (-3.8,-2.6) {IdP\\ {\small 身元確認・当人認証}};
\node[box] (rp) at (3.8,-2.6) {RP\\ {\small サービスの提供}};
\draw[->] (user) -- node[left, align=right, font=\small] {身元確認・\\当人認証を受ける} (idp);
\draw[->] (user) -- node[right, align=left, font=\small] {サービスを\\利用したい} (rp);
\draw[->] (idp) -- node[below, font=\small] {結果（ユーザー情報）を提供} (rp);
\draw[->, dashed] (rp.north west) to[bend right=25] node[above, font=\small] {身元確認・当人認証を委ねる} (idp.north east);
\end{tikzpicture}
\caption{ID連携の登場人物と、IdP・RP が担う ID 管理機能}
\label{fig:ch4-federation-actors}
\end{figure}

図 \ref{fig:ch4-federation-actors} のように、IdP は身元確認・当人認証の結果を RP に提供し、RP はそれらの責務を IdP に委ねる。この役割分担により、IdP と RP のそれぞれで ID 管理機能の性格が変わる。IdP の側の変化を \underline{ID管理機能の強化}、RP の側の変化を \underline{ID管理機能の分業化} と呼ぶ。以下の二つの命題で、それぞれの内容と理由を述べる。

\begin{prop}[IdP：ID管理機能の強化]\label{prop:ch4-idp-strengthening}
IdP では、自己完結型のサービスに比べて、ID 管理機能の強化が不可欠である。すなわち、IdP の ID 管理機能には、より堅牢なセキュリティや高い可用性が求められる。
\end{prop}
\begin{proof}
IdP が管理するアイデンティティの価値が高まることから、その管理に求められる水準も高まることを示す。
まず、IdP は複数の RP から信頼され、それらの RP にユーザー情報を提供する。よって、IdP が管理するアイデンティティは、IdP 自身のサービスだけでなく、自サービス以外の多くのサービスで利用される情報となり、その価値は向上する。
次に、価値が向上した情報を扱う以上、それを管理する ID 管理機能にも、より堅牢なセキュリティや高い可用性が求められる。
したがって、自己完結型のサービスに比べて、IdP では ID 管理機能の強化が当然ながら不可欠である。
\end{proof}

\begin{prop}[RP：ID管理機能の分業化]\label{prop:ch4-rp-division}
RP は、ID 管理機能の一部、特に専門性が求められる身元確認（定義 \ref{dfn:identity-proofing}）や当人認証（定義 \ref{dfn:authentication}）の責務を IdP に委ねることで、自らのサービスの開発・運用コストを軽減できる。一方で、ID 連携を実現するためには、IdP から受け取ったユーザー情報と自社のユーザー情報の紐付けなど、新たに実装が必要となる機能もある。
\end{prop}
\begin{proof}
前半（コストの軽減）と後半（新たな実装の必要）を順に示す。
まず、身元確認や当人認証は専門性が求められる機能であり、自己完結型のサービスではこれを自ら開発・運用しなければならなかった。RP はこの責務を IdP に委ねるので、その分だけ自らのサービスの開発・運用コストを軽減できる。これは、ID 管理機能を IdP と RP の間で分担すること、すなわち ID 管理における「ID 管理機能の分業化」と言える。
次に、RP は IdP から受け取ったユーザー情報を自社のサービスで使うために、それを自社のユーザー情報と紐付けなければならない。この紐付けは自己完結型のサービスには存在しなかった処理であるから、ID 連携を実現するために新たに実装が必要となる。
以上より主張が成り立つ。
\end{proof}

本章では、この ID 管理機能の強化（命題 \ref{prop:ch4-idp-strengthening}）と ID 管理機能の分業化（命題 \ref{prop:ch4-rp-division}）という二つの観点から、ID 連携によって IdP と RP それぞれで変化する ID 管理機能を詳しく見ていく。

%----------------------------------------------------------------------
\subsubsection{ID連携にて行われる各種情報の通知}

ID 連携では、ユーザーの同意に基づく IdP から RP への情報提供や、非同期のリソースアクセスが実現できる。これとは別に、情報通知という観点で、IdP から RP へ、あるいはその逆方向へ情報が送られる仕組みがいくつかある。この項では、そのうちアイデンティティ情報の変更を通知する「アイデンティティプロビジョニング」と、重要な事象を通知する「セキュリティイベントの通知」を導入する。これらは本章の後の二つの項（\ref{subsec:ch4-idp} 項と \ref{subsec:ch4-rp} 項）で IdP・RP の各機能を論じる際に繰り返し用いる。

\underline{アイデンティティ情報の変更を通知するアイデンティティプロビジョニング：}
どのタイミングで IdP から RP へ情報が送られるかは、サービスの種類によって異なる。コンシューマー向けサービスでは、ユーザー自身の操作が起点となり、必要な時に必要な情報だけが連携されることが一般的である。
一方、エンタープライズ向けサービスでは、管理者が IdP にユーザーを登録した時点で、連携先の RP へ自動的に情報を提供することが主流である。

\begin{exa}[エンタープライズ向けサービスにおける新入社員の登録]\label{exa:ch4-enterprise-provisioning}
人事担当者が IdP に新入社員の情報を登録し、アイデンティティライフサイクル（定義 \ref{dfn:identity-lifecycle}）における「仮登録」や「本登録」の状態にすると、その情報は即座に連携先の RP（勤怠管理、経費精算システムなど）へ送られる。これにより、ユーザーが実際に利用する前から、組織情報に基づいた適切なアクセス権を持つアカウントを RP 側に用意でき、組織としてのリアルタイムなアカウント管理が実現する。
\end{exa}

\begin{exa}[属性情報の変更と退職の通知]\label{exa:ch4-provisioning-transfer-retirement}
登録時以外にも、「社員の部署異動」といった属性情報の変更や、社員の退職時に IdP がユーザーを「論理削除」や「未登録」の状態にした際にも、IdP はその状態変更を RP へ通知する。これにより、RP の側でも業務フローを適切に制限できる。
\end{exa}

例 \ref{exa:ch4-enterprise-provisioning} と例 \ref{exa:ch4-provisioning-transfer-retirement} に共通するのは、アイデンティティライフサイクル上の出来事（登録、属性情報の変更、論理削除・削除）を契機として、IdP から RP へ\emph{自動的に}情報が送られ、RP の側のアイデンティティが IdP の側と同期される点である。これを概念として定義する。

\begin{dfn}[アイデンティティプロビジョニング]\label{dfn:identity-provisioning}
アイデンティティのライフサイクル（定義 \ref{dfn:identity-lifecycle}）に応じて、IdP から RP へ自動的に情報を送り、両者の間のアイデンティティ情報の同期を可能とする仕組みを \underline{アイデンティティプロビジョニング}（identity provisioning）と呼ぶ。送られる情報には、ユーザーの登録（「仮登録」「本登録」への遷移）、属性情報の変更、および「論理削除」「未登録」などへの状態変更が含まれる。
\end{dfn}

\begin{figure}[H]
\centering
\begin{tikzpicture}[
  box/.style={draw, rounded corners, minimum width=2.8cm, minimum height=0.9cm, align=center},
  >=Stealth]
\node[box] (hr) at (0,0) {人事担当者\\{\small（管理者）}};
\node[box] (idp) at (4.2,0) {IdP};
\node[box] (rp1) at (9,1.1) {RP\\{\small 勤怠管理}};
\node[box] (rp2) at (9,-1.1) {RP\\{\small 経費精算}};
\draw[->] (hr) -- node[above, font=\small, align=center] {登録・異動・\\退職を反映} (idp);
\draw[->] (idp) -- node[above, sloped, font=\small] {自動で通知} (rp1);
\draw[->] (idp) -- node[below, sloped, font=\small] {自動で通知} (rp2);
\end{tikzpicture}
\caption{エンタープライズ向けサービスのプロビジョニング}
\label{fig:ch4-enterprise-provisioning}
\end{figure}

\underline{セキュリティイベントの通知：}
プロビジョニングが主にライフサイクル上の変化を同期するのに対し、ID 連携ではもう一つ、緊急性のある事象を伝える仕組みが必要になる。その理由を次の命題で述べる。

\begin{prop}[IdP における不正ログインの被害は RP に拡大しうる]\label{prop:ch4-compromise-propagation}
ID 連携において、IdP のユーザーが不正ログインの被害に遭った場合、その情報を利用する RP にも被害が拡大する可能性がある。したがって、IdP は不正ログインを検知した時点で自サービスのセッションを破棄したり、ユーザーを一時的に無効化したりするのが基本的な対策であるが、RP 側でも同様の対応が取れるのが理想である。
\end{prop}
\begin{proof}
被害が拡大する経路を確認し、そこから RP 側の対応が望ましいことを導く。
まず、RP は IdP から提供されるユーザー情報や当人認証の結果に依拠してサービスを提供している（命題 \ref{prop:ch4-rp-division}）。よって、IdP のユーザーが不正ログインの被害に遭えば、その IdP の情報を利用する RP にも被害が拡大する可能性がある。
次に、IdP 自身は、不正ログインを検知した時点で自サービスのセッションを破棄したり、ユーザーを一時的に無効化したりすることで対処できる。しかしこの対処は IdP の中で閉じており、RP の側に既に存在するセッションやアカウントには及ばない。
したがって、被害の拡大を防ぐには、RP 側でも IdP と同様の対応が取れることが理想である。
\end{proof}

RP 側で同様の対応を取るためには、IdP で起きた事象を RP が知る手段が必要である。

\begin{dfn}[セキュリティイベントの通知]\label{dfn:security-event-notification}
不正ログインの検知といった、アイデンティティに関わる重要な事象を IdP から RP へ伝えるための仕組みを \underline{セキュリティイベントの通知}（security event notification）という。この通知は、アイデンティティプロビジョニング（定義 \ref{dfn:identity-provisioning}）と同様のユーザー情報の変更通知に加え、セッション管理に関する通知や、より緊急性のある「イベント」の発生も対象とする。
\end{dfn}

\begin{rem}[プロビジョニングとの重複]\label{rem:ch4-provisioning-overlap}
セキュリティイベントの通知（定義 \ref{dfn:security-event-notification}）の対象には、メールアドレスの変更など、一部アイデンティティプロビジョニング（定義 \ref{dfn:identity-provisioning}）と重複する内容もある。両者は排他的な仕組みではない。
\end{rem}

イベント通知を受け取った RP は、連絡先の同期や、退会者・リスクが検知されたユーザーのアクセス制限など、迅速な対応が可能になる。

\begin{rem}[RP から IdP への通知]\label{rem:ch4-rp-to-idp-notification}
セキュリティイベントの通知は IdP から RP への方向が基本であるが、サービス間の信頼関係によっては、RP から IdP へのセキュリティイベント通知も想定される。
\end{rem}

\begin{figure}[H]
\centering
\begin{tikzpicture}[
  box/.style={draw, rounded corners, minimum width=3.0cm, minimum height=0.9cm, align=center},
  >=Stealth]
\node[box] (idp) at (0,0) {IdP\\{\small 不正ログインを検知}};
\node[box] (rp) at (7,0) {RP\\{\small アクセス制限など}};
\draw[->, thick] ([yshift=4pt]idp.east) -- node[above, font=\small] {セキュリティイベントの通知} ([yshift=4pt]rp.west);
\draw[->, dashed] ([yshift=-4pt]rp.west) -- node[below, font=\small] {（信頼関係によっては逆方向も）} ([yshift=-4pt]idp.east);
\end{tikzpicture}
\caption{IdP から RP へのセキュリティイベントの通知}
\label{fig:ch4-security-event-notification}
\end{figure}

IdP や RP は、アイデンティティプロビジョニングとセキュリティイベントの通知を適切に利用することで、ユーザーにとってシームレスな体験と安全性の向上を実現できる。

%======================================================================
\subsection{IdPのID管理機能強化}\label{subsec:ch4-idp}

命題 \ref{prop:ch4-idp-strengthening} で述べたように、IdP は複数の RP から信頼される立場にあるため、自己完結型のサービスにはなかった機能が求められる。その多くは、IdP 上で起きた変化（新規登録、ログアウト、ユーザー情報の変更、アカウント状態の変更）を RP に伝える通知と、RP からの要求に応える機能である。この項では、第3章で扱った単一サービスの各機能に対応させながら、IdP で強化される ID 管理機能を順に述べる。

%----------------------------------------------------------------------
\subsubsection{新規登録時の情報通知}

第3章の新規登録（定義 \ref{dfn:signup}）は、単一サービスの中でユーザーを「未登録」から「本登録」へ遷移させる機能であった。IdP の場合は、この遷移が IdP の中で閉じない場合がある。
アイデンティティプロビジョニング（定義 \ref{dfn:identity-provisioning}）の説明で述べたとおり、一部のユースケース（例 \ref{exa:ch4-enterprise-provisioning} のようなエンタープライズ向けサービス）では、新規登録のタイミングで IdP から RP への情報通知が行われる。

%----------------------------------------------------------------------
\subsubsection{セッション情報の変更通知}

第3章のログアウト（定義 \ref{dfn:logout}）は、単一サービスのログインセッション（定義 \ref{dfn:login-session}）を破棄すれば完結した。ID 連携では IdP と RP がそれぞれ別のログインセッションを持つため、一方のログアウトが他方に反映されないという新しい問題が生じる。この項ではその問題を定義し、IdP からの通知による解消方法を述べる。

ID 連携において、IdP からユーザー情報を受け取った RP は、独自のログインセッションを作成する。ログイン直後は IdP と RP のユーザーは一致しているが、その後の操作によって両者がずれることがある。

\begin{dfn}[セッションの不整合]\label{dfn:ch4-session-inconsistency}
ID 連携において、IdP でログインしているユーザーと、RP のログインセッションに紐づくユーザーとが異なる状態を \underline{セッションの不整合}という。
\end{dfn}

\begin{prop}[IdP でのログアウトだけではセッションの不整合が生じる]\label{prop:ch4-session-inconsistency}
ユーザーが IdP でログアウトしても RP 側のログインセッションは残ったままになるため、その後別のユーザーが IdP にログインすると、セッションの不整合（定義 \ref{dfn:ch4-session-inconsistency}）が生じる。特に、サービス間を頻繁に行き来する使い方では、この不整合が情報漏洩や意図しない操作の原因となりうる。
\end{prop}
\begin{proof}
状態の推移を順に追う。
まず、ユーザー $u$ が IdP にログインし、ID 連携によって RP にもログインした直後は、IdP のログインユーザーと RP のログインセッションに紐づくユーザーはどちらも $u$ であり、一致している。
次に、$u$ が IdP でログアウトする。RP は独自のログインセッションを持っており、IdP でのログアウトは RP のログインセッションを破棄しないから、RP 側のセッションは $u$ に紐づいたまま残る。
さらに、別のユーザー $u'\neq u$ が IdP にログインすると、IdP のログインユーザーは $u'$、RP のログインセッションのユーザーは $u$ となり、両者は異なる。よって、セッションの不整合が生じる。
最後に、サービス間を頻繁に行き来する使い方では、ユーザー $u'$ は IdP でのログインに続いて RP を利用するつもりで、実際には $u$ のセッションのまま RP を操作しうる。したがって、この不整合は情報漏洩や意図しない操作の原因となりうる。
\end{proof}

この問題を解消するため、IdP でのログアウトやセッション破棄のタイミングで RP へ通知する仕組みがある。通知を受けた RP は自身のセッションを破棄し、この不整合な状態を解消する。具体的な通知方法には、前述のセキュリティイベントの通知（定義 \ref{dfn:security-event-notification}）や、ID 連携プロトコル（定義 \ref{dfn:federation-protocol}）で定義された方法がある。

\begin{figure}[H]
\centering
\begin{tikzpicture}[
  box/.style={draw, rounded corners, minimum width=3.0cm, minimum height=0.9cm, align=center},
  >=Stealth]
\node[box] (idp) at (0,0) {IdP\\{\small ログアウト／セッション破棄}};
\node[box] (rp) at (7,0) {RP\\{\small 自身のセッションを破棄}};
\draw[->, thick] (idp) -- node[above, font=\small] {セッション破棄の通知} (rp);
\end{tikzpicture}
\caption{IdP から RP へのセッション破棄の通知}
\label{fig:ch4-session-revocation}
\end{figure}

%----------------------------------------------------------------------
\subsubsection{RPからの認証要求への対応}

第3章の再認証（定義 \ref{dfn:reauthentication}）では、単一サービスが設定変更や決済といった重要な操作の際に、自らユーザーに再認証を求めた。ID 連携では当人認証の責務が IdP にあるため、再認証も IdP が担うことになる。この項では、その流れと IdP 側の注意点を述べる。

\begin{dfn}[RP からの再認証要求]\label{dfn:ch4-reauth-request}
ID 連携において、RP が ID 連携プロトコル（定義 \ref{dfn:federation-protocol}）で対象ユーザーや認証方式の要件などを指定し、IdP にユーザーの再認証をリクエストすることを、\underline{RP からの再認証要求}という。これを受けた IdP は、要求に応じてユーザーにパスワードの再入力などを促し、その結果を RP に通知する。
\end{dfn}

つまり、単一サービスで「サービスがユーザーに再認証を求める」ことは、ID 連携の場合には「RP が IdP に再認証を要求し、IdP がユーザーに再認証を促し、その結果を RP へ提供する」ことで実現される。

\begin{rem}[IdP と RP のセッションは同期しているとは限らない]\label{rem:ch4-idp-reauth-user-check}
IdP と RP のセッションは必ずしも同期しているわけではない（命題 \ref{prop:ch4-session-inconsistency}）。そのため、RP からの再認証要求（定義 \ref{dfn:ch4-reauth-request}）を受けた IdP は、現在ログイン中のユーザーと RP が指定したユーザーが同一であるかを確認し、セッションの状態に応じた適切な処理を行う必要がある。
\end{rem}

\begin{figure}[H]
\centering
\begin{tikzpicture}[
  box/.style={draw, rounded corners, minimum width=2.6cm, minimum height=0.9cm, align=center},
  >=Stealth]
\node[box] (rp) at (0,0) {RP};
\node[box] (idp) at (5,0) {IdP};
\node[box] (user) at (10,0) {ユーザー};
\draw[->] ([yshift=5pt]rp.east) -- node[above, font=\small, align=center] {(1) 再認証要求\\{\scriptsize 対象ユーザー・認証方式の要件}} ([yshift=5pt]idp.west);
\draw[->] ([yshift=5pt]idp.east) -- node[above, font=\small] {(2) 再認証を促す} ([yshift=5pt]user.west);
\draw[->] ([yshift=-5pt]user.west) -- node[below, font=\small] {(3) パスワード再入力など} ([yshift=-5pt]idp.east);
\draw[->] ([yshift=-5pt]idp.west) -- node[below, font=\small] {(4) 結果を通知} ([yshift=-5pt]rp.east);
\end{tikzpicture}
\caption{RP からの再認証要求を受ける IdP}
\label{fig:ch4-reauth-request}
\end{figure}

%----------------------------------------------------------------------
\subsubsection{ユーザー情報変更時の通知}

第3章の設定変更では、ユーザーは単一サービスの中で自分の情報を更新すればよかった。ID 連携では同じ情報が複数のサービスに保存されているため、更新の手間と更新漏れが問題になる。この項では、IdP からの通知がこれを軽減することを述べる。

\begin{prop}[ユーザー情報の変更通知による同期]\label{prop:ch4-user-info-change-notification}
IdP での情報変更を RP に通知し、RP 側で情報を同期することで、ユーザーが利用中の全サービスで情報を更新する手間と、更新漏れのリスクを軽減できる。実現方法としては、アイデンティティプロビジョニング（定義 \ref{dfn:identity-provisioning}）とセキュリティイベントの通知（定義 \ref{dfn:security-event-notification}）のどちらでも対応可能である。
\end{prop}
\begin{proof}
通知がない場合とある場合を比べる。
まず、通知がない場合、ユーザーがメールアドレスなどを変更した際には、利用中の全サービスで情報を更新しなければならず、これは大変な手間である。また、更新を忘れたサービスには古い情報が残る。
次に、IdP での情報変更が RP に通知され、RP 側で情報が同期されるならば、ユーザーは IdP で一度変更するだけで済み、RP ごとの更新作業も、その更新漏れも生じない。よって、手間と更新漏れのリスクが軽減される。
最後に、実現方法について述べる。アイデンティティプロビジョニングは属性情報の変更を IdP から RP へ自動的に送る仕組みであり（定義 \ref{dfn:identity-provisioning}）、セキュリティイベントの通知もユーザー情報の変更通知を対象に含む（定義 \ref{dfn:security-event-notification}、注意 \ref{rem:ch4-provisioning-overlap}）。したがって、どちらでも対応可能である。
\end{proof}

\begin{exa}[IdP 上でメールアドレスが変更された場合]\label{exa:ch4-email-change-notification}
ユーザーが IdP 上でメールアドレスを変更すると、IdP はその変更を連携先の RP に通知する。通知を受けた RP は、自サービスに保存しているメールアドレスを新しいものに同期する。
\end{exa}

%----------------------------------------------------------------------
\subsubsection{アカウント状態の変更通知（無効化・退会）}

第3章では、単一サービスがユーザーを一時停止（無効化）したり、ユーザーが退会したりする機能を扱った（退会については定義 \ref{dfn:account-deletion}）。これらはアイデンティティライフサイクル（定義 \ref{dfn:identity-lifecycle}）における状態遷移であり、単一サービスでは自サービスの状態を変えれば済んだ。IdP の場合は事情が異なる。

\begin{prop}[アカウント状態の変更通知の必要性]\label{prop:ch4-account-state-notification}
IdP におけるユーザーの無効化や退会は、連携するすべての RP のセキュリティと UX に影響するため、IdP から RP への迅速な通知が不可欠である。
\end{prop}
\begin{proof}
無効化（一時停止）と退会の二つの場合に分けて、通知がないと RP のセキュリティや UX が損なわれ、通知があれば RP が適切に対応できることを示す。

\textbf{一時停止と再有効化}
IdP のユーザーに不正アクセスが疑われた場合、その影響は連携済みの RP にも及ぶ可能性がある（命題 \ref{prop:ch4-compromise-propagation}）。IdP が対象ユーザーを一時停止しつつ状態変更を RP に通知すれば、RP も一時停止や特定機能の利用制限などの自衛措置を講じられる。逆に、通知がなければ RP はこれらの措置を講じる契機を持たない。

\textbf{退会と再登録}
ユーザーが IdP を退会した場合も通知が重要である。通知がなければ、RP には利用実態のないアカウントが残り続け、管理上の負担やセキュリティリスクとなる。IdP からの通知を受け取ることで、RP は自サービス上でもユーザーを退会させたり、ID 連携を解除したりといった適切な対応が可能になる。

いずれの場合も、RP のセキュリティと UX は IdP 上の状態変更に左右され、通知によって初めて RP が適切に対応できる。したがって、IdP から RP への迅速な通知が不可欠である。
\end{proof}

\underline{一時停止と再有効化：}
命題 \ref{prop:ch4-account-state-notification} の証明で述べた IdP から RP への方向に加えて、逆方向の連携も考えられる。

\begin{exa}[RP を起点とする一時停止の連携]\label{exa:ch4-rp-initiated-suspension}
RP 側でユーザーの不正な振る舞いを検知した際に、RP から IdP へ通知し、IdP がユーザーを一時停止して、さらに他の RP へも通知するといった連携も考えられる。これは、注意 \ref{rem:ch4-rp-to-idp-notification} で述べた RP から IdP への通知の一つの使い方である。
\end{exa}

\underline{退会と再登録：}
IdP の退会が RP に通知されない場合、RP には利用実態のないアカウントが残り続ける。通知を受けた RP は、自サービス上でもユーザーを退会させる（定義 \ref{dfn:account-deletion}）か、ID 連携を解除するかを選ぶことになる。どちらを選ぶかは、後の「ID 連携における無効化と退会」で RP の側から論じる。

\begin{figure}[H]
\centering
\begin{tikzpicture}[
  box/.style={draw, rounded corners, minimum width=3.0cm, minimum height=0.9cm, align=center},
  >=Stealth]
\node[box] (idp) at (0,0) {IdP\\{\small 一時停止／退会}};
\node[box] (rp1) at (7,1.0) {RP\\{\small 一時停止・機能制限}};
\node[box] (rp2) at (7,-1.0) {RP\\{\small 退会・ID連携の解除}};
\draw[->] (idp) -- node[above, sloped, font=\small] {無効化の通知} (rp1);
\draw[->] (idp) -- node[below, sloped, font=\small] {退会の通知} (rp2);
\end{tikzpicture}
\caption{無効化、退会処理の通知}
\label{fig:ch4-state-change-notification}
\end{figure}

%======================================================================
\subsection{ID連携におけるRPのID管理機能}\label{subsec:ch4-rp}

命題 \ref{prop:ch4-rp-division} で述べたように、RP は身元確認や当人認証の責務を IdP に委ねる一方で、IdP から受け取った情報と自社のユーザーとの紐付けなどを新たに実装する。この項では、第3章で単一サービスについて定義した各機能（新規登録、ログイン、ログアウト、再認証、セッション管理、アカウントリカバリー、本人確認、無効化と退会、セキュリティイベントログ）が、RP においてどのように変わるかを、第3章のプロセスと対応がわかる形で順に述べる。

%----------------------------------------------------------------------
\subsubsection{新規登録の簡略化}

第3章の新規登録（定義 \ref{dfn:signup}）は、(1) 事前条件の確認、(2) 身元確認、(3) クレデンシャルの設定、(4) 登録・有効化、(5) ログインセッション作成、の五つのプロセスからなっていた。このうち (2) と (3) はユーザーにとって手間のかかる部分である。ID 連携を用いると、この二つが IdP から受け取った情報によって置き換えられる。まず具体例で様子を見る。

\begin{exa}[Google アカウントを IdP とする新規登録]\label{exa:ch4-google-signup}
ユーザーが Google アカウントを IdP として利用し、RP であるコンシューマー向けサービスに新規登録するケースを考える。このとき、RP は Google からユーザー識別子と検証済みのメールアドレスを受け取れる。
RP は IdP である Google を信頼することで、受け取ったメールアドレスを確認済みとみなす（原文はこれに続けて「自サービスで改めて有効性と所持していることを確認する必要があります」と述べている。注意 \ref{rem:ch4-email-reverification} を参照）。受け取った情報を基に RP 内にアカウントを作成し、Google のユーザーと紐付けるだけで、ユーザーはすぐにログインできるようになる。
\end{exa}

\begin{rem}[補足：受け取ったメールアドレスの再確認について]\label{rem:ch4-email-reverification}
例 \ref{exa:ch4-google-signup} の原文は、「確認済みとみなし」と述べた直後に「自サービスで改めて有効性と所持していることを確認する必要があります」と続けており、文意が前後で食い違って読める。本書の他の箇所では、ID 連携が新規登録の身元確認を大幅に簡略化すること（命題 \ref{prop:ch4-federated-signup-simplification}）、およびアカウントリカバリーにおいて検証済みのメールアドレスを受け取った RP は自社でのメールアドレス検証プロセスを省略できること（本項の「ID 連携におけるアカウントリカバリー」）が述べられている。これらと整合させるならば、ここは「改めて確認する必要はない」の意と読むのが自然と考えられる。本ノートでは原文の記述を変更せずに示し、この食い違いを注意として記す。
\end{rem}

\begin{prop}[ID連携による新規登録の簡略化]\label{prop:ch4-federated-signup-simplification}
ID 連携は、第3章の新規登録プロセス（定義 \ref{dfn:signup}）における「身元確認」と「クレデンシャルの設定」という手間のかかる部分を大幅に簡略化する。
\end{prop}
\begin{proof}
二つのプロセスそれぞれについて、ID 連携でどう置き換わるかを示す。
まず身元確認について。単一サービスでは、メールアドレス等の有効性をサービス自身が OTP の送信などで確認していた。ID 連携では、RP は信頼する IdP から検証済みの属性情報（例 \ref{exa:ch4-google-signup} では検証済みのメールアドレス）を受け取り、これを身元確認に用いる。
次にクレデンシャルの設定について。単一サービスでは、パスワードやパスキーなどのクレデンシャルを新たに発行・登録していた。ID 連携では、RP は受け取ったユーザー識別子を用いて、RP 内のアカウントを IdP のユーザーと紐付けるだけでよく、ユーザーはすぐにログインできるようになる。
よって、どちらのプロセスもユーザーにとっての手間が大幅に減る。
\end{proof}

この簡略化された新規登録のプロセスを、第3章の新規登録と同じ五段階で定義する。

\begin{dfn}[ID連携による RP の新規登録]\label{dfn:ch4-federated-signup}
ID 連携を用いた RP の新規登録は、次の五つのプロセスからなる。
\begin{enumerate}
\item 事前条件の確認（IdP から情報を受け取ることへの同意の取得を含む）
\item IdP から受け取った属性情報を用いた身元確認
\item IdP から受け取ったユーザー識別子を用いたクレデンシャル設定
\item 登録、有効化
\item ログインセッション作成
\end{enumerate}
\end{dfn}

第3章の新規登録と対比すると、表 \ref{tab:ch4-signup-comparison} のようになる。

\begin{table}[H]
\centering
\caption{単一サービスの新規登録と ID 連携による RP の新規登録の対比}
\label{tab:ch4-signup-comparison}
\begin{tabular}{c|p{5.2cm}|p{6.2cm}}
\hline
& 単一サービス（定義 \ref{dfn:signup}） & ID 連携による RP（定義 \ref{dfn:ch4-federated-signup}） \\
\hline
(1) & 事前条件の確認（年齢制限、利用規約への同意、招待リンクの検証など） & 事前条件の確認に加え、IdP から情報を受け取ることとその利用方法を規約等に明記し、同意を得る \\
\hline
(2) & 身元確認（メールアドレス等の有効性をサービス自身が確認） & IdP から受け取った確認済みの属性情報を用いた身元確認 \\
\hline
(3) & クレデンシャルの設定（パスワード、パスキー、MFA の設定など） & IdP から受け取ったユーザー識別子を用いたクレデンシャル設定（「どの IdP の、どのユーザーか」との紐付け） \\
\hline
(4) & 登録、有効化（タイミングはサービスの形態による） & 同様にサービスの形態による（即時に有効化、または管理者の承認を経て有効化） \\
\hline
(5) & ログインセッション作成（「新規登録直後に発行されたセッションである」ことを示す情報を含む） & ログインセッション作成（「どの IdP の、どのユーザーと連携しているか」「新規登録で作成された」という情報を含む） \\
\hline
\end{tabular}
\end{table}

以下、定義 \ref{dfn:ch4-federated-signup} の各プロセスを順に述べる。

\underline{(1) 事前条件の確認：}
ID 連携を利用する場合、RP は利用規約やプライバシーポリシーにおいて、IdP から情報を受け取ることやその利用方法について明記し、ユーザーから同意を得る必要がある。

\underline{(2) IdP から受け取った属性情報を用いた身元確認：}
通常、コンシューマー向けサービスでは、連絡先であるメールアドレス等の有効性をサービス自身が確認する。これに対し ID 連携では、信頼する IdP から提供された確認済みの属性情報を身元確認に利用することで、利便性を向上させる。どの情報を身元確認に用いるかは、RP のサービス内容や準拠すべき法制度によって決まる。

\begin{exa}[マイナンバーカードの情報を用いた身元確認]\label{exa:ch4-mynumber-proofing}
デジタル庁の認証アプリを IdP とすれば、マイナンバーカードの情報を身元確認に利用することも可能である。
\end{exa}

\underline{(3) IdP から受け取ったユーザー識別子を用いたクレデンシャル設定：}
RP は、自サービスのユーザーと、「どの IdP の、どのユーザーか」を紐付ける。アイデンティティレジスター（定義 \ref{dfn:identity-register}）や ID 連携のためのデータストア等でこの紐付けを管理することで、ID 連携によるログインを実現できる。
記号で書けば、RP は自サービスのユーザー $u$ に対して、IdP の識別子 $i$ と、その IdP におけるユーザー識別子 $s$ の組 $(i,s)$ を紐付けて保存する。単一サービスにおけるパスワードやパスキーの役割を、この組 $(i,s)$ が担う。

\underline{(4) 登録、有効化：}
アカウントの「登録」と「有効化」のタイミングは、単一サービスと同様にサービスの形態による。ID 連携によって即時に有効化までを行うサービスもあれば、エンタープライズ向けサービスのように、登録後に管理者の承認を経て有効化されるケースもある。

\underline{(5) ログインセッション作成：}
ID 連携による新規登録で作成されるログインセッション（定義 \ref{dfn:login-session}）には、「どの IdP の、どのユーザーと連携しているか」「新規登録で作成された」という情報を含む。

%----------------------------------------------------------------------
\subsubsection{ID連携におけるログイン}

第3章のログイン（定義 \ref{dfn:login}）は、(1) 当人認証、(2) ユーザーがサービスを利用できる状態であることの確認、(3) ログインセッションの作成、の三つのプロセスからなっていた。ID 連携では、このうち (1) の当人認証を IdP への信頼に基づいて行う。

新規登録や後述する設定管理により、IdP ユーザーと RP ユーザーの紐付けが完了している場合、RP は ID 連携を「当人認証」に利用できる。これは、第3章で述べた当人認証の実装パターンのうち、「ユーザーの識別と検証を同時に行う」場合に相当する。実際、RP は IdP から受け取る識別子の組 $(i,s)$ によって、ユーザーが誰であるか（識別）と、そのユーザー本人であること（検証）を一度に判断する。

\begin{dfn}[ID連携による RP のログイン]\label{dfn:ch4-federated-login}
ID 連携を用いた RP のログインは、次の三つのプロセスからなる。
\begin{enumerate}
\item IdP から受け取ったユーザー識別子を用いた当人認証
\item ユーザーがサービスを利用できる状態であることを確認
\item ログインセッションの作成
\end{enumerate}
\end{dfn}

\begin{table}[H]
\centering
\caption{単一サービスのログインと ID 連携による RP のログインの対比}
\label{tab:ch4-login-comparison}
\begin{tabular}{c|p{5.2cm}|p{6.2cm}}
\hline
& 単一サービス（定義 \ref{dfn:login}） & ID 連携による RP（定義 \ref{dfn:ch4-federated-login}） \\
\hline
(1) & 当人認証（識別してから検証する、または識別と検証を同時に行う） & IdP から受け取ったユーザー識別子を用いた当人認証（識別と検証を同時に行う場合に相当）。必要に応じて追加認証 \\
\hline
(2) & ユーザーがサービスを利用できる状態であることを確認 & 同じく必要 \\
\hline
(3) & ログインセッションの作成（認証コンテキストを保持） & ログインセッションの作成（「どの IdP の、どのユーザーと連携しているか」「ログインで作成された」という情報を含む） \\
\hline
\end{tabular}
\end{table}

\underline{(1) IdP から受け取ったユーザー識別子を用いた当人認証：}
ID 連携では IdP への信頼に基づき、IdP の識別子と IdP のユーザー識別子に紐づく RP ユーザーが存在するかを検証する。すなわち、IdP から組 $(i,s)$ を受け取ったとき、$(i,s)$ が紐付けられた RP ユーザー $u$ が存在するかを検証する。RP のポリシーによっては、対象ユーザーが存在すれば当人認証が完了したとみなすこともあれば、さらに別の認証方式による追加認証を要求する場合もある。

\begin{exa}[RP 側の追加認証による MFA の実現]\label{exa:ch4-rp-additional-auth}
決済機能など強度の高い当人認証が求められる RP において、IdP での認証方式がパスワードのみだった場合、RP 側で追加の認証を要求して MFA（多要素認証，定義 \ref{dfn:mfa}）を実現する。
\end{exa}

\underline{(2) ユーザーがサービスを利用できる状態であることを確認：}
ID 連携を利用する場合にも、この確認は必要である。IdP による当人認証が成功しても、RP 上のそのユーザーが「本登録」状態にあるとは限らないからである。

\underline{(3) ログインセッションの作成：}
ID 連携によるログインで作成されるログインセッションには、「どの IdP の、どのユーザーと連携しているか」「ログインで作成された」という情報を含む。

%----------------------------------------------------------------------
\subsubsection{ID連携におけるログアウト}

第3章のログアウト（定義 \ref{dfn:logout}）では、RP は自身のログインセッションを破棄すれば十分であった。ID 連携では、IdP と RP の結びつきの強さによって、これだけでは足りない場合がある。

\begin{prop}[RP から IdP へのログアウト通知]\label{prop:ch4-rp-logout-notification}
IdP と RP が密結合であり「常に IdP と RP でログインユーザーを一致させる」という要件がある場合、RP は自身のセッションを破棄するだけでなく、IdP にログアウト処理が行われたことを通知する必要がある。これにより疎結合なセッション同期が図られる。実現方法には、セキュリティイベントの通知（定義 \ref{dfn:security-event-notification}）や ID 連携プロトコル（定義 \ref{dfn:federation-protocol}）で定義された方法が利用される。
\end{prop}
\begin{proof}
要件と、RP が自身のセッションを破棄しただけの状態とを比べる。
まず、RP が自身のセッションを破棄しただけでは、IdP のログインセッションは残る。このとき RP では未ログイン、IdP ではログイン中となり、両者のログイン状態は一致しない。これは「常に IdP と RP でログインユーザーを一致させる」という要件に反する。
次に、RP が IdP にログアウト処理が行われたことを通知すれば、IdP は自身のセッションを破棄でき、両者のログイン状態は再び一致する。IdP 側から RP へ通知する場合（セッション情報の変更通知）と合わせて、どちらの側でログアウトしても状態を揃えられる。
よって、この要件のもとでは RP から IdP への通知が必要であり、それによって疎結合なセッション同期が図られる。
\end{proof}

\begin{figure}[H]
\centering
\begin{tikzpicture}[
  box/.style={draw, rounded corners, minimum width=3.0cm, minimum height=0.9cm, align=center},
  >=Stealth]
\node[box] (rp) at (0,0) {RP\\{\small 自身のセッションを破棄}};
\node[box] (idp) at (7,0) {IdP\\{\small セッションを同期}};
\draw[->, thick] (rp) -- node[above, font=\small] {ログアウトの通知} (idp);
\end{tikzpicture}
\caption{RP から IdP へのログアウト通知}
\label{fig:ch4-rp-logout-notification}
\end{figure}

%----------------------------------------------------------------------
\subsubsection{ID連携における再認証}

第3章の再認証（定義 \ref{dfn:reauthentication}）は、(1) ログインセッションに紐づいているユーザーに対する当人認証、(2) ログインセッションの情報を更新、の二つのプロセスからなっていた。ID 連携では、(1) を IdP に要求する形で行う。これは IdP の側から見れば、\ref{subsec:ch4-idp} 項で述べた RP からの再認証要求（定義 \ref{dfn:ch4-reauth-request}）にあたる。

RP 上で決済や設定変更といった重要な処理を行う際、RP は IdP に対して再認証を要求する。このとき、対象ユーザーや必要な認証強度などを指定することで、単一サービスの再認証と同様の安全性と UX を担保できる。

\begin{dfn}[ID連携による RP の再認証]\label{dfn:ch4-federated-reauth}
ID 連携を用いた RP の再認証は、第3章の再認証と同じく次の二つのプロセスからなる。
\begin{enumerate}
\item ログインセッションに紐づいているユーザーに対する当人認証（IdP に要求して行う）
\item ログインセッションの情報を更新
\end{enumerate}
\end{dfn}

\underline{(1) ログインセッションに紐づいているユーザーに対する当人認証：}
IdP と RP のログインセッションは必ずしも同期しているわけではない（命題 \ref{prop:ch4-session-inconsistency}）ため、RP は現在ログイン中のユーザーに対して認証を要求する必要がある。その際、紐付けられた IdP のユーザー識別子と、IdP から再度受け取るユーザー識別子の一致を検証する。すなわち、RP のログインセッションのユーザー $u$ に紐付けられた組を $(i,s)$、再認証で IdP $i$ から受け取ったユーザー識別子を $s'$ とするとき、$s'=s$ であることを確かめる。
IdP によっては対象ユーザーを指定できる場合があり、ユーザー体験を損なわずに特定のユーザーに当人認証を要求できる。

\underline{(2) ログインセッションの情報を更新：}
ID 連携による再認証で作成されるログインセッションには、「どの IdP の、どのユーザーに対していつ認証したか」という情報を含む。

%----------------------------------------------------------------------
\subsubsection{ID連携のための設定管理}

第3章の設定変更では、ユーザーが自身の登録情報を変更・更新できるようにした。ID 連携を導入した RP では、それに加えて ID 連携そのものに関する設定をユーザーが管理できる必要がある。

\begin{dfn}[ID連携のための設定管理]\label{dfn:ch4-federation-settings}
RP が管理する、次のような ID 連携に関する設定の機能を、\underline{ID連携のための設定管理}という。
\begin{enumerate}
\item \underline{現在の ID 連携の状態確認}：「どの IdP と連携中か」という基本情報に加え、初回の連携日時や利用環境、IdP から受け取ったメールアドレスなどを表示することで、ユーザーは自身の連携状態をより明確に把握できる。
\item \underline{未連携 IdP との ID 連携}：この管理画面から未連携の IdP と連携設定を行うことで、その IdP をログインに利用可能になる。
\item \underline{既存の ID 連携の解除}：IdP の利用をやめた場合など、ID 連携を解除する必要がある際には、ここから設定を解除できる。
\end{enumerate}
\end{dfn}

(2) の連携設定は、新規登録のプロセス (3)（定義 \ref{dfn:ch4-federated-signup}）と同じく、RP のユーザー $u$ に組 $(i,s)$ を紐付ける操作である。これにより、新規登録時に ID 連携を使わなかったユーザーも、ID 連携によるログイン（定義 \ref{dfn:ch4-federated-login}）を利用できるようになる。

%----------------------------------------------------------------------
\subsubsection{ログインセッションに関する通知の処理}

第3章のセッション管理では、ユーザー自身が有効なログインセッションの一覧を確認し、不要なセッションを遠隔で終了させることができた。ID 連携では、その操作を IdP に伝えることもできる。

ユーザーが RP のセッション管理機能で別環境のログインセッションを破棄した場合でも、IdP がサポートしていれば、IdP への通知が可能である。

%----------------------------------------------------------------------
\subsubsection{ID連携におけるアカウントリカバリー}

第3章のアカウントリカバリー（定義 \ref{dfn:account-recovery}）は、(1) 身元確認、(2) 当人認証のためのクレデンシャル再設定、の二つのステップからなっていた。新規登録と同様に、このうちの身元確認に ID 連携を利用できる。この項では、それに加えて ID 連携がアカウントリカバリーにもたらす効果を述べる。

\begin{prop}[ID連携のフォールバックとしての役割]\label{prop:ch4-federation-fallback}
RP が複数の認証方式をサポートしており、その一つである IdP が幅広い環境からのログインに対応している場合、ID 連携にはフォールバックとしての役割が期待でき、ユーザーが「どうしてもログインできない」という状況に陥る可能性を低減できる。
\end{prop}
\begin{proof}
アカウントリカバリーが必要になる状況から考える。
まず、アカウントリカバリーが必要になるのは、何らかの理由でクレデンシャルが利用できなくなり、利用可能な他の認証方法でもログインできない場合である（定義 \ref{dfn:account-recovery}）。
次に、RP が複数の認証方式をサポートし、その一つとして ID 連携を提供しているならば、他のクレデンシャルが使えなくなっても、ユーザーは IdP によるログインを利用できる。さらに、その IdP が幅広い環境からのログインに対応していれば、どの環境からでもこの代替手段が使える可能性が高い。
よって、ID 連携は他の認証方式が使えないときのフォールバックとして働き、「どうしてもログインできない」という状況に陥る可能性を低減できる。
\end{proof}

さらに、IdP から受け取った属性情報を、本人確認（定義 \ref{dfn:kyc}）やアカウントリカバリー時の身元確認に利用できる。

\begin{exa}[パスキーを用いる RP のアカウントリカバリー]\label{exa:ch4-passkey-recovery}
パスキー（定義 \ref{dfn:passkey}）認証をサポートする RP が、アカウントリカバリー時にはデジタル庁の認証アプリを IdP として利用するなど、低コストで強度の高い身元確認を実現できる。
\end{exa}

第3章のアカウントリカバリーの二つのステップは、ID 連携を用いると次のように変わる。

\underline{(1) IdP から受け取った属性情報を用いた身元確認：}
RP は、IdP から受け取った属性情報と自社で保存している値を比較し、一致することを確認した上で、対象ユーザー自身がアカウントリカバリーを試みていることを確認する。例えば、検証済みのメールアドレスを受け取った場合、RP はそれを信頼し、自社でのメールアドレス検証プロセスを省略できる。また、それ以外の属性情報の組み合わせで、より強度の高い身元確認を実施することもある。

\underline{(2) 当人認証のためのクレデンシャル再設定：}
ID 連携以外の認証方式で利用するクレデンシャルを再設定させる。この時、「デジタル庁の認証アプリを利用した身元確認を行った場合のみパスキーの設定を許可する」というように、身元確認の強度に応じて設定可能な認証方式を制限することもある。また、身元確認を行ったタイミングで、IdP の識別子とユーザー識別子（組 $(i,s)$）を ID 連携によるクレデンシャルとして設定することも可能である。

%----------------------------------------------------------------------
\subsubsection{ID連携における本人確認（eKYC）}

第3章の本人確認（定義 \ref{dfn:kyc}）は、新規登録時以外にも身元確認を行う機能であった。ID 連携では、その身元確認を IdP の結果に依拠させることができる。

IdP が身元確認を行った属性情報を受け取り、RP でも検証済みの情報として扱うことで、本人確認の手間を軽減できる。

\begin{rem}[検証済みとして扱える属性情報と法制度]\label{rem:ch4-ekyc-law}
IdP から受け取った属性情報を検証済みとして扱えるか、また扱える場合はどの属性情報が該当するかなどは、法制度によって定められている場合がある。
\end{rem}

%----------------------------------------------------------------------
\subsubsection{ID連携における無効化と退会}

第3章では、単一サービスがユーザーを一時停止・再有効化する機能と、退会（論理削除・物理削除，定義 \ref{dfn:account-deletion}）の機能を扱った。ID 連携を導入した RP では、RP のユーザーの状態変更に関して、ID 連携の方式によって次のような処理が行われる。

\underline{無効化や退会を IdP に依拠：}
エンタープライズ向けサービスでは、状態変更を IdP に任せる場合がある。RP にて不正な挙動があった際はセキュリティイベントの通知（定義 \ref{dfn:security-event-notification}）などで IdP に通知し、ユーザーの状態変更は IdP からの通知に従うことで、IdP と連動した状態管理が可能である（例 \ref{exa:ch4-rp-initiated-suspension} も参照）。

\underline{退会完了時の連携設定の解除：}
コンシューマー向けサービスのように、ID 連携が任意で利用される場合は、退会完了時に ID 連携の設定を解除するのが一般的である。この場合、ユーザーは再度、新規登録からサービスを利用できる。

\begin{rem}[登録・退会の繰り返しによる不正への対策]\label{rem:ch4-reenrollment-fraud}
退会完了時に ID 連携の設定を解除すると、同じ IdP ユーザーが何度でも新規登録できる。特典目的の登録・退会を繰り返す不正行為への対策として、退会時に過去に連携した IdP のユーザー識別子を記録し、再登録時にチェックするといった方法も考えられる。
\end{rem}

%----------------------------------------------------------------------
\subsubsection{ID連携におけるセキュリティイベントログ}

第3章のセキュリティイベントログ（定義 \ref{dfn:security-event-log}）は、ログイン、設定変更、アカウントリカバリーなどの操作が「いつ、誰によって、何が行われたか」を記録する仕組みであった。ID 連携を利用する RP では、記録すべき内容が増える。

単体サービスのイベントログの内容に加え、ID 連携を利用した処理の際には、どの IdP のどのユーザーによる操作か（組 $(i,s)$）を記録しておくことが重要である。これにより、ある IdP ユーザーに関する操作の追跡が可能になる。

また、IdP との連携が密結合か疎結合かにかかわらず、IdP からのセキュリティイベントの通知（定義 \ref{dfn:security-event-notification}）を適切に受信・処理することで、安全性や利便性、UX の向上が期待できる。

%% =======================================================
%%  part: ch5
%% =======================================================
\section{ID管理機能の設計時に意識したいポイント}

この章では、\cite{Ito2026} の第5章に沿って、ID 管理機能を設計するときに意識したいポイントを扱う。
まず、設計の評価軸となる 4 つの観点（セキュリティ、プライバシー、ユーザビリティ、アクセシビリティ）を定義し、
それらの間に生じる「信頼」と「体験」のトレードオフと、その乗り越え方を論じる。
次に、設計の土台となる基本原則として OECD プライバシー 8 原則とアイデンティティの 7 原則を述べ、
ID 管理に関連するガイドラインとして NIST SP 800-63 とその保証レベル（IAL / AAL / FAL）を紹介する。
さらに、ID 管理機能で利用される標準化仕様（WebAuthn、TOTP、OAuth 2.0、OpenID Connect、SCIM、Shared Signals Framework）を概観し、
最後に、注意したい実装例をいくつか取り上げる。

\subsection{4つの観点}

ID 管理機能（定義 \ref{dfn:identity-management}）を設計するとき、何を基準に「良い設計」を判断すればよいだろうか。
本書は、その評価軸として次の 4 つの観点を挙げる：セキュリティ、プライバシー、ユーザビリティ、アクセシビリティ。
前の二つはユーザーの「信頼」を得るための要件であり、後の二つはユーザーの「体験」を支える要件である。
この項では、それぞれの観点が何を要求するかを定義として固定する。これらの観点の間の関係（トレードオフ）は次の項で扱う。

\subsubsection{セキュリティ}

ID 管理機能はサービスへの入口であり、ここが破られると、サービスが扱う情報が不正に利用されうる。
そこで最初の観点として、セキュリティを置く。

\begin{dfn}[セキュリティの観点]\label{dfn:ch5-security}
ID 管理機能設計における\underline{セキュリティ}（security）の観点とは、
サービスが扱う情報の重要性や、不正利用が発生した際の影響に応じて、ID 管理にも適切なレベルのセキュリティ対策を講じることである。
これには、次のようなものが含まれる。
\begin{enumerate}
\item 多要素認証の導入による不正アクセス防止。
\item パスワードなど機密情報の安全な保管（ハッシュ化・暗号化）。
\item 最小権限の原則に基づいたアクセス制御。
\item 不正ログイン検知の仕組み。
\end{enumerate}
\end{dfn}

ここで (1) の多要素認証は定義 \ref{dfn:mfa}、(3) のアクセス制御は定義 \ref{dfn:access-control} の意味である。
セキュリティ対策の「適切なレベル」は一律ではなく、\emph{情報の重要性と不正利用時の影響}に応じて決まる、という点が定義の要である。
この考え方は、後で述べる NIST SP 800-63 の保証レベル（定義 \ref{dfn:ch5-assurance-level}）にもつながる。

\subsubsection{プライバシー}

ID 管理機能はユーザーの属性情報（定義 \ref{dfn:attribute}）を集め、保管し、利用する。
ユーザーが安心してサービスを利用するには、その扱い方が適切でなければならない。これが第二の観点である。

\begin{dfn}[プライバシーの観点]\label{dfn:ch5-privacy}
ID 管理機能設計における\underline{プライバシー}（privacy）の観点とは、
個人情報保護法などの法令遵守はもちろん、ユーザーが安心してサービスを利用できるよう、プライバシーを保護することである。
具体的には、設計段階からプライバシー保護を組み込む\underline{プライバシー・バイ・デザイン}（privacy by design）の考え方に基づき、
\begin{enumerate}
\item 収集する属性情報を必要最小限に留めること、
\item 利用目的を明確にすること、
\item ユーザー自身が情報をコントロールできる仕組みを提供すること
\end{enumerate}
が求められる。
\end{dfn}

プライバシーの観点は、法令を守るだけでは\emph{足りない}ことに注意する。法令遵守は前提であり、その上でユーザーの安心を設計段階から組み込むことが求められている。
(1)--(3) の考え方は、後の項で述べる OECD プライバシー 8 原則（定義 \ref{dfn:oecd-principles}）やアイデンティティの 7 原則（定義 \ref{dfn:laws-of-identity}）とも対応している。

\subsubsection{ユーザビリティ}

セキュリティとプライバシーを確保しても、使いにくければユーザーは離れていく。そこで第三の観点として、ユーザビリティを置く。

\begin{dfn}[ユーザビリティの観点]\label{dfn:ch5-usability}
ID 管理機能設計における\underline{ユーザビリティ}（usability）の観点とは、
セキュリティやプライバシーを確保した上で、ユーザーがストレスなく使える体験を提供することである。
煩雑な登録・ログインプロセスはユーザー離脱の原因になるため、
パスキー認証やソーシャルログインの導入で入力を簡略化したり、エラーメッセージを分かりやすくしたりと、
ユーザーが直感的に操作できる設計を心がける必要がある。
\end{dfn}

パスキー認証は定義 \ref{dfn:passkey}、ソーシャルログインは ID 連携（定義 \ref{dfn:identity-federation}）を用いたログインである。
定義の冒頭の「セキュリティやプライバシーを確保した\emph{上で}」という限定が重要であり、ユーザビリティのために前二者を犠牲にしてよいという意味ではない。

\subsubsection{アクセシビリティ}

ユーザビリティが「使いやすさ」の観点であるのに対し、アクセシビリティは「誰もが使えること」の観点である。

\begin{dfn}[アクセシビリティの観点]\label{dfn:ch5-accessibility}
ID 管理機能設計における\underline{アクセシビリティ}（accessibility）の観点とは、
ID 管理機能を、年齢や身体的な制約、利用環境に関わらず、サービスのユーザーであれば誰もが利用できるようにすることである。
具体的には、スクリーンリーダーへの対応や、キーボードのみで操作できる設計、十分な文字サイズや色のコントラストを確保するといった、
ウェブアクセシビリティの標準（例：WCAG）に準拠することが求められる。
\end{dfn}

ID 管理機能はサービスのすべての利用に先立つ入口であるから、ここが利用できないユーザーはサービス全体を利用できない。
その意味で、アクセシビリティの要求は「サービスのユーザーであれば\emph{誰もが}」という全称的な形で述べられている。

\subsection{「信頼」と「体験」のトレードオフを乗り越える}

前の項で定義した 4 つの観点は、どれも単独では当然の要求に見える。しかし、実際の設計ではこれらが互いに衝突する。
この項では、まずその衝突（トレードオフ）がなぜ生じるかを述べ、次に二つの典型的な状況、
すなわち画一的な高セキュリティ要件がユーザビリティを下げる状況と、セキュリティ強化策がかえって新たなリスクを生む状況とを取り上げ、
それぞれの解決のアプローチを論じる。

ID 管理機能を設計する上で、4 つの観点のバランスを考慮することは不可欠である。
ユーザーの\underline{「信頼」}を得るためのセキュリティとプライバシー（定義 \ref{dfn:ch5-security}, 定義 \ref{dfn:ch5-privacy}）の要件は、
ID が利用されるサービス（例：金融、行政、SNS）の性質によって決まる。
その上で\underline{「体験」}を支えるユーザビリティやアクセシビリティ（定義 \ref{dfn:ch5-usability}, 定義 \ref{dfn:ch5-accessibility}）を高めるのが理想である。

\begin{prop}[信頼と体験のトレードオフ]\label{prop:ch5-trust-experience-tradeoff}
「信頼」を支えるセキュリティ・プライバシーと、「体験」を支えるユーザビリティ・アクセシビリティとを高いレベルで両立させることは難しく、
多くの場合トレードオフが発生する。
\end{prop}

\begin{proof}
信頼の側の要件はサービスの性質から外的に決まるが、体験の側はそれを所与として高めるしかない、という構造から説明する。
まず、セキュリティとプライバシーの要件は、ID が利用されるサービス（金融、行政、SNS など）の性質によって決まり、設計者が体験のために自由に緩めることはできない。
次に、ユーザビリティやアクセシビリティは、その要件を満たした\emph{上で}高めるものである（定義 \ref{dfn:ch5-usability}）。
ところが、高い信頼の要件を満たすための対策（設定項目や確認手順、追加の認証など）は、ユーザーの操作や負担を増やす。
よって、信頼の要件が高いほど体験を高める余地は狭まり、両者を同時に高いレベルにすることは難しい。
具体的な現れ方は、以下の命題 \ref{prop:ch5-uniform-security-usability} と命題 \ref{prop:ch5-security-new-risk} で見る。
\end{proof}

\subsubsection{画一的な高セキュリティ要件によるユーザビリティの低下}

トレードオフの第一の典型は、もっとも厳しいサービスに合わせたセキュリティ要件を全員に課してしまう場合である。
ここでは、SNS から金融系まで、幅広いサービスで利用される ID 管理機能を想定する。

\begin{exa}[SNS から金融系まで共通の ID 管理機能]\label{exa:ch5-sns-to-finance}
一つの ID 管理機能が、SNS のような気軽なサービスと、金融サービスのようなセキュリティ要件が高いサービスの両方で利用されているとする。
このとき、金融サービスレベルの高いセキュリティ要件をすべてのユーザーに画一的に適用すると、新規登録（定義 \ref{dfn:signup}）時に設定すべき項目が多くなる。
\end{exa}

\begin{prop}[画一的な高セキュリティ要件はユーザビリティを下げる]\label{prop:ch5-uniform-security-usability}
例 \ref{exa:ch5-sns-to-finance} の状況で、金融サービスレベルの高いセキュリティ要件をすべてのユーザーに画一的に適用すると、
SNS のような気軽なサービスから使い始めたいユーザーが新規登録を断念してしまう。
すなわち、\underline{セキュリティとユーザビリティのトレードオフ}が発生する。
\end{prop}

\begin{proof}
要件の画一的な適用が、登録時の負担を経由して離脱につながることを示す。
まず、金融サービスレベルの要件を全員に適用すると、新規登録時に設定すべき項目が多くなり、ユーザビリティは低下する。
次に、定義 \ref{dfn:ch5-usability} で述べたとおり、煩雑な登録プロセスはユーザー離脱の原因になる。
とくに、SNS のような気軽なサービスだけを使いたいユーザーにとって、金融サービス向けの負担は本来不要なものである。
したがって、こうしたユーザーは新規登録を断念してしまう。セキュリティを高めた結果ユーザビリティが失われているので、これはセキュリティとユーザビリティのトレードオフである。
\end{proof}

この課題は、ユーザーの行動に応じてセキュリティレベルを段階的に引き上げるアプローチで解決できる。

\begin{dfn}[段階的なセキュリティ強化]\label{dfn:ch5-step-up-security}
\underline{段階的なセキュリティ強化}とは、ユーザーの行動に応じてセキュリティレベルを段階的に引き上げるアプローチである。
すなわち、
\begin{enumerate}
\item 新規登録時は、メールアドレスの確認やソーシャルログインなど、可能な限りシンプルなユーザー体験を意識する。
\item その後、ユーザーが金融サービスのようなセキュリティ要件が高い機能を利用する際に、
「安全のため、電話番号による追加認証を設定してください」といった形で、セキュリティ強化を要求する。
\end{enumerate}
\end{dfn}

\begin{prop}[段階的なセキュリティ強化によるトレードオフの解消]\label{prop:ch5-step-up-resolves}
段階的なセキュリティ強化（定義 \ref{dfn:ch5-step-up-security}）により、すべてのユーザーに最初から高い負担を強いることなく、
必要な場面で必要なセキュリティを確保し、ユーザビリティとの両立を図ることができる。
\end{prop}

\begin{proof}
登録時と高セキュリティ機能の利用時とに分けて考える。
まず、新規登録時には定義 \ref{dfn:ch5-step-up-security} (1) によりシンプルな体験だけを求めるので、
命題 \ref{prop:ch5-uniform-security-usability} の証明で見た「登録時の設定項目の多さ」による離脱は起きない。
次に、セキュリティ要件が高い機能を利用するユーザーには、その時点で (2) により追加認証の設定を要求するので、その機能に必要なセキュリティは確保される。
SNS だけを使うユーザーは (2) の負担を負わず、金融サービスを使うユーザーだけが必要な負担を負う。
以上より、負担は必要な場面に限られ、セキュリティとユーザビリティの両立が図られる。
\end{proof}

\subsubsection{セキュリティ強化が引き起こす新たなリスク}

トレードオフの第二の典型は、より見えにくい。セキュリティを強化するつもりの対策が、ユーザーの行動を変えることで、かえってセキュリティを損なう場合である。
当人認証（定義 \ref{dfn:authentication}）を用いたセキュリティの強化策として、安易にパスワード（定義 \ref{dfn:password-auth}）や
MFA の OTP（定義 \ref{dfn:otp-auth}）による追加認証の機会を増やすことにも注意が必要である。

\begin{prop}[認証要求の増加が中継型フィッシングのリスクを高める]\label{prop:ch5-security-new-risk}
パスワードや OTP のように、ユーザーが認証情報を手動で入力する認証方式において、追加認証の機会を安易に増やすと、
中継型のフィッシング攻撃の被害に遭うリスクが高まる。
\end{prop}

\begin{proof}
認証情報の入力が習慣化することで、本来の前提が崩れることを示す。
まず、パスワードや OTP による認証は、本来「正規のサービスであることを確認した上で認証情報を入力する」ことが求められるものである。
次に、ユーザーが頻繁にこれらの認証情報を求められると、入力は次第に「何も考えずに行う作業」になってしまい、正規のサービスであることの確認が行われなくなる。
その結果、攻撃者が用意した偽サイトへの誘導があった際にも、ユーザーは正規サービスの認証情報をためらいなく入力してしまうこととなり、中継型のフィッシング攻撃の被害に遭うリスクが高まる。
\end{proof}

\begin{rem}[認証方式と強化策の組み合わせによる危険]\label{rem:ch5-manual-input-risk}
命題 \ref{prop:ch5-security-new-risk} の事態は、「ユーザーが認証情報を手動で入力する」という認証方式の特徴と、
セキュリティ強化策（認証機会の増加）の組み合わせが、ユーザーの警戒心を解き、かえって危険に晒すというものである。
どちらか一方だけを見ていては気づきにくい点に注意する。
\end{rem}

この事態を回避するためには、「単純に当人認証を要求する機会を増やす」ことが本当にセキュリティ強化につながるかを、より具体的に検討することが必要である。
そのために、当人認証を要求するにあたって次の概念を導入する。

\begin{dfn}[リスクベース判定]\label{dfn:ch5-risk-based}
\underline{リスクベース判定}（risk-based assessment）とは、当人認証を要求するかどうかを、
ログインする場所、時間、デバイスなどの状況（\underline{コンテキスト}, context）から判断したリスクに基づいて決めることである。
典型的には、普段と違う環境からのアクセス時のみ追加認証を要求する。
\end{dfn}

\begin{prop}[リスクベース判定の効果]\label{prop:ch5-risk-based-effect}
リスクベース判定（定義 \ref{dfn:ch5-risk-based}）を導入し、普段と違う環境からのアクセス時のみ追加認証を要求することで、
ユーザーの負担を減らし、認証要求の重要性を保つことができる。
\end{prop}

\begin{proof}
認証要求の頻度と、その要求が持つ意味の二点から示す。
まず、普段どおりの環境からのアクセスでは追加認証を要求しないので、認証を求められる回数が減り、ユーザーの負担は減る。
次に、追加認証が求められるのは普段と違う環境からのアクセスという例外的な場面に限られるので、
認証要求は命題 \ref{prop:ch5-security-new-risk} の証明で述べた「何も考えずに行う作業」にはなりにくく、ユーザーにとって注意を払うべき出来事であり続ける。
よって、認証要求の重要性が保たれる。
\end{proof}

\subsection{ID管理における基本原則}

ID 管理機能を設計する上で、サービス毎の要件に応じて前述の 4 つの観点（定義 \ref{dfn:ch5-security}--定義 \ref{dfn:ch5-accessibility}）を組み合わせることが重要である。
そして、その組み合わせ方を決める土台となる考え方や哲学が\underline{基本原則}である。
この項では、代表的な基本原則として、ユーザー情報をどう扱うべきかの国際的な基準である「OECD プライバシー 8 原則」と、
ユーザー自身が情報をコントロールできるようにするための「アイデンティティの 7 原則」を取り上げる。

\subsubsection{OECDプライバシー8原則}

最初に、成立の背景を述べる。
\underline{OECD 8 原則}は、経済協力開発機構（OECD）が 1980 年に採択した「プライバシー保護と個人データの国際流通に関するガイドライン」の中で示された、
プライバシー保護に関する 8 つの基本的な考え方である。
コンピュータと通信技術の発展により、国境を越えて個人データがやり取りされるようになったことを背景に、
個人のプライバシーを保護しつつ、自由な情報流通も妨げないための国際的な共通ルールとして策定された。
この原則は、特定の法律や規制\emph{そのものではなく}、各国の政府・立法機関、企業や組織から一般の個人向けまで、様々なサービスに活用されている。

OECD 8 原則は、以下の 8 つの項目から成り立っている。

\begin{dfn}[OECD プライバシー 8 原則]\label{dfn:oecd-principles}
\underline{OECD プライバシー 8 原則}とは、個人データの取り扱いに関する次の 8 つの原則である。
\begin{enumerate}
\item \underline{収集制限の原則}（Collection Limitation Principle）：
個人データの収集は、適法かつ公正な手段で行うべきである。また、原則として本人に通知するか、同意を得た上で収集しなければならない。
\item \underline{データ内容の原則}（Data Quality Principle）：
収集した個人データは、利用目的に必要な範囲内で、正確、完全、かつ最新の状態に保つべきである。
\item \underline{目的明確化の原則}（Purpose Specification Principle）：
個人データを何のために集めるのか、その利用目的を収集する時点で明確にしなければならない。
\item \underline{利用制限の原則}（Use Limitation Principle）：
本人の同意がある場合や法律で定められている場合を除き、明確にした目的以外のために個人データを開示、提供、利用してはならない。
\item \underline{安全保護の原則}（Security Safeguards Principle）：
個人データが漏えい、紛失、破壊、改ざんされたり、不正に利用されたりしないよう、合理的な安全対策を講じなければならない。
\item \underline{公開の原則}（Openness Principle）：
個人データの取り扱いに関する方針や実務（どんなデータを、何のために、誰が管理しているかなど）は公開し、誰でも容易に知ることができるようにすべきである。
\item \underline{個人参加の原則}（Individual Participation Principle）：
データを提供した本人は、自分に関するデータがどこに存在し、どのように利用されているかを確認できる。また、その内容に誤りがあれば、訂正や削除を求める権利が保障されるべきである。
\item \underline{責任の原則}（Accountability Principle）：
個人データを取り扱う管理者（データ管理者）は、上記 (1)--(7) の原則を遵守する責任を負わなければならない。
\end{enumerate}
\end{dfn}

各原則は ID 管理機能の設計にも直接対応する。本書が挙げている対応を、原則の番号ごとにまとめておく。

\begin{rem}[OECD 8 原則と ID 管理機能]\label{rem:ch5-oecd-in-idm}
定義 \ref{dfn:oecd-principles} の各原則について、ID 管理における意味は次のとおりである。
\underline{データ内容の原則（(2)）：}ID 管理における例としては、メールアドレスや電話番号が最新のものかどうかの確認を求められる場合がある。

\underline{安全保護の原則（(5)）：}ID 管理機能で扱う属性情報に、漏えい、紛失、破壊、改ざん、不正利用といった事態が発生すると、
ユーザー、他サービス含めて大きな影響が出ることが考えられる。

\underline{公開の原則（(6)）：}ID 管理機能に限らず、様々なデータを適切に扱っているかどうかがサービス全体の信頼に繋がる。

\underline{個人参加の原則（(7)）：}ID 管理機能の設定変更、退会（定義 \ref{dfn:account-deletion}）時のデータの扱いで、これを意識する必要がある。
\end{rem}

プライバシーの観点（定義 \ref{dfn:ch5-privacy}）の「収集する属性情報を必要最小限に留める」「利用目的を明確にする」「ユーザー自身が情報をコントロールできる」は、
それぞれ (1)(2)、(3)(4)、(7) の原則と同じ方向を向いている。

\subsubsection{アイデンティティの7原則}

次に、もう一つの基本原則である \underline{アイデンティティの 7 原則}（The 7 Laws of Identity）について、まず背景を述べる。
アイデンティティの 7 原則は、マイクロソフトのチーフアーキテクトであったキム・キャメロン氏が 2005 年に発表した、デジタルアイデンティティ（定義 \ref{dfn:digital-identity}）に関する設計思想である。
これは国際機関が採択したもの\emph{ではなく}、デジタル世界における個人のアイデンティティはどうあるべきか、またそれを実現するシステムはどう設計されるべきかを示した、
影響力の大きい技術的なビジョン・提言として知られている。

この原則は、テクノロジーではなく\emph{ユーザーの視点}からデジタルアイデンティティがどうあるべきかを定義しており、サービスの企画・開発におけるガイドラインとして広く活用されている。
また、OpenID Foundation のような団体による標準化仕様の策定においても、この原則の思想が影響を与えている。

\begin{dfn}[アイデンティティの 7 原則]\label{dfn:laws-of-identity}
\underline{アイデンティティの 7 原則}とは、デジタルアイデンティティに関する次の 7 つの原則である。
\begin{enumerate}
\item \underline{ユーザーによる制御と同意}（User Control and Consent）：
デジタルアイデンティティはユーザー自身が管理・制御するものであるべきである。個人に関する情報が本人の明確な同意なしに使用されたり、共有されたりしてはならない。
\item \underline{限定的な利用のための最小限の開示}（Minimal Disclosure for a Constrained Use）：
あるサービスを利用する際、相手に開示する個人情報は、その目的に必要な最小限であるべきである。
\item \underline{正当な関係者}（Justifiable Parties）：
個人情報を受け取る相手（サービス提供者など）は、その情報を受け取るに値する正当な関係者でなければならない。ユーザーは、誰に、なぜ自分の情報が渡されるのかを理解できる必要がある。
\item \underline{指向性アイデンティティ}（Directed Identity）：
デジタルアイデンティティは、特定の相手との関係性において使われる一対一（指向性）のものであるべきである。
\item \underline{事業者と技術の多元性}（Pluralism of Operators and Technologies）：
デジタルアイデンティティの仕組みは、特定の企業や特定の技術に縛られるものであってはならない。様々な技術やサービスが相互に連携できるオープンなエコシステムを目指すべきである。
\item \underline{人間との統合}（Human Integration）：
システムは、人間にとって直感的で理解しやすく、使いやすいものでなければならない。複雑な技術を意識することなく、ユーザーが自然にアイデンティティを管理できるインターフェースが求められる。
\item \underline{コンテキストを越えた一貫した経験}（Consistent Experience Across Contexts）：
ユーザーが様々なサービスやシステムを利用する際に、その体験（どう情報を開示し、どう認証するかなど）に一貫性があるべきである。
\end{enumerate}
\end{dfn}

定義 \ref{dfn:laws-of-identity} のいくつかの原則には、本書による具体例や趣旨の説明が添えられている。これを以下にまとめる。

\begin{exa}[年齢確認における最小限の開示]\label{exa:ch5-age-minimal-disclosure}
定義 \ref{dfn:laws-of-identity} (2)（限定的な利用のための最小限の開示）の例として、年齢確認を考える。
年齢確認のために生年月日そのものを渡すのではなく、「20 歳以上である」という事実だけを伝えるべき、というのがこの原則の考え方である。
\end{exa}

属性情報の集合としてのアイデンティティ（定義 \ref{dfn:identity}）の言葉で言えば、相手に渡す属性情報を、目的の達成に必要な部分に限るということである。
生年月日は「20 歳以上か」より多くの情報を含むので、年齢確認という目的に対しては過剰な開示になる。

\begin{rem}[指向性アイデンティティの趣旨]\label{rem:ch5-directed-identity}
定義 \ref{dfn:laws-of-identity} (4)（指向性アイデンティティ）は、あるサービスで利用した識別情報（ID）が、本人の知らないところで
別のサービスでの活動と勝手に紐づけられる（\underline{名寄せ}される）ことを防ぐ思想である。
\end{rem}

\begin{rem}[一貫した経験の効果]\label{rem:ch5-consistent-experience}
定義 \ref{dfn:laws-of-identity} (7)（コンテキストを越えた一貫した経験）が満たされると、ユーザーは混乱することなく、安心してシステムを利用できる。
\end{rem}

(1)--(3) はプライバシーの観点（定義 \ref{dfn:ch5-privacy}）と、(6)(7) はユーザビリティの観点（定義 \ref{dfn:ch5-usability}）と、それぞれ重なる内容をユーザーの視点から述べたものと読むことができる。

\subsection{ID管理に関連するガイドライン}

基本原則は設計の哲学を与えるが、具体的に「どの程度の対策を講じればよいか」までは決めない。
そこで、より具体的な技術的要件を与えるガイドラインが参照される。
この項では、ID 管理に関連するガイドラインの代表として NIST SP 800-63 を取り上げ、その構成と、身元確認・当人認証・ID 連携のそれぞれについての保証レベル（IAL / AAL / FAL）を述べる。

\subsubsection{NIST SP 800-63 Digital Identity Guidelines}

\underline{NIST SP 800-63}（Digital Identity Guidelines）は、米国国立標準技術研究所（NIST）が発行する、デジタルアイデンティティに関するガイドラインである。
元々は米国政府機関向けに策定されたが、その網羅性と信頼性の高さから、現在では民間企業も含め、
世界中の ID 管理システムを設計する際の事実上のグローバルスタンダードとして広く参照されている。

このガイドラインは、パスキー認証（定義 \ref{dfn:passkey}）の登場など、世の中の動向に合わせて継続的に更新されている。
最新版は 2025 年 7 月に発行された第 4 版であり、本書の内容も NIST SP 800-63 の内容を強く意識している。

\subsubsection{NIST SP 800-63の概要と構成}

このガイドラインの主な目的は、デジタルアイデンティティのリスクを効果的に管理するための技術的な要件と推奨事項を定義することである。
厳格なコンプライアンスモデルではなく、各組織が自身のリスク評価に基づいてセキュリティ対策を調整する\underline{リスクベースのアプローチ}を推進している。

NIST SP 800-63 は、次の 4 つのドキュメントから構成されている。

\begin{table}[H]
\centering
\caption{NIST SP 800-63 を構成するドキュメント}\label{tab:ch5-sp800-63-documents}
\begin{tabular}{ll}
\hline
ドキュメント & 内容 \\
\hline
SP 800-63 & 全体のフレームワークと概要 \\
SP 800-63A & 「身元確認」に関する詳細な要件 \\
SP 800-63B & 「当人認証」に関する詳細な要件 \\
SP 800-63C & 「ID 連携」に関する詳細な要件 \\
\hline
\end{tabular}
\end{table}

NIST SP 800-63 は、デジタルアイデンティティのライフサイクルを「身元確認」「当人認証」「ID 連携」の 3 つの領域
（第1章の定義 \ref{dfn:identity-proofing}, 定義 \ref{dfn:authentication}, 定義 \ref{dfn:identity-federation}）に分け、
それぞれの信頼性の度合いを保証レベルで定義している。

\begin{dfn}[保証レベル]\label{dfn:ch5-assurance-level}
NIST SP 800-63 において、身元確認・当人認証・ID 連携の各領域の信頼性の度合いを表す段階を\underline{保証レベル}（Assurance Level）という。
身元確認については IAL（定義 \ref{dfn:ial}）、当人認証については AAL（定義 \ref{dfn:aal}）、ID 連携については FAL（定義 \ref{dfn:fal}）が定められている。
\end{dfn}

\begin{prop}[保証レベルによる過不足のない対策]\label{prop:ch5-assurance-level-proportionate}
保証レベルを用いることで、システムで扱う情報の重要性やリスクに応じた、過不足のない適切なセキュリティ対策を講じることが可能になる。
\end{prop}

\begin{proof}
保証レベルが信頼性の度合いを段階化していることから従う。
まず、身元確認・当人認証・ID 連携のそれぞれについて、信頼性の度合いが段階（保証レベル）として定義されている。
次に、システムの側では、扱う情報の重要性やリスクを評価し、それに見合う段階を選ぶことができる。
低い段階を選べば不必要に重い対策を避けられ、高い段階を選べばリスクに見合う強い対策が要求される。
よって、一律の対策ではなく、情報の重要性やリスクに応じて過不足のない対策を講じることができる。
これはセキュリティの観点（定義 \ref{dfn:ch5-security}）の「適切なレベル」を具体化するものであり、リスクベースのアプローチの現れでもある。
\end{proof}

\subsubsection{SP 800-63A：身元確認とアイデンティティの登録}

SP 800-63A では、様々な保証レベルでユーザーのデジタルアイデンティティを確立・管理する方法を規定している。
第1章で概要を定義した身元確認（定義 \ref{dfn:identity-proofing}）を、SP 800-63A は次の 3 つのステップに分けて精密化する。

\begin{dfn}[SP 800-63A における身元確認の 3 ステップ]\label{dfn:ch5-proofing-steps}
SP 800-63A において、身元確認は以下の 3 つのステップで構成される。
\begin{enumerate}
\item \underline{アイデンティティ解決}（Identity Resolution）：申請者が実在する唯一の個人であることを確認するため、証拠を収集するプロセス。
\item \underline{アイデンティティ検証}（Identity Validation）：収集した証拠が本物であり、正確であることを確認するプロセス。
\item \underline{ID 検証}（Identity Verification）：申請者が、提示された証拠の正当な所有者本人であることを確認するプロセス。
\end{enumerate}
\end{dfn}

(1) は「誰についての話か」を一意に定め、(2) は「証拠そのもの」の真正性を、(3) は「証拠と申請者の結びつき」を確かめる。
証拠が本物であっても、提示しているのがその所有者本人でなければ身元確認は成り立たないので、(2) と (3) は別のステップとして必要になる。

さらに SP 800-63A は、身元確認の信頼性を次の保証レベルで規定している。

\begin{dfn}[身元確認の保証レベル（IAL）]\label{dfn:ial}
\underline{身元確認の保証レベル}（Identity Assurance Levels : IALs）は、次の 3 段階である。
\begin{enumerate}
\item \underline{IAL1}（基本レベル）：主張されたアイデンティティが実在し、申請者がそれに関連していることを保証するレベル。大規模な自動攻撃を防ぐ。
\item \underline{IAL2}（高レベル）：より強力な証拠と厳格な検証プロセスを要求し、標的型の攻撃から保護する。
\item \underline{IAL3}（最高レベル）：訓練を受けた担当者との直接対話と生体情報の収集を義務付け、最も高度な攻撃から保護する。
\end{enumerate}
\end{dfn}

各 IAL には、使用できる証拠の種類や、検証・本人確認の方法に関する具体的な要件が詳述されている。

\subsubsection{SP 800-63B：当人認証と認証器の管理}

SP 800-63B では、当人認証（定義 \ref{dfn:authentication}）のプロセスと、認証器のライフサイクル管理に関する要件を定めている
（認証器の登録から回復までの局面については、第2章の定義 \ref{dfn:authenticator-phases} を参照）。

\begin{dfn}[認証器]\label{dfn:ch5-authenticator}
\underline{認証器}（Authenticator）とは、パスワードやセキュリティキーなど、ユーザーが本人であることを証明する手段を指す。
\end{dfn}

認証器の信頼性は、次の保証レベルで評価される。

\begin{dfn}[当人認証の保証レベル（AAL）]\label{dfn:aal}
\underline{当人認証の保証レベル}（Authenticator Assurance Levels : AALs）は、次の 3 段階である。
\begin{enumerate}
\item \underline{AAL1}（基本レベル）：単一の認証要素で許可される。リプレイ攻撃への耐性が必要。
\item \underline{AAL2}（高レベル）：2 つの異なる認証要素（多要素認証）を要求する。フィッシング耐性のある認証手段の提供が必須。
\item \underline{AAL3}（最高レベル）：AAL2 に加え、フィッシング耐性が必須となる。さらに、認証器が持つ秘密鍵はハードウェアで保護されている必要がある。
\end{enumerate}
\end{dfn}

AAL2 の「2 つの異なる認証要素」は、第2章の認証の 3 要素（定義 \ref{dfn:authentication-factors}）のうち異なるものを 2 つ組み合わせること、すなわち多要素認証（定義 \ref{dfn:mfa}）である。
AAL2 ではフィッシング耐性のある認証手段を\emph{提供する}ことが必須であるのに対し、AAL3 ではフィッシング耐性そのものが\emph{必須}となる、という違いに注意する。
各 AAL には、認証プロセスの各ステップで満たすべき具体的な要件が詳述されている。

\subsubsection{SP 800-63C：ID連携とその中でやり取りされる情報}

SP 800-63C は、ID 連携（Federation、定義 \ref{dfn:identity-federation}）、いわゆる\underline{シングルサインオン}（SSO）に関する技術要件を定めている。
ID 連携におけるトランザクションのセキュリティは、次の保証レベルで分類される。

\begin{dfn}[ID 連携の保証レベル（FAL）]\label{dfn:fal}
\underline{ID 連携の保証レベル}（Federation Assurance Levels : FALs）は、次の 3 段階である。
\begin{enumerate}
\item \underline{FAL1}（基本レベル）：署名付きアサーションを利用する基本的なレベル。
\item \underline{FAL2}（高レベル）：各アサーションは単一の RP（Relying Party、定義 \ref{dfn:rp}）向けとし、不正利用からの保護が必須。
\item \underline{FAL3}（最高レベル）：FAL2 に加え、「Holder-of-Key」アサーション等を必須とし、トランザクション中にユーザーが特定の鍵を所持していることを保証する。
\end{enumerate}
\end{dfn}

SP 800-63C には次のような特徴がある。

\begin{enumerate}
\item \underline{信頼関係の確立と鍵管理：}IdP（定義 \ref{dfn:idp}）と RP の間の信頼関係を確立し、暗号鍵を管理するプロセス（Discovery and Registration）について、手動または動的な方法が FAL に応じて規定されている。
\item \underline{プライバシーへの配慮：}ユーザー情報の伝達を最小限に抑えることや、データの二次利用に対する同意取得の重要性が強調されている。
\item \underline{セキュリティへの配慮：}アサーションインジェクション攻撃からの保護や、設定情報の安全な伝達など、セキュリティに関する考慮事項が全体を通して詳述されている。
\end{enumerate}

(2) は OECD 8 原則の収集制限・利用制限（定義 \ref{dfn:oecd-principles} (1)(4)）や、アイデンティティの 7 原則の最小限の開示（定義 \ref{dfn:laws-of-identity} (2)）と同じ方向の要求である。

以上の 3 種類の保証レベルを、本書の記述の範囲で表にまとめる。

\begin{table}[H]
\centering
\caption{NIST SP 800-63 の保証レベル（IAL / AAL / FAL）}\label{tab:ch5-assurance-levels}
\small
\begin{tabular}{|p{0.13\textwidth}|p{0.25\textwidth}|p{0.25\textwidth}|p{0.25\textwidth}|}
\hline
 & IAL（身元確認, SP 800-63A） & AAL（当人認証, SP 800-63B） & FAL（ID 連携, SP 800-63C） \\
\hline
レベル 1（基本レベル） & 主張されたアイデンティティが実在し、申請者がそれに関連していることを保証。大規模な自動攻撃を防ぐ。 & 単一の認証要素で許可。リプレイ攻撃への耐性が必要。 & 署名付きアサーションを利用。 \\
\hline
レベル 2（高レベル） & より強力な証拠と厳格な検証プロセスを要求。標的型の攻撃から保護。 & 2 つの異なる認証要素（多要素認証）を要求。フィッシング耐性のある認証手段の提供が必須。 & 各アサーションは単一の RP 向け。不正利用からの保護が必須。 \\
\hline
レベル 3（最高レベル） & 訓練を受けた担当者との直接対話と生体情報の収集を義務付け。最も高度な攻撃から保護。 & AAL2 に加え、フィッシング耐性が必須。認証器の秘密鍵はハードウェアで保護。 & FAL2 に加え、「Holder-of-Key」アサーション等を必須とし、トランザクション中のユーザーによる特定の鍵の所持を保証。 \\
\hline
\end{tabular}
\end{table}

\subsection{ID管理と標準化仕様}

ガイドラインが「どの程度の対策が必要か」を与えるのに対し、標準化仕様は「それをどう実現するか」の具体的な設計図を与える。
この項では、まず標準化仕様をサービス独自の仕様より優先すべき理由を述べ、
次に ID 管理機能の中で利用されている代表的な標準化仕様について、その概要、ID 管理機能への適用範囲、策定団体を順に述べる。

\subsubsection{標準化仕様の優位性}

\begin{dfn}[標準化仕様]\label{dfn:ch5-standard-spec}
IETF や W3C などの団体が策定した仕様を\underline{標準化仕様}という。
\end{dfn}

\begin{prop}[標準化仕様の優先]\label{prop:ch5-prefer-standards}
ID 管理機能を構築する際、標準化仕様がある場合は、サービス独自の仕様よりも優先して採用すべきである。
\end{prop}

\begin{proof}
標準化仕様が安全性の面と相互運用性の面の両方で利点を持つことを示す。
まず、標準化仕様は「有識者によるレビュー済みの安全な設計図」である。
したがって開発者は、安全性や利便性の検証といった「\underline{車輪の再発明}」を避け、実績のある堅牢な機能を実装できる。
独自の仕様では、この検証を自ら行わなければならない。
次に、とくに ID 連携（定義 \ref{dfn:identity-federation}）では、関係者全員が標準に対応したライブラリや製品を利用することで、迅速かつ安全なシステム連携が実現する。
すなわち、標準化仕様の採用により相互運用性の確保と実装コストの抑制という大きな利点が得られる。
以上より、標準化仕様がある場合は独自仕様よりも優先して採用すべきである。
\end{proof}

ID 管理機能の中では、いくつかの標準化仕様が利用されている。
各仕様の概要、ID 管理機能への適用範囲、そして策定団体などを認識しておくことは、システムの設計や実装において重要である。
以下ではこれらを順に述べ、最後に表 \ref{tab:ch5-standards} にまとめる。

\subsubsection{Web Authentication (WebAuthn): An API for accessing Public Key Credentials}

パスキー認証（定義 \ref{dfn:passkey}）や FIDO 認証（定義 \ref{dfn:fido-auth}）を Web で実現するための仕様が WebAuthn である。

\begin{dfn}[WebAuthn]\label{dfn:ch5-webauthn}
\underline{Web Authentication}（\underline{WebAuthn}）: An API for accessing Public Key Credentials は、
W3C（World Wide Web Consortium）と FIDO Alliance が FIDO2 プロジェクトの一部として共同で策定した、パスキー認証を実現するための Web ブラウザ標準 API であり、
内部で管理されるデータモデルなども定義されている。
WebAuthn は、公開鍵暗号方式を利用して安全な認証を実現する、Web ブラウザ向けの API 仕様である。
\end{dfn}

WebAuthn の利点は、利便性と安全性の両面にある。

\begin{prop}[WebAuthn の利便性と安全性]\label{prop:ch5-webauthn-merits}
WebAuthn による認証には、次の利点がある。
\begin{enumerate}
\item 利便性が高まる。
\item サービスが侵害されて漏洩した認証情報を悪用するリスクを根本から排除する。
\item フィッシングのリスクを大幅に軽減できる。
\end{enumerate}
\end{prop}

\begin{proof}
各利点を、WebAuthn の仕組みから順に示す。

\underline{(1) について：}ユーザーの操作は、デバイスの画面ロック解除（生体認証や PIN 入力など）や PIN 入力で済むため、利便性が高まる。

\underline{(2) について：}パスキー認証を要求するサービスには公開鍵のみが保存され、秘密鍵はユーザーのデバイスおよびパスワードマネージャーに保存される。
したがって、サービスが侵害されて保存データが漏洩しても、そこに含まれるのは公開鍵だけであり、漏洩した認証情報を悪用するリスクは根本から排除される。

\underline{(3) について：}クレデンシャルがサービスのドメインに紐づくこと、およびブラウザやプラットフォームの機能が認証を仲介することにより、フィッシングのリスクを大幅に軽減できる。
\end{proof}

ID 管理機能においては、WebAuthn は、フィッシング耐性が求められる高セキュリティな当人認証を利用する機能（ログイン（定義 \ref{dfn:login}）、再認証（定義 \ref{dfn:reauthentication}））で利用される。
命題 \ref{prop:ch5-webauthn-merits} (3) により、WebAuthn は命題 \ref{prop:ch5-security-new-risk} で述べた中継型フィッシングのリスクへの対策にもなる。

\subsubsection{RFC 6238 TOTP: Time-Based One-Time Password Algorithm}

第2章で扱った TOTP 認証（定義 \ref{dfn:totp-auth}）の基礎となる仕様が RFC 6238 である。

\begin{dfn}[RFC 6238 TOTP]\label{dfn:ch5-rfc6238}
\underline{TOTP}（Time-Based One-Time Password Algorithm）は、IETF（Internet Engineering Task Force）により RFC 6238 として標準化されている。
TOTP は、時刻を基に、一定時間（通常 30 秒）ごとに変化する使い捨てのパスワード（ワンタイムパスワード）を生成するアルゴリズムである。
Google Authenticator などの TOTP 対応アプリと認証サービスが共有した秘密鍵（\underline{シード}）を使い、それぞれが時刻から生成したパスワードを比較することで認証を行う。
\end{dfn}

ID 管理機能では、TOTP はパスワード認証と組み合わせた多要素認証（MFA、定義 \ref{dfn:mfa}）において、「所持情報」を用いた認証を実現する仕組みとして広く利用されている。

\subsubsection{OAuth 2.0}

ID 管理機能では、認証だけでなく、サービス間でユーザーのデータをやり取りする場面もある。そのための仕様が OAuth 2.0 である。

\begin{dfn}[OAuth 2.0]\label{dfn:ch5-oauth2}
\underline{OAuth 2.0} は、IETF（Internet Engineering Task Force）によって、RFC 6749 など一連のドキュメントとして標準化されている。
OAuth 2.0 は、ユーザーの同意に基づき、あるアプリケーションが別のサービス上にあるユーザーのデータ（例：プロフィール、写真リストなど）へ限定的にアクセスするための
\underline{アクセストークン}を発行し、それを用いてリソースへアクセスする方法を定めている。
\end{dfn}

ID 管理機能においては、サービス間のデータのやり取りを行う標準化仕様で OAuth 2.0 が利用されている。

\subsubsection{OpenID Connect (OIDC)}

ID 連携のフェデレーションプロトコル（第2章の定義 \ref{dfn:federation-protocol}）のうち、OAuth 2.0 の上に構築されたものが OpenID Connect である。

\begin{dfn}[OpenID Connect]\label{dfn:ch5-oidc}
\underline{OpenID Connect}（\underline{OIDC}）は、OpenID Foundation が策定した、異なるサービス間で認証情報とユーザープロフィールを安全に連携させるためのプロトコルである。
OIDC では、ユーザー情報を利用する依拠当事者（RP、定義 \ref{dfn:rp}）が、アイデンティティプロバイダー（IdP、定義 \ref{dfn:idp}）に認証を要求し、
IdP はユーザーの同意に基づき、認証結果を含む\underline{ID トークン}を RP に提供する。
\end{dfn}

\begin{rem}[OIDC と OAuth 2.0・SAML の関係]\label{rem:ch5-oidc-oauth-saml}
OIDC は、OAuth 2.0（定義 \ref{dfn:ch5-oauth2}）のリソースアクセスの仕組みをベースにして、ID トークンのやり取りを定義する形となっており、
OAuth 2.0 のアクセストークンを用いた非同期なユーザー情報の取得についても定義している。
また、SAML が XML ベースであるのに対し、OIDC は REST/JSON ベースのため、モダンな Web アプリケーションやモバイルアプリと高い親和性をもつ。
\end{rem}

\subsubsection{System for Cross-domain Identity Management (SCIM)}

第4章で扱うアイデンティティプロビジョニング（定義 \ref{dfn:identity-provisioning}）を、組織やサービスの境界を越えて自動化するための仕様が SCIM である。

\begin{dfn}[SCIM]\label{dfn:ch5-scim}
\underline{System for Cross-domain Identity Management}（\underline{SCIM}）は、IETF（Internet Engineering Task Force）によって、RFC 7643、7644 として標準化されている。
SCIM は、異なるドメイン（組織やサービス）間で、ユーザー ID 情報のプロビジョニング（作成、更新、削除）を自動化するためのプロトコルである。
REST と JSON をベースに、ユーザーやグループに対する操作（CRUD）を標準化している。
\end{dfn}

ID 管理機能においては、SCIM はアイデンティティプロビジョニングで利用されている。

\subsubsection{Shared Signals Framework (SSF)}

第4章で扱うセキュリティイベントの通知（定義 \ref{dfn:security-event-notification}）を、システム間でリアルタイムに行うための仕様群が SSF である。

\begin{dfn}[Shared Signals Framework]\label{dfn:ch5-ssf}
\underline{Shared Signals Framework}（\underline{SSF}）は、OpenID Foundation 内のワーキンググループで仕様策定中の仕様群で、
システム間でセキュリティイベントをリアルタイムに共有するためのフレームワークである。
\end{dfn}

\begin{exa}[SSF によるセキュリティイベントの共有]\label{exa:ch5-ssf-usage}
あるサービスで不審なアクティビティが検知されたり、ユーザーがセッションを自らログアウト（定義 \ref{dfn:logout}）したりした場合、
SSF を用いてその情報を即座に他の連携サービスに伝達することで、ID 連携による安全性を高める。
また、OAuth 2.0 の各種トークンの有効期限を待たず、継続的なアクセス評価を行い、そのリスクに応じてトークンの即時無効化が可能になる。
\end{exa}

ID 管理機能においては、SSF はセキュリティイベント通知で利用されている。

以上の標準化仕様を、策定団体と ID 管理機能での主な適用範囲とともにまとめる。

\begin{table}[H]
\centering
\caption{ID 管理機能で利用される標準化仕様}\label{tab:ch5-standards}
\small
\begin{tabular}{|p{0.2\textwidth}|p{0.3\textwidth}|p{0.4\textwidth}|}
\hline
仕様 & 策定団体（文書） & ID 管理機能での主な適用範囲 \\
\hline
WebAuthn & W3C と FIDO Alliance（FIDO2 プロジェクトの一部） & フィッシング耐性が求められる高セキュリティな当人認証を利用する機能（ログイン、再認証） \\
\hline
TOTP & IETF（RFC 6238） & パスワード認証と組み合わせた MFA における「所持情報」を用いた認証 \\
\hline
OAuth 2.0 & IETF（RFC 6749 など一連のドキュメント） & サービス間のデータのやり取り \\
\hline
OIDC & OpenID Foundation & 異なるサービス間での認証情報とユーザープロフィールの連携（ID 連携） \\
\hline
SCIM & IETF（RFC 7643、7644） & アイデンティティプロビジョニング \\
\hline
SSF & OpenID Foundation 内のワーキンググループ（策定中） & セキュリティイベント通知 \\
\hline
\end{tabular}
\end{table}

\subsection{注意したい実装例}

ここまでの観点・原則・ガイドライン・標準化仕様を踏まえても、実装の細部で問題が生じることがある。
この項では、本書が挙げる注意したい実装例を 4 つ取り上げ、それぞれ何が問題で、なぜ問題なのかを論じる。

\subsubsection{可変なユーザー識別子の内部利用}

第3章で紹介したアイデンティティレジスター（定義 \ref{dfn:identity-register}）では、属性情報をユーザーに紐づけて管理する。
その紐づけの「キー」に何を使うかが、最初の注意点である。

\begin{prop}[属性情報は不変なユーザー識別子に紐づける]\label{prop:ch5-immutable-identifier}
アイデンティティレジスターにて属性情報をユーザーに紐づける際は、不変であるユーザー識別子に対して属性情報を紐づける形で管理される必要がある。
とくに、ログイン画面でユーザーが入力するメールアドレスや電話番号、サービス内で利用しているユーザー名といったユーザー識別子は、
属性情報を紐づけるためのユーザー識別子としての利用は不適切である。
\end{prop}

\begin{proof}
ログイン時に入力される識別子が可変であることから示す。
まず、ログイン画面において、ユーザーはメールアドレスや電話番号、サービス内で利用しているユーザー名といったユーザー識別子を入力して、自身が誰なのかを主張する。
次に、これらの値は時間と共に変わる可能性がある。
もしこれらを属性情報の紐づけ先の識別子にすると、値が変わるたびに、アイデンティティレジスターとそれに紐づく全てのデータの識別子を更新しなければならない（例 \ref{exa:ch5-email-as-identifier}）。
よって、これらは紐づけ用の識別子として不適切であり、紐づけ先は不変であるユーザー識別子でなければならない。
\end{proof}

\begin{exa}[メールアドレスをユーザー識別子とした場合]\label{exa:ch5-email-as-identifier}
メールアドレスをアイデンティティレジスターのユーザー識別子として利用している状況を考える。
ここでメールアドレスが変更された場合、アイデンティティレジスターとそれに紐づく全てのデータの識別子を更新する必要がある。
そのため、処理が複雑になったり、メールアドレス変更を許可できない事態が考えられる。
\end{exa}

\subsubsection{攻撃者に親切なエラー表示}

エラーメッセージを分かりやすくすることはユーザビリティの観点（定義 \ref{dfn:ch5-usability}）から望ましいが、
分かりやすさが攻撃者にとっての手がかりになる場合がある。

\begin{exa}[ログイン時のエラー表示]\label{exa:ch5-login-error-message}
パスワード認証（定義 \ref{dfn:password-auth}）を用いたログインにおいて、ユーザーが入力するメールアドレスが登録済み、かつパスワードが正しくなかったとする。
このとき、「メールアドレスは登録済みですがパスワードが正しくありません」ではなく、
「メールアドレスもしくはパスワードが正しくありません」というエラーを表示することが一般的である。
\end{exa}

\begin{prop}[認証前に登録状況を明かすエラー表示の危険]\label{prop:ch5-helpful-error-risk}
パスワード認証が成功していない段階で、メールアドレスが登録済みかどうかを区別するエラー表示を行うと、
悪意のある第三者に有益な情報を与えてしまう可能性がある。
\end{prop}

\begin{proof}
認証前の操作者が誰であるか分からないことと、登録状況という情報の攻撃上の価値の二点から示す。
まず、パスワード認証が成功していない段階では、この操作を行っているのが正規のユーザーであるかがわからない。
したがって、親切なエラー表示は、正規のユーザーだけでなく悪意のある第三者に対しても同じ情報を与える。
次に、あるメールアドレスが登録済みかどうかを知られると、不正ログインを試みる攻撃の対象が登録済みユーザーに絞られたり、
登録済みメールアドレスへのメールを起点としたフィッシング攻撃などのリスクが伴う。
よって、そのようなエラー表示は悪意のある第三者に有益な情報を与えてしまう可能性があり、例 \ref{exa:ch5-login-error-message} の一般的なエラー表示はこれを避けるためのものである。
\end{proof}

\begin{rem}[新規登録・アカウントリカバリーでのエラー表示]\label{rem:ch5-signup-recovery-error}
一方で、新規登録（定義 \ref{dfn:signup}）やアカウントリカバリー（定義 \ref{dfn:account-recovery}）の文脈では、
メールアドレスや電話番号といった連絡先が登録されていないことをエラーとして表示するサービスが見られる。
このような実装は、膨大なメールアドレスとパスワードのリストの中から有効なメールアドレスのみに絞る処理などを容易にする可能性がある。
メールアドレスの状態を正規のユーザーのみに伝えたい場合は、エラー表示ではなく、実際のメールの内容を利用するといった対策がある。
\end{rem}

\subsubsection{セキュリティ重視によるユーザビリティ低下}

次の例は、命題 \ref{prop:ch5-trust-experience-tradeoff} のトレードオフが、単一のパラメータの設定に現れる場合である。

\begin{prop}[OTP の有効期限の短縮とユーザビリティ]\label{prop:ch5-otp-expiry}
メールや SMS によるワンタイムパスワード（OTP）認証（定義 \ref{dfn:otp-auth}）では、OTP の有効期限が短いほど安全性が高まる。
しかし、有効期限を短くしすぎると、多数のユーザーがログインできなくなるという事態が生じうるので、これは避けたい。
\end{prop}

\begin{proof}
有効期限内に OTP を入力し終えるまでにかかる時間が、ユーザーの側で必ずしも短くないことから示す。
まず、メールや SMS では配送に遅延もありうる。
次に、ユーザーによっては、受け取った OTP の入力に時間がかかってしまうこともありうる。
有効期限を短くすると、これらの遅れを含めた所要時間が有効期限を超えるユーザーが増え、そのユーザーは正しい OTP を持っていてもログインできない。
よって、安全性のために有効期限を短くした結果、多数のユーザーがログインできなくなる事態が生じうる。
\end{proof}

\subsubsection{当人認証、身元確認の手段を同時に失う“詰み”につながる実装}

最後の例は、個々には合理的な技術の変化が組み合わさることで、ユーザーが回復不能な状況に陥る場合である。

\begin{rem}[当人認証と身元確認の手段の変化]\label{rem:ch5-means-shift-mobile}
当人認証の手段としては、知識情報から所持情報を用いたものへの移り変わりが見られる（認証の 3 要素については定義 \ref{dfn:authentication-factors}）。
また、身元確認（定義 \ref{dfn:identity-proofing}）で使える身分証明書は、モバイル端末内で管理されるようになるといった変化が見られる。
\end{rem}

\begin{prop}[同一端末への集中による“詰み”]\label{prop:ch5-single-device-deadlock}
当人認証に利用するクレデンシャルと、身元確認で利用する身分証明書が同一のモバイル端末で管理されている場合、
そのモバイル端末が故障、紛失した際に、当人認証と身元確認の両方の手段を一気に失ってしまう\underline{“詰み”}の状況がある。
\end{prop}

\begin{proof}
二つの手段の依存先が一つの端末に重なることから示す。
まず、当人認証のクレデンシャルがその端末で管理されているので、端末が故障・紛失すると当人認証の手段が失われる。
次に、身元確認で利用する身分証明書も同じ端末で管理されているので、同じ出来事によって身元確認の手段も失われる。
当人認証の手段を失ったときには、身元確認によって本人であることを確かめてアカウントを回復する（定義 \ref{dfn:account-recovery}）ことが考えられるが、その身元確認の手段も同時に失われている。
よって、一つの端末の故障・紛失によって両方の手段を一気に失い、回復の道が断たれた“詰み”の状況になる。
\end{proof}

このような事態を避けるため、サービスは自身のユーザーのクレデンシャルと身分証明書の管理方法を意識し、
ユーザーの環境に何か問題が起こった際にもサービスを継続して利用できるような状況になるような設計と実装が求められる。
これは、当人認証の 4 つの局面（定義 \ref{dfn:authenticator-phases}）のうち、盗難・紛失等への対応とアカウントの回復の局面に関わる設計上の要請である。

%% =======================================================
%%  part: ch6
%% =======================================================
\section{ID管理機能を支える技術と開発スタイル}

この章では、\cite{Ito2026} の第6章に沿って、ID 管理機能を実際に作るうえで土台となる技術と、ID 管理機能の開発スタイルを扱う。
前半（ID 管理機能を支える技術）では、まず ID 管理機能を Web アプリケーションとして提供するという構成の利点を述べ、
次に HTTP を用いたフロントチャネルとバックチャネルの通信、Base64URL エンコードと CBOR によるエンコード・デコード、
ID 管理で利用する暗号技術（暗号化と復号、署名の生成と検証、共通鍵暗号と公開鍵暗号）、それらを組み合わせた仕様群である JOSE と COSE、
リプレイ攻撃を防ぐチャレンジレスポンス認証、Web ブラウザの果たす役割、デスクトップアプリ・ネイティブアプリとの連動を順に説明する。
後半（ID 管理機能の開発スタイル）では、スクラッチ実装、既存のライブラリやフレームワークの利用、ID 管理の製品・IDaaS の利用という 3 つのスタイルを、
自由度とコストの観点から比較する。

%======================================================================
\subsection{ID管理機能を支える技術}\label{subsec:ch6-technologies}

この節では、これまでの章で述べてきた身元確認・当人認証・ID 連携・アクセス制御・セッション管理といった ID 管理機能を、
実際のオンラインサービスの上で実現するための技術的な基盤を整理する。
個々のプロトコルの詳細に入る前に、どのような形態で ID 管理機能を提供するのが有利か、
そしてその上で通信の経路、データの表現、暗号技術、Web ブラウザがそれぞれどのような役割を果たすかを順に見ていく。

%----------------------------------------------------------------------
\subsubsection{「Webアプリケーションとして ID 管理機能を提供する」という選択肢}

オンラインサービスの提供形態はさまざまであり、それに応じて ID 管理機能をどこに実装するかにも選択の余地がある。
この項では、複数の環境でサービスを提供する場合に、ID 管理機能を Web アプリケーションとして実装し、
各プラットフォームのアプリケーションからそれを利用する構成が有利であることを、本書の理由づけに沿って述べる。

オンラインサービスの提供形態には、Web アプリケーションのみ、ネイティブアプリケーションのみ、あるいはその両方、
さらには複数のプラットフォームでの提供など、さまざまなパターンが存在する。それに伴い、ID 管理機能の提供方法もいくつか考えられる。

ネイティブアプリケーションのみで提供されるサービスの場合、ID 管理機能もアプリケーションの一部として実装されることが一般的である。
しかし、PC ユーザー向けに Web アプリケーションを提供したり、他のプラットフォームでもネイティブアプリケーションを展開したりする際には、
それぞれで ID 管理機能を実装する必要が生じる。

このように複数の環境でサービスを提供するうえで、ID 管理機能の実装に適しているのが Web アプリケーションである。
サービスが Web アプリケーションで提供されている場合はもちろん、ネイティブアプリケーションの場合でも、
OS が提供する\underline{セキュアなブラウザコンポーネント}（Android では AuthTabs、iOS では WebAuthenticationSession）を呼び出すことで、
Web 上の ID 管理機能を利用できる。これにより、ネイティブアプリケーションと遜色のない UX を実現できる。

\begin{prop}[Web アプリケーションへの ID 管理機能の集約の利点]\label{prop:ch6-web-idm-aggregation}
複数の環境（Web、iOS、Android など）でサービスを提供する場合、ID 管理機能を一つの Web アプリケーションに集約し、
各プラットフォームのアプリケーションからそれを利用する構成をとると、次の利点が得られる。
\begin{enumerate}
\item ID 管理機能を集約し、システム全体で UX の共通化と向上を図ることができる。
\item プラットフォームごとに ID 管理機能を実装するための初期実装のコストを削減できる。
\item 将来的な機能追加も Web アプリケーションの改修のみで済むため、開発・運用コストを低く抑えられる。
\end{enumerate}
\end{prop}

\begin{proof}
この構成が有効である背景には、OpenID Connect のような ID 連携の標準プロトコル（第2章のフェデレーションプロトコル、定義 \ref{dfn:federation-protocol}）を適用しやすいという事情がある。
これを軸に各項目を示す。

まず (1) について。ID 管理機能を提供する Web アプリケーションが IdP（Identity Provider、定義 \ref{dfn:idp}）として振る舞い、
各プラットフォームのアプリケーション（Web、iOS、Android など）が RP（Relying Party、定義 \ref{dfn:rp}）として振る舞うとする。
すると、どのプラットフォームから利用しても当人認証などの ID 管理機能は同じ IdP が担うことになるので、ID 管理機能が一か所に集約され、
システム全体で UX を共通化でき、その改善も全体に及ぶ。また、ネイティブアプリケーションからは OS のセキュアなブラウザコンポーネントを通して
この Web 上の機能を呼び出せるので、ネイティブアプリケーションと遜色のない UX を保ったままこの構成をとれる。

次に (2) について。ID 管理機能の実体は IdP である Web アプリケーション一つだけであり、各プラットフォームのアプリケーションは RP として
標準プロトコルに従って IdP を利用するだけでよい。したがって、プラットフォームごとに ID 管理機能を実装する必要がなくなり、その初期実装のコストを削減できる。

最後に (3) について。ID 管理機能が IdP である Web アプリケーションに集約されているので、新しい機能を追加するときも、
その Web アプリケーションを改修すれば全プラットフォームに反映される。よって、開発・運用コストを低く抑えられる。
\end{proof}

この構成を図示すると図 \ref{fig:ch6-web-idm} のようになる（本書の図「Web アプリケーションに ID 管理機能を集約するメリット」に対応する）。

\begin{figure}[H]
\centering
\begin{tikzpicture}[
  box/.style={draw, rounded corners, align=center, minimum width=26mm, minimum height=9mm},
  >=Stealth]
\node[box, minimum width=60mm] (idp) at (0,2.4) {ID 管理機能を提供する Web アプリケーション\\（IdP）};
\node[box] (web) at (-4,0) {Web アプリケーション\\（RP）};
\node[box] (ios) at (0,0) {iOS アプリ\\（RP）};
\node[box] (and) at (4,0) {Android アプリ\\（RP）};
\draw[<->] (web) -- node[left, font=\small, xshift=-2mm] {OpenID Connect} (idp);
\draw[<->] (ios) -- (idp);
\draw[<->] (and) -- (idp);
\node[font=\small, align=center] at (0,-1.2) {iOS・Android からは OS のセキュアなブラウザコンポーネントを通して利用する};
\end{tikzpicture}
\caption{Web アプリケーションに ID 管理機能を集約する構成}\label{fig:ch6-web-idm}
\end{figure}

%----------------------------------------------------------------------
\subsubsection{HTTPを用いたフロントチャネル、バックチャネルのリクエストとレスポンス}

ID 管理機能を Web アプリケーションで実装する際には、HTTP プロトコルを用いた、異なるサービス・アプリケーション間の通信が発生する。
その通信においては、\underline{フロントチャネル}（Front-Channel）と\underline{バックチャネル}（Back-Channel）という 2 つの経路の概念がある。
この 2 つは、通信の主体と通信が行われる場所が異なる概念であり、ID 連携のための標準化プロトコルにおいても重要な意味合いを持つ。
この項では、それぞれの経路を定義し、HTTP のリクエストとレスポンスの具体例で動きを確かめたうえで、両者を比較する。

\underline{フロントチャネル：}
まず、ユーザーの Web ブラウザを経由する経路を定義する。

\begin{dfn}[フロントチャネル]\label{dfn:front-channel}
ユーザーの Web ブラウザを介して行われるサーバー間の通信で利用される経路を\underline{フロントチャネル}（Front-Channel）という。
フロントチャネルでは、サーバー同士が直接やり取りするのではなく、ブラウザのリダイレクト機能などを利用して情報を伝達する。
\end{dfn}

\begin{prop}[フロントチャネルではログイン中のユーザーと紐付けられる]\label{prop:ch6-front-channel-session}
フロントチャネルでアクセスを受けたサービスは、Web ブラウザのセッション管理の仕組みを利用できる。
このため、フロントチャネルは「ログイン状態のユーザーに何かを求める」ユースケースで使われている。
\end{prop}

\begin{proof}
フロントチャネルの通信は、各サービスへのリクエストを Web ブラウザ自身が送信する形で進む。
Web ブラウザはサービスへのリクエストに、そのサービスに対して保存している Cookie の値を含めて送る。
よって、アクセス先のサービスはその Cookie からログインセッション（定義 \ref{dfn:login-session}）に保存されている情報、
すなわちログイン中のユーザーなどの情報にアクセスでき、受け取ったパラメータをそのユーザーと紐付けて処理できる。
したがって、ログイン状態のユーザーに何かを求めるユースケースに適している。この流れを次の例で具体的に確かめる。
\end{proof}

\begin{exa}[フロントチャネルによるサービス間のデータ送信]\label{exa:ch6-front-channel}
あるサービス（service1.example.com、以下サービス 1）から、フロントチャネルで別のサービス（service2.example.net、以下サービス 2）にデータを送る例を考える。
通信は次の 3 段階で進む（図 \ref{fig:ch6-front-channel}）。
\begin{enumerate}
\item 初めに Web ブラウザがサービス 1 にアクセスする。このときのリクエストには、Web ブラウザがサービス 1 に対して保存している Cookie の値も含まれるため、
サービス 1 はログイン中のユーザーなどの情報にアクセスできる。
\item サービス 1 は何らかの理由で、ログイン中のユーザーをサービス 2 に送るためのレスポンスを返す。ここで、サービス 1 はサービス 2 に渡したいデータをクエリパラメータに指定する。
\item レスポンスを受け取った Web ブラウザはサービス 2 にアクセスする。これがいわゆる\underline{リダイレクト}と呼ばれる処理であり、
ユーザー目線では Web ブラウザに表示されているアクセス先がサービス 1 からサービス 2 に変わる。
最初のリクエストと同様に、このときのリクエストには Web ブラウザがサービス 2 に対して保存している Cookie の値が含まれる。
サービス 2 はここからログインセッションに保存されている情報にアクセスし、指定されたパラメータと紐付けるといった処理が可能になる。
\end{enumerate}
各段階の HTTP メッセージは、この例の直後に示す。
\end{exa}

段階 (1) のサービス 1 へのリクエストは次のとおりである（ヘッダーの一部は省略している）。最終行の \texttt{Cookie} ヘッダーに、
Web ブラウザがサービス 1 に対して保存している値 \texttt{session\_id} が含まれている。

\begin{verbatim}
GET /path HTTP/1.1
Host: service1.example.com
...(中略)...
Cookie: session_id=a1b2c3d4e5f67890
\end{verbatim}

段階 (2) のサービス 1 のレスポンスは次のとおりである。ステータス \texttt{302 Found} と \texttt{Location} ヘッダーにより、
Web ブラウザをサービス 2 へ向かわせ、サービス 2 に渡したいデータ（\texttt{param1}, \texttt{param2}）をクエリパラメータに載せている。

\begin{verbatim}
HTTP/1.1 302 Found
Location: https://service2.example.net/path?param1=value1&param2=value2
\end{verbatim}

段階 (3) の、リダイレクトによるサービス 2 へのリクエストは次のとおりである。サービス 2 に対して保存されている、
サービス 1 のものとは別の \texttt{session\_id} が送られている。

\begin{verbatim}
GET /path?param1=value1&param2=value2 HTTP/1.1
Host: service2.example.net
...(中略)...
Cookie: session_id=z9y8x7w6v5f43210
\end{verbatim}

\begin{figure}[H]
\centering
\begin{tikzpicture}[>=Stealth, font=\small]
\node[draw, rounded corners, minimum width=26mm, minimum height=8mm] (s1) at (0,0) {サービス 1};
\node[draw, rounded corners, minimum width=26mm, minimum height=8mm] (br) at (5,0) {Web ブラウザ};
\node[draw, rounded corners, minimum width=26mm, minimum height=8mm] (s2) at (10,0) {サービス 2};
\draw (s1) -- (0,-4.2); \draw (br) -- (5,-4.2); \draw (s2) -- (10,-4.2);
\draw[->] (5,-1.0) -- node[above, align=center] {(1) GET /path\\Cookie（サービス 1 用）} (0,-1.0);
\draw[->] (0,-2.2) -- node[above, align=center] {(2) 302 Found\\Location: サービス 2（クエリパラメータ）} (5,-2.2);
\draw[->] (5,-3.4) -- node[above, align=center] {(3) GET /path?param1=...\\Cookie（サービス 2 用）} (10,-3.4);
\end{tikzpicture}
\caption{フロントチャネルを用いた通信}\label{fig:ch6-front-channel}
\end{figure}

\underline{バックチャネル：}
次に、Web ブラウザを経由しない経路を定義する。

\begin{dfn}[バックチャネル]\label{dfn:back-channel}
サーバー間で直接行われる通信で利用される経路を\underline{バックチャネル}（Back-Channel）という。
バックチャネルでは、ユーザーのブラウザを介さずに、プログラムが裏側で HTTP リクエストを送信する。
\end{dfn}

\begin{figure}[H]
\centering
\begin{tikzpicture}[>=Stealth, font=\small]
\node[draw, rounded corners, minimum width=26mm, minimum height=8mm] (s1) at (0,0) {サービス 1};
\node[draw, rounded corners, minimum width=26mm, minimum height=8mm] (s2) at (8,0) {サービス 2};
\node[draw, rounded corners, dashed, minimum width=26mm, minimum height=8mm, text=gray] (br) at (4,1.6) {Web ブラウザ（介在しない）};
\draw[->] ([yshift=1.5mm]s1.east) -- node[above] {HTTP リクエスト} ([yshift=1.5mm]s2.west);
\draw[<-] ([yshift=-1.5mm]s1.east) -- node[below] {HTTP レスポンス} ([yshift=-1.5mm]s2.west);
\end{tikzpicture}
\caption{バックチャネルを用いた通信}\label{fig:ch6-back-channel}
\end{figure}

\begin{exa}[バックチャネルの最もシンプルなリクエスト]\label{exa:ch6-back-channel-simple}
バックチャネルを利用する最もシンプルなリクエストは、次の HTTP リクエストのようになる（この例の直後に示す）。
フロントチャネルのリクエストの例（例 \ref{exa:ch6-front-channel}）にあった Cookie の値などはなく、
このリクエストだけでは、リクエストの送信元を検証してのアクセス制限などはできない。
\end{exa}

\begin{verbatim}
GET /api/v1/products HTTP/1.1
Host: service2.example.net
\end{verbatim}

\begin{exa}[OAuth 2.0 のリソースアクセス]\label{exa:ch6-back-channel-bearer}
RFC 6749 The OAuth 2.0 Authorization Framework で定義されているリソースアクセスでは、Authorization Header に値を指定する
（次の HTTP リクエストでは \texttt{Authorization: Bearer ...}）。リクエストを受け取ったサービスは、この値を検証することで、
アクセス元の検証や対象ユーザーの特定などを行う。
\end{exa}

\begin{verbatim}
GET /resource/1 HTTP/1.1
Host: service2.example.net
Authorization: Bearer mF_9.B5f-4.1JqM
\end{verbatim}

\begin{prop}[バックチャネルでは機密情報も安全に送信できる]\label{prop:ch6-back-channel-confidential}
バックチャネルでは、機密情報なども安全に送信できる。
\end{prop}

\begin{proof}
本書は 2 つの理由を挙げる。
まず、例 \ref{exa:ch6-back-channel-bearer} の Authorization Header のほかにも、バックチャネルのリクエストを保護・強化する仕組みが存在する。
次に、バックチャネルの通信はサーバー間で直接行われるので（定義 \ref{dfn:back-channel}）、フロントチャネルに比べて Web ブラウザの介入がない。
よって、ブラウザを経由することによる漏洩のリスクを避けつつリクエストを保護できるので、機密情報なども安全に送信できる。
\end{proof}

以上の 2 つの経路の特性をまとめると、表 \ref{tab:ch6-channels} のようになる。表中の記号は、セキュリティの観点での相対的な評価であり、
△ は機密情報を扱えないこと、◎ は機密情報を扱えることを表す。

\begin{table}[H]
\centering
\caption{フロントチャネルとバックチャネルの比較}\label{tab:ch6-channels}
\begin{tabular}{|p{20mm}|p{58mm}|p{48mm}|}
\hline
特性 & フロントチャネル & バックチャネル \\ \hline
通信経路 & サーバー $\to$ Web ブラウザ $\to$ サーバー & サーバー $\to$ サーバー \\ \hline
ユーザーの関与 & あり（ログイン中のユーザーとの紐付けなど） & なし（裏側で実行） \\ \hline
セキュリティ & △（安全な保護、漏洩リスクの観点から機密情報を扱えない） & ◎（機密情報を扱える） \\ \hline
主な用途 & サービスと ID 管理機能の遷移；ID 連携のユーザー同意取得 & ID 連携のトークン取得；リソースアクセス \\ \hline
\end{tabular}
\end{table}

\begin{rem}[フロントチャネルとバックチャネルの使い分け]\label{rem:ch6-channel-usage}
表 \ref{tab:ch6-channels} が示すように、フロントチャネルはログイン中のユーザーとの紐付けができる代わりに、安全な保護と漏洩リスクの観点から機密情報を扱えない。
そのため、主な用途はサービスと ID 管理機能の間の遷移や、ID 連携におけるユーザーの同意取得である。
一方、バックチャネルはユーザーの関与がなく裏側で実行されるが、機密情報を扱えるので、ID 連携のトークン取得やリソースアクセスに用いられる。
\end{rem}

%----------------------------------------------------------------------
\subsubsection{エンコード、デコード：Base64URL エンコード、CBOR}

ID 管理機能では、同一サービス内や他サービスとの間でさまざまな形式のデータがやり取りされる。
ところが、その際に利用される HTTP などのプロトコルでは、データは文字列やバイナリフォーマットとして扱われることが一般的である。
数値や文字列から構造化されたデータまで、さまざまなフォーマットのデータを扱うために、データを文字列やバイナリデータへ変換する処理と、
元に戻す処理が必要になる。この項では、まずこの 2 つの処理を定義し、主な方式である Base64URL エンコードと CBOR を説明する。

\begin{dfn}[エンコードとデコード]\label{dfn:ch6-encode-decode}
さまざまなフォーマットのデータ（数値、文字列、構造化されたデータなど）を、プロトコルで扱える文字列やバイナリデータに変換する処理を
\underline{エンコード}（encode）という。エンコードされたデータを元のデータフォーマットに戻す処理を\underline{デコード}（decode）という。
元のデータ全体の集合を $D$、変換後の文字列またはバイナリデータ全体の集合を $S$ とすれば、エンコードは写像 $e:D\to S$、デコードは写像 $d:S\to D$ であり、
デコードが元のデータフォーマットに戻す処理であることは
\[
d\bigl(e(x)\bigr) = x \qquad (x\in D)
\]
と表せる。
\end{dfn}

主なエンコードの方式に、Base64URL エンコードと CBOR がある。

\underline{Base64URL エンコード：}

\begin{dfn}[Base64URL エンコード]\label{dfn:base64url}
バイナリデータや JSON テキストなど、任意のデータを、URL に含めても安全に扱える 64 文字の英数字からなる文字列に変換する方式を
\underline{Base64URL エンコード}という。
\end{dfn}

Base64URL エンコードを使うことで、画像のようなバイナリデータから文字列への変換や、JSON 文字列で表現した構造化されたデータの文字列への変換が実現できる
（本書の図「Base64URL エンコードによる文字列への変換」は、Base64URL エンコードによる文字列への変換を図示したものである）。

\begin{rem}[Base64 エンコードとの違い]\label{rem:ch6-base64-vs-base64url}
名前が似ている方式として \underline{Base64 エンコード}があるが、Base64URL エンコードでは「\texttt{+}」や「\texttt{/}」のように
URL として特別な意味を持つ文字を\emph{使わない}という特徴がある。これが、Base64URL エンコードの結果を URL に含めても安全に扱える理由である。
\end{rem}

\underline{CBOR（Concise Binary Object Representation）：}

\begin{dfn}[CBOR]\label{dfn:cbor}
さまざまなフォーマットのデータや構造化されたデータを、よりコンパクトなバイナリ形式で表現するための方式を
\underline{CBOR}（Concise Binary Object Representation）という。CBOR は RFC 7049 として定められている。
\end{dfn}

\begin{prop}[CBOR によるサイズの削減]\label{prop:ch6-cbor-compact}
データの送受信の観点では、CBOR を用いると、JSON で表現された文字列を Base64URL エンコードした文字列よりも、サイズを抑えたバイナリデータとしてデータを表現できる。
そのため CBOR は、外部のデバイスや Web ブラウザなどのソフトウェアが介入するプロトコルなどで利用されている。
\end{prop}

\begin{proof}
定義 \ref{dfn:cbor} のとおり、CBOR はデータをよりコンパクトな\emph{バイナリ}形式で表現する方式である。
一方、JSON 文字列を Base64URL エンコードする方法では、データはまず JSON というテキストで表現され、さらに URL で安全に扱える文字列（定義 \ref{dfn:base64url}）へ変換される。
よって、同じデータを送受信する場合、CBOR で表現したほうがサイズを抑えられる。このサイズの小ささが、外部のデバイスや Web ブラウザなどが介入するプロトコルで CBOR が使われる理由である。
\end{proof}

本書の図「CBOR によるバイナリデータへの変換」は、CBOR によるこのバイナリデータへの変換を図示したものである。

%----------------------------------------------------------------------
\subsubsection{ID管理で利用する暗号技術}

ID 管理機能には、暗号技術を利用するものがいくつかある。
この項では、暗号技術の主な用途である「暗号化、復号」と「署名の生成、検証」、そして「共通鍵暗号」と「公開鍵暗号」という鍵の形式による違いを整理する。
これらを理解しておくことが、後で述べる JOSE・COSE やパスキー認証の仕組みを理解するうえで重要である。

以後、この項では次の記号を用いる。

\begin{nota}[暗号技術の記号]\label{nota:ch6-crypto}
送りたいデータ（平文）を $m$、鍵を $k$ と書く。鍵 $k$ による暗号化を $\mathrm{Enc}_k$、鍵 $k'$ による復号を $\mathrm{Dec}_{k'}$ と書く。
公開鍵暗号の秘密鍵を $sk$、公開鍵を $pk$ と書き、秘密鍵 $sk$ による署名の生成を $\mathrm{Sign}_{sk}$、公開鍵 $pk$ による署名の検証を $\mathrm{Verify}_{pk}$ と書く。
\end{nota}

\underline{暗号化、復号：}

\begin{dfn}[暗号化と復号]\label{dfn:ch6-encryption}
送信者が鍵 $k$ を使って、送りたいデータ（\underline{平文}）$m$ を、鍵なしでは読めないデータ（\underline{暗号文}）$c=\mathrm{Enc}_k(m)$ に変換する処理を
\underline{暗号化}（encryption）という。暗号文の受信者が、同一もしくは対となる鍵 $k'$ を使って、暗号文を元の平文に戻す処理を
\underline{復号}（decryption）という。すなわち
\[
\mathrm{Dec}_{k'}\bigl(\mathrm{Enc}_k(m)\bigr) = m
\]
であり、$k'$ は $k$ と同一の鍵であるか、$k$ と対になる鍵である。
\end{dfn}

暗号化と復号は、正しい鍵を持つ受信者だけがデータにアクセスできるようにすること、すなわち第三者から中身を隠してデータの\underline{機密性}を守ることを目的とした技術である
（本書の図「暗号化、復号機能の利用」は、この暗号化と復号の利用を図示したものである）。

\underline{署名の生成、検証：}

\begin{dfn}[署名の生成と検証]\label{dfn:ch6-signature}
送信者が鍵を使って、データ $m$ に対する\underline{デジタル署名} $\sigma$ を作成する処理を\underline{署名の生成}という。
受信者が対になる鍵を使って、その署名が正当なものか、またデータが改ざんされていないかを確認する処理を\underline{署名の検証}という。
公開鍵暗号の記号（記号 \ref{nota:ch6-crypto}）では、$\sigma=\mathrm{Sign}_{sk}(m)$ が署名の生成、$\mathrm{Verify}_{pk}(m,\sigma)$ が
署名 $\sigma$ とデータ $m$ の組み合わせの検証にあたる。
\end{dfn}

署名の生成と検証は、送受信におけるデータの\underline{完全性}（改ざんされていないこと）、送信者の確認、\underline{否認防止}（作成した事実を後から否定できないこと）の保証を目的とした技術である
（本書の図「署名の生成、検証機能の利用」は、この署名の生成と検証の利用を図示したものである）。

\begin{rem}[暗号化と署名の目的の違い]\label{rem:ch6-encryption-vs-signature}
定義 \ref{dfn:ch6-encryption} と定義 \ref{dfn:ch6-signature} は、どちらも鍵を使う処理であるが目的が異なる。
暗号化・復号の目的は\emph{機密性}（第三者から中身を隠すこと）であり、署名の生成・検証の目的は\emph{完全性}・送信者の確認・\emph{否認防止}である。
\end{rem}

\underline{共通鍵暗号と公開鍵暗号：}
暗号化・復号と署名の生成・検証は、どのような鍵を使うかによって 2 つの形式に分かれる。

\begin{dfn}[共通鍵暗号]\label{dfn:ch6-symmetric-key}
1 つの\underline{共通鍵} $k$ を利用して「暗号化、復号」または「署名の生成、検証」を行う方法を\underline{共通鍵暗号}という。
暗号化・復号の場合、定義 \ref{dfn:ch6-encryption} の $k'$ は $k$ と同一であり、$\mathrm{Dec}_{k}\bigl(\mathrm{Enc}_k(m)\bigr)=m$ となる。
\end{dfn}

\begin{rem}[共通鍵暗号の長所と課題]\label{rem:ch6-symmetric-tradeoff}
共通鍵暗号は処理が高速である一方、送信者と受信者が同じ鍵を持つ必要があるので、安全に鍵を共有・管理する仕組みが必要である。
\end{rem}

\begin{dfn}[公開鍵暗号]\label{dfn:ch6-public-key}
\underline{秘密鍵} $sk$ と\underline{公開鍵} $pk$ のペアを使用して「暗号化、復号」または「署名の生成、検証」を行う方法を\underline{公開鍵暗号}という。
\end{dfn}

公開鍵暗号の 2 つの使い方を、本書の記述に沿って命題として述べる。

\begin{prop}[公開鍵暗号による暗号化と復号]\label{prop:ch6-public-key-encryption}
公開鍵暗号による「暗号化、復号」では、送信者が受信者の公開鍵 $pk$ を利用して送りたいデータを暗号化して $c=\mathrm{Enc}_{pk}(m)$ を送信し、
受信者は秘密鍵 $sk$ を用いて $m=\mathrm{Dec}_{sk}(c)$ と復号してそのデータを処理することが一般的である。
この使い方では、秘密鍵の持ち主のみがデータを復号して処理できる。
\end{prop}

\begin{proof}
暗号文 $c$ を元の平文に戻すには、定義 \ref{dfn:ch6-encryption} より、暗号化に使った鍵 $pk$ と対になる鍵、すなわち秘密鍵 $sk$ が必要である。
公開鍵 $pk$ は誰が持っていてもよいが、秘密鍵 $sk$ は持ち主だけが持つ。よって、$pk$ で暗号化されたデータを復号して処理できるのは秘密鍵の持ち主のみである。
\end{proof}

この流れは本書の図「公開鍵暗号を用いた暗号化と復号」に示されており、図 \ref{fig:ch6-pk-enc} のように描ける。

\begin{figure}[H]
\centering
\begin{tikzpicture}[>=Stealth, font=\small,
  box/.style={draw, rounded corners, align=center, minimum width=24mm, minimum height=8mm}]
\node[box] (snd) at (0,0) {送信者\\（公開鍵 $pk$）};
\node[box] (rcv) at (8,0) {受信者\\（秘密鍵 $sk$）};
\draw[->] (snd) -- node[above] {暗号文 $c=\mathrm{Enc}_{pk}(m)$} node[below] {$m=\mathrm{Dec}_{sk}(c)$ で復号} (rcv);
\end{tikzpicture}
\caption{公開鍵暗号を用いた暗号化と復号}\label{fig:ch6-pk-enc}
\end{figure}

\begin{prop}[公開鍵暗号による署名の生成と検証]\label{prop:ch6-public-key-signature}
公開鍵暗号による「署名の生成、検証」では、送信者は送りたいデータ $m$ に対して自身の秘密鍵 $sk$ で署名 $\sigma=\mathrm{Sign}_{sk}(m)$ を作成して $m$ と一緒に送信し、
受信者は事前に共有されている送信者の公開鍵 $pk$ で、その署名とデータの組み合わせを $\mathrm{Verify}_{pk}(m,\sigma)$ により検証するという使い方が一般的である。
この使い方では、秘密鍵の持ち主のみが署名を生成でき、検証者は公開鍵があればその署名を検証できる。
\end{prop}

\begin{proof}
署名の生成には秘密鍵 $sk$ を使い、秘密鍵は持ち主だけが持つ。よって、正当な署名を生成できるのは秘密鍵の持ち主のみである。
一方、検証には定義 \ref{dfn:ch6-signature} より生成に使った鍵と対になる鍵、すなわち公開鍵 $pk$ を使う。
公開鍵は事前に共有されているので、検証者は公開鍵さえあればその署名を検証できる。
\end{proof}

この流れは本書の図「公開鍵暗号を用いた署名生成と検証」に示されており、図 \ref{fig:ch6-pk-sig} のように描ける。

\begin{figure}[H]
\centering
\begin{tikzpicture}[>=Stealth, font=\small,
  box/.style={draw, rounded corners, align=center, minimum width=24mm, minimum height=8mm}]
\node[box] (snd) at (0,0) {送信者\\（秘密鍵 $sk$）};
\node[box] (rcv) at (8,0) {受信者\\（送信者の公開鍵 $pk$）};
\draw[->] (snd) -- node[above] {$m$ と $\sigma=\mathrm{Sign}_{sk}(m)$} node[below] {$\mathrm{Verify}_{pk}(m,\sigma)$ で検証} (rcv);
\end{tikzpicture}
\caption{公開鍵暗号を用いた署名生成と検証}\label{fig:ch6-pk-sig}
\end{figure}

\begin{rem}[公開鍵暗号と署名の利用先]\label{rem:ch6-public-key-signature-usage}
公開鍵暗号方式と「署名の生成、検証」の組み合わせは、ID 連携のプロトコル（定義 \ref{dfn:federation-protocol}）やパスキー認証（定義 \ref{dfn:passkey}）で使われている。
後者については、チャレンジレスポンス認証の項（命題 \ref{prop:ch6-passkey-challenge-response}）で具体的な流れを述べる。
\end{rem}

%----------------------------------------------------------------------
\subsubsection{JOSEとCOSE}

これまでに紹介したエンコード／デコード方式と暗号方式を組み合わせ、さまざまな形式のデータに署名や暗号化を施しつつサービス間でやり取りするための仕様群が、JOSE や COSE である。
この項では、JSON をベースとする JOSE と、CBOR をベースとする COSE をそれぞれ定義し、構成する仕様と利用例を見ていく。

\underline{JOSE（JSON Object Signing and Encryption）：}

\begin{dfn}[JOSE]\label{dfn:jose}
JSON をベースとしたデータ（特に JWT）に署名や暗号化を行うための標準仕様群を\underline{JOSE}（JSON Object Signing and Encryption）という。
JOSE は複数の RFC により定義されており、その中でも次の仕様は Web API 連携における認証・認可で広く利用されている。
\begin{enumerate}
\item RFC 7519 JSON Web Token（JWT）（定義 \ref{dfn:jwt}）
\item RFC 7515 JSON Web Signature（JWS）（定義 \ref{dfn:ch6-jws}）
\item RFC 7516 JSON Web Encryption（JWE）（定義 \ref{dfn:ch6-jwe}）
\item RFC 7517 JSON Web Key（JWK）（定義 \ref{dfn:ch6-jwk}）
\item RFC 7518 JSON Web Algorithms（JWA）（定義 \ref{dfn:ch6-jwa}）
\end{enumerate}
\end{dfn}

上の各仕様の内容を、個別に定義として述べる。

\begin{dfn}[JWT]\label{dfn:jwt}
RFC 7519 で定められた、2 つのサービス間で構造化されたデータをやり取りするために使える文字列表現を\underline{JWT}（JSON Web Token）という。
\end{dfn}

\begin{dfn}[JWS]\label{dfn:ch6-jws}
RFC 7515 で定められた、署名付きの JWT の生成と検証をするための仕様を\underline{JWS}（JSON Web Signature）という。
\end{dfn}

\begin{dfn}[JWE]\label{dfn:ch6-jwe}
RFC 7516 で定められた、ペイロードの暗号化を伴う JWT の生成と復号のための仕様を\underline{JWE}（JSON Web Encryption）という。
\end{dfn}

\begin{dfn}[JWK]\label{dfn:ch6-jwk}
RFC 7517 で定められた、署名生成と検証、暗号化と復号で利用するための鍵表現についての仕様を\underline{JWK}（JSON Web Key）という。
\end{dfn}

\begin{dfn}[JWA]\label{dfn:ch6-jwa}
RFC 7518 で定められた、署名生成と検証、暗号化と復号で利用する際の暗号化アルゴリズムについての仕様を\underline{JWA}（JSON Web Algorithms）という。
\end{dfn}

JOSE の一番最初のユースケースとして、ID 連携のためのプロトコルである OpenID Connect の \underline{ID Token} がある。

\begin{exa}[OpenID Connect の ID Token]\label{exa:ch6-id-token}
OpenID Connect の ID Token では、複数サービス間でユーザー情報を含む構造化されたデータをトークン文字列としてやり取りするために、
JWT、JWS、JWK、JWA の組み合わせ（定義 \ref{dfn:jwt}、定義 \ref{dfn:ch6-jws}、定義 \ref{dfn:ch6-jwk}、定義 \ref{dfn:ch6-jwa}）が採用されている。
ID Token の文字列は、JWS で定義されている \underline{JWS Compact Serialization} というフォーマットである。
その具体例と、それを Base64URL デコード（定義 \ref{dfn:base64url}）した内容を、この例の直後に示す。
\end{exa}

ID Token の文字列の例は次のとおりである。実際には改行を含まない 1 本の文字列であり、紙面に収めるために途中で折り返している。
この文字列は「\texttt{.}」で区切られた 3 つのパートからなる。

\begin{verbatim}
eyJhbGci0iJSUzI1NiIsImtpZCI6IjFl0WdKazcifQ.ewogImlzcyI6ICJodHRw0i8vc2Vy
dmVyLmV4YW1wbGUuY29tIiwKICJzdWIiAiMjQ4Mjg5NzYxMDAxIiwKICJhdWQiOiAicZCAa
GRSa3F0MyIsCiAibm9uY2UiOiAibi0wUzZfV3pBMk1qIiwKICJleHAiOiAxMzExMjgxOTcw
LAogImlhdCI6IDEzMTEyODA5NzAKfQ.ggW8hZ1EuVLuxNuuIJKX_V8a_0MXzR0EHR9R6jgd
qr00F4daGU96Sr_P6qJp6IcmD3HP990bi1PRs-cwh3L0-p146waJ8IhehcwL7F09JdijmBq
kvPeB2T9CJNqeGpe-gccMg4vfKjkM8FcGvnzZUN4_KSP0aAp1t0J1zZvgjxq6ByKHi0tX7T
pdQyHE5lcMiKPxfEIQILVq0pc_E2DzL7emopWoaoZTF_m0_N0YzFC6g6Ejb0EoRoSK5hoDa
lrcvRYLSrQAZZKflyuVCyixEoV9GfNQC3_osjzw2PAithfubEEBLuVVk4XUVrWOLrLl0nx7
RkkU8NXNHq-rvKMzgg
\end{verbatim}

「\texttt{.}」で区切られた部分を Base64URL デコードすると、最初の 2 パートはそれぞれ構造化されたデータ（JSON）となる。
最初のパートに含まれるのは\underline{メタデータ}であり、次のようになる。\texttt{alg} は JWA で定義されている署名アルゴリズム、
\texttt{kid} は JWK で表現される鍵情報の識別子である。

\begin{verbatim}
{
  "alg": "RS256",   # JWAで定義されている署名アルゴリズム
  "kid": "1e9gdk7"  # JWKで表現される鍵情報の識別子
}
\end{verbatim}

次のパートには、JWT で定義されている\underline{標準クレーム}による、ユーザーと ID 連携フローの情報が含まれる。

\begin{verbatim}
{
  "iss": "http://server.example.com",
  "sub": "248289761001",
  "aud": "s6BhdRkqt3",
  "nonce": "n-0S6_WzA2Mj",
  "exp": 1311281970,
  "iat": 1311280970
}
\end{verbatim}

\begin{rem}[ID Token の構成のまとめ]\label{rem:ch6-id-token-structure}
例 \ref{exa:ch6-id-token} の ID Token は、JOSE の各仕様を次のように組み合わせている。
文字列全体の形式（JWS Compact Serialization）は JWS で定められ、各パートは Base64URL エンコードされている。
第 1 パート（メタデータ）の \texttt{alg} の値は JWA で定義された署名アルゴリズムを指し、\texttt{kid} は JWK で表現される鍵情報を指す。
第 2 パートは JWT で定義された標準クレーム（\texttt{iss}, \texttt{sub}, \texttt{aud}, \texttt{nonce}, \texttt{exp}, \texttt{iat}）からなる。
本書がデコード結果を示しているのは最初の 2 パートであり、Base64URL デコードすると構造化されたデータになるのもこの 2 パートである。
\end{rem}

\underline{COSE（CBOR Object Signing and Encryption）：}

\begin{dfn}[COSE]\label{dfn:cose}
データ表現に CBOR（定義 \ref{dfn:cbor}）を利用している、JOSE（定義 \ref{dfn:jose}）の CBOR 版と位置づけられる仕様群を
\underline{COSE}（CBOR Object Signing and Encryption）という。
\end{dfn}

COSE の仕様の内容は JOSE とほぼ同じものが揃っており、JOSE よりもデータ効率が求められる状況において採用されるケースが増えている。
COSE は、最初に仕様策定された RFC 8152 CBOR Object Signing and Encryption (COSE) に代わり、次のような仕様が標準化されている。
\begin{enumerate}
\item RFC 9052 CBOR Object Signing and Encryption (COSE): Structures and Process\\
COSE のデータ構造、処理モデル、および鍵の表現方法を詳細に定義しており、セキュリティプロファイルと拡張性のための基盤を提供する。
\item RFC 9053 CBOR Object Signing and Encryption (COSE): Hash Algorithms\\
COSE で使用されるハッシュアルゴリズムの仕様を記述し、署名および MAC 計算におけるアルゴリズムの選択と利用方法を提供する。
\item RFC 9054 CBOR Object Signing and Encryption (COSE): Algorithms\\
COSE で使用されるさまざまな署名、MAC、および暗号化アルゴリズム（例えば、AES、ECDH、EdDSA など）の詳細な仕様。
\item RFC 9055 CBOR Object Signing and Encryption (COSE): Header Parameters\\
COSE メッセージのヘッダーパラメータ（\texttt{kid}、\texttt{alg}、\texttt{crit} など）の定義と使用方法を規定し、
メッセージの受信者が署名や暗号化を検証するために必要な情報を効率的に伝達する方法。
\end{enumerate}

\begin{exa}[パスキー認証における CBOR と COSE]\label{exa:ch6-passkey-cbor}
パスキー認証（定義 \ref{dfn:passkey}）では、認証を利用するサービス、Web ブラウザ、認証器が登場する。
デバイス内で行われる認証器の処理において作成される認証情報には、CBOR 形式によるデータ表現と、COSE で定義されている鍵情報の表現が利用されている
（本書の図「パスキー認証における CBOR の使われ方」は、この CBOR の使われ方を図示したものである）。
\end{exa}

%----------------------------------------------------------------------
\subsubsection{チャレンジレスポンス認証}

暗号技術を当人認証や ID 連携に用いるとき、問題になる攻撃の一つが、通信を盗聴してその結果を再利用する攻撃である。
この項では、その攻撃を定義し、それを防ぐ方式であるチャレンジレスポンス認証と、パスキー認証における公開鍵暗号との組み合わせを述べる。

\begin{dfn}[リプレイ攻撃]\label{dfn:ch6-replay-attack}
一度盗聴した通信結果を利用して不正ログインを試みる攻撃を\underline{リプレイ攻撃}（replay attack）という。
\end{dfn}

\begin{dfn}[チャレンジレスポンス認証]\label{dfn:challenge-response}
ID 連携を行うサービス間や、当人認証の際のサービスと当人認証の認証器の間のやり取りにおいて、
リクエストの送信者が\underline{チャレンジ}という一意の値を送り、受信者はそれに対して、あらかじめ決められている署名生成や計算ロジックに基づいた出力を
\underline{レスポンス}として返す方式を\underline{チャレンジレスポンス認証}という。このとき、チャレンジには毎回\emph{異なる}値を利用する。
\end{dfn}

\begin{prop}[チャレンジレスポンス認証はリプレイ攻撃を防ぐ]\label{prop:ch6-challenge-response-replay}
チャレンジレスポンス認証では、万が一レスポンスが盗聴されても、そのレスポンスを再利用できない。したがって、チャレンジレスポンス認証はリプレイ攻撃（定義 \ref{dfn:ch6-replay-attack}）を防ぐ方式である。
\end{prop}

\begin{proof}
レスポンスは、チャレンジに対して、あらかじめ決められている署名生成や計算ロジックに基づいて作られた出力である。
そして、チャレンジレスポンス認証では毎回異なるチャレンジを利用する。よって、ある回のレスポンスを盗聴しても、それはその回のチャレンジに対する出力であり、
次の回には別のチャレンジが送られるので、盗聴したレスポンスは新しいチャレンジに対するものとして通用しない。
したがって、盗聴したレスポンスは再利用できず、一度盗聴した通信結果を利用する不正ログインは成功しない。
\end{proof}

パスキー認証では、このチャレンジレスポンス認証と、公開鍵暗号の署名生成・検証（命題 \ref{prop:ch6-public-key-signature}）を組み合わせたものが利用されている。
この観点で簡略化した流れを次の命題で述べる。

\begin{prop}[公開鍵暗号とチャレンジレスポンス認証を組み合わせたパスキー認証]\label{prop:ch6-passkey-challenge-response}
パスキー認証（定義 \ref{dfn:passkey}）は、この観点で簡略化すると次の流れで行われる。
\begin{enumerate}
\item サービスは、\underline{チャレンジ}と呼ばれる一度しか使えないランダムな文字列 $r$ を生成し、セッションに保持しつつパスキー認証を要求する。
\item 認証器は、そのチャレンジ $r$ を含むデータに対して、自身の秘密鍵 $sk$ で署名 $\sigma=\mathrm{Sign}_{sk}(r,\dots)$ を生成する。
チャレンジの値や署名を含む認証レスポンスを生成してサーバーに返す。
\item サーバーは、保持している公開鍵 $pk$ で認証レスポンスの署名を検証する。署名が正しければ、セッションに保存されているチャレンジの値と比較し、一致したら認証成功と判断する。
\end{enumerate}
この流れにより、リプレイ攻撃を防ぎつつ、秘密鍵を持つ認証器による当人認証が実現される。
\end{prop}

\begin{proof}
流れ (1)--(3) は本書の記述そのものである。最後の主張を示す。
まず、署名を生成できるのは秘密鍵 $sk$ の持ち主のみであり、サーバーは保持している公開鍵 $pk$ でそれを検証できる（命題 \ref{prop:ch6-public-key-signature}）。
よって (3) で署名が正しいことが確認できれば、認証レスポンスはその秘密鍵を持つ認証器が作ったものである。
次に、チャレンジは一度しか使えないランダムな文字列であり、サーバーはセッションに保存したチャレンジと認証レスポンス中のチャレンジが一致することを確かめる。
したがって、過去の認証レスポンスを盗聴して送り直しても、チャレンジが一致せず認証は成功しない（命題 \ref{prop:ch6-challenge-response-replay}）。
以上より、この流れはチャレンジレスポンス認証と公開鍵暗号の署名生成・検証を組み合わせたものとして機能する。
\end{proof}

この流れを図示すると図 \ref{fig:ch6-passkey-cr} のようになる（本書の図「公開鍵暗号とチャレンジレスポンス認証を組み合わせたパスキー認証」に対応する）。

\begin{figure}[H]
\centering
\begin{tikzpicture}[>=Stealth, font=\small]
\node[draw, rounded corners, minimum width=30mm, minimum height=8mm, align=center] (sv) at (0,0) {サービス（サーバー）\\公開鍵 $pk$ を保持};
\node[draw, rounded corners, minimum width=30mm, minimum height=8mm, align=center] (au) at (8,0) {認証器\\秘密鍵 $sk$ を保持};
\draw (sv) -- (0,-4.4); \draw (au) -- (8,-4.4);
\node[draw, fill=white, font=\footnotesize, align=center] at (0,-1.1) {(1) チャレンジ $r$ を生成し\\セッションに保持};
\draw[->] (1.6,-1.1) -- node[above] {パスキー認証を要求（$r$）} (8,-1.1);
\node[draw, fill=white, font=\footnotesize, align=center] at (8,-2.3) {(2) $r$ を含むデータに\\$sk$ で署名};
\draw[->] (6.5,-2.3) -- node[above] {認証レスポンス（$r$、署名）} (0,-2.3);
\node[draw, fill=white, font=\footnotesize, align=center] at (0,-3.6) {(3) $pk$ で署名を検証し\\保存した $r$ と比較};
\end{tikzpicture}
\caption{公開鍵暗号とチャレンジレスポンス認証を組み合わせたパスキー認証}\label{fig:ch6-passkey-cr}
\end{figure}

%----------------------------------------------------------------------
\subsubsection{Webブラウザの果たす役割}

ここまで、フロントチャネルの通信やパスキー認証の説明に Web ブラウザが繰り返し登場した。
この項では、Web ブラウザが ID 管理機能を支える重要な構成要素として担っている役割を整理する。

\begin{prop}[ID 管理機能における Web ブラウザの役割]\label{prop:ch6-browser-roles}
Web ブラウザは、ID 管理機能を支える重要な構成要素として、次の役割を担っている。
\begin{enumerate}
\item \underline{Cookie/Web Storage：}セッショントークンやログイン状態のユーザー情報などを安全に保持する。
\item \underline{セキュリティポリシーの強制：}\underline{CORS}（Cross-Origin Resource Sharing）や \underline{CSP}（Content Security Policy）といった仕組みにより、
安全でない通信や意図しないスクリプトの実行を防ぐ。
\item \underline{ブラウザ API：}サービスが呼び出すことで、Web ブラウザ/OS と連動した処理を実行できる。
パスキー認証では \underline{Web Authentication API} というブラウザ API を呼び出すことで、認証器と連動した安全な当人認証が実現できる。
\end{enumerate}
\end{prop}

\begin{proof}
各項目について、それがどのように ID 管理機能を支えるかを述べる。
まず (1) について。フロントチャネルの通信（命題 \ref{prop:ch6-front-channel-session}）で見たように、サービスは Web ブラウザが保持する Cookie からログインセッションの情報にアクセスする。
Web ブラウザがセッショントークンやログイン状態のユーザー情報を安全に保持することで、セッション管理（定義 \ref{dfn:session-management}）が成り立つ。
次に (2) について。CORS や CSP によって安全でない通信や意図しないスクリプトの実行が防がれるので、ID 管理機能が扱う情報が不正な通信やスクリプトにさらされることを防げる。
最後に (3) について。パスキー認証で用いられる認証器はデバイス側にあり、サービスは Web Authentication API を呼び出すことで、Web ブラウザ/OS を通して認証器と連動した当人認証を行える。
以上より、Web ブラウザはこれらの役割を通して ID 管理機能を支えている。
\end{proof}

\begin{rem}[ブラウザごとの差異]\label{rem:ch6-browser-differences}
命題 \ref{prop:ch6-browser-roles} の各機能の導入状況は、Web ブラウザの種類やバージョンによって異なる。
そのような状況の中で、ID 管理機能はさまざまなプラットフォーム／OS において安定した動作が求められる。
\end{rem}

%----------------------------------------------------------------------
\subsubsection{デスクトップアプリ、ネイティブアプリとの連動}

この節の冒頭（命題 \ref{prop:ch6-web-idm-aggregation}）で、ネイティブアプリケーションからも Web 上の ID 管理機能を利用できることを述べた。
この項では、デスクトップアプリやスマートフォンアプリから Web 上の ID 管理機能を利用する際に、どのようにブラウザを呼び出すのが推奨されるかを述べる。

\begin{dfn}[AppAuth パターンと WebView]\label{dfn:ch6-appauth}
デスクトップアプリやスマートフォンアプリから Web 上の ID 管理機能を利用する際に、OS の標準ブラウザを呼び出す方式を\underline{AppAuth パターン}という。
これに対し、アプリ内に埋め込まれたブラウザを\underline{WebView}（埋め込みブラウザ）という。
\end{dfn}

\begin{prop}[AppAuth パターンの推奨]\label{prop:ch6-appauth-recommended}
デスクトップアプリやスマートフォンアプリから Web 上の ID 管理機能を利用する場合、アプリ内に WebView を呼び出す方法よりも、AppAuth パターンが推奨されている。
\end{prop}

\begin{proof}
2 つの方法を比較する。
まず、アプリ内に WebView を呼び出す方法では、一部の機能が制限されたり、ログインセッションの管理が個別に行われたりすることがある。
次に、AppAuth パターンでは、ユーザーは普段使い慣れたブラウザで操作を行う。そのため、利便性の向上や安全性の担保が期待できる。
よって、AppAuth パターンが推奨される。
\end{proof}

\begin{rem}[外部 IdP による外部ブラウザの推奨・強制]\label{rem:ch6-external-idp-browser}
特に、外部 IdP（定義 \ref{dfn:idp}）を利用した ID 連携（定義 \ref{dfn:identity-federation}）によるログインにおいては、外部 IdP の目線で見ると、
他社アプリ内の WebView でログイン処理が行われることによるセキュリティリスクが懸念される。
そのため、外部 IdP が外部ブラウザやセキュアなブラウザコンポーネントの利用を推奨、ないし強制している場合がある。
\end{rem}

%======================================================================
\subsection{ID管理機能の開発スタイル}\label{subsec:ch6-development-styles}

前節で述べた技術の上に、実際に ID 管理機能をどう作るかには、いくつかのスタイルがある。
この節では、スタイルを選ぶ際の判断基準を述べたうえで、スクラッチ実装、既存のライブラリやフレームワークの利用、ID 管理の製品・IDaaS の利用の 3 つを、
それぞれのメリット・デメリットとともに説明し、最後に比較表にまとめる。

%----------------------------------------------------------------------
\subsubsection{汎用性と自由度のバランス}

この項では、開発スタイルを選ぶ際の判断基準と、それぞれのスタイルがその基準のもとでどのような位置にあるかの概観を述べる。

ID 管理機能の開発スタイルはいくつかある。どのスタイルを選ぶかの判断基準として、次の 2 つがある。

\begin{dfn}[開発スタイルの判断基準]\label{dfn:ch6-style-criteria}
ID 管理機能の開発スタイルを選ぶ際の判断基準として、次の 2 つを考える。
\begin{enumerate}
\item \underline{機能やデザイン、UX などの自由度}
\item \underline{実装コストと開発期間}
\end{enumerate}
\end{dfn}

\begin{prop}[自由度とコストのトレードオフ]\label{prop:ch6-freedom-cost-tradeoff}
定義 \ref{dfn:ch6-style-criteria} の 2 つの基準について、各スタイルは次のような位置にある。
\begin{enumerate}
\item すべてを自社で開発すると、自由度は最大化されるが、実装コストと開発期間は大きなものになる。
\item ID 管理機能のための各プログラミング言語のライブラリを利用すると、細かい機能単位で実装コストを下げて開発期間を短縮できる。
\item 単体で Web アプリケーションとして動作するプロダクトや外部サービスを利用すると、実装コストや開発期間を大きく下げられる。
その分、その製品やサービスの仕様に合わせなければならないという制約や、既存機能との調整が必要になったり、運用コストも発生したりする場合がある。
\end{enumerate}
\end{prop}

\begin{proof}
(1) について。すべてを自社で開発すれば、機能やデザイン、UX のすべてを自社で決められるので、自由度は最大になる。
一方で、すべての機能を自社で実装する必要があるので、実装コストと開発期間は大きくなる。

(2) について。実際には、第三者が提供する各プログラミング言語のライブラリを利用すれば、そのライブラリが担う機能の単位で自社の実装を省ける。
よって、細かい機能単位で実装コストを下げ、開発期間を短縮できる。

(3) について。単体で Web アプリケーションとして動作するプロダクトや外部サービスは、ID 管理機能をまとめて提供するので、実装コストや開発期間を大きく下げられる。
その分、機能はその製品やサービスの仕様によって決まるので、それに合わせなければならないという制約が生じ、既存機能との調整が必要になったり、
その製品やサービスを動かし続けるための運用コストが発生したりする場合がある。
\end{proof}

\begin{rem}[開発スタイルの選択]\label{rem:ch6-style-choice}
したがって、サービスの規模、ID 管理機能への要件、開発期間、開発/運用リソースといったさまざまな要素を考慮して、どのスタイルを選択するかを決める必要がある。
\end{rem}

%----------------------------------------------------------------------
\subsubsection{スクラッチ実装}

この項では、最も自由度の高いスタイルであるスクラッチ実装を定義し、その歴史的な位置づけとメリット・デメリットを述べる。

\begin{dfn}[スクラッチ実装]\label{dfn:ch6-scratch}
ID 管理機能を、ライブラリや外部サービスに頼らずにゼロから開発するスタイルを\underline{スクラッチ実装}という。
\end{dfn}

\begin{rem}[スクラッチ実装の減少]\label{rem:ch6-scratch-history}
オンラインサービスに会員機能が登場したころは、ほとんどのサービスがスクラッチ実装であった。
しかし、標準化仕様の採用や、各プログラミング言語の Web アプリケーションフレームワークの登場によって、すべてをスクラッチ実装するスタイルは減った。
\end{rem}

\underline{スクラッチ実装のメリット：}
自社のサービスの要件に合わせて、データ構造から UI/UX まで、すべてを完全に自由に設計できる。

\underline{スクラッチ実装のデメリット：}

\begin{prop}[スクラッチ実装のコスト]\label{prop:ch6-scratch-cost}
スクラッチ実装では、開発と運用のコストはそれなりに大きくなる。
\end{prop}

\begin{proof}
まず、OpenID Connect や FIDO2 といった複雑な標準仕様（定義 \ref{dfn:federation-protocol}、定義 \ref{dfn:fido-auth}）への準拠や、常に変化する脅威への対策などを、
すべて自社で行わなければならず、そのためには高度な専門知識が不可欠である。
次に、それらを自社でゼロから実装するので、脆弱性を生み出すリスクも高まる。
よって、専門知識の確保と脆弱性への対応のために、開発と運用のコストはそれなりに大きくなる。
\end{proof}

%----------------------------------------------------------------------
\subsubsection{既存のライブラリやフレームワークの利用}

この項では、スクラッチ実装と製品・サービスの利用との中間に位置する、既存のライブラリやフレームワークを利用するスタイルを述べる。

\begin{dfn}[既存のライブラリやフレームワークの利用]\label{dfn:ch6-library}
各プログラミング言語で提供されている、認証・認可処理をサポートするライブラリやフレームワークを活用して ID 管理機能を開発するスタイルを、
\underline{既存のライブラリやフレームワークの利用}という。
\end{dfn}

\begin{exa}[認証・認可のライブラリやフレームワーク]\label{exa:ch6-libraries}
例として、Java の Spring Security、.NET の IdentityModel、Node.js の Passport.js などがある。
\end{exa}

\begin{prop}[既存ライブラリの利用は多くの場合で最もバランスがよい]\label{prop:ch6-library-balance}
既存のライブラリやフレームワークの利用は、多くの Web アプリケーション開発において、最もバランスの取れた選択肢となるケースが多い（というのが本書の著者の印象である）。
\end{prop}

\begin{proof}
まず、標準化仕様やベストプラクティスがまとめられているような機能については、すでに第三者が開発してさまざまなサービスで利用されているライブラリを利用することで、
開発／運用コストを小さくできる。
次に、ライブラリを適用するのは主にバックエンドの処理であり、それ以外の部分は自社で自由に作れるので、自由度は保たれる。
よって、自由度を保ったまま開発コストを下げられるので、定義 \ref{dfn:ch6-style-criteria} の 2 つの基準の間で最もバランスの取れた選択肢となるケースが多い。
\end{proof}

\underline{既存ライブラリやフレームワークのメリット：}
スクラッチ実装に比べて開発工数を大幅に削減できる。標準プロトコルに準拠した実装が容易になり、コミュニティによって発見された脆弱性への対応も迅速に行える。

\underline{既存ライブラリやフレームワークのデメリット：}
ライブラリの仕様や設計思想に沿った実装が求められるため、一定の制約は発生する。また、ライブラリ自体の学習コストも必要である。

%----------------------------------------------------------------------
\subsubsection{ID管理の製品、IDaaSの利用}

この項では、実装コストと開発期間を最も大きく下げられるスタイルである、ID 管理の製品や IDaaS の利用を述べる。

\begin{dfn}[IDaaS と ID 管理の製品の利用]\label{dfn:ch6-idaas}
ID 管理機能を提供するクラウドサービスを\underline{IDaaS}（Identity as a Service）という。
Keycloak のような ID 管理機能のための OSS 製品や、Okta や Amazon Cognito といった IDaaS を利用して ID 管理機能を実現するスタイルを、
\underline{ID 管理の製品、IDaaS の利用}という。
\end{dfn}

\begin{rem}[ID 管理の製品・IDaaS の普及]\label{rem:ch6-idaas-adoption}
ID 管理の製品や IDaaS は、一通りの機能をサポートするものから、既存の ID 管理システムと連携して認証認可の機能を実装できるものまで、
さまざまな要件に対応できるようになった。このことから、エンタープライズ向け、コンシューマー向けの両方で利用が進んでいる。
\end{rem}

\underline{ID 管理の製品や IDaaS のメリット：}
ID 管理機能で提供する一通りの機能を、手動の設定や API 呼び出しにより迅速に組み込める。
セキュリティインフラの運用や法規制への対応（GDPR など）、最新仕様のサポートをベンダーに任せることができる。
そのため、自社の開発チームは本来のビジネスロジック開発に集中できる。

\underline{ID 管理の製品や IDaaS のデメリット：}
月額利用料などのランニングコストが発生する場合がある。UI などは提供される機能の範囲内でしかカスタマイズできず、
他のクラウドサービスとの関係によっては、特定のベンダーに依存する（\underline{ベンダーロックイン}）可能性もある。

%----------------------------------------------------------------------
\subsubsection{開発スタイルの比較}

最後に、この節で述べた 3 つの開発スタイルのメリットとデメリットを、表 \ref{tab:ch6-dev-styles} にまとめる。
表の内容はすべて上の各項の記述によるものである。

\begin{table}[H]
\centering
\caption{ID 管理機能の開発スタイルの比較}\label{tab:ch6-dev-styles}
\begin{tabular}{|p{26mm}|p{24mm}|p{40mm}|p{40mm}|}
\hline
スタイル & 例 & メリット & デメリット \\ \hline
スクラッチ実装 & （ライブラリや外部サービスに頼らない） &
データ構造から UI/UX まですべてを完全に自由に設計できる（自由度が最大） &
複雑な標準仕様への準拠や変化する脅威への対策に高度な専門知識が不可欠；脆弱性を生むリスクが高く、開発・運用コストが大きい \\ \hline
既存のライブラリやフレームワークの利用 & Spring Security（Java）、IdentityModel（.NET）、Passport.js（Node.js） &
開発工数を大幅に削減；標準プロトコル準拠の実装が容易；コミュニティが発見した脆弱性に迅速に対応できる；自由度を保ったまま開発コストを下げられる &
ライブラリの仕様や設計思想に沿う必要があり一定の制約が発生；ライブラリ自体の学習コスト \\ \hline
ID 管理の製品、IDaaS の利用 & Keycloak（OSS）、Okta、Amazon Cognito（IDaaS） &
一通りの機能を設定や API 呼び出しで迅速に組み込める；セキュリティインフラ運用・法規制対応（GDPR など）・最新仕様のサポートをベンダーに任せられる；ビジネスロジック開発に集中できる &
月額利用料などのランニングコスト；UI などは提供範囲内でしかカスタマイズできない；ベンダーロックインの可能性 \\ \hline
\end{tabular}
\end{table}

\begin{thebibliography}{9}
\bibitem{Ito2026}
いとう りょう，『はじめてのデジタルアイデンティティ―Webサービスに欠かせない認証・認可・ID管理』，技術評論社，2026．
\end{thebibliography}

\end{document}
